\documentclass[11pt,a4paper]{article}
\usepackage[hmargin=26mm,vmargin=25mm]{geometry}
\usepackage{iftex}
\usepackage{amsmath,amsthm,mathtools,booktabs,array}
\ifPDFTeX
  \usepackage[T1]{fontenc}
  \usepackage{lmodern,amssymb}
\else
  \usepackage{amssymb,fontspec,unicode-math}
  \setmainfont{Latin Modern Roman}
  \setmathfont{Latin Modern Math}
  \AtBeginDocument{\renewcommand{\Join}{\mathbin{\bowtie}}}
\fi
\usepackage{enumitem}
\SetEnumitemKey{titledsteps}{
    label=\arabic*., leftmargin=1.85em, labelsep=.55em,
    topsep=4pt plus 1pt minus 1pt, partopsep=0pt,
    itemsep=3pt plus 1pt minus 1pt, parsep=2pt plus 1pt,
    beginpenalty=10000
}
\usepackage{xcolor}
\usepackage{tikz-cd}
\usetikzlibrary{arrows.meta,calc}
\usepackage[section]{placeins}
\tikzcdset{diagrams={labels={font=\small}}}
\setlength{\abovecaptionskip}{6pt}
\tikzset{
    figwire/.style={-{Latex[length=1.6mm]},line width=.55pt},
    figbox/.style={draw,rounded corners=1.5pt,fill=black!3,
        inner xsep=6pt,inner ysep=4pt,align=center},
    figlabel/.style={font=\footnotesize,inner sep=2pt,fill=white}
}
\tikzset{
    figline/.style={line width=.55pt},
    cutmark/.style={line width=1pt},
    dot/.style={circle,fill=black,inner sep=1.35pt}
}


\usepackage[colorlinks=true,linkcolor=black,citecolor=blue!44!black,urlcolor=blue!44!black]{hyperref}
\usepackage{aliascnt}
\usepackage[nameinlink,capitalise,noabbrev]{cleveref}
\newtheorem{theorem}{Theorem}[section]
\newaliascnt{proposition}{theorem}
\newtheorem{proposition}[proposition]{Proposition}
\aliascntresetthe{proposition}
\newaliascnt{lemma}{theorem}
\newtheorem{lemma}[lemma]{Lemma}
\aliascntresetthe{lemma}
\newaliascnt{corollary}{theorem}
\newtheorem{corollary}[corollary]{Corollary}
\aliascntresetthe{corollary}
\newtheorem{comparison}{Proposition}
\renewcommand{\thecomparison}{\Alph{comparison}}
\crefname{comparison}{Proposition}{Propositions}
\Crefname{comparison}{Proposition}{Propositions}
\theoremstyle{definition}
\newaliascnt{definition}{theorem}
\newtheorem{definition}[definition]{Definition}
\aliascntresetthe{definition}
\newaliascnt{construction}{theorem}
\newtheorem{construction}[construction]{Construction}
\aliascntresetthe{construction}
\newaliascnt{example}{theorem}
\newtheorem{example}[example]{Example}
\aliascntresetthe{example}
\theoremstyle{remark}\newaliascnt{remark}{theorem}
\newtheorem{remark}[remark]{Remark}
\aliascntresetthe{remark}
\crefname{construction}{Construction}{Constructions}
\Crefname{construction}{Construction}{Constructions}
\providecommand{\llbracket}{\mathopen{[\mkern-3mu[}}
\providecommand{\rrbracket}{\mathclose{]\mkern-3mu]}}
\newcommand{\E}{\mathbf{Eff}}
\newcommand{\N}{\mathbb N}
\newcommand{\Q}{\mathbb Q}
\newcommand{\R}{\mathbb R}
\newcommand{\Fin}{\operatorname{Fin}}
\newcommand{\Idl}{\operatorname{Idl}}
\newcommand{\V}{\mathsf V_{\le1}}
\newcommand{\Hom}{\operatorname{Hom}}
\newcommand{\Sym}{\operatorname{Sym}}
\newcommand{\id}{\mathrm{id}}
\newcommand{\ev}{\operatorname{ev}}
\newcommand{\enc}{\operatorname{enc}}
\newcommand{\para}{\mathsf{Para}}
\newcommand{\lens}{\mathsf{Lens}}
\newcommand{\NN}{\mathcal N}
\newcommand{\dif}{\mathrm d}
\newcommand{\att}{\operatorname{Att}}
\newcommand{\op}{\mathrm{op}}
\newcommand{\Sq}{\operatorname{Sq}}
\newcommand{\pair}[2]{\langle #1,#2\rangle}
\newcommand{\norm}[1]{\lVert #1\rVert}
\newcommand{\onto}{\twoheadrightarrow}
\newcommand{\tensor}{\mathbin{\odot}}
\newcommand{\ang}[2]{\langle #1,#2\rangle}
\newcommand{\STF}{\operatorname{STF}}
\newcommand{\Spec}{\operatorname{Spec}}
\setlength{\parskip}{1pt plus .5pt minus .5pt}
\setlength{\textfloatsep}{14pt plus 3pt minus 3pt}
\setlength{\intextsep}{10pt plus 2pt minus 2pt}
\clubpenalty=5000
\widowpenalty=5000
\displaywidowpenalty=5000
\setlength{\emergencystretch}{2em}
\allowdisplaybreaks[1]
\setcounter{tocdepth}{1}
\raggedbottom
\numberwithin{equation}{section}
\hypersetup{bookmarksdepth=2,pdfauthor={Conor Sheridan},
  pdftitle={Reading neural representations as logical queries: Compilation, readback and operational grounding over algebraic databases}}
\title{Reading neural representations as logical queries\\\large Compilation, readback and operational grounding over algebraic databases}
\author{Conor Sheridan}
\date{}
\begin{document}

\maketitle
\begin{abstract}
Which logical queries can be recovered from a neural program, and what can a hidden register tell us about the environment in which the program acts? For a rational sum--product program on finite instances of a finite entity category, we prove a unique expansion in decorated embedding queries. Exact evaluations at finitely many shapes and inputs recover its coefficients given bounds on degree and witness size. The expansion retains rows used only as witnesses and composes by substitution. When the register list itself is available, a separate structural readback also recovers gate occurrences, shared parameter addresses and wiring.

For a deterministic interaction, we construct a common quotient of reached port values and environmental responses to action words. Its predicates are precisely the properties readable through both observations. Closing the defining identifications under continuation gives the largest family of common predicates stable under action inverse images. On a finite transition table we compute a port formula and a sensor-experiment formula for each quotient class.

Opening a register to independent replacements gives a response quotient, which identifies replacements with equal downstream responses in every bypass context. For frozen rational affine--ReLU scores on compact rational boxes, we express the normalized supremum of the response discrepancy as a rational \L{}ukasiewicz term. A two-unit network illustrates the distinction between recovering all decisions on reachable inputs and preserving scores under intervention. The same keyed presentations give compositional cotangent pullbacks at shared parameters and a target-indexed Kan extension characterizing all parameter tables that attain least error.
\end{abstract}
\clearpage
\begingroup
\makeatletter
\small
\setlength{\parskip}{0pt}
\renewcommand*{\l@section}{\@dottedtocline{1}{0em}{2.4em}}
\renewcommand*{\l@part}[2]{\addpenalty{-200}\addvspace{.65\baselineskip}%
  {\bfseries\noindent #1\hfill #2\par}\nobreak}
\tableofcontents
\makeatother
\endgroup
\clearpage

\section{Which neural representations admit a logical readback?}\label{sec:intro}
A keyed neural program stores a field of values indexed by database rows at each register. Reading a register back as a query means recovering an expression for that field in the input fields and foreign keys. A hidden register presents a morphism \(h:X\to Z\). Cutting its outgoing wires leaves a continuation \(C:X\times Z\to O\), where the first factor retains everything that bypasses the register. Ordinary execution is the restriction of \(C\) along the graph of \(h\). For a proposed decoder \(d:Z\to D\), corestrict to its image \(D_0\). On the left, $F$ is the ordinary output. On the right, a factor through the decoder exists precisely when its code suffices to compute the response in every bypass context. Such a factor is unique.
\[
\begin{tikzcd}[column sep=large]
 X\arrow[r,"{\langle 1,h\rangle}"]\arrow[dr,"F"']&X\times Z\arrow[d,"C"]\\
 &O
\end{tikzcd}
\qquad
\begin{tikzcd}[column sep=small]
 X\times Z\arrow[rr,two heads,"1\times d"]\arrow[dr,"C"']&&X\times D_0\arrow[dl,dashed,"\bar C"]\\
 &O
\end{tikzcd}
\]
Learning a predicate on reachable values uses the first triangle. Testing whether the response to every admissible replacement factors through the decoder uses the second. The different domains already matter for a register with two rectifiers (\cref{ex:trained-discrepancy}).

On a database shape \(A\), write \(h_A\) and \(C_A\) for the two maps, \(x\) for the independent input cells and \(\bar b\) for an output key. A vocabulary of hidden attributes \(E_m\) determines which predicates are available at the port. For an interaction, let $p$ be a reached history, $h(p)$ its port value, and $q(p)$ its environmental response, with $q(p)(w)$ the sensor observation after the action word $w$. The interpretation questions are
\[
\begin{aligned}
\text{query equality:}\quad&
  \forall A,\bar b,x\quad F_A(x)_{\bar b}=G_A(x)_{\bar b},\\
\text{explanation of a target:}\quad&
  \forall x\in X\quad P(x)\ \longleftrightarrow\ B\big((E_m(h_A(x)))_{m\in M}\big),\\
\text{interventional equivalence:}\quad&
  z\sim z'\ \longleftrightarrow\ \forall x\in X\quad C_A(x,z)=C_A(x,z'),\\
\text{operational grounding:}\quad&
  \text{which predicates }P\text{ admit }P=h^*U=q^*V\,?
\end{aligned}
\]
\Cref{thm:support-readback} answers the first question in the rational sum--product language. Its inverse recovers a unique expansion in decorated embeddings from finitely many exact evaluations at the supplied bounds. A pattern can contain a row that supplies no numerical factor: the row still affects the number of witnesses, and varying the database shape reveals its coefficient. The structural inverse of \cref{thm:compiler-struct} instead starts with a register list and reads off gate occurrences and parameter-owner keys. The target and intervention tests start by naming a field or opening one of those occurrences. A reachable explanation factors through the chosen profile (\cref{thm:concept-expressibility}); a sufficient intervention decoder determines the response image (\cref{prop:response-quotient}). For bounded rational rectifier scores we also compute the greatest discrepancy between responses (\cref{cor:dmv-response}).

\paragraph{Which concepts does an interaction make available?}
For the fourth question the data are the port, the actions and the observed sensor responses. The target $P$ ranges over the predicates to be discovered. We call an interpretation \emph{operationally grounded} when its two descriptions agree on every reached history. \Cref{thm:ground-common} classifies all such predicates, and \cref{thm:ground-stable} identifies the ones whose pullbacks along every action word still admit port descriptions. In \cref{ex:ground-pulse}, two histories have the same port value and current sensor token. After one more action their sensor tokens differ. This is a finite witness to a distinction the environment makes and the port has lost.

\paragraph{Two backward maps for a parameterized program.}
Gavranovi\'c develops parameter wires and bidirectional information flow through parametric weighted optics \cite{gavranovic-thesis}. Here we use the parameterized architectures of Gavranovi\'c et al.\ \cite{cdl} and the bimorphic lenses of Cruttwell, Gavranovi\'c, Ghani, Wilson and Zanasi \cite{lens:deep}. For objects \(A,A',B,B'\) in a cartesian category, a lens \(L:(A,A')\to(B,B')\) has forward map \(l:A\to B\) and backward map \(l^\flat:A\times B'\to A'\). Here a prime denotes the type of feedback. For a second lens \(K\), putting \(b=l(a)\) gives
\[
 (KL)^+(a)=k(b),\qquad (KL)^\flat(a,c')=l^\flat(a,k^\flat(b,c')) .
\]
Parallel composition pairs both forward maps and pairs their backward maps at their respective inputs. Its unit is \((1,1)\). Retaining \(b\) gives associativity by substitution, and componentwise pairing gives interchange. A parametric lens has type
\((P,P')\otimes(X,X')\to(Y,Y')\): its backward map returns feedback at both the parameter and the input. Freezing \(p\) gives a forward map \(f_p:X\to Y\). A predicate on \(Y\) pulls back along \(f_p\) (\cref{sec:hidden-concepts}). At a differentiable execution an output covector pulls back jointly to \(P\) and \(X\) (\cref{thm:para-lens}). These choices of feedback give different comparison maps for parallel composition, as \eqref{eq:predicate-tensor} and \cref{neural:thm:cotangent} show.
\paragraph{Earlier translations between logic and networks.}
Garcez and Zaverucha's C-IL$^2$P translates a propositional program into a trainable network and extracts revised rules \cite{cil2p}. TensorLog compiles probabilistic logic queries into differentiable computations \cite{tensorlog}; DeepProbLog embeds learned neural predicates in probabilistic programs \cite{deepproblog}. Pugh et al.\ present a dense floating-point architecture by a coherent category and obtain inference by restriction along the input--output interface, including for networks with fixed or tied weights \cite[Theorems 4.13--4.14]{pugh}. \Cref{prop:coherent-port} constructs the comparison with our keyed field theory on models and their transformations. Varying the concrete finite models gives the evaluation family in \cref{thm:support-readback}. Soeteman et al.'s neuro-relational programs are simultaneously relational queries and neural architectures. They give sum--product semantics to joins and aggregation and characterize the expressiveness of ReLU-based fragments using counting logic \cite{nrp}. Their conjunction rule records both whether an answer tuple exists and its embedding. A weighted sum at an output can be zero both for an empty fibre and for a fibre with one zero-valued contribution. To compare the two results one can return the witness count alongside the sum; the count is positive exactly when the join has a witness. The schema collage, provenance semirings and polynomial substitution supply the structural semantics used here \cite{algebraic-databases,green-provenance,poly:gambino-kock}.

\paragraph{Which Kan extensions occur?}
Along the inclusion of the object of the action category, left extension propagates identifications through all continuations and right extension assembles the sensor readings of all future experiments (\cref{prop:ground-action-kan,thm:ground-stable}). For fitting, Pugh et al.\ relate a universal inference extension to least error \cite[Theorem 3.15]{pugh-kan}. Our inference map on a sample has a target-indexed category of loss comparisons. Its universal arrows classify every attaining table, including ties (\cref{thm:fit-kan}). Even equality of the entire predicted function leaves room for two hidden ports with different experimentally grounded concepts (\cref{ex:fit-ground-distinction}).

\paragraph{What an exact correspondence observes.}
Tena Cucala et al.\ compute an equivalent Datalog program for the predicted facts of a monotonic max--sum GNN \cite{gnn}. Yang, Chiang and Angluin characterize a strictly masked hard-attention transformer class without positional embeddings by linear temporal logic \cite{hard-attn-ltl}. Both give exact comparisons on their stated observables. Zhang and B\"olcskei use substitution graphs and many-valued logic to identify ReLU networks by the functions they compute on a cube \cite{zhang-bolcskei-mv}. Their rational score calculus applies to the frozen open term in \cref{cor:dmv-response}. Our canonical inverse observes numerical values at typed keys and permits the database shape to vary. Puny et al.\ give a multigraph basis for bounded-degree equivariant graph polynomials at fixed vertex count \cite{graph-poly-basis}; the relation to our embedding basis is computed in \cref{sec:semantic}. Homomorphism bases, equivalence of aggregate queries and exact identification of conjunctive queries from examples have independent precedents \cite{motif-basis,aggregate-equivalence,cq-learning}. Our result combines foreign-key-closed witnesses, finite evaluation bounds and an executable normal form for this rational fragment. Exhaustive reconstruction gives a decision procedure; its cost grows with the shapes and monomials enumerated.

\paragraph{Concepts in a hidden representation.}
Hirth and Hanika construct many-valued and symbolic conceptual views of neural networks, compare representations, and extract rules from neuron activations \cite{concept-neural-views}. Yeh et al.\ quantify how completely a concept vocabulary explains predictions and propose a method for discovering concepts \cite{concept-completeness}. For a supplied target we begin with a specified vocabulary of observations at a typed register. Formal concept analysis supplies the extent--intent correspondence \cite{fca-foundations}. Here the compiler constructs its incidence columns as formulas in the original database attributes. It can certify a description throughout a declared domain, find inputs that defeat the vocabulary, and compute all thresholds realizing a specified Boolean template along fixed directions. A semantic name such as \emph{large discrepancy} enters through its target predicate. Agreement with that predicate is a mathematical claim about the given network. In \cref{sec:ground-operational} the environment supplies another observation language. A pair of formulas can then identify a concept before it has been given a task label.

\paragraph{What makes an explanation mechanistic.}
Tracr compiles programs into transformers with known computations, and InterpBench tests interpretability methods on models with known circuits \cite{tracr,interpbench}. Friedman, Wettig and Chen train a modified transformer that can be converted to a discrete program \cite{transformer-programs}. Huang et al.\ decompile trained transformers into RASP and use interventions to identify sufficient subprograms \cite{huang-rasp}. Geiger et al.\ formalize agreement between mechanisms under interventions \cite{causal-abstraction}. Our question is what can be certified at proposed, aligned ports. In the polynomial fragment the intervened operation has a unique symbolic query across shapes. In the fixed-shape rational affine--ReLU fragment agreement of a proposed representation or intervened prediction is decidable by linear inequalities. This uses established verification methods for rectifier networks \cite{reluplex,relu-geometry}. At the same port we can therefore compare two descriptions with different claims to fidelity. One reproduces decisions on the reachable states. The other identifies precisely those replacement states having equal downstream responses. The trained example computes both.

\paragraph{Universal properties used in the readback.}
The choice of what to expose at a port determines which universal construction answers the interpretation question.
\begin{center}
\small
\begin{tabular}{@{}p{.29\linewidth}p{.46\linewidth}p{.17\linewidth}@{}}
\toprule
What is given & What determines the answer & Result\\
\midrule
Evaluations across shapes & Unique finite expansion in decorated embedding queries at the supplied bounds & \cref{thm:support-readback}\\[2pt]

A presented architecture & Models over a field interpretation are parameter tables; coherent restriction selects a port & \cref{prop:coherent-port}\\[2pt]
A target and its profile & Left and right adjoints to inverse image give its two optimal profile bounds & \cref{prop:profile-adjoints}\\[2pt]
A cut and its contexts & Reflection into pairings with distinguishable replacement states & \cref{thm:fibred-reflection}\\[2pt]
A port and an environmental response & Common quotient of observations and its quotient under word actions & \cref{thm:ground-common,thm:ground-stable}\\[2pt]
A cell and its parameter & Unique algebra morphism from free syntax, evaluated with the parameter shared & \cref{poly:thm:parametric}\\[2pt]
A first variation & The representing module of derivations and its pullback along the program & \cref{neural:lem:kahler}\\[2pt]
A sample and its target & Universal arrows to evaluation classify all tables attaining least error & \cref{thm:fit-kan}\\[2pt]

A solved equation & The solution graph projects isomorphically onto its execution domain & \cref{thm:solved-composition}\\
\bottomrule
\end{tabular}
\end{center}
Part~\ref{part:functors} gives structural and semantic readback. Part~\ref{part:ports} moves from deciding whether a specified target is readable to discovering the predicates shared by a port and an environment. It then compiles the implications of finite observation vocabularies. Parts~\ref{part:algebras} and~\ref{part:feedback} construct the underlying syntax and its differential semantics. Part~\ref{part:solutions} identifies the observables of concrete architectures. Dependent value vectors, for example, can expose the same output under different attention weights \cite{attn-identifiability}. Part~\ref{part:learners} takes the resulting feedback into an optimizer state.

\part{Functors and normal forms for readback}\label{part:functors}

\section{Presented architectures and structural readback}\label{sec:compiler}
A field is a scalar value at each row of a specified type. Its type determines which joins are legal and where a sum lands. We first compare two ways of writing a computation with these fields. We then ask the stronger question of recovering a query from its values on finite databases (\cref{sec:semantic}). Both constructions keep the schema in the interface.

\subsection{Keyed fields and concrete models}\label{sec:field-language}
The entity category \(C\) has row sorts as objects and foreign keys as arrows. A finite shape \(A:C\to\mathsf{FinSet}\) supplies their carriers and address maps. A type theory \(\mathcal T\) supplies the scalar sorts and operations. An observable of sort \(t\) at \(c\) has a value in \(S(t)\) at every row of \(A(c)\), for a type algebra \(S\). Reading through a foreign key and then applying a scalar operation gives the mixed terms of the schema collage (\cref{thm:schema-presentation,lem:collage}).

Build key contexts \(U\) from entity sorts, a singleton, tuples, finite tags and restrictions by decidable structural key tests. Key terms are foreign-key paths, projections, tuples and their induced maps. A field sort \([U:t]\) has concrete interpretation \(M_A[U:t]=S(t)^{U(A)}\). The finite-product theory \(\mathbb F\) has finite ordered records of field sorts as objects. Its terms read fields, form constant fields, apply the declared pointwise operations and use the following keyed operations for a key map \(f:U\to V\) and arithmetic sort \(R\).
\[
 f^*:[V:R]\to[U:R],\qquad
 \Sigma_f,\Pi_f:[U:R]\to[V:R].
\]
For \(u\in U(A)\), \(v\in V(A)\), their coordinates in \(M_A\) are
\[
 (f^*b)_u=b_{f_A(u)},\qquad
 (\Sigma_f a)_v=\sum_{f_A(u)=v}a_u,\qquad
 (\Pi_f a)_v=\prod_{f_A(u)=v}a_u .
\]
The empty fibres yield zero and one, respectively. A guarded scope takes its eligible table as input (\cref{thm:elaboration}). Composition of terms is substitution, pairing assembles records, and the equations on keys, observables and scalar operations are those of the selected presentation. Independent data fields are variables. A shape with prescribed background attributes uses their given values; the \hyperref[sec:legal-attribute-tables]{equalizer of observable tables} gives the object of compatible assignments when the schema relates them. A parameter owner is a designated input port, and \(\rho^*\theta\) reads it through its key map \(\rho\). Enumerating the rows specializes these operations to finite scalar terms (\cref{app:field-presentations}).

\subsection{Gate occurrences, wiring and the two translations}
An ordered register presentation records each gate occurrence, its ordered operands and all output references. A term in the finite-product theory records the result of substituting its defining gate terms. To recover the wiring we first work with the occurrence presentation.

Fix typed input ports \(\Gamma\) and an ordered output record \(\Delta\). For an instruction list of length \(n\), put \(V_i=\Gamma+\{r_j\mid 1\le j\le i\}\), with the type of \(r_j\) given by instruction \(j\). An instruction has a typed gate label \(g_i\) (including its key map, if any) and an ordered operand list \(\vec s_i\) in \(V_{i-1}\). Repeated operands and unused instruction results are allowed. The output is a record of references \(\vec o\) in \(V_n\) of types \(\Delta\). Define the canonical administrative terms at this interface by
\[
 \mathsf{let}\ r_1=g_1(\vec s_1)\ \mathsf{in}\ \cdots\
 \mathsf{let}\ r_n=g_n(\vec s_n)\ \mathsf{in}\ \vec o .
\]
Only instructions introduce names in this grammar. Aliases, record construction and projections on records are resolved as operations on references. For example, \(\mathsf{let}\ y=x\ \mathsf{in}\ y\) normalizes to \(x\). In \(\mathsf{let}\ y=g(x)\ \mathsf{in}\ (y,y)\) the one instruction and the two output references both survive.

\begin{theorem}[Structural correspondence]\label{thm:compiler-struct}
At each interface \((\Gamma,\Delta)\), compilation \(C\) and readback \(D\) give inverse bijections between canonical administrative terms modulo renaming and ordered keyed register lists modulo renaming. On unrestricted expressions, \(DC\) is administrative normalization. Forgetting topological order identifies lists related by exchanges of independent instructions. The translations preserve port substitution, shared parameter addresses and interpretation on every presented instance.
\end{theorem}
\begin{proof}\leavevmode
\begin{enumerate}[titledsteps]
    \item \textbf{Compile the instruction list.} For a prefix \(L\) and an environment \(e\) of typed references, write \(\mathcal C(t;e,L)=(L',s)\). Variables return \((L,e(x))\). A gate compiles its operands in socket order and appends one instruction with their returned references. A record compiles its components in order, threading the instruction list through them, and returns their tuple of references. A projection selects from that tuple and retains the emitted list. Thus an unused result and two references to a shared result each keep their original gate occurrences. The two binding rules are
\[
\begin{aligned}
 \mathcal C(\mathsf{let}\ x=t\ \mathsf{in}\ u;e,L)
   &=\mathcal C(u;e[x\mapsto s],L')\quad
        \text{where }(L',s)=\mathcal C(t;e,L),\\
 D(L\mathbin{+\!\!+}[r=g(\vec s)],\vec o)
   &=D(L,\ \mathsf{let}\ r=g(\vec s)\ \mathsf{in}\ \vec o).
\end{aligned}
\]

    \item \textbf{Recover the wiring.} In the second line \(D(L,-)\) wraps a continuation in the lets of \(L\). Thus \(CD(L,\vec o)=(L,\vec o)\) by prefix induction. Conversely the compiler emits one instruction for each canonical let and reproduces its operand references, giving \(DC(t)=_\alpha t\) on the grammar above. Connecting two lists concatenates their instructions after translating each downstream external reference through the upstream output record. This is precisely the binding rule for port substitution. Two linear extensions of the same finite dependency order differ by adjacent exchanges of incomparable instructions: move the first desired instruction left across its incomparable predecessors and induct on the suffix.


    \item \textbf{Identify the interpreted fields.} Fix an input instance. Induction on emitted registers identifies each stored field with the value of its defining expression at every row. For a lookup the row read is $f(u)$ on both sides. For a reduction the enumeration on both sides is $f^{-1}(v)$ with the stated empty convention. The scalar and let cases follow by application and substitution. The computation of a typed register can be expanded on a fixed finite instance to scalar gates and wires. In particular sums are affine, and multiplication over $\Q$ or $\R$ is realized by $uv=((u+v)^2-(u-v)^2)/4$. Thus a polynomial field program yields an actual finite affine--square neural circuit at every presented shape. The schema-level instructions retain the information needed for structural readback across shapes.\qedhere
\end{enumerate}
\end{proof}

The theorem operates on the presented gate syntax. To impose an equation between observable terms on the symbolic side, compile both sides and generate the corresponding congruence on programs under the same typed substitutions. Readback carries each generating pair back to its original pair up to administrative normalization; induction on congruence derivations proves the same for every consequence. Thus the two translations also descend to inverse maps of these presented quotients. For an interpretation on a restricted numerical locus, both sides are evaluated on that locus. Different readbacks can have equal values there, as the Boolean example below illustrates.
\paragraph{Sharing can be observed by a cut.}
One program computes \(r=g(x)\) and returns \((r,r)\). A second computes \(r=g(x)\) and \(s=g(x)\) and returns \((r,s)\). Both interpret as the cartesian arrow \(\langle g,g\rangle\). Replacing the first register result by an independent \(z\) gives \((z,z)\) and \((z,g(x))\), respectively. At \(g(x)=x+1,\ x=0,\ z=2\), the answers are \((2,2)\) and \((2,1)\). The designated vertex and its outgoing uses carry precisely the information needed for this operation. The relation between such graphs and cartesian terms is the sharing distinction in the term-graph semantics of Corradini and Gadducci \cite{corradini-gadducci}. Figure~\ref{fig:shared-gate-cut} shows the two cuts.

\begin{figure}[!htbp]
    \centering\small
    \begin{tikzpicture}[x=1cm,y=1cm]
        \begin{scope}
            \node at (1.9,1.30) {One shared gate};
            \node (x) at (0,0) {$x$};
            \node[figbox] (gate) at (1,0) {$g$};
            \draw[figwire] (x) -- (gate);
            \draw[figline] (gate) -- (2.35,0);
            \draw[cutmark] (1.72,-.21) -- (1.86,.21);
            \node[above] at (1.79,.18) {$z$};
            \node[dot] (copy) at (2.35,0) {};
            \draw[figwire] (copy) -- ++(.45,.46) -- (3.45,.46);
            \draw[figwire] (copy) -- ++(.45,-.46) -- (3.45,-.46);
            \node[right] at (3.45,.46) {$z$};
            \node[right] at (3.45,-.46) {$z$};
        \end{scope}
        \begin{scope}[xshift=6.2cm]
            \node at (1.9,1.30) {Two gate occurrences};
            \node (x) at (0,0) {$x$};
            \node[dot] (copy) at (.6,0) {};
            \draw[figline] (x) -- (copy);
            \node[figbox] (first) at (1.92,.46) {$g$};
            \node[figbox] (second) at (1.92,-.46) {$g$};
            \draw[figwire] (copy) -- (1.12,.46) -- (first);
            \draw[figwire] (copy) -- (1.12,-.46) -- (second);
            \draw[figwire] (first) -- (3.45,.46);
            \draw[figwire] (second) -- (3.45,-.46);
            \draw[cutmark] (2.60,.25) -- (2.74,.67);
            \node[above] at (2.67,.63) {$z$};
            \node[right] at (3.45,.46) {$z$};
            \node[right] at (3.45,-.46) {$g(x)$};
        \end{scope}
    \end{tikzpicture}
    \caption{A dot copies its input. Replace the value at the slash by $z$. In the left circuit the cut precedes the copy, so both outputs change. In the right circuit the second gate still receives $x$.}
    \label{fig:shared-gate-cut}
\end{figure}


\paragraph{Substitution at a port.}
Fix an interface and keep the order of the register list. Write \(\mathsf{Term}\) for the canonical terms and \(\mathsf{Reg}\) for these lists, both up to renaming. For a composable pair the comparison is a square of sets of arrows (typed interfaces are suppressed):
\[
\begin{tikzcd}[column sep=large]
 \mathsf{Term}(Y,Z)\times\mathsf{Term}(X,Y)\arrow[r,"\mathrm{subst}"]\arrow[d,"C\times C"']&
 \mathsf{Term}(X,Z)\arrow[d,"C"]\\
 \mathsf{Reg}(Y,Z)\times\mathsf{Reg}(X,Y)\arrow[r,"\mathrm{wire}"']&
 \mathsf{Reg}(X,Z).
\end{tikzcd}
\]
The binding rule proves commutation, including two references to the same register. The empty instruction list implements the identity. Reading the square upwards with \(D\) therefore gives the inverse functor. When two terms run in parallel, exchanging independent instructions identifies the two orders in which their lists can be emitted.

For edge keys \(s,t:E\to V\) and owner key \(\rho:E\to P\), the four instructions and their readback are
\[
 r_1=\rho^*\theta,\quad r_2=s^*x,\quad r_3=r_1r_2,\quad r_4=\Sigma_t r_3,\qquad
 (r_4)_j=\sum_{i\in V}\underbrace{\sum_{e:s(e)=i,\,t(e)=j}\theta_{\rho(e)}}_{K_{ji}}x_i.
\]
A unit feature probe at \(i\) recovers \(K_{ji}\). Two independent parallel weights \((p,q)\) and \((p+\delta,q-\delta)\) give the same column. In the register list the two occurrences and their owner paths are still present. Varying the shape will give the additional observations used in \cref{thm:support-readback}.

\subsection{Parameter tables and interpretations of a port}\label{sec:model-interfaces}
Let \(X,\Theta,Y\) be the data, parameter and output records. Adjoining a parameter table to \(\mathbb F\) gives a finite-product category \(\mathbb F_\Theta\) on the same objects, with
\[
 \mathbb F_\Theta(U,V)=\mathbb F(\Theta\times U,V),\qquad
 g\star f=g\langle\pi_1,f\rangle,\qquad 1_U=\pi_2 .
\]
Thus successive calls both read \(p\) in \(g(p,f(p,u))\). Pairing and terminal maps are pointwise in \(p\). Write \(J(a)=a\pi_2\) for the canonical functor and \(\tau:1\to\Theta\) for the generic table, represented by \(\pi_1:\Theta\times1\to\Theta\). This is the cartesian parameter construction of Lambek and Scott \cite[Part I, Proposition 6.1]{lambek-scott}.

For a list \(N\), form a presentation \(P_N\) over \(\mathbb F_\Theta\) by naming each register \(r_i:X\to V_i\) and giving it its displayed gate definition. In that context the data references are projections, the parameter references are components of \(\tau!\), and earlier names can be reused in any socket. Pair the designated output references to obtain \(F_N:X\to Y\). Write \(|P_N|\) for the category presented by these definitions.

\begin{proposition}[Models and port interpretations]\label{prop:relative-models}
For a finite-product target \(\mathcal D\), models of \(|P_N|\) are equivalent, over the category of product-preserving \(M:\mathbb F\to\mathcal D\), to pairs \((M,p:1\to M\Theta)\). A transformation \(a:M\Rightarrow M'\) lifts precisely when \(a_\Theta p=p'\), and its lift is unique. Given \(t:X\to V\) in \(|P_N|\), the interpretation
\[
 i_t:\mathbb F[o:X\to V]\longrightarrow |P_N|,\qquad o\longmapsto t
\]
selects the port. Restricting a model along \(i_t\) returns its map \(MX\to MV\).
\end{proposition}
\begin{proof}
Eliminate the register names in list order. Each \(r_i\) becomes its gate applied to the interpretations of its operands. This defines an inverse to the map \(\mathbb F_\Theta\to |P_N|\): one composite fixes the original terms, and induction using each defining equality shows that the other fixes every \(r_i\). For \(f:\Theta\times U\to V\), extend \(M\) by
\[
 M_p(f)=M(f)\langle p!,1_{MU}\rangle .
\]
The composition formula above proves functoriality, and pairing proves product preservation. Every arrow of the extension has the form \(J(f)\langle\tau!,1_U\rangle\), so its value is forced by \(M\) and \(p\). Naturality at \(\tau\) gives \(a_\Theta p=p'\), which suffices on every arrow by the same factorization. The free-arrow interface \(\mathbb F[o:X\to V]\) interprets \(o\) as any map of the prescribed type. Substitution of \(t\) therefore gives \(i_t\), and precomposition gives the stated map on models.
\end{proof}

Over a concrete set-valued \(M_A\), the models with that fixed interpretation and identity comparison form the discrete set \(M_A\Theta\). Selecting \(F_N\) returns the whole input--output function at the chosen table; selecting \(r_i\) returns its hidden field. Elimination identifies \(|P_N|\) with \(\mathbb F_\Theta\) even for programs with different gate lists. Specifying \(i_{F_N}\) records the computed arrow, and designating an occurrence in the presentation specifies its cut.

\begin{proposition}[Coherent reduct at a named port]\label{prop:coherent-port}
Put \(T=|P_N|\), \(I=\mathbb F[o:X\to V]\), and write \(C(T)\) for the coherent syntactic category of the equational theory of \(T\). For every coherent target \(\mathcal D\) there is a square commuting up to canonical natural isomorphism
\[
\begin{tikzcd}[column sep=large]
 \mathsf{Coh}(C(T),\mathcal D)\arrow[r,"C(i_t)^*"]\arrow[d,"\simeq"']&
 \mathsf{Coh}(C(I),\mathcal D)\arrow[d,"\simeq"]\\
 \mathsf{FP}(T,\mathcal D)\arrow[r,"i_t^*"']&
 \mathsf{FP}(I,\mathcal D).
\end{tikzcd}
\]
The vertical functors evaluate on sorts and terms. Here \(\mathsf{FP}\) denotes finite-product-preserving functors, and both model categories use natural transformations. At the model \((M,p)\) of \cref{prop:relative-models} the distinguished arrow in the reduct is \(M_p(t)\). These comparisons compose under interpretations and restrict to the maximal subgroupoids of models.
\end{proposition}
\begin{proof}
Send a sort context to its product and an equation to the equalizer of its two terms. Interpret conjunction by intersection, disjunction by finite union, and existential quantification by regular image. These operations interpret coherent deduction. The graph of a total functional relation has a first projection which is both regular epic and monic. Its inverse followed by the second projection interprets the corresponding arrow. Existential composition of graphs interprets composition of arrows. Thus each product model extends to a coherent functor. Conversely a coherent functor preserves each of the operations used in this construction. Induction on formulas identifies its value on every formula with the interpretation in its restricted model, and the graph description gives naturality on arrows.

A homomorphism on sorts carries tuples satisfying each equation to tuples satisfying its image. Restriction to intersections and finite unions preserves this implication. For an existential formula the map on witness tuples induces a map on their regular images. Induction thus restricts the homomorphism to every formula subobject, and preservation of the functional graphs proves naturality. The inclusions into contexts force each component, giving full faithfulness. This proves the vertical equivalences, as in the coherent syntactic-category theorem \cite[Theorem 8.1.3]{makkai-reyes}. Substitution along \(i_t\) on terms sends \(o\) to \(t\). Its translation on formulas commutes with all the interpreting operations, so the two routes in the square agree on sorts, graphs and transformations. For successive substitutions the comparison is fixed by its components on sorts, proving the composition and identity laws. Passing to invertible transformations on both sides gives the last assertion.
\end{proof}
The common interface with \cite[Definition 4.9 and Theorem 4.13]{pugh} is the coherent theory of a span. The interpretation \(\gamma:\mathsf{Span}_{\rm syn}\to C(I)\) sends its generating span to \(X\xleftarrow{1}X\xrightarrow{o}V\). It exists by translation of the three sorts and two arrows. Consequently \(C(i_t)\gamma\) sends that span to the graph of \(t\), and restricting the coherent model gives \(MX\xleftarrow{1}MX\xrightarrow{M_p(t)}MV\). Composing two such graph spans computes the graph of the composite, since the pullback of the first output along the second identity is the first input. This specifies the comparison on the common span signature while retaining the scalar theory of each presentation. Varying-fibre aggregation enters through the field operation \(\Sigma_f\). The one-child and two-child models in \hyperref[sec:coherent-row-language]{the row-language comparison} distinguish the signature needed for this operation.

\paragraph{Opening a named occurrence.}\phantomsection\label{sec:cut-substitution}
In context \(X\times Z\) of \(\mathbb F_\Theta\), pull data references back from \(X\) and use the second projection for every reference to the chosen result of type \(Z\), including any output reference. Evaluate the other gates in list order. Write \(h:X\to Z\) for the original register and \(C_N:X\times Z\to Y\) for the open output. For any \(a:U\to X,\ w:U\to Z\), filling the open list gives \(C_N\langle a,w\rangle\). Indeed the selected reference becomes \(w\), and precomposition with \(\langle a,w\rangle\) commutes with each gate's operand tuple. This induction includes all bypass paths. Taking \(U=X\times Z,\ (a,w)=(\pi_X,\pi_Z)\) proves uniqueness of the arrow representing all fillings, and taking \((a,w)=(1_X,h)\) gives \(F_N=C_N\langle1_X,h\rangle\). The interpretation \(i_{C_N}\) exposes the independently variable port and bypass inputs.

\section{Semantic normal forms and orbit quotients}\label{sec:semantic}
Here a shape is a finite-valued functor $A:C\to\mathsf{FinSet}$. Begin with a finite entity category given by its arrow and composition tables. This covers, for example, graphs with two endpoint keys and a separate parameter-owner key. A finite list of input field tags specifies the sort of each coordinate in its ordered key tuple; outputs have the same kind of interface. We give the argument for one output tag, componentwise for a record. Input field values are independent rational coordinates; derived observables are evaluated by substituting their definitions. Finite tagged contexts split into cases. A context restricted by structural key tests retains only its eligible tuples. Denote the set of legal input cells by $X_A$; a cell includes its field tag. A row size $|A|$ counts all entity rows once, summing across sorts.

A finite composition table can itself be recovered from a finite presentation by generators and path relations whenever the presented category has finitely many arrows. Enumerate finite category tables $D$ on the given objects, images of the generators satisfying the relations, and a typed representative path $w_d$ for each arrow of $D$ evaluating to $d$. Dovetail finite equational derivations of $w_{1_i}=1_i$ and $g w_d=w_{\bar g d}$ for every composable generator $g$ with image $\bar g$. A completed certificate gives a functor from the presented category to $D$. Induction on path length proves $v=w_{\bar v}$ in the presentation. Thus $d\mapsto[w_d]$ is inverse on each hom-set, and composition is transported by evaluation. For a finite presented quotient its own table, path representatives and finitely many derivations occur in the search. The procedure therefore halts throughout this domain and supplies the table used below.

The grammar here forms rational constants, reads at typed key tuples, addition, subtraction, multiplication, lets, and sums binding rows of a named sort. A structural guard is a Boolean combination of equalities between parallel typed key terms. In particular an inclusion \(i:B\hookrightarrow A\) of closed subshapes satisfies
\[
                    G_B(\bar u)\ \Longleftrightarrow\ G_A(i\bar u).
\]
For an equality this follows from commutation with the keys and injectivity at its target sort. Boolean connectives preserve the equivalence. A fibre sum uses an equality to the target key, and a product of row sorts uses multiple binders. A product of a fixed list of factors belongs to the grammar. The variable-length product \(\prod_i(1+x_i)\) in \cref{thm:compiler-struct} has degree equal to the fibre size. A decidable test inspecting the entire host, such as whether a sort is inhabited, also has a different support behavior, computed below.
The guard syntax also captures every decidable isomorphism-invariant predicate on a fixed tuple of row slots that is preserved and reflected by embeddings. Indeed let \(T=\coprod_j C(c_j,-)\) be the finite free shape on the slots. For each sorted congruence \(\pi\) on \(T\), form the conjunction specifying \emph{all} equalities and disequalities between same-sort row terms in \(T\) according to \(\pi\). Every assignment \(T\to A\) satisfies exactly the conjunction for its kernel. It factors as \(T\twoheadrightarrow T/\pi\hookrightarrow A\). Test the predicate on the generators in \(T/\pi\), then disjoin the passing conjunctions. Embedding invariance identifies this formula with the original predicate on every host. This supplies an explicit translation from semantic local key tests to the displayed grammar.
\paragraph{Which changes of shape are functorial?}
Let \(\mathcal G\) be the groupoid of finite shapes and isomorphisms for this schema. Transport of the legal keyed cells along an isomorphism gives forward transport of each field by reindexing along its inverse. An equivariant family is therefore a natural transformation between the input and output field functors on \(\mathcal G\). Extending this assertion to every inclusion of shapes would exclude even the scalar query \(A\mapsto |A|\) on one entity sort. Its input is terminal and its scalar output has identity transport, so naturality for the inclusion of one row into two would require \(1=2\). In the inverse formula below, closed subshape inclusions index the coefficient subtraction. Equivariance under isomorphisms ensures that the resulting coefficients depend only on the decorated isomorphism class.

\subsection{The inverse formula}
Write $F_A(x)_{\bar b}$ for the output at legal tuple $\bar b$. On each shape this is a formal rational polynomial in $x_c$, $c\in X_A$. Reindexing by a shape isomorphism transports its coordinates and outputs. For a monomial $m$ (a multiset of tagged cells), let $a(A,\bar b,m)=[x^m]F_A(x)_{\bar b}$. Let $K_A(\bar b,m)$ be the smallest subfunctor containing all row coordinates of $\bar b$ and $m$. A subshape $B\le A$ is a family of subsets closed under all foreign keys, with the restricted structure. Membership of a tuple and all structural equality tests are preserved and reflected by inclusion. Define
\begin{equation}\label{eq:support-residual}
 r(A,\bar b,m)=a(A,\bar b,m)-
 \sum_{K_A(\bar b,m)\le B\lneq A}r(B,\bar b,m).
\end{equation}
This is M\"obius inversion on the finite interval of closed subshapes containing $K_A(\bar b,m)$ \cite{rota}. Proper inclusion decreases row size, so the recursion terminates. All subsets and their closure tests are finite. It commutes with isomorphisms since they transport every summand bijectively.

A pattern $H=(P,\bar u,M)$ is a finite shape with a legal ordered output and a multiset of legal tagged input cells. Other rows of $P$ are allowed, even when no input field reads them. Write $a_H$ for the number of shape automorphisms fixing $\bar u$ and preserving $M$ as a multiset. Define the query
\begin{equation}\label{eq:embedding-query}
 B_H(A,x)_{\bar b}=\frac{1}{a_H}
 \sum_{\substack{i:P\hookrightarrow A\\i\bar u=\bar b}}\prod_{c\in M}x_{i c}.
\end{equation}
The product is over occurrences with multiplicities. The sum binds an image row for each row of $P$, requires preservation of all foreign keys, distinct images within each sort, and the stated output equations. Thus \eqref{eq:embedding-query} is executable query syntax. A repeated factor does not add a row automorphism. In signed arithmetic a test of distinctness can also be eliminated by multiplying by one minus the equality indicator.

\begin{theorem}[Canonical semantic correspondence]\label{thm:support-readback}
For the finite entity category, independent rational fields and structural guards specified above, the following three presentations describe the same arrows at each field interface. The first two use finite codes. A name for an arrow in the third consists of the indicated evaluation procedure and bounds.
\begin{enumerate}
\item rational sum-comprehension queries, or their keyed sum--product programs, modulo equality on all rational inputs at all finite shapes;
\item finite rational linear combinations of the \(B_H\), collected by decorated isomorphism class;
\item an exact evaluator for an isomorphism-equivariant polynomial family and integers \(k,d\ge0\) witnessing \(r(A,\bar b,m)=0\) whenever \(|A|>k\) or \(|m|>d\). Names of the same family represent the same arrow.
\end{enumerate}
Writing $[H]$ for a decorated isomorphism class, the inverse is
\begin{equation}\label{eq:canonical-semantic-readback}
 F=\sum_{[H]} r(H)\,B_H .
\end{equation}
For finite syntax the construction computes both bounds. Given exact evaluation and valid bounds $k,d$, finitely many probes recover the unique coefficients in \eqref{eq:canonical-semantic-readback}. The translations preserve composition, identities and the chosen finite products of field records. Thus the three presentations are isomorphic as many-sorted finite-product theories. In particular semantic readback is independent of the architecture used to compute the family.
\end{theorem}
Equivalently, let \(\mathsf{Qry}\) be the first category in the statement, \(\mathsf{Pat}\) the second with the substitution below, and \(\mathsf{Beh}\) the third. All three have the same field interfaces. Their hom-set bijections assemble as
\[
\begin{tikzcd}[column sep=large]
 \mathsf{Qry}\arrow[r,"\mathrm{expand}","\sim"']\arrow[dr,"\mathrm{eval}"']&
 \mathsf{Pat}\arrow[d,"\mathrm{eval}_B","\sim"']\\
 &\mathsf{Beh}.
\end{tikzcd}
\]
In this diagram \(\mathsf{Beh}\) already identifies equal evaluations. Bounds are data for computing the inverse functor on names. They do not label distinct semantic arrows.

\begin{proof}

The chosen product of two interfaces is concatenation of their output records. In \(\mathsf{Qry}\) projection selects components and pairing assembles a record; these operations descend through equality on all finite models. In \(\mathsf{Pat}\) tag each pattern by its output component. Combining two expansions with their respective tags gives their pair, and extracting one tag block gives the projection of that pair. Uniqueness follows componentwise. The empty output record has the unique empty expansion. In \(\mathsf{Beh}\) these operations are pointwise, with bounds \(\max(k_1,k_2),\max(d_1,d_2)\) for a pair of evaluators. The coefficient calculation for each tag uses only that component. Consequently both the expansion and its inverse preserve these products.

\begin{enumerate}[titledsteps]
    \item \textbf{Invert the sum over images.} Fix an output \(\bar b\), monomial \(m\) and closure \(K=K_A(\bar b,m)\). For an embedding contributing to \([x^m]B_H(A,x)_{\bar b}\), take its image with the induced output and multiset. The fibre over a concrete image \(B\) is
\[
 \operatorname{Iso}\bigl(H,(B,\bar b,m)\bigr)
 \ \cong\ \operatorname{Aut}(H)
\]
after choosing any one such isomorphism. Thus this finite fibre has size \(a_H\), independently of the choice. On the finite interval of closed subshapes above \(K\), write \((\zeta f)(A)=\sum_{K\le B\le A}f(B)\). Then
\[
 [x^m]B_H(A,x)_{\bar b}
   =(\zeta\,\mathbf1_{[H]})(A),\qquad
 \zeta^{-1}[x^m]B_H=\mathbf1_{[H]}.
\]
Here \(\mathbf1_{[H]}(B)\) tests the decorated type of \((B,\bar b,m)\), and \eqref{eq:support-residual} computes the inverse of this unitriangular operator. For an arbitrary family the same recursion gives
\[
 a(A,\bar b,m)=\sum_{K\le B\le A}r(B,\bar b,m).
\]
The coefficient of $x^m$ in $\sum_H r(H)B_H$ is precisely this sum. This proves existence and uniqueness coefficient by coefficient. With finite ranges there are finitely many candidate types, since there are finitely many choices of finite carriers, structure tables and multisets within the bounds.


    \item \textbf{Expand the syntax and bound the patterns.} Inline lets for the coefficient calculation. With the output slots fixed, encode a clause as \((c,\vec v,G,M)\) with value \(\sum_{\vec v} c\,\mathbf1_G\prod_{a\in M}x_a\). Write \(L(t)\) for the finite list of clauses of \(t\). Constants, reads and structural indicators give singleton lists. The induction steps, extended over lists by distributivity, are
\[
\begin{aligned}
 L(t+t')&=L(t)\mathbin{+\!\!+}L(t'),\\
 L(-t)&=\{(-c,\vec v,G,M):(c,\vec v,G,M)\in L(t)\}_{\rm list},\\
 (c,\vec v,G,M)(c',\vec w,H,N)
   &=(cc',(\vec v,\vec w),G\wedge H,M\uplus N),\\
 \sum_{u}(c,\vec v,G,M)&=(c,(u,\vec v),G,M).
\end{aligned}
\]
Rename bound slots apart before multiplication. This rule keeps two independent witness choices even if their read addresses coincide. Finite tags give separate cases and structural restrictions contribute guards.

For a clause with typed free and bound row slots $v$, form the free shape $T=\coprod_v C(\mathrm{sort}(v),-)$. A tuple of assignments in $A$ extends uniquely to a homomorphism $T\to A$ by following paths. Enumerate the sorted equivalence relations on $T$ stable under all entity arrows. For each such congruence $\pi$, the maps with that kernel are exactly the embeddings $T/\pi\hookrightarrow A$ after the quotient map. Evaluate the clause's guard on the named quotient rows. Injectivity implies that its truth there is exactly its truth on every such assignment. Keep the passing quotient and transport the ordered output and tagged factor multiset to it. Its unnormalized embedding sum contributes the clause coefficient times $a_H B_H$. This partitions all assignments, including those with collisions. Put $L=\max_c\sum_t|C(c,t)|$ (use $0$ for an empty entity category). A clause with $g$ slots contributes row size at most $gL$ and degree at most its factor count. Taking the maxima over the lists gives $k,d$, with both bounds zero for an empty list. Combining equal finite decorated types is decidable by enumerating sorted bijections. This proves the translation from syntax and gives an exact decision procedure for uniform equality in this fragment.


    \item \textbf{Recover coefficients by finite probes.} Let \(S_k\) be the finite enumeration of legal shape representatives of size \(\le k\). At each \(A\in S_k\) use the grid \(G_A=\{0,\ldots,d\}^{X_A}\). With \(\ell_j(T)=\prod_{0\le i\le d,\,i\ne j}(T-i)/(j-i)\), interpolation gives the executable coefficient formula
\[
 a(A,\bar b,m)=[x^m]\sum_{v\in G_A} F_A(v)_{\bar b}
                                      \prod_{c\in X_A}\ell_{v_c}(x_c).
\]
Indeed \(a=\zeta r\) already gives \(a(A,\bar b,m)=0\) for \(|m|>d\) on every host, so interpolation in each variable is exact. For a single monomial one may set variables outside its support to zero. Compute \(r\) on \emph{concrete} subshapes in increasing size, including multiplicities when several subshapes have the same type. For \(d=0\) or \(X_A=\varnothing\) the empty products give the constant evaluation. Compiling the resulting finite combination proves the inverse law on behaviors.


    \item \textbf{Compose by identifying shared rows.} For composition, fix an outer pattern \(G\), and for each occurrence \(e\) of its multiset choose an inner pattern \(F_e\) of the appropriate output type. Tag all copies separately. In the disjoint union of their carriers impose the boundary identifications \(\bar u_{F_e}=\mathrm{key}_G(e)\) componentwise. Enumerate sorted congruences \(\pi\) extending these identifications whose restriction to each original pattern is equality. Denote by \(H_\pi\) the quotient with the outer output and the union of the transported inner multisets. For each host \(A\) there is a bijection
\[
 \left\{(i_G,(i_e)_e):
 \begin{array}{l} i_G:G\hookrightarrow A,\ i_e:F_e\hookrightarrow A,\\
 i_e\bar u_{F_e}=i_G\mathrm{key}_G(e)\end{array}\right\}
 \ \cong\ \coprod_{\pi}\operatorname{Emb}(H_\pi,A).
\]
Here embeddings act on the underlying shapes and the outer output is fixed at the requested key. The forward map takes the simultaneous kernel and the induced injection. Its inverse restricts along the quotient maps. Thus with inner coefficients \(\lambda_{F_e}\) and outer coefficient \(\lambda_G\), the contribution at each \(\pi\) is
\[
 \frac{\lambda_G}{a_G}\left(\prod_e\frac{\lambda_{F_e}}{a_{F_e}}\right)
                 a_{H_\pi} B_{H_\pi},\qquad
 |H_\pi|\le k_G+d_Gk_F,\quad \deg H_\pi\le d_Gd_F .
\]
Sum over choices and congruences and collect decorated types. Repeated occurrences at one cell still have independent inner choices. Nullary multisets use the empty product. Identities are direct reads. The bijection evaluates the formula as ordinary substitution, so uniqueness of the coefficients transfers the associative and identity laws.\qedhere
\end{enumerate}
\end{proof}

The boundary gluing in this composition is a colimit in the category of finite shapes. More explicitly, write \(U\) for the underlying outer shape, \(V=\coprod_e V_e\) for the inner shapes and \(L=\coprod_e L_e\), where \(L_e\) is free on the typed components of the boundary tuple at occurrence \(e\). The tuple gives both legs from \(L\). For compatible embeddings into a host \(A\) one has
\[
\begin{tikzcd}[column sep=large,row sep=normal]
 L\arrow[r]\arrow[d]&V\arrow[d]\arrow[ddr,bend left=13]\\
 U\arrow[r]\arrow[drr,bend right=13]&U\amalg_L V\arrow[dr,dashed,"u"]\\
 &&A .
\end{tikzcd}
\]
The right-hand cocone uses the copairing of the inner embeddings. The unique \(u\) factors as an epimorphism onto its image followed by a monomorphism into \(A\), computed sortwise. The epi part may identify rows in different copies. Its restriction to each constituent shape is injective because that constituent already embeds into \(A\). Conversely a further quotient of the pushout with this restriction, followed by an embedding, recovers precisely one compatible family. Taking the kernel congruence indexes these factorizations once. This accounts for the full sum over collisions in the substitution formula.
\begin{corollary}[Differentiating the canonical query]\label{cor:normal-form-derivative}
For a pattern $H=(P,\bar b,M)$, introduce a fresh direction field for each input field. For a distinct tagged cell $c$ of multiplicity $m_c$ in $M$, let $H_c$ replace one occurrence at $c$ by the corresponding direction-field occurrence. Then
\[
 D B_H(x;h)=
 \sum_{[c]\in\operatorname{supp}(M)/\operatorname{Aut}(H)}
                 m_c\,B_{H_c}(x,h).
\]
Thus semantic reconstruction of the derivative agrees with differentiating the reconstructed query and collecting equal decorated types.
\end{corollary}
\begin{proof}
Fix \(G=\operatorname{Aut}(H)\), \(G_c=\operatorname{Stab}_G(c)\), and let \(S_c\) be the unnormalized embedding sum for a mark at \(c\). Restoring the fresh direction tag to its original tag recovers \(H\), so \(\operatorname{Aut}(H_c)=G_c\). Extracting \([\epsilon]\) at \(x+\epsilon h\) gives
\[
 DB_H=\frac1{|G|}\sum_c m_c S_c
    =\sum_{[c]}m_c\,\frac{|G\cdot c|}{|G|}S_c
    =\sum_{[c]}m_c\,\frac{S_c}{|G_c|}.
\]
The last equality is orbit--stabilizer. The sum is empty for degree zero. Uniqueness in \cref{thm:support-readback} identifies the collected coefficients with semantic reconstruction of the derivative.
\end{proof}

\paragraph{Readback at a named interface.}
Write \(U=\Theta\times X\) to make the shared parameters explicit. For a program in the grammar of \cref{thm:support-readback}, designating a hidden output \(h:U\to Z\) selects a record of terms; opening the register gives \(C_N:U\times Z\to Y\) by the substitution above. Both are queries in the same grammar. Denote the expansion into decorated patterns by \(\mathsf{expand}\). The substitution law gives
\[
 \mathsf{expand}(F_N)
 =\mathsf{expand}(C_N)\circ_{\mathsf{Pat}}
       \langle 1_U,\mathsf{expand}(h)\rangle .
\]
Thus an evaluator for the open interface with valid bounds determines its canonical continuation, including its response away from the graph of \(h\). An evaluator restricted to that graph supplies the composite on the left.

\subsection{What the bounds detect}
On a single sort $V$ with no nonidentity arrows the constant field $v\mapsto1$ followed by summation gives $|V|$. Its canonical scalar query is the one-row pattern with no input factors. With two independent row binders the count is $|V|^2$. Equal and unequal assignments give coefficients $1$ on the one-row pattern and $2$ on the unordered two-row pattern. On $n$ rows the reconstruction reads $n+2\binom n2=n^2$. A requirement that every witness row be generated by an output or an input cell would lose both counts. More precisely, invariance on old outputs under extending inputs by zero is equivalent to $r(A,\bar b,m)=0$ for $A\ne K_A(\bar b,m)$. Indeed setting new variables to zero identifies exactly the coefficients of monomials in old cells. Their equality with the coefficient on $K$ is equivalent, by \eqref{eq:support-residual} and induction, to the stated vanishing. Retaining the other residuals gives compositional readback directly on the original input interface. Figure~\ref{fig:two-row-interval} shows the concrete summands when $n=2$.

\begin{figure}[!htbp]
    \centering\small
    \begin{tikzcd}[row sep=4mm,column sep=13mm]
        & \{a,b\}\quad(r=2) & \\
        \{a\}\quad(r=1) \arrow[ur,hook] && \{b\}\quad(r=1) \arrow[ul,hook'] \\
        & \varnothing\quad(r=0) \arrow[ul,hook] \arrow[ur,hook'] &
    \end{tikzcd}
    \caption{Concrete subshapes of the two-row host for the query $|V|^2$. The two singleton subshapes have the same decorated type and each contributes one. The top residual counts the two orders of a pair of distinct rows. Summing the four displayed residuals gives $4$.}
    \label{fig:two-row-interval}
\end{figure}


Degree alone cannot certify finite support. For \(a(n)=\mathbf1_{n>0}\),
\[
 r(0)=0,\qquad r(n)=(-1)^{n+1}\ (n>0),\qquad
 a(n)=\sum_{j=0}^n\binom nj r(j).
\]
The last equality is the binomial theorem. Even an adaptive evaluator-only inverse cannot uniformly discover a size bound on the finite-query families. Run it on the zero query. If it terminates, let \(N\ge1\) exceed the size of every host it inspected. The query \(B_N(n)=\binom nN\) counts images of an \(N\)-row pattern with empty multiset. It returns the same answer to every probe in that run, forcing the same subsequent probes and final output. Yet \(B_N(N)=1\). This proves that a bound in the evaluator presentation of \cref{thm:support-readback} carries information.
Even a partial procedure on evaluator codes that is required to halt only on the promised families cannot extract a size bound. From an index $i$ compile a total evaluator $g(i)$ by inserting $i$ into the following finite program. On a host of size $n$ simulate $i(i)$ for $n$ steps. Until a result appears return zero. If it returns $k$ at step $t$, put $N=\max(t,k+1)$ and return $\binom nN$. Suppose $B$ extracts a valid size bound at degree zero, and take an index $a$ for the program $i\mapsto B(g(i))$. If $a(a)$ returns $k$ at step $t$, then $g(a)$ denotes exactly the $N$-row pattern: at sizes below $t$ its binomial coefficient is zero as well. Its residual at size $N$ forces $k\ge N$, a contradiction. Hence $a(a)$ cannot return. But then $g(a)$ is the zero family, with the explicit bound zero, and $B(g(a))=a(a)$ must return. This argument uses only bounded simulation in the evaluator.

Degree carries independent information even at row bound zero. Use a nullary scalar field, enumerate the rationals without repetitions as $(q_j)_{j\ge0}$, and on $x=q_j$ simulate $i(i)$ for $j$ steps. Until a result appears return zero. With $k,t,N$ as above return $\prod_{h<N}(x-q_h)$. If the simulation eventually returns, every earlier test input is a root, so the evaluator denotes exactly this monic degree-$N$ polynomial on every host. If it never returns, the evaluator is zero. With $g(i)$ now coding this program and $a$ indexing $i\mapsto B(g(i))$ for a proposed degree extractor, termination again contradicts $N>k$, and impossibility of termination makes $g(a)$ the zero polynomial, forcing termination. Neither argument requires deciding membership in the promise.
 
This same residual quantifies the effect of freezing a parameter separately on each shape. The program \(F(p)=p\) is a direct scalar read. Choosing \(p_A=\mathbf1_{\{|A(V)|>0\}}\) gives exactly the family above, of degree zero in the remaining inputs and unbounded residual support. Keeping \(p\) as an input, or substituting a parameter query in the bounded grammar, preserves the finite expansion by the composition law.


The same inversion extends to finitely presented entity schemas with infinitely many paths. For every finite instance, closed subshapes are still computed by scanning just the generating arrows. The image-count proof and residual identity therefore hold verbatim. Every finite pattern can be queried by checking preservation of only those arrows. An evaluator supplied with computable per-shape degree bounds gives a computable coefficient table indexed by all finite decorated patterns: enumerate finite action tables on the generators, test the finitely many relations, interpolate and subtract. At a fixed host one needs only its finitely many closed subshapes. For degree enumeration on that host use the maximum of the supplied bounds on its subshapes; recursion can import a higher degree from a subshape even when it cancels in the host output. This gives an effective, locally finite semantic readback even on a schema with cycles.

The finite-arrow step in the finite normal form has specific content. On the schema with one unconstrained endomorphism, take hosts consisting of a single cycle of length $n$. The only closed subshapes are the empty shape and the entire cycle. The ordinary scalar row count has residual $n$ on every such cycle. It has a one-binder source program and a one-reduction network. Reconstructing it by \emph{total-subshape} patterns needs infinitely many pattern sizes. Structural readback remains finite on this schema. Raw attributes tied by schema equations are likewise evaluated by substitution or restriction to their model locus. For example, on the locus $x_1=x_2$ separate coefficients of the two ambient coordinates cannot be distinguished by probes. With one independent owner field and two reads the construction computes the unique coefficient in the owner coordinate.

The embedding/homomorphism change of basis for graph counts is developed by Curticapean, Dell and Marx \cite{motif-basis}. Puny et al.\ construct a basis of equivariant graph polynomials \cite{graph-poly-basis}. Cohen, Nutt and Sagiv characterize equivalences for classes of conjunctive aggregate queries \cite{aggregate-equivalence}. Here the inversion is on foreign-key-closed subinstances, indexed by a tagged input monomial and an ordered output; the substitution rule also accounts for unlabelled witness rows and tied input cells.

For a concrete comparison with the graph basis of Puny et al.\ \cite[Theorem 2.1 and Equation (7)]{graph-poly-basis}, take one vertex sort and a binary field $X_{ij}$. Let $H_0$ have $r$ vertices, all occurring in an output or an input cell, and let $H$ adjoin $k$ further vertices with no occurrences. At any host size $n\ge r$,
\[
                    B_H=\binom{n-r}{k}B_{H_0}.
\]
For each embedding of \(H_0\) there are \((n-r)(n-r-1)\cdots(n-r-k+1)\) extensions, and \(a_H=k!\,a_{H_0}\). This proves the binomial factor, also when the complement has fewer than \(k\) rows. With the notation of \cite[Equations (4) and (7)]{graph-poly-basis}, the injective index sum is \(Q_{H_0}=a_{H_0}B_{H_0}\). For the unrestricted sum, partition vertex assignments by their kernel to obtain
\[
             P_{H_0}=\sum_{\pi\text{ a vertex partition}}
                       a_{H_0/\pi}\, B_{H_0/\pi}.
\]
Carry the output pair and all input-edge occurrences through the quotient. At a fixed output, an incompatible quotient contributes zero. These equations distinguish the change of basis at fixed host size from the binomial dependence of unused witnesses as the size varies.

\begin{proposition}[A boundary for validity across all shapes]\label{prop:shape-concept-limit}
There is no algorithm that, given a finite discrete entity schema and a polynomial equality in its row counts, decides whether the equality implies false on every finite instance.
\end{proposition}
\begin{proof}
Given $p\in\mathbb Z[t_1,\ldots,t_s]$, take independent sorts $V_1,\ldots,V_s$. The scalar query $n_i=\sum_{v:V_i}1$ reads the $i$th cardinality, so the sum--product language computes $p(n_1,\ldots,n_s)$. Every tuple in $\N^s$, including zero entries, occurs as the sizes of a finite instance. Thus a procedure for this implication would decide whether a given Diophantine equation has a solution in natural numbers, by reversing its answer \cite{diophantine-undecidability}. The canonical readback in \cref{thm:support-readback} decides equality of bounded polynomial operations. Existence of a shape at which a polynomial vanishes is a different quantification.
\end{proof}

\paragraph{Approximation on a register.}
On a fixed finite vector of independent real inputs with access to its scalar coordinates, affine--square circuits can form every rational polynomial, by the product realization in \cref{thm:compiler-struct}. On a compact subset of $\mathbb R^d$ these polynomials are uniformly dense in the continuous real functions by Stone--Weierstrass: real polynomials separate points, and their finitely many coefficients can be approximated by rationals uniformly on the compact set. The comparison in \cref{thm:support-readback} uses exact equality across database shapes. Uniform approximation on one register also leaves derivative data undetermined. On $[-1,1]$ take $p_n(x)=x(1-x^2)^n$ for $n\ge1$. Its maximum absolute value is attained at $|x|=1/\sqrt{2n+1}$, as differentiation shows, and is at most $1/\sqrt{2n+1}$. Yet $p_n'(0)=1$ for every $n$. Thus even a sequence of exact polynomial circuits converging uniformly to zero can have a persistent first derivative at an interior point.

\paragraph{Executable finite probes.}
The ancillary file \texttt{support\_readback.py} takes a finite schema presentation, an exact equivariant evaluator, and valid bounds on witness size and degree. It interpolates on labelled hosts, subtracts contributions of concrete closed subshapes, and evaluates the recovered embedding queries. From exact values of $(|V|+\sum_v x_v)^2$ it recovers six decorated types at row and degree bounds $(2,2)$. For an idempotent key $p$ it recovers $\sum_{p(e)=j}(1+x_e)$ at bounds $(2,1)$: each witness uses only $e$ and $p(e)$, since $p^2=p$. For every output on all 56 labelled idempotent maps on at most four rows, evaluation of the recovered embedding queries agrees with the direct fibre sum on the Boolean grid. These functions are affine in the input fields, so the grid check also verifies equality on all rational inputs at these hosts. For the unrestricted unary key the residual of the row count is the cycle length on cycles of sizes one through three, as predicted above.

\subsection{Orbit quotients at a fixed shape}\label{sec:fixed-shape-orbits}
For the fixed-shape comparison take finite address sets and a group acting on them. Write \(F_R(I)\) for the free \(R\)-module on the addresses \(I\). The following result uses the orbit coequalizer of the simultaneous action on output and input addresses.
\begin{theorem}[Orbit keys classify equivariant linear layers]\label{thm:orbit}
Let a presented finite group $G$ act on finite decidable input and output address sets $I$ and $J$. On their permutation modules over $R$ with trivial scalar action, the orbit map $\rho:J\times I\to O=(J\times I)/G$ gives
\[
 R^O \ \cong\ \Hom_{R,G}(R^I,R^J),\qquad
 (L_\theta x)_j=\sum_{i\in I}\theta_{\rho(j,i)}x_i .
\]
\end{theorem}
\begin{proof}\leavevmode
\begin{enumerate}[titledsteps]
    \item \textbf{Identify the coefficient table.} Write \(F_R(I)\) and \(F_R(J)\) for the two permutation modules. The coefficient pairing induces the canonical identification
\[
 \Hom_R(F_R(I),F_R(J))\cong F_R(J)\otimes_R F_R(I)^*\cong R^{J\times I}.
\]
The first map uses duality for the finite free module, and the second uses its distinguished address basis. This identification intertwines conjugation on maps with transport on coefficient fields: on an elementary map from address \(i\) to address \(j\), conjugation by \(g\) gives the elementary map from \(gi\) to \(gj\). Equivariant maps are exactly the fixed fields for this action. The orbit map is the coequalizer of the action and projection \(G\times(J\times I)\rightrightarrows J\times I\). Applying \(R^{(-)}\) identifies its equalizer with \(R^O\). The two comparisons are linear and mutually inverse. The coordinate formula in the statement evaluates the resulting map on an input field. Enumerating the finite action computes the orbit key.


    \item \textbf{Factor through the orbit key.} The triangle exhibits the factorization that determines the parameter key. The square states the intertwining equation for the layer obtained from that key.
\[
\begin{tikzcd}[column sep=small]
 J\times I\arrow[r,"\rho"]\arrow[dr,"K"']&O\arrow[d,dashed,"\theta"]\\
 &R
\end{tikzcd}
\qquad
\begin{tikzcd}[column sep=small]
 R^I\arrow[r,"L_\theta"]\arrow[d,"g"']&R^J\arrow[d,"g"]\\
 R^I\arrow[r,"L_\theta"']&R^J .
\end{tikzcd}
\]
Thus the parameter table has one address per orbit \cite{orbit:sharing}. The quotient construction uses no division by the order of the group.\qedhere
\end{enumerate}
\end{proof}
For $G=I=J=\mathbb Z/n$ with translation, $d=j-i$ gives a complete orbit label. The layer is cyclic convolution $y_j=\sum_i\theta_{j-i}x_i$. A local kernel selects the displacement orbits on which to place edges. Equivariance determines the allowed sharing even when the chosen window is the whole group.

\paragraph{Degree-bounded polynomial fields.}
At a fixed finite action on input and output keys extend \cref{thm:orbit} from linear gates to all rational polynomials of degree at most $d$. Index monomials by multisets $M$ of at most $d$ input addresses. The group acts on pairs $(j,M)$ by transport of both coordinates. Write
\[
 F_j(x)=\sum_{|M|\le d}a(j,M)x^M.
\]
Taking coefficients in \(F(gx)=gF(x)\) gives
\[
 [x^M]F_j=[x^{gM}]F_{gj},\qquad
 \{F:\deg F\le d,\ F(gx)=gF(x)\}
    \ \cong\ \Q^{(J\times\operatorname{Mult}_{\le d}(I))/G}
\]
for rational coefficients. Here \(\operatorname{Mult}_{\le d}(I)\) denotes the finite set of multisets of size at most \(d\). The inverse reads one orbit coefficient per monomial and forms the displayed polynomial. For evaluator input use the grid of \cref{thm:support-readback}. For a linear map just evaluate \(F(e_i)_j\). Distinct orbits with equal numerical coefficients still have distinct keys in the orbit table. No division by a stabilizer order occurs.


\clearpage
\part{Fibrations, adjunctions and port interpretations}\label{part:ports}

\section{Families, predicates and their adjoints}\label{sec:indexed-readback}

There are two changes of input at a hidden port. A parameter substitution transports the entire family of experiments. Restricting the available inputs within one experiment can change which predicates the experiment distinguishes. To compute these operations, put \(\mathcal B=\mathbf{Set}\) for external semantics or the effective topos \(\mathcal B=\E\) for internal semantics. The effective presentations for the latter are constructed in \cref{sec:eff}. We use the codomain and subobject fibrations in the sense of \cite{jacobs-fibrations}.

\subsection{Changing the base of an experiment}
A family over \(S\) is an arrow \(p:A\to S\). An arrow of families over \(r:T\to S\) is a commuting square. Pullback gives the cartesian lift on the left below. Direct image of an entire family is composition, with its opcartesian lift on the right for \(p':A'\to T\).
\[
\begin{tikzcd}[column sep=large]
 r^*A\arrow[r,"\bar r"]\arrow[d]&A\arrow[d,"p"]\\
 T\arrow[r,"r"']&S
\end{tikzcd}
\qquad
\begin{tikzcd}[column sep=large]
 A'\arrow[r,equal]\arrow[d,"p'"']&A'\arrow[d,"rp'"]\\
 T\arrow[r,"r"']&S .
\end{tikzcd}
\]
Any square into \(p\) over \(ru\) factors uniquely through the first square over \(u\), by the pullback property. Any square out of \(p'\) over \(vr\) factors uniquely through the second over \(v\), using its original top arrow. This exhibits the bifibration \(\operatorname{cod}:\mathcal B^{\to}\to\mathcal B\). Pasting cartesian squares gives the canonical comparison \(u^*r^*\cong(ru)^*\). Uniqueness makes the two comparisons for three substitutions equal.

In the fibre of families the two adjoints of \(r^*\) have the hom-set formulas
\[
\begin{aligned}
 \Hom_{/S}(\Sigma_r A',B)&\cong\Hom_{/T}(A',r^*B),\\
 \Hom_{/T}(r^*B,A')&\cong\Hom_{/S}(B,\Pi_r A').
\end{aligned}
\]
The first bijection is the cartesian factorization just described. For the second, a map on \(T\) assigns, for each \(b\in B_s\), a value in \(A'_t\) at every \(t\) over \(s\). Currying this assignment gives a map into \((\Pi_r A')_s=\prod_{t\in T_s}A'_t\), and evaluation is its inverse. This is an argument in the internal language of either base. For a pullback square with sides \(r:T\to S,\ r':T'\to S'\) and horizontal maps \(v:T'\to T,\ u:S'\to S\), the resulting comparisons are
\[
 u^*\Sigma_r\cong\Sigma_{r'}v^*,\qquad
 u^*\Pi_r\cong\Pi_{r'}v^* .
\]
Over \(s'\) both sides use the canonical isomorphism \(T'_{s'}\cong T_{u(s')}\) with the same family in each summand or factor. These Beck--Chevalley isomorphisms transport the witness objects of \cref{thm:elaboration}. The numerical aggregation identity in \cref{lem:bc} interprets the corresponding finite fibres in a semiring.

\subsection{Possible and certain explanations on the profile image}
In the fibre \(\operatorname{Sub}(S)\), the two adjoints of inverse image are the quantifiers \(\exists_r\dashv r^*\dashv\forall_r\). Existential image takes the image of the composite of a mono with \(r\). Universal image classifies the assertion that every point in the indicated fibre satisfies the predicate. They obey the same two base-change comparisons, since a pullback identifies the quantified fibres. Quantifying a family of witnesses and quantifying its support have different types. Support takes the regular image of the family's projection.

Let \(b:X\to 2^M\) record the observed Boolean attributes. Factor it as \(X\xrightarrow{e}B\hookrightarrow 2^M\), with \(B\) the reachable profiles and \(e\) regular epic. For a target predicate \(P\hookrightarrow X\), write \(P_{\rm may}=\exists_eP\) and \(P_{\rm must}=\forall_eP\). Both are predicates on \(B\). Let \(\mathsf{Sat}_e(X)\) consist of the predicates on \(X\) pulled back from \(B\), ordered by inclusion.

\begin{proposition}[Optimal profile bounds and exact descent]\label{prop:profile-adjoints}
Inverse image identifies \(\operatorname{Sub}(B)\) with \(\mathsf{Sat}_e(X)\). For its inclusion \(i:\mathsf{Sat}_e(X)\hookrightarrow\operatorname{Sub}(X)\), the left and right adjoints are
\[
 U(P)=e^*P_{\rm may},\qquad L(P)=e^*P_{\rm must},\qquad U\dashv i\dashv L .
\]
In particular the least profile-definable upper bound and the greatest profile-definable lower bound coincide exactly when \(P\) descends. Equivalently,
\[
 P_{\rm must}=P_{\rm may}
 \quad\Longleftrightarrow\quad P=e^*D\text{ for a unique }D\hookrightarrow B
 \quad\Longleftrightarrow\quad p_1^*P=p_2^*P
\]
for the kernel pair \(p_1,p_2:X\times_B X\rightrightarrows X\). Here a profile-definable predicate means a subobject pulled back from \(B\). The finite Boolean formula realizing a decidable such predicate on the effective presentations is constructed in \cref{thm:concept-expressibility}.
\end{proposition}
\begin{proof}
The internal fibre of \(e\) over a profile is inhabited in the existential sense. Hence \(\exists_e e^*D=D=\forall_e e^*D\) by eliminating the fibre witness into the membership proposition. This also shows \(e^*\) reflects inclusion. The two adjunctions give
\[
 U(P)\le e^*D\ \Longleftrightarrow\ P\le e^*D,\qquad
 e^*D\le L(P)\ \Longleftrightarrow\ e^*D\le P .
\]
Their unit and counit give \(L(P)\le P\le U(P)\). Equality of the two profile predicates therefore gives \(P=e^*D\), and the inverse direction follows from the two identities above. A pulled-back predicate has equal inverse images along \(p_1,p_2\). Conversely, a witness \(x'\) for \((\exists_eP)(e x)\) has \(e x'=e x\). Equality on the kernel pair gives \(P(x')\Rightarrow P(x)\), which proves \(U(P)\le P\). This constructs the descended predicate without choosing representatives. On the full Boolean cube, an unreachable profile has false existential image and true universal image. Corestriction to \(B\) is essential to the equality test.
\end{proof}
Equivalently, on the discrete categories of inputs and reachable profiles, a predicate is a diagram valued in truth values ordered by entailment. At \(b\in B\) the pointwise extensions of \(P\) along \(e\) are
\[
 (\operatorname{Lan}_e P)(b)=\exists_{x:e(x)=b}P(x),\qquad
 (\operatorname{Ran}_e P)(b)=\forall_{x:e(x)=b}P(x).
\]
Indeed both comma categories are the discrete fibre of \(e\) at \(b\). Its join and meet have precisely the cocone and cone properties in the truth-value order. Internally these are joins and meets in \(\Omega\) indexed by that fibre. Thus exact descent in \cref{prop:profile-adjoints} means that the two extensions agree on the reachable profiles.

On the finite and rational-linear presentations, a reachable profile in \(P_{\rm may}\) outside \(P_{\rm must}\) admits a positive and a negative realizing input. The searches in \cref{thm:concept-expressibility} return this pair. Internal universal image is available as a semantic operation beyond those presentations. For instance, bounded simulation decides whether program \(e\) is still running after \(n\) steps. Its universal image over \(n\in\N\) is nonhalting. A uniform positive test for that image, run in parallel with the halting test, would decide halting. This locates the computational boundary of the right adjoint.



\paragraph{Restriction of a population and change of parameter.}
Let \(f:X'\to X\) restrict the inputs inside one experiment. Let \(e':X'\twoheadrightarrow B'\) be the profile image after restriction, with \(j:B'\hookrightarrow B\). The square commutes since \(je'=ef\). For \(P'=f^*P\) one always has
\[
 \exists_{e'}P'\ \le j^*\exists_eP,\qquad
 j^*\forall_eP\ \le\ \forall_{e'}P'.
\]
A realizing input in the smaller population maps to one in the larger population, which gives the first comparison. A predicate true throughout an old fibre remains true throughout the retained inputs, which gives the second. Both comparisons are equalities if \(X'\to B'\times_B X\) is regular epic: every old input in each retained profile then lifts locally to the smaller population, and existential elimination proves both reverse inclusions. Pullback of the entire experiment along an external parameter map gives the Beck--Chevalley equalities in its slice. Mere restriction can be strict even when \(B'=B\). Take two inputs of the same profile with opposite decisions and retain only the positive one. The universal image of the target changes from false to true, so a classification can become exact through this loss of inputs.

\section{Adjunctions and quotients at a hidden port}\label{sec:hidden-concepts}
Fix a port of a frozen program, and let $X$ be the inputs to be explained. There are three computations in this section. First recover the value at the port as a query on $X$. Then use a vocabulary of observations at that port to express a target on $X$. Finally cut the port and ask which observations distinguish the responses to independent replacements. Each computation retains the key of the field being observed. \Cref{sec:ground-operational} uses the profile and response constructions together when a history carries both a port observation and an environmental observation.

\subsection{From a frozen network to formulas for its hidden fields}\label{sec:representation-readback}
At a finite presented shape flatten each exposed scalar cell to a variable named by its field tag and key tuple. An affine layer, a foreign-key lookup and a finite sum over a known fibre then have rational affine coordinate formulas when weights and biases are rational. Work here over exact rational values with the ordered scalar language of affine terms, conjunction, disjunction and equality or order comparisons. The corresponding sets of ordinary real solutions are the same rational polyhedral descriptions interpreted over $\mathbb R$. On located real names, numerical rectification is executable as $[s]_+=(s+|s|)/2$; deciding the truth of an arbitrary real threshold on its equality boundary is a separate question. Rounded machine operations have their own semantics. A finite binary floating-point parameter can be read as an exact dyadic coefficient, which specifies the exact-arithmetic realization used below.

\begin{theorem}[Exact hidden fields and decision predicates]\label{thm:relu-formula}
For a finite acyclic network of rational affine and ReLU gates, construct, for any selected tuple of internal registers $h$, a quantifier-free formula $\Phi_h(x,v)$ of rational linear arithmetic. For all rational inputs, $\Phi_h(x,v)$ holds exactly when $v$ is the actual tuple of hidden values at that input. The same construction returns an input-only formula for any Boolean combination of affine tests on these registers or on the output. It applies to a named collection of channels or a rational affine probe of a layer.
\end{theorem}
\begin{proof}\leavevmode
\begin{enumerate}[titledsteps]
    \item \textbf{Propagate affine formulas.} Give each selected register its external output name and take its ancestors in topological order. For each activity pattern $\sigma$ on their rectifier gates, propagate affine expressions in $x$. At an affine gate substitute the expressions for its operands. At a rectifier with preactivation $u_\sigma(x)$, record $u_\sigma(x)\ge0$ and propagate $u_\sigma(x)$ on the active branch, or record $u_\sigma(x)\le0$ and propagate zero on the other branch. Let $P_\sigma$ be the conjunction of all recorded tests and $L_{\sigma,h}$ the tuple at the requested registers. Return
\[
 \Phi_h(x,v)=\bigvee_\sigma\big(P_\sigma(x)\ \wedge\ v=L_{\sigma,h}(x)\big).
\]

    \item \textbf{Verify the graph on rational inputs.} For a rational input, decidable comparison selects an activity pattern. Induction on the sorted gates gives \(x\in P_\sigma\Rightarrow h(x)=L_{\sigma,h}(x)\). When the preactivation is zero the two branch expressions agree. For an output test \(T\), substitute \(T(L_{\sigma,h}(x))\) for the output equation and include the desired tie convention. Eliminate empty regions by rational feasibility. An alternative of linear size in the wiring keeps each gate's local graph equation and existentially binds its private result. The disjunction above can have exponentially many patterns.


    \item \textbf{Interpret the graph over the reals.} Over classical reals the same activity-pattern argument applies. For internal located reals the value has a constructive graph as well: with \(u\) the preactivation,
\[
 y=[u]_+\quad\Longleftrightarrow\quad
 0\le y\ \wedge\ u\le y\ \wedge\ \neg(0<y\ \wedge\ u<y).
\]
Indeed \(y>\max(0,u)\) is equivalent to both strict inequalities, and \(y\le\max(0,u)\) is the negation of this strict comparison. Together the two upper-bound inequalities and this last condition force equality. This graph evaluates the value without choosing a sign. A uniform choice of a valid activity bit would give \(u\ge0\lor u\le0\). The argument in \cref{prop:profile-risk} exhibits the obstruction.\qedhere
\end{enumerate}
\end{proof}

\paragraph{A term for the score.}
At a fixed finite shape, freeze the rational parameter table and take a compact box with rational endpoints for the exposed inputs. Substitute singleton coordinates as constants and rescale each remaining interval affinely to \([0,1]\). The activity cells in the proof of \cref{thm:relu-formula} are rational polytopes; every scalar score is continuous and affine with rational coefficients on each cell. Rational score bounds \(l,u\) can be computed by evaluating its affine pieces at the vertices of their nonempty cells. For \(l<u\) put \(s^\sharp=(s-l)/(u-l)\). A constant score can instead be recorded by its value and the normalized term \(0\).

In the standard divisible MV-algebra (DMV), \(a\oplus b=\min(1,a+b)\), \(\neg a=1-a\), and \(\delta_n a=a/n\) for \(n\ge1\). Write \(a\ominus b=\neg(\neg a\oplus b)\) and \(a\vee b=(a\ominus b)\oplus b\). Gerla identifies the term functions on the cube with the continuous \([0,1]\)-valued rational piecewise-affine functions. Equality there is exactly equational derivability in DMV \cite[Theorems 3.15 and 4.5]{gerla-dmv}.

Here is a direct translation of a frozen score. First, every clipped rational affine function \(\operatorname{clip}(L)=\min(1,\max(0,L))\) on the cube has an explicit term. Split \(L=P-N\) by collecting the positive and negative coefficients separately, including its intercept. Both \(P,N\) have nonnegative coefficients. Choose an integer \(m\ge1\) bounding both coefficient sums and use \(a^{\oplus j}\) for the \(j\)-fold truncated sum, with \(a^{\oplus0}=0\). For a rational coefficient \(r=k/d\in[0,1]\), the term \((\delta_d v)^{\oplus k}\) realizes \(rv\). Use \(v=1=\neg0\) for the constant coordinate. These terms assemble \(P/m,N/m\) by \(\oplus\) exactly, since each total is at most one. Hence
\[
 \operatorname{clip}(L)=\bigl((P/m)\ominus(N/m)\bigr)^{\oplus m}.
\]
Next, propagate rational interval bounds on the finite register list using affine interval arithmetic and rectification. A register known to be constant can be substituted. For an affine gate with nonconstant bounds \([l_i,u_i]\), express its raw affine value \(L_i\) by rescaling its normalized predecessors. The clipped affine term for \((L_i-l_i)/(u_i-l_i)\) represents its normalized value exactly. For a ReLU with positive upper bound \(u_i\), use output bounds \([0,u_i]\); the term for \(\operatorname{clip}(L_i/u_i)\) is its normalized output. A nonpositive preactivation upper bound gives the zero term. Induction constructs a term for every selected register. Clipping the final affine rescaling gives \(s^\sharp\) in the chosen score units. This construction also applies componentwise to a field record.

Zhang and B\"olcskei give a layerwise extraction and a reconstruction from substitution graphs for their non-degenerate networks \cite[Proposition 5 and Section 6.2]{zhang-bolcskei-mv}. Their round trip starts with the presented network. The score alone identifies its equational class, while the register presentation of \cref{thm:compiler-struct} specifies which occurrence can be cut. A threshold at an interior crossing can be discontinuous and is read by the order formula of \cref{thm:relu-formula}. Freezing a weight also matters here: already the joint operation \((p,x)\mapsto px\) restricts on \(p=x\) to \(x^2\), which is affine on no interval of positive length. The parameterized field interpretation still evaluates this operation.

For true label \(Y\) and predictions \(f,g\), the pointwise inequality
\[
 |\mathbf1_{f=Y}-\mathbf1_{g=Y}|\le\mathbf1_{f\ne g}
\]
gives \(|\Pr(f=Y)-\Pr(g=Y)|\le\Pr(f\ne g)\) by integration. Exact readback gives zero on the right. Decisions alone can leave the hidden code undetermined (\cref{ex:trained-discrepancy}); Theorems \ref{thm:relu-formula} and \ref{thm:support-readback} specify the different data used for their inverses.

\subsection{From a named field to a formal context}
\paragraph{Predicate readback as a lens.}
\begin{comparison}[Predicate readback and parallel composition]\label{comp:predicate}
In sets and internally in \(\E\), inverse image of predicates defines a symmetric colax monoidal functor into bimorphic lenses. Its tensor comparison forms the rectangle of two predicates and can fail to be invertible.
\end{comparison}
\begin{proof}\leavevmode
\begin{enumerate}[titledsteps]
    \item \textbf{Pull predicates back along a map.} For the set semantics of this section put \(\mathcal OX=(X,\mathcal P X)\). A map \(f:X\to Y\) induces the lens \(\mathcal Of=(f,(x,B)\mapsto f^{-1}B)\). The unit law is \(1^{-1}B=B\), and the lens composition law is \((gf)^{-1}B=f^{-1}(g^{-1}B)\). For a family \(f:P\times X\to Y\) this construction applies in each parameter fibre: \(r:P'\to P\) gives \((f\circ(r\times 1))_{p'}^{-1}=f_{r(p')}^{-1}\). An attribute at a register is thus an output observation transported backwards along an actual map of the network.


    \item \textbf{Form the parallel comparison.} The comparison for parallel observations is a lens
\(\kappa_{X,Y}:\mathcal O(X\times Y)\to\mathcal OX\otimes\mathcal OY\).
Its forward component is the identity. Its backward component sends \((A,B)\) to \(A\times B\), independently of the retained point. For \(f:X\to X'\) and \(g:Y\to Y'\) its naturality square is
\begin{equation}\label{eq:predicate-tensor}
\begin{tikzcd}[column sep=large]
 \mathcal O(X\times Y)\arrow[r,"\kappa_{X,Y}"]\arrow[d,"{\mathcal O(f\times g)}"']&
 \mathcal OX\otimes\mathcal OY\arrow[d,"{\mathcal Of\otimes\mathcal Og}"]\\
 \mathcal O(X'\times Y')\arrow[r,"\kappa_{X',Y'}"']&
 \mathcal OX'\otimes\mathcal OY' .
\end{tikzcd}
\end{equation}
Indeed \((f\times g)^{-1}(A\times B)=f^{-1}A\times g^{-1}B\). The unit comparison \(\mathcal O1\to(1,1)\) returns the full predicate on \(1\). The unitors, associator and symmetry identify the corresponding rectangular subsets, proving the symmetric colax monoidal laws. The direction of \(\kappa\) matters. On the two-point set \(2\), the diagonal in \(2\times2\) is a joint test. A rectangle containing both diagonal points contains the off-diagonal points as well. Thus this comparison is not invertible. Internally replace \(\mathcal P X\) by \(\Omega^X\), where \(\Omega\) classifies subobjects. The backward map is precomposition, and a rectangular predicate has characteristic map \((x,y)\mapsto A(x)\wedge B(y)\). The same identities follow from these maps and the meet and unit laws in \(\Omega\). A decidable attribute factors through \(2\to\Omega\). The finite Boolean synthesis below uses the stated decidable or rational presentations.\qedhere
\end{enumerate}
\end{proof}
A register joining two features can therefore carry observations that require a joint predicate.
Let $h:X\to H$ be the tuple of values at the chosen port. An attribute vocabulary consists of names $m\in M$ and predicates $E_m$ on $H$. For example, one name can denote the firing of a keyed channel and another a threshold on the difference between two channels. Write $E\subseteq H\times M$ for their incidence relation. In the equipment of sets, maps and relations, the companion of $h$ is its graph $\Gamma_h$. Horizontal composition with this graph gives
\begin{equation}\label{eq:concept-incidence}
 I_h=E\odot\Gamma_h,\qquad
 I_h(x,m)\ \Longleftrightarrow\
 \exists z\,[z=h(x)\wedge E_m(z)]\ \Longleftrightarrow\ E_m(h(x)).
\end{equation}
Thus the two presentations compute the same column. On the neural side, observe the register and test its value. On the database side, substitute its readback in the attribute predicate. If $h=v\circ u$, then $\Gamma_h=\Gamma_v\odot\Gamma_u$. Restricting $E$ along $v$ and then along $u$ therefore gives exactly the composite column. At fixed shape, \cref{thm:relu-formula} writes every column for a rational linear test as a formula in the input fields. In the other direction, a Boolean expression in the column names instructs the program to evaluate the register once, apply the declared comparison gates, and combine their bits. Gate-by-gate substitution proves agreement. The order tests belong to the observation language at the port.

For $A\subseteq X$ and $T\subseteq M$ form the derivations
\[
 A^\uparrow=\{m\mid\forall x\in A\ I_h(x,m)\},\qquad
 T^\downarrow=\{x\mid\forall m\in T\ I_h(x,m)\}.
\]
Both reverse inclusion. The equivalence $A\subseteq T^\downarrow\iff T\subseteq A^\uparrow$ is the Isbell adjunction for the relation. A \emph{formal concept} is a pair $(A,T)$ with $A=T^\downarrow$ (its extent) and $T=A^\uparrow$ (its intent). This is the classical formal context construction \cite{fca-foundations}, and the two derivations are the Boolean instance of the adjunction for an enriched distributor \cite{isbell-distributors}. Extents are ordered by inclusion. Meets are intersections, and the join of a family of extents is the closure of its union under $A\mapsto A^{\uparrow\downarrow}$. Indeed the adjunction makes this composite extensive, monotone and idempotent, hence its value on a set is the least closed superset.

\begin{theorem}[Concepts at a port and on its input domain]\label{thm:reachable-concepts}
Let $R=h(X)\subseteq H$ carry the restriction of $E$, and factor $h$ as $X\xrightarrow{q}R\hookrightarrow H$. Inverse image induces an isomorphism of complete concept lattices
\[
  (B,T)\longmapsto(q^{-1}B,T)
  \quad\text{from the context on $R$ to the context on $X$.}
\]
Its inverse is $(A,T)\mapsto(q(A),T)$. In this set semantics, for a finite vocabulary the intent reads back as the database formula $\bigwedge_{m\in T}E_m(h(x))$ and compiles to the same test at the neural port.
\end{theorem}
\begin{proof}
Keeping track of the context in the subscripts gives
$A^{\uparrow_X}=q(A)^{\uparrow_R}$ and
$T^{\downarrow_X}=q^{-1}(T^{\downarrow_R})$, directly from \eqref{eq:concept-incidence}. Since $q$ is onto, $q(q^{-1}B)=B$. It follows that the intents of both contexts have the same closure $T\mapsto T^{\downarrow\uparrow}$. Every input extent is $T^{\downarrow_X}$ and thus contains entire fibres of $q$. The two displayed maps are therefore inverse, preserve inclusion and carry the same intent. The last assertion is column substitution followed by conjunction, including the empty conjunction.
\end{proof}
The incidence itself is a pullback. Let \(\lvert E\rvert\hookrightarrow H\times M\) and \(\lvert I_h\rvert\hookrightarrow X\times M\) denote the two incidence subobjects. Write \(\lvert E\rvert_R\) for the restriction over \(R\). Factoring \(h\) through its image gives two cartesian squares:
\[
\begin{tikzcd}[column sep=large]
 \lvert I_h\rvert\arrow[r]\arrow[d,hook]&\lvert E\rvert_{R}\arrow[r]\arrow[d,hook]&\lvert E\rvert\arrow[d,hook]\\
 X\times M\arrow[r,"q\times 1"']&R\times M\arrow[r,hook]&H\times M .
\end{tikzcd}
\]
Pullback pasting identifies the outer rectangle with the graph-composition formula \eqref{eq:concept-incidence}. The bottom-left arrow is a regular epimorphism. Its inverse image reflects equality of subobjects, so equality of input columns is exactly equality on \(R\).

\paragraph{Restricting a population of inputs.}
For a map $f:X'\to X$, each attribute column restricts by $I_{h f}(x',m)=I_h(f(x'),m)$, the graph-composition calculation of \eqref{eq:concept-incidence}. For any premise set $T$, its extent on $X'$ is $f^{-1}(T^{\downarrow_X})$. Consequently
\[
 \operatorname{cl}_{X'}(T)=
 \{m\mid \forall x\in f(X')\cap T^{\downarrow_X}\ I_h(x,m)\}
 \ \supseteq\ \operatorname{cl}_X(T).
\]
Surjectivity gives equality. More generally, equality for this $T$ holds exactly when every attribute refuted somewhere on $T^{\downarrow_X}$ is refuted on that extent inside $f(X')$. For example, restrict the discrepancy grid to its lower tail. The implication $\varnothing\Rightarrow L$ becomes valid there, although it fails on the original grid. This restriction is a change in the population on which the compiled fields are observed. When recomputing an aggregate on a changed database, the comparison of the fields themselves is instead supplied by the pullback square in \cref{lem:bc} when that square applies.

\subsection{Closure and descent of target predicates}
Take a target $P\subseteq X$. It may be a specified task or a decision of the frozen model. Write $b(x)=\{m\in M\mid I_h(x,m)\}$ for the profile of an input. There are two useful logical fragments even when all observations are Boolean.

\begin{theorem}[Exact explanations in a finite vocabulary]\label{thm:concept-expressibility}
In classical sets, for finite $M$, the strongest positive conjunction valid on $P$ has intent $T_P=P^\uparrow$ and denotes $T_P^\downarrow$. A conjunction defines $P$ exactly if and only if $T_P^\downarrow=P$. A Boolean expression in the same attributes defines $P$ exactly if and only if
\[
 b(x)=b(x')\quad\Longrightarrow\quad (x\in P\ \Longleftrightarrow\ x'\in P)
 \qquad(x,x'\in X).
\]
Both criteria yield a formula or a certificate of failure on an explicitly presented finite table. They do so as well on a fixed database shape over $\Q$ when $X$, $P$ and the attribute predicates are defined by rational linear arithmetic and $h$ is a rational affine--ReLU register tuple.
\end{theorem}
\begin{proof}\leavevmode
\begin{enumerate}[titledsteps]
    \item \textbf{Construct the strongest conjunction.} For each \(m\in M\) compute \(A_m=\{x\in P:\neg I_h(x,m)\}\). Then
\[
 T_P=\{m:A_m=\varnothing\},\qquad
 P\subseteq T_P^\downarrow\subseteq T^\downarrow
                 \quad\text{whenever }P\subseteq T^\downarrow .
\]
The second inclusion follows from \(T\subseteq T_P\). A point of \(T_P^\downarrow\setminus P\) is therefore a simultaneous counterexample to every positive conjunction true on \(P\). The empty set of names denotes the empty conjunction.
A failure certificate can itself be checked using only finite data. Give \(x_-\notin P\). For each \(m\notin b(x_-)\), supply \(x_m\in P\) with \(m\notin b(x_m)\). Every name true throughout \(P\) then belongs to \(b(x_-)\). Conversely any \(x_-\in T_P^\downarrow\setminus P\) admits exactly these witness searches: each missing name was removed when \(A_m\) was found inhabited. Thus at most \(1+|M|\) inputs certify failure, even when \(X\) is infinite on the effective rational domain.

    \item \textbf{Test agreement on complete profiles.} Every Boolean expression depends only on the profile. Conversely, for $S\subseteq M$ form the complete literal test
\[
 L_S=\bigwedge_{m\in S}m\ \wedge\!
       \bigwedge_{m\in M\setminus S}\neg m .
\]
Let \(F_S=\{x:b(x)=S\}\), \(F_S^+=F_S\cap P\) and \(F_S^-=F_S\setminus P\). Form the disagreement witnesses
\[
 W=\coprod_{S\subseteq M}F_S^+\times F_S^- ,\qquad
 \mathcal S_+=\{S:\exists x\in F_S^+\}.
\]
Each \((S,x_+,x_-)\in W\) refutes every Boolean expression in the vocabulary: its truth agrees on the two inputs. If \(W=\varnothing\), then \(P=\bigcup_{S\in\mathcal S_+}F_S\). For the reverse inclusion, a negative input in this union would pair with a positive witness of its profile. The algorithm returns \(\bigvee_{S\in\mathcal S_+}L_S\), with empty disjunction false. On the two effective domains in the theorem the tests decide inhabitation and return witnesses for inhabited sets.


    \item \textbf{Compute the witnesses.} The literal formula uses the original attribute names. Substituting their columns gives a query over input fields, while connecting the named tests to the exposed port gives its neural evaluation by \eqref{eq:concept-incidence}. On finite tables these tests scan rows. In rational linear arithmetic substitute the piecewise affine graphs of \cref{thm:relu-formula}, split Boolean combinations into finite systems of weak and strict affine inequalities, and test feasibility. Elimination pairs each lower bound of a variable with each upper bound, marking the comparison strict when either bound is strict. Equality can first be split into two weak inequalities. Back substitution chooses an allowed rational endpoint, a rational interior point, or a point past a one-sided bound. This gives rational counterexamples whenever a system is feasible. Enumeration can grow exponentially with vocabulary size or gate count.\qedhere
\end{enumerate}
\end{proof}
\paragraph{Descent of a decision.}
Corestrict the profile to \(e:X\twoheadrightarrow B=\operatorname{im}(b)\). In sets the kernel pair \(K=X\times_B X\) has coequalizer \(e\). Thus the Boolean part of the theorem asks for the dashed arrow:
\[
\begin{tikzcd}[column sep=large]
 K\arrow[r,shift left=.6ex,"p_1"]\arrow[r,shift right=.6ex,"p_2"']&
 X\arrow[r,two heads,"e"]\arrow[dr,"\chi_P"']&
 B\arrow[d,dashed,"\overline{\chi}_P"]\\
 &&2 .
\end{tikzcd}
\]
If \(\chi_Pp_1=\chi_Pp_2\), the value assigned to each profile is independent of its representative. This constructs the unique predicate on \(B\). Two predicates on the full cube \(2^M\) can still differ off \(B\). The disjunction in the proof extends it by false there. Conversely a pair in \(K\) with unequal decisions witnesses failure of descent. The positive-conjunction part asks in addition whether \(P\) is fixed by the closure \((-)^{\uparrow\downarrow}\). These categorical tests explain the different failure certificates.

Internally the kernel-pair test always constructs the factor on the reachable image. A Boolean formula in $M$ gives a decidable extension to the full cube. The distinction already occurs for a uniformly indexed context with empty vocabulary. Let $U_i$ and $V_i$ be the disjoint semidecidable events that program $i$ on $i$ returns $0$ and $1$, respectively. Take the input fibre $U_i+V_i$ with target bit $1$ on the first summand and $0$ on the second. The target is constant on each inhabited fibre, so descends to its support. A uniform formula on the whole cube would give a total recursive bit $b(i)$ with $b=1$ on $U$ and $b=0$ on $V$. At an index for $b$ itself either answer contradicts that specification. On the two execution domains in the theorem the tests for $F_S^+$ construct the extension explicitly.

\begin{corollary}[Synthesis of threshold observations]\label{cor:concept-thresholds}
On the rational linear domains of \cref{thm:concept-expressibility}, fix rational affine forms $\ell_j$ on $H$ and a Boolean formula $B$ in $k$ named tests. The set of rational thresholds $\tau\in\Q^k$ for which $B$ explains $P$ is an effectively computed finite union of regions defined by weak or strict rational affine inequalities.
\end{corollary}
\begin{proof}
Set $B_\tau(z)=B((\ell_j(z)\ge\tau_j)_{j=1}^k)$. Construct the set of bad thresholds by projecting
\[
 \{(x,z,\tau)\mid x\in X,\ \Phi_h(x,z),\
                 P(x)\mathbin{\mathrm{xor}}B_\tau(z)\}
\]
onto $\tau$. All comparisons are affine in the joint variables because the coefficients of $\ell_j$ are fixed. The elimination in the theorem computes the projection. Its complement is the required region. Feasibility and back substitution return a threshold tuple from any nonempty region. A weak upper threshold uses the negated direction and threshold. Allowing the directions themselves to vary would introduce products of unknown coefficients and hidden values.
\end{proof}

\subsection{Open ports and their response images}
The two preceding tests admit different universal descriptions. Positive conjunctions are the extents fixed by the incidence closure. For Boolean formulas, \cref{prop:profile-adjoints} gives the largest lower approximation and least upper approximation in the inclusion order; explicitly, \(L(P)\) implies \(P\), and \(P\) implies \(U(P)\). Equality produces the descended decision on reachable profiles.

\paragraph{An open port as substitution.}
In the category of a given query fragment, let \(C:X\times Z\to O\) be the computation downstream of the cut. A way of filling its input and port in a common context \(T\) is a pair \(u:T\to X,\ w:T\to Z\). Wiring the pair into \(C\) gives a family natural in the testing context:
\[
\begin{tikzcd}[column sep=large]
 \Hom(T,X\times Z)\arrow[r,"C\circ(-)"]\arrow[d,"(-)\circ v"']&
 \Hom(T,O)\arrow[d,"(-)\circ v"]\\
 \Hom(T',X\times Z)\arrow[r,"C\circ(-)"']&\Hom(T',O) .
\end{tikzcd}
\]
Conversely a natural transformation of these representables determines \(C\) by its value on \(1_{X\times Z}\). Naturality with respect to \(\langle u,w\rangle\) then recovers every filling. The \hyperref[sec:cut-substitution]{register substitution} of \cref{sec:model-interfaces} constructs this representing arrow; the application following \cref{cor:normal-form-derivative} computes its decorated expansion in the bounded grammar. Ordinary runs are the fillings \(\langle1,h\rangle u\), which test its restriction to the graph of \(h\). The identity filling tests the whole open arrow in \cref{prop:port-interpretation}.
Write \(C_N(x,z)\) for this open continuation. The open list constructs its maximal execution domain. Start with $D_0=X\times Z$ on the legal external interface. At instruction $i$ form its operand tuple $v_i:D_{i-1}\to U_i$ from retained fields. Pull the primitive domain $U_i^{\rm def}\hookrightarrow U_i$ back along $v_i$ to obtain $D_i$, restrict the retained fields, and append the new result. A total gate leaves this domain intact. Equations and branch constraints imposed by the program are pulled back where they are tested. Induction on the list proves that a family of fillings factors through $D_n$ exactly when each executed instruction is defined. A protocol can further restrict the experiment by intersecting its replacement relation with $D_n$. For $X=1$, $Z=2$, $h=0$ and $C(z)=z$, the singleton protocol and the full protocol have the same uncut execution and different available responses. Thus the choice of experiment has independent observable content. Ordinary execution is $C_N(x,h(x))$; replacing $z$ by $h(x')$ gives an interchange intervention from a second input. Multiple aligned ports can be cut simultaneously. These operations implement the declared interventions in the sense of \cite{causal-abstraction}.
\begin{proposition}[Checking an interpretation at a named port]\label{prop:port-interpretation}
In the rational sum-comprehension fragment, the cut program has a canonical query across database shapes. Two programs with aligned intervention ports agree under every replacement on their full rational field interfaces exactly when these canonical queries agree. At fixed shape in the rational affine--ReLU fragment, a proposed polyhedral predicate about the hidden fields can be checked for agreement with an input predicate of rational linear arithmetic. One can also test whether a proposed rational affine decoder $d(z)$ retains all information used by the downstream observable under the admitted replacements.
\end{proposition}
\begin{proof}\leavevmode
\begin{enumerate}[titledsteps]
    \item \textbf{Open the polynomial program.} The polynomial cut program has one additional typed field and uses the same term formers. Apply \cref{thm:support-readback} at the enlarged interface; syntactic bounds are obtained from that program. 
    \item \textbf{Test the rational conditions.} For the rectifier case, substitute the piecewise affine fields from \cref{thm:relu-formula} into the proposed hidden predicate and seek an input where its truth differs from the proposed input predicate. For the sufficiency test, seek a solution to
\[
      d(z)=d(z'),\qquad C_N(x,z)\ne C_N(x,z')
\]
with the same $x$ in both executions and the declared rational polyhedral domain constraints. The observable may be a score, a tuple or a decision, with its corresponding equality test. Enumerate the finitely many linear cases of both executions and split disequalities into strict inequalities. Rational linear elimination decides feasibility and recovers a rational witness by back substitution (for an open interval between finite bounds use its midpoint). 
    \item \textbf{Construct the factored observable.} If the test is infeasible the observable is constant on every admitted decoder fibre. Its graph in the reduced vocabulary is
\[
       G(x,c,y)\ \Longleftrightarrow\
       \exists z\,[\,z\text{ admissible at }x\ \wedge\
                         d(z)=c\ \wedge\ C_N(x,z)=y\,].
\]
It has a unique output on the image, and elimination gives a formula there. Conversely any factorization through $(x,d(z))$ forces equality on those fibres. This certifies a proposed decoder at the named port even if some external fields bypass it.\qedhere
\end{enumerate}
\end{proof}

\begin{proposition}[The coarsest sufficient description under intervention]\label{prop:response-quotient}
Let $Z$ be a common domain of replacements, $X$ the bypass inputs, and $C:X\times Z\to O$ the observable of the cut program. Define
\[
 z\sim_C z'\quad\Longleftrightarrow\quad
           \forall x\in X\ C(x,z)=C(x,z').
\]
The quotient map $q_C:Z\to Z/{\sim_C}$ retains exactly the distinctions required by $C$. A decoder $d:Z\to D$ is sufficient, meaning that $C(x,z)$ factors through $(x,d(z))$ on $X\times d(Z)$, precisely when there is a map $r:d(Z)\to Z/{\sim_C}$ with $q_C=r\circ d$ (corestricting $d$ to its image). This map is unique. In the fixed-shape rational affine--ReLU case with rational linear domains and a score or decision observable, $\sim_C$ admits an effective formula on pairs of replacements.
\end{proposition}
\begin{proof}
Transpose \(C\) to \(\widehat C:Z\to O^X\) and put \(Q=\operatorname{im}\widehat C\). The two maps in the image factorization are \(q(z)=\widehat C(z)\) and inclusion \(j:Q\hookrightarrow O^X\). Then \(z\sim_C z'\iff q(z)=q(z')\), and
\[
\begin{tikzcd}[column sep=large]
 X\times Z\arrow[r,"1\times q"]\arrow[dr,"C"']&
 X\times Q\arrow[d,"{\mathrm{ev}\circ(1\times j)}"]\\
 &O .
\end{tikzcd}
\]
For \(d:Z\twoheadrightarrow d(Z)\), sufficiency is equivalent to \(\ker d\subseteq\ker q\). The graph of the resulting factor \(r\) consists of \((d(z),q(z))\). The kernel inclusion gives functionality and surjectivity gives totality and uniqueness. For rational affine--ReLU programs eliminate \(x\) from the formula for \(C(x,z)\ne C(x,z')\), then complement the projection.

\end{proof}
\begin{comparison}[Change of observation at a common port]\label{comp:response}
Restricting the bypass contexts and postcomposing the observable induce quotient maps of response images. Pairing two observations of a common replacement intersects their kernel pairs. Formation of the response image in a parameter slice commutes with reindexing in the parameter.
\end{comparison}
\begin{proof}\leavevmode
\begin{enumerate}[titledsteps]
    \item \textbf{Change the observation.} Write \(q_C:Z\twoheadrightarrow Q_C\) for the response image in \cref{prop:response-quotient}. Changing the available bypass contexts by \(a:X'\to X\) and observing the output through \(v:O\to O'\) gives \(C'=vC(a\times1_Z)\). Define \(t:O^X\to (O')^{X'}\) by \(t(s)=vsa\). Transposition and image factorization give the diagram
\begin{equation}\label{eq:response-change}
\begin{tikzcd}[column sep=large]
 Z\arrow[r,two heads,"q_C"]\arrow[d,equal]&Q_C\arrow[r,hook,"j_C"]\arrow[d,two heads,"\bar t"]&
 O^X\arrow[d,"t"]\\
 Z\arrow[r,two heads,"q_{C'}"']&Q_{C'}\arrow[r,hook,"j_{C'}"']&(O')^{X'} .
\end{tikzcd}
\end{equation}
The middle arrow sends \(q_C(z)\) to \(q_{C'}(z)\). Equality of two \(q_C\)-values implies equality after applying \(t\), and surjectivity follows from \(q_{C'}=\bar t q_C\). The epi \(q_C\) proves identity and composition for these induced arrows by cancellation. Consequently a witness separating two states after the change also separates them before the change.


    \item \textbf{Pair two outputs.} There is a distinct law for observing two outputs at the same port. For \(C_i:X\times Z\to O_i\) put \(C=\langle C_1,C_2\rangle\). The canonical isomorphism \((O_1\times O_2)^X\cong O_1^X\times O_2^X\) and uniqueness of image factorization give
\[
 Q_C\ \cong\ \operatorname{im}\bigl(Z\xrightarrow{\langle q_{C_1},q_{C_2}\rangle}Q_{C_1}\times Q_{C_2}\bigr),
 \qquad \ker q_C=\ker q_{C_1}\cap\ker q_{C_2}.
\]
Here kernels denote kernel pairs as subobjects of \(Z\times Z\). Indeed two replacements have the same paired response exactly when both component responses agree. The image in this formula retains correlations caused by observing the same replacement twice.


    \item \textbf{Reindex the response image.} For the base-change assertion, first consider sets.
For a parameter object \(S\), do the construction in \(\mathbf{Set}/S\). Write \(X,Z,O\) for families over \(S\) and take \(C:X\times_S Z\to O\) over \(S\). The exponential and its image have fibres \((O_s)^{X_s}\) and \(Q_s=\operatorname{im}(Z_s\to (O_s)^{X_s})\). For \(r:S'\to S\), the square of total spaces
\[
\begin{tikzcd}[column sep=large]
 r^*Q\arrow[r]\arrow[d]&Q\arrow[d]\\
 S'\arrow[r,"r"']&S
\end{tikzcd}
\]
is cartesian, and \(r^*Q\) is the response quotient for \(r^*C\): its fibre at \(s'\) is precisely \(Q_{r(s')}\) with the same quotient map. In a slice of \(\E\) one obtains the same result without choosing global points. Write \(\Sigma_r\dashv r^*\) for composition with \(r\) and pullback. For \(T\to S'\) there are natural bijections
\[
\begin{aligned}
 \Hom_{/S'}(T,r^*(O^X))
 &\cong \Hom_{/S}(\Sigma_rT\times_S X,O)\\
 &\cong \Hom_{/S'}(T\times_{S'}r^*X,r^*O).
\end{aligned}
\]
The second uses \(\Sigma_r(T\times_{S'}r^*X)\cong\Sigma_r T\times_S X\), by pullback pasting. Yoneda identifies \(r^*(O^X)\) with \((r^*O)^{r^*X}\) in the new slice. Pullback preserves both parts of the regular epi--mono image factorization in a topos. Applying it to the transposed response gives the claimed base-change isomorphism. Uniqueness of image comparison gives its coherence under successive base changes.\qedhere
\end{enumerate}
\end{proof}
\subsection{The response quotient as a fibred reflection}\label{sec:pairing-reflection}
A change of interface may also map the replacement states themselves. Over \(S\), form a category \(\mathsf{Port}_S\) in the slice \(\mathcal B/S\) of \cref{sec:indexed-readback}. Its objects are pairings \(C:X\times_S Z\to O\). For \(D:Y\times_S W\to V\) an arrow \(C\to D\) consists of
\[
 a:Y\to X,\qquad d:Z\to W,\qquad v:O\to V,\qquad
 D(y,dz)=v(C(ay,z)).
\]
The equation holds as maps \(Y\times_S Z\to V\). Replacement states move forwards and a test context for the target pulls back to a test for the source. For fixed outcomes and \(v=1\), this is the Chu morphism law, with replacements as points and contexts as tests \cite{pratt-chu}. Two successive triples compose as \((a_1a_2,d_2d_1,v_2v_1)\), by substitution in the equation. Identities use identity maps. Call \(C\) \emph{separable} when its transpose \(\widehat C:Z\to O^X\) is monic in the slice. Write \(i_S:\mathsf{Sep}_S\hookrightarrow\mathsf{Port}_S\) for the full subcategory of such pairings. In matrix language this is separation of the rows.

\begin{theorem}[Reflection of port pairings over the parameter base]\label{thm:fibred-reflection}
The response image defines a left adjoint \(L_S\dashv i_S\). For substitution \(r:T\to S\), the canonical isomorphism \(r^*L_S\cong L_T r^*\) respects the units and composes coherently. Consequently these adjunctions give a reflection over \(\mathcal B\) whose reflector preserves cartesian lifts.
\end{theorem}
\begin{proof}\leavevmode
\begin{enumerate}[titledsteps]
    \item \textbf{Reflect a port pairing.} Factor \(\widehat C=j_Cq_C\) as a regular epimorphism followed by a monomorphism. Define \(L_SC\) by transposing \(j_C:Q_C\hookrightarrow O^X\). Its unit is \((1_X,q_C,1_O)\). For a triple \(C\to D\), let \([a,v]\) take \(s\) to \(vsa\). The outer rectangle commutes by the equation for the triple:
\[
\begin{tikzcd}[column sep=large]
 Z\arrow[r,two heads,"q_C"]\arrow[d,"d"']&Q_C\arrow[r,hook,"j_C"]\arrow[d,dashed,"\bar d"]&O^X\arrow[d,"{[a,v]}"]\\
 W\arrow[r,two heads,"q_D"']&Q_D\arrow[r,hook,"j_D"']&V^Y .
\end{tikzcd}
\]
On the kernel pair of \(q_C\), both composites into \(V^Y\) agree. Monicity of \(j_D\) gives equality after \(q_Dd\), so the coequalizer property yields the unique \(\bar d\) making the left square commute. Cancelling the epi \(q_C\) gives the right square, functoriality and naturality of the unit. When \(D\) is separable, \(q_D\) is invertible. The inverse to precomposition with the unit sends \((a,d,v)\) to \((a,q_D^{-1}\bar d,v)\). Its two inverse laws follow from \(\bar d q_C=q_D d\) and epimorphic cancellation. This proves the natural hom-set bijection.


    \item \textbf{Transport the reflection across the base.} Pulling back all three objects and the pairing along \(r:T\to S\) reindexes \(\mathsf{Port}_S\) into \(\mathsf{Port}_T\). Pullback pasting supplies the pseudofunctor coherence. In its Grothendieck construction an arrow over \(r\) from \(A\) over \(T\) to \(D\) over \(S\) is precisely a vertical arrow \(A\to r^*D\). The arrow \(r^*D\to D\) represented by the identity is cartesian, since a further arrow over \(ru\) is exactly a vertical arrow into \(u^*r^*D\). By \cref{comp:response}, reindexing preserves the exponential and the response image. It carries \(q_C\) to the unit at \(r^*C\) up to the unique image isomorphism. These isomorphisms agree on identities and on composites by uniqueness. For separable \(D\), the bijection over \(r\) is now the fibre bijection with target \(r^*D\), which is again separable. This proves the assertion about total categories and cartesian lifts.\qedhere
\end{enumerate}
\end{proof}
The variance matters even for finite tables. Start with one context in which \(z_0,z_1\) have the same output. Include that context into a larger family with a second context distinguishing them. The identity on states and outputs respects the smaller table. It cannot induce a map from the old one-point quotient to the two-point quotient compatible with both state projections. Restriction from the larger family to the smaller one has exactly the context map required by \cref{thm:fibred-reflection}.

\begin{corollary}[Least sufficient refinement of a proposed decoder]\label{cor:repair-decoder}
For a decoder \(d:Z\to H\) on a common replacement domain, the image of \(\langle d,q_C\rangle\) gives the coarsest sufficient refinement of \(d\). In the poset of equivalence relations on \(Z\) ordered by inclusion, the inclusion of sufficient relations into all relations of this poset has right adjoint \(R\mapsto R\cap\ker q_C\).
\end{corollary}
\begin{proof}
Corestrict decoders to their images. A refinement \(e\) of \(d\) is sufficient precisely when both \(d\) and \(q_C\) factor through \(e\). Pair the two factors. Their pairing lands in \(\operatorname{im}\langle d,q_C\rangle\) because \(e\) is regular epic. This gives the unique factor from \(e\)'s image to the proposed refinement and proves minimality. In kernel pairs the displayed image has relation \(\ker d\cap\ker q_C\). More generally a sufficient relation \(K\) satisfies \(K\le\ker q_C\), so \(K\le R\) is equivalent to \(K\le R\cap\ker q_C\). Every internal equivalence relation here is effective. This proves the adjunction without choosing a representative of any class.
\end{proof}

There is also an effective description at the level of pairs. For the rational-linear interfaces of \cref{prop:port-interpretation}, equality in the refined quotient is the conjunction
\[
 d(z)=d(z')\quad\wedge\quad
 \neg\exists x\,\bigl(C(x,z)\ne C(x,z')\bigr).
\]
The elimination procedure of that proposition computes this relation. Its computation requires no choice of coordinates for the quotient object \(Q_C\).

\paragraph{Recovering the global quotient from a cover of contexts.}
For a finite jointly epimorphic family \(a_i:X_i\to X\) over the same base and \(C_i=C(a_i\times 1_Z)\), put \(q_i=q_{C_i}\). Then
\[
 \ker q_C=\bigcap_i\ker q_i,\qquad
 Q_C\ \cong\ \operatorname{im}\!\left(Z\xrightarrow{\langle q_i\rangle_i}\prod_i Q_{C_i}\right).
\]
Indeed \(\coprod_i X_i\to X\) is regular epic and remains so after product with any test object in the slice. Two generalized responses equal on every \(X_i\) are therefore equal on \(X\) by epimorphic cancellation. The intersection formula follows. Both maps onto the stated images coequalize the same kernel pair, giving the canonical isomorphism. Even for a two-piece cover the image can be proper: if each context returns the same Boolean replacement itself, the paired image is the diagonal in \(2\times2\). The comparison uses one replacement shared across all pieces of the cover.

A family of decoders indexed by a base and a single decoder shared across different admissibility domains have different universal properties.
\paragraph{Varying intervention domains.}
If admissibility varies with \(x\), write \(A=\{(x,z):z\in Z_x\}\) and let \(C:A\to O\). A fixed decoder \(d\) suffices exactly when
\[
 \forall x,\ z,z'\in Z_x\quad d(z)=d(z')\Longrightarrow C(x,z)=C(x,z').
\]
The coarsest quotient in each fibre is \(\operatorname{im}(C_x)\). A single coarsest decoder on \(Z\) can fail to exist. Take three fibres \(\{a,c\},\{a,b\},\{b,c\}\), with outputs \((0,1),(0,0),(0,0)\) in the displayed orders. The partitions \(ab\mid c\) and \(a\mid bc\) both suffice. A quotient coarser than both identifies \(a=b=c\), contradicting the first fibre. Thus the product domain in the response-function construction has mathematical content. The paired test in \cref{prop:port-interpretation} still decides sufficiency of each proposed decoder.
On a finite presented table there is an exact existence test. Put an edge between \(z,z'\) when some common fibre admits both with different outputs. A decoder is sufficient exactly when each of its fibres is an independent set. Write \(z\,R\,z'\) when there is no such edge. Every nonedge can be merged by a sufficient partition (merge just that pair). A coarsest sufficient partition must therefore contain \(R\), while sufficiency forces its equivalence relation to be contained in \(R\). Thus one exists exactly when \(R\) is transitive. In that case the decoder is \(Z\to Z/R\). In the three-fibre example \(aRb\) and \(bRc\), while \(aRc\) fails.

\subsection{How far apart are two intervention responses?}\label{sec:response-distance}
For the response image in \cref{prop:response-quotient}, a metric on the observed values measures the size of the distinction made by all bypass inputs. We first work with sets and ordinary bounded real-valued metrics.

\begin{proposition}[Distance on the response image]\label{prop:response-distance}
For a common replacement domain \(Z\), a set \(X\) of bypass inputs and \(C:X\times Z\to(O,d_O)\), put
\[
 D_C(z,z')=\sup_{x\in X}d_O(C(x,z),C(x,z')),\qquad \sup\varnothing=0
\]
for this supremum in the nonnegative reals. It is a pseudometric with zero relation \(\ker q_C\). It induces a metric on the response image \(Q_C\). For any decoder \(e:Z\to E\) and tolerance \(\varepsilon\ge0\), every two replacements with the same code differ in output by at most \(\varepsilon\) in every context if and only if \(D_C\le\varepsilon\) on \(\ker e\). At zero this is the sufficient-decoder criterion.
\end{proposition}
\begin{proof}
Equip \(O^X\) with the supremum metric. Boundedness gives an upper bound, and the triangle inequality follows by bounding each pointwise sum by the sum of the two suprema. Supremum zero is equivalent to pointwise equality, including when \(X\) is empty, in which case \(O^X\) is a singleton. Pull this metric back along the transposed response \(\widehat C\). Its restriction to the subset \(Q_C=\operatorname{im}\widehat C\) is a metric, and its pullback along \(q_C\) is exactly \(D_C\). The tolerance test follows from the defining upper-bound property on each pair in \(\ker e\); at zero apply \cref{prop:response-quotient}.
\end{proof}
For an internal experiment with located pointwise distances one can always take the supremum in the nonnegative lower reals $Q$ of \cref{sec:eff}. Its cut contains every negative rational and contains $r\ge0$ exactly when $\exists x\, [r<d_O(C(x,z),C(x,z'))]$. The pointwise inequalities imply the triangle inequality by monotonicity of lower-real addition, and vanishing is equivalent to equality at every context by separatedness of the output metric. A located value for every such experiment would decide arbitrary truth values: take a subterminal bypass $X=p$, two replacement bits, and $C(x,z)=z$ with discrete distance. Locatedness at $1/3<2/3$ gives a lower bound implying $p$ or an upper bound implying $\neg p$. The rational program below constructs the located distance on its boxes.

Restricting the bypass contexts decreases \(D_C\), and taking their union takes the supremum of the restricted distances. Postcomposition by a nonexpansive map decreases \(D_C\). For two observations on the same contexts and replacements, give the paired output the maximum metric. The resulting distance is \(\max(D_{C_1},D_{C_2})\), since each pointwise maximum is bounded by this number and dominates each component pointwise. These operations recover the kernel comparisons of \cref{comp:response} at distance zero.

\begin{corollary}[A rational term for the response distance]\label{cor:dmv-response}
Take a fixed finite rational affine--ReLU program with frozen parameters and a scalar score at its open port. On nonempty compact rational boxes for the bypass and replacement coordinates, normalize the score and the coordinates as above. From the open program one can compute a DMV term in the two replacement tuples alone whose value is \(D_{C^\sharp}(z,z')\) for the absolute-value metric. For rational replacements the construction also yields a rational maximizing context. In the original units \(D_C=(u-l)D_{C^\sharp}\) for the chosen nonconstant score scale; a constant score has zero distance.
\end{corollary}
\begin{proof}\leavevmode
\begin{enumerate}[titledsteps]
    \item \textbf{Construct the activity cells.} Let \(c(x,z)\) be the normalized score term. Form the term
\[
 b(x,z,z')=
 (c(x,z)\ominus c(x,z'))\vee(c(x,z')\ominus c(x,z)).
\]
It evaluates to the absolute score difference. For any such term \(b(x,w)\) with \(x\in[0,1]^n\), \(w=(z,z')\), build its closed rational activity cells by splitting at each piecewise-affine operation. On a cell the value is affine, and the domain has a presentation
\[
 A x+B w\le c
\]
including the bounds of the cube. 
    \item \textbf{Reach a vertex by rational moves.} At rational $w,x$ use rational comparisons to select a cell containing $(x,w)$. If the tight $x$-rows have rank less than $n$, elimination finds a nonzero rational vector orthogonal to them. Choose its sign to make the affine objective nondecreasing. Move to the nearest blocking row in this direction by taking the minimum of the finitely many positive rational slacks divided by positive directional coefficients. A box row blocks every nonzero direction. The new tight row is independent of the old tight rows and the point stays feasible. After at most $n$ moves this gives a vertex with objective at least the starting value. 
    \item \textbf{Enumerate the clipped candidates.} Thus it suffices to enumerate invertible $n$-row submatrices $A_I$ of all cells. Form the terms
\[
 v_I(w)=\operatorname{clip}\bigl(A_I^{-1}(c_I-B_I w)\bigr)
\]
by the clipped-affine construction, coordinatewise. The index \(I\) here includes its cell. These terms are defined on the entire parameter cube. When \(n=0\), take the single empty tuple. The promised elimination is the finite join
\[
 m(w)=\bigvee_I b\bigl(v_I(w),w\bigr).
\]

    \item \textbf{Prove the least-upper-bound property.} Each clipped tuple lies in the box, so evaluating the whole term there gives an actual score difference. The rational vertex construction proves $b(x,w)\le m(w)$ at rational arguments. Both sides are continuous with computable Lipschitz bounds obtained by adding and multiplying rational bounds along the affine and lattice gates. Rational approximation in the box extends the inequality to located arguments. Every upper bound at $w$ bounds each actual candidate value, hence their finite lattice join. Thus $m(w)$ is the least upper bound. In classical sets this finite join is attained by a candidate even at irrational $w$. At rational $w$ the finite list has decidable comparisons, so its largest entry supplies a rational maximizing bypass tuple. Undoing the coordinate and score changes gives the assertions.\qedhere
\end{enumerate}
\end{proof}

For the coordinate projection \(\pi:X\times W\to W\), write \(\exists_\pi b=m\). In the pointwise order on DMV term functions, \(\exists_\pi b\le v\) if and only if \(b\le v\pi\). For a tuple \(g:W'\to W\) of term functions,
\[
 (\exists_\pi b)\circ g=\exists_{\pi'}\bigl(b\circ(1_X\times g)\bigr),
 \qquad \pi':X\times W'\to W'.
\]
Both sides take the least upper bound over the same \(x\) at \(g(w')\). This is a numerical instance of the projection adjunction used for predicates in \cref{prop:profile-adjoints}. Divisible propositional \L{}ukasiewicz logic also admits this quantifier elimination \cite{pkhakadze-qe}. The polyhedral construction here operates on the presented boxes and returns a term for this particular intervention observable.

\paragraph{Division and the bypass variable both matter.}
Using the score in the units \([0,1]\), let \(\rho(t)=\max(0,t)\), take \(x,z\in[0,1]\), and open the score \(C(x,z)=\rho(x-\rho(2x-z))\). From left to right in \(x\) its values are \(x\) on \([0,z/2]\), \(z-x\) on \([z/2,z]\), and \(0\) on \([z,1]\). In fact the distance on the entire replacement interval is
\[
 D_C(z,z')=|z-z'|\wedge\delta_2(z\vee z'),
 \qquad a\wedge b=\neg(\neg a\vee\neg b).
\]
To see this, order the replacements as \(a\le b\). At each \(x\), \(C(x,\cdot)\) is nondecreasing and \(1\)-Lipschitz, by the two rectifications. Hence the difference is at most \(b-a\), and it is at most the height \(b/2\) of the larger tent. At \(x=b/2\) the smaller tent has height \(\max(0,a-b/2)\), attaining the lesser of these bounds. Absolute difference here is the join of the two truncated differences, so the display is a DMV term. In particular \(D_C(z,0)=z/2=\delta_2 z\). At \(z=1\) this is \(1/2\); any ordinary MV term in \(z\) with only \(0,1\) as constants returns a Boolean value at that input, by induction on its operations. Nor can a fixed finite list of bypass probes compute this distance for all \(z\). If the list has a positive probe, choose \(0<z\) smaller than its least positive one; otherwise take \(z=1\). All listed differences vanish, while the true distance is positive. Figure~\ref{fig:tent-response} shows an attaining context for two replacements.

\begin{figure}[!htbp]
    \centering\small
    \begin{tikzpicture}[x=6cm,y=4cm]
        \draw[figwire] (0,0) -- (1.08,0) node[right] {$x$};
        \draw[figwire] (0,0) -- (0,.61) node[above] {score};
        \foreach \x/\t in {0/0,.5/{1/2},.75/{3/4},1/1}
            \draw (\x,0) -- ++(0,-.018) node[below] {$\t$};
        \foreach \y/\t in {.25/{1/4},.5/{1/2}}
            \draw (0,\y) -- ++(-.014,0) node[left] {$\t$};
        \draw[black!65,line width=1.3pt] (0,0) -- (.5,.5) -- (1,0);
        \draw[black,dash pattern=on 3pt off 2pt,line width=.6pt] (0,0) -- (.375,.375) -- (.75,0) -- (1,0);
        \draw[{Latex[length=1.35mm]}-{Latex[length=1.35mm]},thin]
            (.5,.255) -- (.5,.495);
        \draw[black!65,line width=1.3pt] (1.14,.40) -- (1.30,.40) node[black,right=3pt] {$C(x,1)$};
        \draw[black,dash pattern=on 3pt off 2pt] (1.14,.28) -- (1.30,.28) node[right=3pt] {$C(x,3/4)$};
    \end{tikzpicture}
    \caption{Two replacements in the tent example. The rising sides coincide up to $x=3/8$. At $x=1/2$ the scores are $1/2$ and $1/4$, so the vertical segment realizes the response distance $D_C(1,3/4)=1/4$.}
    \label{fig:tent-response}
\end{figure}


\subsection{Two learned units, from observations to interventions}
\begin{example}[What two trained hidden units encode]\label{ex:trained-discrepancy}\leavevmode
\begin{enumerate}[titledsteps]
    \item \textbf{Freeze and evaluate the trained model.} A record has two integer readings $a,b\in\{-3,\ldots,3\}$. The target is $\mathbf1\{|a-b|\ge2\}$. We trained a dense network with two inputs, two rectified hidden units and one affine score thresholded at zero. It received only the two readings and the task label. Training used the 38 pairs with $(a+2b)\bmod4\ne0$; the remaining 11 pairs were held out. For reproducibility, in the order $(w_{11},w_{12},w_{21},w_{22},b_1,b_2,v_1,v_2,c)$ initialize from NumPy's PCG64 generator with seed 13 using nine independent normal draws of standard deviation $1/2$. Run 1,000 full-batch Adam updates on mean binary cross entropy with logits, step size $0.03$, decay factors $0.9,0.999$, bias correction and denominator offset $10^{-8}$; take derivative zero at a zero rectifier input. To fix the model independently of training-library arithmetic, the following integers divided by $2^{53}$ specify its frozen binary64 parameters exactly:
\[
\begin{aligned}
 (w_{11},w_{12})\,2^{53}&=(45120229354350536,-45097680132202112),\\
 (w_{21},w_{22})\,2^{53}&=(21862191961228660,-21859388735148040),\\
 (b_1,b_2)\,2^{53}&=(-3922626985195417,43734115255595168),\\
 (v_1,v_2)\,2^{53}&=(46467639698804896,-44964886573674144),\\
 c\,2^{53}&=53097338566859432 .
\end{aligned}
\]
Write $u_i=w_{i1}a+w_{i2}b+b_i$, $h_i=[u_i]_+$ and $k=-v_2>0$. The positive decision has the following input-only formula, valid on the whole rational plane:
\[
             k u_2\le c \quad\lor\quad k u_2\le c+v_1u_1 .
\]
Indeed $v_1,c>0$, so $k[u_2]_+\le c+v_1[u_1]_+$ is equivalent to $ku_2\le\max(c,c+v_1u_1)$. Splitting this comparison gives the two indicated affine half-spaces. Each hidden value has graph $(u_i\le0\wedge h_i=0)\lor(u_i\ge0\wedge h_i=u_i)$. These formulas read the frozen parameters and require no examples.

On the entire integer task domain the hidden \emph{activity predicates} simplify further:
\[
 h_1>0\ \Longleftrightarrow\ a-b\ge1,\qquad
 h_2>0\ \Longleftrightarrow\ a-b\ge-2 .
\]
There is a short integer certificate for the two signs. Set \(D=2^{53}\), \(U_i=Du_i\), and write \(K_1=1,K_2=-2\). For each \(t=a-b\), substitute \(a=t+b\) and take the endpoints of
\(\max(-3,-3-t)\le b\le\min(3,3-t)\), since \(U_i(t+b,b)\) is affine in \(b\). Taking extrema over the thirteen integer differences gives
\[
\begin{array}{c|r|r}
 i&\max_{t<K_i} U_i&\min_{t\ge K_i} U_i\\ \hline
 1&-3854979318750145&41129954702709847\\
 2&-21844050949848952&6928107057228
\end{array}
\]
For the score, let \(V_i=Dv_i\), \(C=Dc\), and \(S=DC+V_1[U_1]_++V_2[U_2]_+=D^2s\). Enumerating \(a,b=-3,\ldots,3\) gives \(S<-6D^2\) on the 19 negative labels and \(S>5D^2\) on the 30 positive labels. Splitting the same comparisons by \((a+2b)\bmod4\) gives 38 correct training and 11 correct held-out decisions. At \((1,0),(2,0)\) both activity bits are one, while the decisions are zero and one.


    \item \textbf{Round the parameters and compare firing profiles.} A second frozen model is obtained by rounding each parameter to the nearest quarter. With $t=a-b$, its exact hidden fields and score are
\[
 h'_1=[5t-\tfrac12]_+,\qquad h'_2=[\tfrac52t+\tfrac{19}4]_+,\qquad
 s'=\tfrac{21}4h'_1-5h'_2+6 .
\]
The same half-space elimination gives $s'\ge0$ exactly when $t\le-71/50$ or $t\ge163/110$. Thus the rounded model gives the same 49 task decisions. Its second unit is active on the grid exactly when $t\ge-1$. There are five pairs with $t=-2$; on every one the original unit was active and the rounded one is zero. The rounded model therefore preserves every task decision while changing a hidden activity predicate.

For the rounded network, the two ordinary firing bits partition the grid as follows. Each difference $t$ has multiplicity $7-|t|$.
\[
\begin{array}{cclrr}
\toprule
 \mathbf1_{h'_1>0} & \mathbf1_{h'_2>0} &\text{differences }t
 &\text{negative decisions}&\text{positive decisions}\\
\midrule
 0&0&t\le -2 &0&15\\
 0&1&t=-1,0  &13&0\\
 1&1&t\ge 1  &6&15\\
\bottomrule
\end{array}
\]
Every Boolean expression in those two bits is constant on each row. Choosing the majority decision in each row is achievable by a disjunction of complete profiles and yields exactly $15+13+15=43$ correct decisions out of 49. Figure~\ref{fig:firing-decisions} locates the six errors on the input grid.

\begin{figure}[!htbp]
    \centering\small
    \begin{tikzpicture}[x=.45cm,y=.45cm]
        \foreach \panel/\shift/\heading in {0/0/{Task decisions},1/13/{Best firing-bit decoder}}{
            \begin{scope}[shift={(\shift,0)}]
                \node at (0,5.18) {\heading};
                \draw[figwire] (-3.75,-3.75) -- (3.9,-3.75) node[right] {$a$};
                \draw[figwire] (-3.75,-3.75) -- (-3.75,3.8) node[above] {$b$};
                \foreach \k in {-3,0,3}{
                    \draw (\k,-3.75) -- ++(0,-.14) node[below,font=\scriptsize] {$\k$};
                    \draw (-3.75,\k) -- ++(-.14,0) node[left,font=\scriptsize] {$\k$};
                }
                \foreach \aval in {-3,...,3}{
                    \foreach \bval in {-3,...,3}{
                        \pgfmathtruncatemacro{\tval}{\aval-\bval}
                        \pgfmathtruncatemacro{\positive}{abs(\tval)>=2 || (\panel==1 && \tval==1)}
                        \ifnum\positive=1
                            \fill (\aval,\bval) circle[radius=1.6pt];
                        \else
                            \draw (\aval,\bval) circle[radius=1.6pt];
                        \fi
                        \ifnum\panel=1
                            \ifnum\tval=1 \draw (\aval,\bval) circle[radius=3.2pt];\fi
                        \fi
                    }
                }
            \end{scope}
        }
    \end{tikzpicture}
    \caption{Each point is one input pair. Filled points mean a positive decision. In the rounded network the firing profile $(1,1)$ groups the six inputs with $a-b=1$ together with the fifteen inputs with $a-b\ge2$. Choosing its majority label gives the six circled errors on the right.}
    \label{fig:firing-decisions}
\end{figure}



    \item \textbf{Recover the grid decision from one coordinate.} Now keep only the second coordinate but change its observation language. On the grid $h'_2$ is zero for $t\le-2$, takes the successive values $9/4,19/4,29/4$ for $t=-1,0,1$, and is at least $39/4$ for $t\ge2$. The complete threshold region for
$(h'_2\le\alpha)\vee(h'_2\ge\beta)$ is therefore
\[
       0\le\alpha<9/4,\qquad 29/4<\beta\le39/4 .
\]
Indeed, rejecting the three middle values forces $\alpha<9/4$ and $\beta>29/4$. Accepting zero and $39/4$ then forces the other bounds. Monotonicity on the upper tail proves sufficiency. Thus the two names $L=(h'_2=0)$ and $U=(h'_2\ge39/4)$ give $L\vee U$ as an exact explanation of the task on $X$. In the formal context with precisely these two attributes, each individual tail is closed. Both tails are nonempty and disjoint, so their union has empty intent, whose extent is all 49 records. Their join in the concept lattice is the whole grid. Their logical disjunction denotes the 30 positive records. This is a target that requires the Boolean stage of \cref{thm:concept-expressibility} in the stated vocabulary.

The reachable hidden pairs on the whole rational plane have an explicit parametrization. Write \(v=h'_2\). If \(v>0\), the second preactivation gives \(5t-\tfrac12=2v-10\). If \(v=0\), then \(t\le-19/10\) and \(h'_1=0\). Conversely each \(v\ge0\) occurs by taking \(t=\tfrac25(v-19/4)\). Hence
\[
 \operatorname{im}h'=\{([2v-10]_+,v):v\in\Q_{\ge0}\},\qquad
 s'=\begin{cases}6-5v&0\le v\le5,\\(11v-93)/2&v\ge5.\end{cases}
\]
On this entire image the score factors through \(v\), and the decision is \(v\le6/5\) or \(v\ge93/11\). The threshold region above concerns the 49-point grid.


    \item \textbf{Cut the hidden pair.} Cut the hidden pair and admit all of \(\Q_{\ge0}^2\). At \((0,2)\) and \((1,2)\) the cut scores are \(-4\) and \(5/4\). The first pair lies on the displayed graph and the second lies outside it. With \(d(z)=21z_1-20z_2\), the two response quotients are determined by
\[
 C_{\rm score}(z)=6+d(z)/4,\qquad C_{\rm decision}(z)=\mathbf1_{d(z)\ge-24}.
\]
Injectivity of \(u\mapsto6+u/4\) gives the score kernel pair \(\ker d\). The decision quotient has the two nonempty threshold classes witnessed above.

On the common replacement box \([0,2]^2\), the raw score ranges from \(-4\) to \(33/2\). There is no remaining bypass dependence, so for the displayed pair the response distance is \(21/4\) score units. Its normalized value is \((21/4)/(41/2)=21/82\). In unit coordinates \(u=z_1/2,\ v=z_2/2\), the normalized open score has the explicit term
\[
 c(u,v)=(\delta_{41}u)^{\oplus21}\oplus(\delta_{41}\neg v)^{\oplus20}
       =\frac{21u+20(1-v)}{41}.
\]
Its absolute difference on \((0,1)\) and \((1/2,1)\) is \(21/82\). Thus \cref{cor:dmv-response} gives both the term metric on this response image and the raw-unit separation of the two opposite decisions.
\end{enumerate}
\end{example}

\subsection{Graded observations at the same port}\label{sec:graded-concepts}
Some observations come with a specified membership scale. For instance, an attribute may return one of three ordered grades with two named cut levels. Fix a complete integral commutative residuated lattice $L$ with product $\otimes$, unit $\top$ and residual $\Rightarrow$, so $a\otimes b\le c$ exactly when $b\le a\Rightarrow c$. Give each name a membership map $E_m:H\to L$. In the equipment of $L$-valued relations on sets, the graph of $h$ has value $\top$ on its graph and $\bot$ elsewhere. Composing by a join of products again gives $I(x,m)=E_m(h(x))$. Residuated Galois connections and their enriched distributor formulations are established in \cite{fuzzy-galois,isbell-distributors,multi-concept-lattices}.

\begin{proposition}[Graded extents and reachable states]\label{prop:graded-port}
For $u:X\to L$ and $v:M\to L$, put
\[
 u^\uparrow(m)=\bigwedge_{x\in X}(u(x)\Rightarrow I(x,m)),\qquad
 v^\downarrow(x)=\bigwedge_{m\in M}(v(m)\Rightarrow I(x,m)).
\]
These derivations give an antitone adjunction. Its fixed pairs on inputs and on reachable states are order-isomorphic by pulling the extent back along $q:X\onto R$. The two-element scale recovers \cref{thm:reachable-concepts}.
\end{proposition}
\begin{proof}
The adjunction and the two change-of-domain identities are
\[
\begin{aligned}
 v\le u^\uparrow
  &\iff \forall x,m\quad u(x)\otimes v(m)\le I(x,m)
   \iff u\le v^\downarrow,\\
 v^{\downarrow_X}&=q^*(v^{\downarrow_R}),\qquad
 (q^*w)^{\uparrow_X}=w^{\uparrow_R}.
\end{aligned}
\]
Residuation and commutativity give the first line. For the second, \(I(x,m)=E_m(qx)\) and taking a meet over a surjective reindexing leaves it unchanged. Consequently \(v^{\downarrow_X\uparrow_X}=v^{\downarrow_R\uparrow_R}\). Each fixed intent determines its extent by the downward derivation, giving inverse order-preserving maps.
\end{proof}
With finite $X,M,L$ and presented operation tables these meets are finite computations. On the rational domains of \cref{thm:concept-expressibility}, a finite chain of grades with linear-definable level tests also permits symbolic evaluation. For a given finite intent $v$, compute its extent by combining the finitely many grade tests with the tables of $\Rightarrow$ and meet. To recover the common grades, test each proposed bound $\lambda\le u^\uparrow(m)$ by seeking an $x$ at which $\lambda\otimes u(x)>I(x,m)$, for $u$ given by such level tests. Linear feasibility settles each query. For a general scale and infinite domains the displayed meets give the semantics. A membership grade belongs to an individual observation. The frequency with which a vocabulary predicts correctly requires a law on inputs, introduced in \cref{prop:profile-risk}.

\section{Operational grounding at a hidden port}\label{sec:ground-operational}

In \cref{sec:hidden-concepts}, descent decides whether a supplied target can be expressed at a port. To discover which predicates have an environmental interpretation, equip each reached history with two observations. One reads a named neural port at the current history. The other continues the interaction by a word of actions and reads the environment. A port predicate and an environmental predicate express the same property when their inverse images on histories agree. Which of these common predicates can still be read at the port after any specified sequence of actions? The answer uses the response image of the environmental readout. A learned forecast can be compared with that readout on the same histories.

\subsection{Response images and common predicates}
Work in a fibre $\mathcal C=\mathcal B/S$ of \cref{sec:indexed-readback}, with $\mathcal B=\mathbf{Set}$ or $\mathbf{Eff}$. Products, exponentials and subobjects below are in this fibre. Let $A$ be the object of actions and $M=\operatorname{List}(A)$ its free monoid. A deterministic interaction is a coalgebra $X\to O\times X^A$. Its components give the environmental readout $o:X\to O$ and the update for one action. Iteration gives a right action $\tau:X\times M\to X$, written $xw=\tau(x,w)$. We use the reached carrier. From a chosen reset state it is the image of the orbit map $w\mapsto x_0w$, with its induced action. One may equally use the histories themselves as carrier, acting by concatenation.

Transpose the experiment $(w,x)\mapsto o(xw)$ and take its image, as in \cref{prop:response-quotient}. Corestrict the named port to its image as well. This gives the span and the inclusion
\[
 H\xleftarrow{\ h\ } X\xrightarrow{\ q\ } Q
       \xhookrightarrow{\ j\ } O^M,
 \qquad (jq)(x)(w)=o(xw).
\]
Both $h$ and $q$ are regular epimorphisms. If a vocabulary is specified at the port, $h$ is its reached profile. Its definition requires only the value at the reached execution. Action-conditional predictions also give coordinates for predictive representations of state \cite{ground-psr}.

\begin{proposition}[Kan extensions along the action interface]\label{prop:ground-action-kan}
Let \(\mathbb BM\) be the internal one-object category of the word monoid, with composition \(v\circ u=vu\). Write \(k:1\to(\mathbb BM)^{\rm op}\) for the functor selecting its object. Restriction of a right action along \(k\) is \(U:\mathsf{Act}_M(\mathcal C)\to\mathcal C\). Its pointwise left and right Kan extensions are the free and cofree actions
\[
 F=\operatorname{Lan}_k \ \dashv\ U=k^*\ \dashv\
 R=\operatorname{Ran}_k,\qquad F(V)=V\times M,\qquad R(O)=O^M .
\]
These are indexed constructions in the slice \(\mathcal C\). Both commute with substitution of parameters.
\end{proposition}
\begin{proof}
Internally a right action \(Z\) is a discrete opfibration over \((\mathbb BM)^{\rm op}\). Its arrow above \(u\) with source \(z\) has target \(z\cdot u\), and the unique lifts compose by the action law. Conversely unique lifts give this action on the object of objects. Pullback to \(1\) gives the discrete family \(UZ\). On the free action put \((v,u)\cdot w=(v,uw)\), and on the cofree action put \((s\cdot u)(w)=s(uw)\). Associativity and the empty word give both action laws. For any action \(Z\) the two transpositions are
\[
\begin{aligned}
 b:V\to UZ &\longmapsto b^\sharp:F(V)\to Z,
     & b^\sharp(v,w)&=b(v)\cdot w ,\\
 o:UZ\to O &\longmapsto o^\flat:Z\to R(O),
     & o^\flat(z)(w)&=o(z\cdot w).
\end{aligned}
\]
Equivariance follows by associating the same product \(uw\). The inverse in the first line restricts to \((v,\epsilon)\), and the inverse in the second evaluates at \(\epsilon\). Indeed an equivariant \(f\) satisfies \(f(v,w)=f(v,\epsilon)\cdot w\), and an equivariant \(g\) satisfies \(g(z)(w)=g(z\cdot w)(\epsilon)\). The formulas commute with precomposition and postcomposition, giving the two adjunctions. The comma categories for the extensions at the sole object are discrete families indexed by the arrows \(w:M\), since the source of \(k\) has only identities. Their indexed coproduct of \(V\) is \(V\times M\), and their indexed product of \(O\) is \(O^M\). Composition gives exactly the displayed actions. This proves pointwiseness for internal diagrams, with the words ranging over the object \(M\) in the slice. Pullback preserves these products and slice exponentials. The length and finite tuple description of words in \cref{cor:ground-output} identifies the reindexed word object. The formulas for both transpositions then reindex verbatim.
\end{proof}

\begin{lemma}[Environmental response to finite plays]\label{lem:ground-finite-play}
In the parameter fibre the final $O\times(-)^A$-coalgebra is $O^M$. Its readout is evaluation at $\epsilon$, its action is residual, and the behaviour map of $X$ is $b=jq$. The image $Q$ is closed under residuals. The kernel pair of $q$ is the largest action-stable equivalence contained in $\ker o$.
\end{lemma}
\begin{proof}
The cofree transpose \(o^\flat\) in \cref{prop:ground-action-kan} gives the displayed response map. A coalgebra morphism $b:X\to O^M$ must satisfy
\[
 b(x)(\epsilon)=o(x),\qquad b(x)(aw)=b(xa)(w).
\]
List recursion gives $b(x)(w)=o(xw)$, and induction on $w$ gives uniqueness in the fibre. In particular $b(xv)=b(x)\cdot v$, so the regular image inherits the action and $q$ is equivariant. Equality of behaviours implies equality at the empty word and is stable under every action. Conversely, if a related pair has equal readouts and all its successors remain related, word induction gives equality of all its observations. This proves the kernel assertion and identifies the response image with the behavioural quotient (cf.\ \cite{rutten-systems}).
\end{proof}

\begin{proposition}[The common quotient of two reached observations]\label{thm:ground-common}
Write $p_1,p_2:R_h=X\times_H X\rightrightarrows X$ for the kernel pair of $h$. Let $e_0$ be the coequalizer of $qp_1,qp_2$. It makes the displayed square a pushout. The subobjects descending along its common composite $g_0$ are precisely those descending along both legs.
\[
\begin{tikzcd}[column sep=large,row sep=small]
 X\arrow[r,two heads,"q"]\arrow[d,two heads,"h"']&
 Q\arrow[d,two heads,"e_0"]\\
 H\arrow[r,two heads,"d_0"']&G_0 .
\end{tikzcd}
\qquad
\begin{gathered}
 g_0=d_0h=e_0q,\\
 \mathsf{Sat}_{g_0}(X)=
 \mathsf{Sat}_h(X)\cap\mathsf{Sat}_q(X).
\end{gathered}
\]
Here $\mathsf{Sat}_e(X)$ denotes the subobjects of $X$ pulled back along $e$, as in \cref{prop:profile-adjoints}. Complete recovery of the environmental behaviour at the port is equivalent to $qp_1=qp_2$. In that case the arrow $H\to Q$ recovering it is unique.
\end{proposition}
\begin{proof}
The composite $e_0q$ equalizes $p_1,p_2$, so it descends uniquely along $h$ to $d_0$. Its image is all of $G_0$, whence $d_0$ is regular epic. Given $u:H\to Y$ and $v:Q\to Y$ with $uh=vq$, the arrow $v$ equalizes $qp_1,qp_2$. Its unique factor through $e_0$ agrees with $u$ after $d_0$, by cancellation of $h$. This proves the pushout property. The kernel of $g_0$ contains $\ker h$ and $\ker q$. Conversely, the quotient by any equivalence containing both kernels gives a common cocone and factors through $g_0$. Thus $\ker g_0=\ker h\vee\ker q$. Taking $Y$ to be the subobject classifier gives the predicate identity. The choice $Y=2$ gives the decidable predicates. Applying effective descent along $h$ directly to $q$ gives the recovery criterion.
\end{proof}

\subsection{Descent under continuation}
Extend each \(qp_i\) freely using \cref{prop:ground-action-kan}. The resulting equivariant arrows impose the observed port identifications after every continuation, using only the transition on the environmental response.

\begin{theorem}[The common quotient under word actions]\label{thm:ground-stable}
Let $e$ be the coequalizer in $\mathsf{Act}_M(\mathcal C)$ displayed below.
\[
 F(R_h)\underset{(qp_2)^\sharp}{\overset{(qp_1)^\sharp}{\rightrightarrows}}
 Q\xrightarrow{\ e\ }G_{\mathrm{act}}.
\]
There is a unique $d:H\to UG_{\mathrm{act}}$ with $dh=U(e)q$. For every $M$-action $K$, composition with $d$ and $e$ gives a natural bijection
\[
 \operatorname{Hom}_{M}(G_{\mathrm{act}},K)
 \cong
 \{(u,v)\in\operatorname{Hom}_{\mathcal C}(H,UK)\times
              \operatorname{Hom}_{M}(Q,K)\mid uh=U(v)q\}.
\]
Write $g=dh$. For a subobject $P\hookrightarrow X$ the predicate comparison is
\begin{equation}\label{eq:ground-action-predicates}
 P\in\mathsf{Sat}_g(X)
 \quad\Longleftrightarrow\quad
 \tau^*P\in\mathsf{Sat}_{g_0\times 1_M}(X\times M).
\end{equation}
In sets these are precisely the members of the largest Boolean subalgebra of $\mathsf{Sat}_{g_0}(X)$ closed under inverse image by each action. The displayed comparison also holds for subobjects internally.
\end{theorem}
\begin{proof}
Take the underlying coequalizer of the displayed pair in $\mathcal C$. Product with $M$ preserves this coequalizer, since it has right adjoint $(-)^M$. Equivariance of the generating pair therefore descends the action of $Q$ to its quotient. Cancellation of the quotient, also after product with $M$ or $M\times M$, proves the unit and multiplication laws and equivariance of the induced factors. Thus $U$ creates this coequalizer. Restriction along the empty word gives $U(e)qp_1=U(e)qp_2$ and hence $d$. An equivariant $v:Q\to K$ coequalizes the two free extensions exactly when $U(v)qp_1=U(v)qp_2$, by the free-action bijection. Descent along $h$ identifies this condition with the existence of the unique $u$ in the statement. These factorizations commute with postcomposition, giving naturality. The pushout also gives the comparison $G_0\twoheadrightarrow UG_{\mathrm{act}}$. In particular $\ker g$ is the least action-stable equivalence containing both kernels. Quotienting $X$ by any such equivalence gives a common cocone with an induced action, to which the bijection applies.

For the predicate calculation let $\Omega$ denote the classifier in the fibre and let $\chi_P:X\to\Omega$ classify $P$. The transpose
\[
 b_P:X\longrightarrow\Omega^M,\qquad b_P(x)(w)=\chi_P(xw)
\]
is equivariant for the residual action. The right side of \eqref{eq:ground-action-predicates} says exactly that $b_P$ factors through $g_0$, by transposition. Such a factor gives compatible arrows from $H$ and $Q$ to $\Omega^M$. The arrow from $Q$ is equivariant since $b_P$ and $q$ are equivariant and $q\times1_M$ is epic. The hom-set bijection now factors $b_P$ through $g$. Evaluation at the empty word descends $\chi_P$. Conversely, given $\chi_P=\chi g$, the arrow $(s,w)\mapsto\chi(s\cdot w)$ on $G_{\mathrm{act}}\times M$ pulls back to $\chi_P\tau$. Composing with the comparison from $G_0$ proves the reverse implication.

In sets the right side says that $x\mapsto\chi_P(xw)$ factors through $g_0$ for every word $w$. Closure under one-action inverse images implies this condition by word induction. The predicates pulled back from $G_{\mathrm{act}}$ themselves have this closure by equivariance of $g$. This proves maximality. Internally the product with $M$ in \eqref{eq:ground-action-predicates} tests words in every parameter context. Replacing $\Omega$ by $2$ gives the corresponding comparison for decidable predicates.
\end{proof}

\begin{corollary}[Parameter substitution and the environmental readout]\label{cor:ground-output}
Both quotient constructions commute with pullback in the parameter $S$, up to canonical isomorphism. The readout $s\mapsto j(s)(\epsilon)$ on $Q$ descends along $e$ exactly when $q$ descends along $h$. In that case $Q\cong G_0\cong UG_{\mathrm{act}}$. If the kernel pair of $h$ is stable under the word action, then $G_0\cong UG_{\mathrm{act}}$ even for a partial common description.
\end{corollary}
\begin{proof}
For $r:T\to S$, pullback preserves kernel pairs, images and slice exponentials. It preserves the coequalizers since $r^*\dashv\Pi_r$. A finite word consists of a length $n$ and a tuple of letters indexed by $\operatorname{Fin}(n)$. The length object is the pullback of the natural numbers object. The dependent sums and products in this description reindex by the comparisons of \cref{sec:indexed-readback}. Hence $r^*(V\times M)$ is the reindexed free action, and both coequalizer diagrams transport. The isomorphisms respect identities and compose by uniqueness of the induced factors. Bijections of action, sensor and port names likewise conjugate the diagrams, and the corresponding experiment formulas transport by renaming their letters and tokens.

If the readout on $Q$ descends to $G_{\mathrm{act}}$, equivariance gives its value after every word there as well. The resulting behaviour map $G_{\mathrm{act}}\to O^M$ composes with $e$ to give the mono $j$. Thus $e$ is both monic and regular epic, hence invertible. The map from $H$ then recovers $q$. Conversely, if $q$ descends along $h$, the generators $qp_1,qp_2$ agree, as do their free extensions. Both coequalizers are isomorphisms out of $Q$. Finally, when $R_h$ is stable, it carries the componentwise action and $qp_i$ are equivariant. Their underlying coequalizer $e_0$ inherits the action by the product argument above and has the same universal property as $e$.
\end{proof}

\begin{proposition}[Change of the observed port]\label{prop:ground-port-change}
Fix the reached interaction \(X\), its actions, and its response quotient \(q:X\twoheadrightarrow Q\). Order the reached ports \(h:X\twoheadrightarrow H\) by arrows \(h\to h'\) given by factorizations \(h'=u h\). Each of the two grounding quotients is functorial in this preorder. The induced maps are regular epimorphisms, equivariant for the quotient under word actions. Pullback along them includes the predicates grounded at \(h'\) among those grounded at \(h\), with their same interpretations on \(X\).
\end{proposition}
\begin{proof}
Since \(h\) is epic, \(u\) is unique. Factoring \(h'\) through \(h\) gives \(R_h\hookrightarrow R_{h'}\) over \(X\times X\). Hence the quotient of \(Q\) for \(h'\) coequalizes every generating pair for \(h\), also after free extension in the action case. Write \(e_h,e_{h'}\) for the quotient maps of either kind. The unique factor \(v\) with \(v e_h=e_{h'}\) is regular epic since its image contains the image of \(e_{h'}\). In the action case the factor is equivariant by the coequalizer there. Its underlying regular epimorphism is also regular epic in actions, since the kernel pair carries the componentwise action and its coequalizer is created by \(U\). Cancellation of \(e_h\) proves the identity and composition laws. Finally, inverse image along \(v\) reflects inclusion by regular epimorphy, and \(v e_h q=e_{h'}q\) identifies the interpretations on histories. In particular ports with the same kernel pair have canonically isomorphic grounding quotients, even when their reached values have different names.
\end{proof}

\paragraph{Comparison with a port cut.}
Either quotient supplies a code $d':H\twoheadrightarrow G$ on reached port values. For a specified downstream continuation $C:B\times H\to Y$, \cref{prop:response-quotient} gives the additional test
\[
 C=\bar C(1_B\times d')\quad\Longleftrightarrow\quad
                 \ker d'\ \subseteq\ \ker\widehat C,
 \qquad \widehat C:H\to Y^B .
\]
This tests sufficiency under the replacements in $H$. To exhibit an influence on the output, two replacements in a single bypass context must give different responses. A larger replacement domain requires a code on that domain.

\subsection{Finite readback from action and sensor observations}
For a word $w$ and a sensor value $t$, let $[w:t]$ denote the pullback of $t$ along $x\mapsto o(xw)$. On a transition table this atom follows the foreign keys for $w$ and tests the resulting sensor field. Inverse image under action $a$ takes $[w:t]$ to $[aw:t]$. At the port, equality to a reached tuple of rational register values is a conjunction of coordinate equalities. Equality to a Boolean profile is its complete literal test.

\begin{theorem}[Paired formulas on a finite presentation]\label{thm:ground-finite}
On a finite deterministic joint transition table with a reset row, a finite action alphabet with decidable equality and decidable equality on the observed sensor and port values, both quotients have computable presentations. Each quotient class has a Boolean formula in the experiment atoms and one in the port tests, with equal inverse images on reachable rows. Disjoining these formulas presents every subset of classes in classical sets. Internally, every quotient predicate has a finite expansion in the same class tests with its membership truth values as coefficients. Those admitting ordinary Boolean formulas in these tests are precisely the decidable quotient predicates. Failure of complete recovery returns two access words with the same port value and a shortest separating continuation for that pair.
\end{theorem}
\begin{proof}\leavevmode
\begin{enumerate}[titledsteps]
\item \textbf{Compute the response image.} Enumerate reachable rows from reset and store an access word for each. Write $T_a(x)=xa$ and let $D_n$ be the kernel of the response on words of length at most $n$. This is the state-distinguishability refinement for a deterministic Moore machine \cite{moore}. Splitting a nonempty word at its first letter gives
\[
 D_0=\ker o,\qquad D_{n+1}=D_0\cap\bigcap_{a\in A}(T_a\times T_a)^*D_n.
\]
Each strict refinement increases the number of blocks, and equality at any step is permanent. Thus for $N\ge1$ reachable rows, $D_{N-1}=\ker q$. Its classes present $Q$ and their successors give its residuals. Two rows first separated at depth zero differ on $\epsilon$. For a first split at positive depth choose an action whose successors split at the preceding depth, and prepend it to their separator. Induction gives a shortest word.
\item \textbf{Compute the two coequalizers.} For each $z\in H$ list the $Q$-states paired with $z$. Identifying the entries of each list and taking equivalence closure computes $e_0$. For $e$, also merge the two successors under $a$ of any related pair, for every $a$, until stable. Each effective merge reduces the number of classes, and every congruence containing the generating pairs contains each merge. These are therefore the coequalizers specified above. A list with two distinct entries gives the access words and the separating experiment. Singleton lists give the unique recovery map $H\to Q$. Equivalently $G_0$ is the set of connected components of the bipartite incidence on $H\amalg Q$. This is the common-function construction on a finite support \cite{ground-common-information}.
\item \textbf{Read both formulas.} Collect one separating word for each unordered pair of distinct elements of $Q$ in a finite set $W$. Then $\chi_s=\bigwedge_{w\in W}[w:j(s)(w)]$ selects exactly $q^{-1}(s)$. For a quotient class $c$, let $Q_c,H_c$ denote its two fibres. The formulas
\[
             \Phi_c=\bigvee_{s\in Q_c}\chi_s
          \qquad\text{and}\qquad
             \Psi_c(u)=\bigvee_{z\in H_c}(u=z)
\]
give the same predicate on $X$ when the port variable $u$ is filled by $h$. Empty meets and joins give true and false. For a predicate $D$ on either quotient $G$, let $d:H\to G$ and $e:Q\to G$ denote its two quotient maps. In the internal language,
\[
 \begin{aligned}
 D(d(z))&\ \longleftrightarrow\ \bigvee_{c\in G}\bigl(D(c)\wedge\Psi_c(z)\bigr),\\
 D(e(q(x)))&\ \longleftrightarrow\ \bigvee_{c\in G}\bigl(D(c)\wedge\Phi_c(x)\bigr).
 \end{aligned}
\]
Indeed the class tests partition the reached carriers, and $dh=eq$. For bit-valued membership, scan the classes and retain precisely the terms with bit one. Conversely a Boolean formula in the decidable observations defines a bit on the reached rows. If it represents a quotient predicate, this bit is constant on the kernel pair and descends to $G$ since the quotient is regular epic. The distinction already occurs on a one-class table over the parameter $i\in\N$: taking $D(*)$ to be diagonal halting gives an internal predicate, and a uniform Boolean class formula for it would decide halting. At a fixed finite unrolling, \cref{thm:relu-formula} substitutes the graph of a rational affine--ReLU port into its test. The same register comparisons implement the formula at the port. The positive-conjunction fragment uses the separate extent criterion of \cref{thm:concept-expressibility}.
\end{enumerate}
\end{proof}

\begin{proposition}[Reconstruction from reset experiments]\label{thm:ground-reset}
For a deterministic reset presentation with the finite observables above and a supplied bound $N\ge1$ on the number of reachable joint states, observations after words of length at most $2N-1$ suffice for both quotient presentations and their paired formulas. For the environmental predictor alone an environmental state bound suffices.
\end{proposition}
\begin{proof}
Put $L=A^{\le N-1}$. Every reached joint row has an access word in $L$, since deleting a loop from an access path preserves its endpoint. For $p\in L$ read the port and form the signature $\sigma(p)=(o(pw))_{w\in L}$ together with $\sigma(pa)$ for each action. Each entry is obtained by replay from reset. The longest experiment has length $(N-1)+1+(N-1)$. By \cref{thm:ground-finite}, signature equality is precisely equality in $Q$. Every successor signature occurs at some short access word, so $\sigma(p)\cdot a=\sigma(pa)$ presents its action. The pairs $(h(p),\sigma(p))$ range over the entire reached incidence. They present the two coequalizers by the same computation. Omitting the port and applying the argument to the environmental machine gives the last assertion.
\end{proof}

The recovered transition table $\delta:Q\times A\to Q$ and sensor map $r:Q\to O$ have an exact keyed recurrent realization. On one-hot state and action fields $y,u$ form
\[
 v_{ia}=[y_i+u_a-1]_+,\qquad
 y'_j=\sum_{\delta(i,a)=j}v_{ia},\qquad
 s_t(y)=\sum_{r(i)=t}y_i .
\]
The unique active incidence gives its successor and the active state gives its sensor. Induction from $\sigma(\epsilon)$ identifies every finite rollout. To test a separately specified autonomous predictor, form the synchronous product of its update with the environmental update. From the specified initial pairs enumerate the reached product. When this presentation is finite, equality of the two readouts there is decidable and is equivalent by word induction to agreement of all rollouts. Breadth-first search returns a distinguishing continuation on failure.

\subsection{A finite sensor and the limits of probing}
\begin{example}[A common predicate lost under one action]\label{ex:ground-pulse}
A device accepts actions $a,f,k$ and emits tokens $2,7,9$. Its environmental state is $(b,i)\in 2\times\mathbb Z/3$, initially $(0,0)$. The action $a$ advances $i$, $f$ flips $b$, and $k$ fixes both coordinates. At phase zero the sensor returns $7$ for $b=0$ and $9$ for $b=1$. At the other two phases it returns $2$. All six states are reachable. Distinct phases are separated by advancing until one first reaches zero, and distinct bits at the same phase are separated on reaching zero. Thus $Q$ has six states.

With environmental bound six, \cref{thm:ground-reset} recovers the following table using only sensor tokens. Taking the first access word in length order, with $a,f,k$ in that order, names each distinct signature.
\begin{center}
\begin{tabular}{@{}lcccc@{}}
\toprule
Access word & Sensor & After $a$ & After $f$ & After $k$\\
\midrule
$\epsilon$ & 7 & $a$ & $f$ & $\epsilon$\\
$a$        & 2 & $aa$ & $af$ & $a$\\
$f$        & 9 & $af$ & $\epsilon$ & $f$\\
$aa$       & 2 & $\epsilon$ & $aaf$ & $aa$\\
$af$       & 2 & $aaf$ & $a$ & $af$\\
$aaf$      & 2 & $f$ & $aa$ & $aaf$\\
\bottomrule
\end{tabular}
\end{center}
Compile this table into the recurrent realization above. Adjoin a neural memory bit $j$, initially zero, with update $j'=j+u_k-2[j+u_k-1]_+$. Expose the field tuple
\[
 z_1=y_f+y_{af}+y_{aaf},\qquad z_2=y_\epsilon+y_f,\qquad z_3=j.
\]
Thus $h(b,i,j)=(b,\mathbf1_{i=0},j)$. Action $k$ toggles $j$ independently of the environment. There are twelve reached joint rows and eight reached port values.

For each $b$, the pair $qp_1,qp_2$ identifies $(b,1)$ with $(b,2)$ in $Q$. These are the only distinct states identified at a fixed port value, so $G_0$ has the four classes $(b,\mathbf1_{i=0})$. Applying $a$ to the same pair also identifies $(b,2)$ with $(b,0)$. Equality of $b$ is preserved by all three actions, so $G_{\mathrm{act}}$ has exactly the two classes indexed by $b$. On this quotient $a,k$ act by the identity and $f$ interchanges the two classes.

The two generators of the static predicate algebra have the paired descriptions
\[
\begin{aligned}
 z_1=1 &\quad\longleftrightarrow\quad
       [\epsilon:9]\lor[a:9]\lor[aa:9],\\
 z_2=1 &\quad\longleftrightarrow\quad [\epsilon:7]\lor[\epsilon:9].
\end{aligned}
\]
The first formula tests one full phase cycle while preserving $b$, so it generates the predicate algebra of $G_{\mathrm{act}}$. The second tests phase zero. Its inverse image under $a$ tests phase two, which the port identifies with phase one. Indeed, the histories $a$ and $aa$ both give port $(0,0,0)$ and sensor $2$. One more $a$ yields $2$ and $7$, respectively. This is a shortest separating continuation. Varying $j$ preserves the complete environmental response, so its value supplies no additional common predicate.

The finite check compares the compiled predictor with the independent simulator on all twelve reached product states and all $36$ transitions. Enumeration of the $64$ predicates on $Q$ gives $16$ with a port interpretation. Closing each predicate under action inverse images and testing the same condition leaves four. These are the predicate counts of the four-class and two-class quotients.
\end{example}

\begin{example}[The bound in a reset presentation]\label{ex:ground-delayed}
Consider a constant port and one action $a$. On the finite machine with constant sensor $0$, a terminating probe procedure inspects finitely many powers of $a$. Choose $m$ greater than every inspected length. A second finite machine emits $0$ at lengths below $m$ and $1$ from length $m$ onwards. All answers, including those to adaptively chosen probes, agree along the first run. The first response image is a singleton. In the second, $\epsilon$ and $a^m$ have distinct present observations at the same port. Consequently no procedure terminating on all finite reset machines with unspecified size can decide complete recovery from probes alone.
\end{example}

\begin{example}[A common subobject beyond the finite experiment language]\label{ex:ground-infinite}
Take the two-action history object and expose the full word at the port. The sensor is $1$ exactly on words $a^nba^n$ for $n\ge0$. After $bb$ the response is identically zero. After $a^nb$ its unique continuation with value one is $a^n$. The singleton of the zero response in $Q$ therefore defines a common subobject. Any finite Boolean formula of experiment atoms mentions finitely many continuations. Choose $n$ with $a^n$ outside this list. The two histories then agree on every atom in the formula. Hence the zero-response subobject has no formula in this language. For finite tables the separating family in \cref{thm:ground-finite} gives class tests, and the class membership bits select the tests for a decidable predicate.
\end{example}

\section{Closure operators and their Horn presentations}\label{sec:early-horn}
The paired formulas in \cref{thm:ground-finite} give predicates on the same reached rows. Choose a finite collection as attributes. Their incidence on those rows determines which conjunctions imply which further attributes. \Cref{cor:concept-horn} compiles this implication closure. We begin by comparing logical and numerical evaluation on the same unquotiented gate syntax. Let \(\mathcal G\) be the free cartesian category on a sort \(B\), constants \(0,1\) and binary gate symbols \(\wedge,\vee\). Arrows between finite contexts are tuples of terms, composed by substitution. Only the cartesian wiring laws enter their syntax.
\begin{theorem}[Boolean restriction as a natural comparison]\label{neural:thm:boolean-comparison}
There are product-preserving interpretations $\mathcal B:\mathcal G\to\E$ and $\mathcal P:\mathcal G\to\mathsf{Poly}_{\R}$, and a monic natural comparison $\iota:\mathcal B\Rightarrow U\mathcal P$, with $U$ the underlying-carrier functor in the numerical realization, defined by
\[
 \mathcal B(B)=2,\qquad \mathcal P(B)=\R,\qquad
 \mathcal P(\wedge)(x,y)=xy,\quad
 \mathcal P(\vee)(x,y)=x+y-xy.
\]
For each circuit $g:B^I\to B^J$ the square
\[
\begin{tikzcd}
2^I \arrow[r,"\mathcal B(g)"] \arrow[d,hook,"\iota_I"'] & 2^J \arrow[d,hook,"\iota_J"]\\
\R^I \arrow[r,"\mathcal P(g)"'] & \R^J
\end{tikzcd}
\]
commutes. Every polynomial circuit so constructed has an exact finite realization using affine wiring and the square activation $s(x)=x^2$.
\end{theorem}
\begin{proof}
Interpret constants and cartesian wiring by their canonical maps. Both gate polynomials restrict to the Boolean truth tables and preserve the embedded Boolean cube. The free cartesian universal property extends the interpretations uniquely. Induction under composition and pairing proves naturality. To realize multiplication, use
\[
 xy=\frac{(x+y)^2-(x-y)^2}{4}.
\]
Every occurrence of multiplication is therefore a finite affine--square--affine circuit. Addition, subtraction, copying, and rational constants are affine wiring. The construction is effective by finite syntax recursion and commutes with relabeling of the presented ports.
\end{proof}
Let $\NN_\square$ be the cartesian circuit category generated by rational affine maps and square gates. Replacing each gate of $\mathcal G$ by its circuit defines a compiler $C:\mathcal G\to\NN_\square$. Agreement on generators extends to products and composites, so $\ev C=\mathcal P$. On Boolean inputs the two routes in the restriction square agree as well. For example, $x\wedge x$ and $x$ define the same Boolean function, but their specified polynomial extensions $x^2$ and $x$ have different derivatives.

In a finite positive Horn database, \(T_P(E)\) denotes the set of heads of rules whose bodies lie in \(E\). Write \(K_P(E)\) for the least rule-closed set containing \(E\), obtained by inflationary iteration on the finite atom set. Let \(\enc(E)\) be the characteristic vector of the input set \(E\).
\begin{theorem}[Exact Horn network and certificates]\label{neural:thm:horn-network}
For a finite ground positive Horn program $P$ on a presented atom set $A$, construct a polynomial circuit $H_P:\R^A\to\R^A$ by
\begin{equation}\label{neural:eq:horn-network}
 (H_P(x))_a=1-(1-x_a)
   \prod_{r:\operatorname{head}(r)=a}
      \left(1-\prod_{b\in\operatorname{body}(r)}x_b\right).
\end{equation}
The empty body product is $1$ and the empty rule product is $1$. On encoded finite databases it realizes inflationary forward chaining. Iteration from $\enc(E)$ reaches $\enc(K_P(E))$ after at most $|A|$ rounds. The compiled iteration has a tracked extension that retains a finite Horn proof for each activated atom.
\end{theorem}
\begin{proof}
For Boolean $x$, a rule body product equals $1$ exactly when every premise holds. The outer expression is the disjunction of the old head activity and the body activities. Hence
\[
H_P\enc=\enc\circ(E\mapsto E\cup T_P(E)).
\]
Put \(E_0=E\), \(E_{n+1}=E_n\cup T_P(E_n)\). Then \(H_P^n\enc(E)=\enc(E_n)\). If \(E_{n+1}\ne E_n\), then \(|E_{n+1}|\ge |E_n|+1\). If \(D\) is closed and contains \(E\), induction gives \(E_n\subseteq D\). Hence \(E_{|A|}=K_P(E)\).

Define certificates simultaneously for all \(a\in E_n\), initializing \(p_0(a)=\mathsf{fact}(a)\). In the finite rule order, put
\[
 p_{n+1}(a)=
 \begin{cases}
 p_n(a),&a\in E_n,\\
 \mathsf{rule}\bigl(r,(p_n(b))_{b\in\mathrm{body}(r)}\bigr),
       &a\in E_{n+1}\setminus E_n,
 \end{cases}
\]
where \(r\) is the first rule with head \(a\) and body contained in \(E_n\). In the second case membership in \(T_P(E_n)\) ensures that the finite search succeeds. Induction gives conclusion \(a\) and proof height at most \(n+1\), including the nullary-rule case.
\end{proof}
\paragraph{Closed theories as a reflective poset.}
Let \(K\) be forward chaining to saturation and write \(\mathsf{Mod}\) for the subsets of the finite vocabulary satisfying the rules. The least-model property gives \(K(T)\subseteq S\Longleftrightarrow T\subseteq S\) for \(S\in\mathsf{Mod}\). Thus \(K\) is left adjoint to the inclusion of models and the induced monad is idempotent. Attribute exploration obtains valid implications from queries and counterexamples \cite{ganter-exploration}. For the complete implication system of a context, the next corollary identifies this reflector with the closure on attribute \emph{intents}. The adjoints in \cref{prop:profile-adjoints} instead act on predicates of \emph{inputs}. Their carriers distinguish closure under logical consequences from saturation under equality of observed profiles.

\begin{corollary}[Compiling the implications of a hidden context]\label{cor:concept-horn}
In classical sets every context on a finite vocabulary $M$ has a finite Horn presentation of $\operatorname{cl}(T)=T^{\downarrow\uparrow}$. Its models are the intersections of families of realized profiles, including the empty intersection $M$. On the finite-table and rational-linear domains of \cref{thm:concept-expressibility} this presentation and its circuit are uniformly computable. More generally, for an internal family of contexts over $S$ there is a largest subobject of $S$ on which the closure of every bit-encoded premise has a uniform Boolean circuit.
\end{corollary}
\begin{proof}\leavevmode
\begin{enumerate}[titledsteps]
    \item \textbf{Compile the valid implications.} Write $K=2^M\times M$ for the finite list of premise--head pairs. In each fibre form the validity table as an arrow $v:S\to\Omega^K$ by
\[
 v_s(T,m)=\forall x:X_s\;(T\subseteq b_s(x)\Rightarrow m\in b_s(x)).
\]
Pull back $2^K\hookrightarrow\Omega^K$ along $v$ and call the resulting subobject $S_H$. There the unique bit table selects the rules $T\Rightarrow m$ by scanning $K$. On input $T$, every selected rule with body contained in $T$ has head in $\operatorname{cl}(T)$ by monotonicity. Every $m\in\operatorname{cl}(T)$ is obtained from the rule with body $T$. Thus a single inflationary round returns the closure. Idempotence follows from the derivation adjunction. 
    \item \textbf{Identify the largest circuit domain.} Conversely an exact Boolean circuit over any test base $Y\to S$ returns the bit $v(T,m)$ by evaluation at $T$ and projection to $m$. Its bits give the unique factor through $S_H$. Pullback pasting gives the base-change law. In classical sets $S_H=S$. On the two computational domains, scan rows or eliminate the rational linear constraints of a counterexample to each implication as in \cref{thm:concept-expressibility}. These tests construct the factorization. Finally $\operatorname{cl}(T)=\bigcap_{x:T\subseteq b(x)} b(x)$, each profile is closed, and every closed set is this intersection for its own premise. This identifies the models.


    \item \textbf{Test the limit of bit-valued compilation.} For an internal truth-valued premise the derivation adjunction still constructs closure on $\Omega^M$. Even a one-row context with one always-false attribute sends a premise of truth value $p$ to $\neg\neg p$: the row satisfies the premise just when $\neg p$, and its attribute is false. The finite bit table of this context gives the identity on $2$. Also a decidable incidence by itself need not give a computable validity table on bit premises. Set $M=\{m\}$, $X=\N$, and let $m$ hold at $n$ precisely when machine $e$ on $e$ has not halted within $n$ steps. Validity of $\varnothing\Rightarrow m$ is nonhalting. A uniform compiler for these descriptions would decide it by evaluating its output at $\varnothing$.\qedhere
\end{enumerate}
\end{proof}

Consider the rule $q\wedge r\Rightarrow p$. Its Boolean certificate records the firing rule and its premises. The forward circuit for the head is $p'=p+(1-p)qr$. Applying the marking rule to this particular polynomial gives its arithmetic variation:


\[
 \dot p'=(1-qr)\dot p+(1-p)r\,\dot q+(1-p)q\,\dot r.
\]
A certificate of entailment carries no information about the chosen polynomial extension. Even for that extension, first variation and finite change differ. For example, the finite change in a join includes the cross term $\dot R\Join\dot S$. In the square-zero calculation its coefficient is $\varepsilon^2=0$. This is why the backward signal alone does not determine the effect of a finite step.

Wilson and Zanasi give a presentation of polynomial circuits with a reverse derivative \cite{wz}. The present compiler preserves the chosen polynomial. The square below compares differentiating that polynomial with first compiling it to an affine--square circuit and then differentiating the circuit.
\[
\begin{tikzcd}[column sep=large,row sep=large]
\mathcal G
 \arrow[r,"\mathcal P"] \arrow[d,"C"']
 & \mathsf{Poly}_{\R}
 \arrow[d,"\mathsf D^*"]\\
\NN_\square
 \arrow[r,"\mathsf D^*\ev"']
 & \lens.
\end{tikzcd}
\]
Its commutativity follows from $\ev C=\mathcal P$ on the gate generators and functoriality under composition. The Boolean restriction in \cref{neural:thm:boolean-comparison} checks the values on encoded evidence. This square also checks derivatives at numerical executions. For the interpretation of backpropagation as a learner, see \cite{neural:fong}.

For instance, $x$ and $x^2$ give the same unary Boolean map. At the encoded point $0$ their classes in the pointed second-jet module are respectively $d$ and $d^2$. The Boolean comparison supplies the value at the encoded points; the polynomial extension determines which loss jets pull back along them.

\subsection{An exact extensional readback for finite Boolean control}
There is a canonical repair when the source identifies two proofs by their entire Boolean truth table. For a finite context of $n$ Boolean inputs set $X=2^n$. Use the joint-state register $\Q^X$ with one coordinate per valuation. For a Boolean map $f:X\to Y$ compile the incidence matrix
\[
 (T_f p)_y=\sum_{x:f(x)=y}p_x.
\]
This is an affine neural layer. Evaluating it at $\delta_x$ gives $\delta_{f(x)}$. Reading the unique unit entry in each rational basis column recovers the truth table. Conversely a matrix with one unit entry and all other entries zero in every column is exactly the compilation of that recovered function. Any output bit is recovered as the disjunction, over its true input valuations, of their complete literal tests; on positive finite Horn inputs the forward-chaining circuit above continues to provide proof certificates. Regrouping the finite sum gives $T_gT_f=T_{gf}$, and identity maps give identity matrices. Parallel independent maps give Kronecker products. Copying a shared state sends $\delta_x$ to $\delta_{(x,x)}$, so on a general input joint law $p$ it produces a distribution supported on the diagonal with entries $p_x$. This explicitly preserves correlation under fan-out. The columns can be computed by running the finite Horn closure once per Boolean valuation. This behavioral compiler is exponential in the number of atoms and applies to the finite Boolean function. Its linear derivative on joint-state coordinates is $T_f$ itself. The occurrence-level polynomial extension on the original scalar ports computes the gradients of the actual gate realization.

\clearpage
\part{Indexed witnesses and free algebras}\label{part:algebras}

\section{Effective families, schema collages and finite evidence}\label{sec:eff}

A join compares keys, whereas a numerical layer multiplies or sums attributes stored at those keys. The keys in a presented finite table have decidable equality. An attribute may be known only through observations that arrive over time. To treat both in one instance, we first specify what a program can observe and then show how a query builds its witnesses from those observations.

\subsection{Realizers, internal families and positive observations}
Kleene's first partial combinatory algebra \(\mathcal K_1\) has carrier \(\mathbb N\) and application \(e\cdot n=\varphi_e(n)\), defined when the \(e\)th partial recursive function halts on \(n\). An \emph{assembly} is a set \(X\) together with a nonempty set \(E_X(x)\) of natural-number realizers for each \(x\). A map \(f:X\to Y\) is tracked by an index \(a\) if \(a\cdot n\) is defined and belongs to \(E_Y(fx)\) whenever \(n\in E_X(x)\). Identity and composition are tracked by identity and program composition. Pairing and projection of numbers track products. A partitioned assembly has a singleton realizer set $E_X(x)$ at each element. The exact completion of the regular category of assemblies gives \(\E\) \cite{van-oosten}. In particular, equivalence relations and quotients can be treated internally even where equality of denotations is undecidable.

The natural numbers object has $E_{\N}(n)=\{n\}$, so finite port sets $[k]=\{0,\ldots,k-1\}$, lists and tables admit numerical encodings. We treat an $I$-indexed family of instances in the slice $\E/I$. For $f:A\to B$, reindexing a family gives $\Delta_f$, with adjoints $\Sigma_f\dashv\Delta_f\dashv\Pi_f$. These operations will construct the witnesses and input positions of a query. To execute a join we compare the finite port IDs, which requires no decision about equality of their attribute values.

Let \(\Omega\) denote the subobject classifier. For computations returning positive evidence we use the subobject
\[
 \Sigma=\{p:\Omega\mid\exists b:2^{\N}\
          (p\ \Leftrightarrow\ \exists n\ b(n)=1)\}.
\]
This is the semidecidable truth object \cite[\S3.2]{bauer}. The definition is internal. On named inputs a witnessing stream is implemented by a total bit-producing program. Conversely, bounded simulation turns a halting test into such a stream. A semantic predicate \(X\to\Omega\) and a positive test \(X\to\Sigma\) have distinct types.

\begin{lemma}[Search operations used by finite and countable queries]\label{lem:eff-search}
Finite conjunction and disjunction preserve semidecidable truth. The existential quantifier over \(\N\) takes a uniformly semidecidable family to a semidecidable result. If \(p\) is semidecidable and \(q\) is semidecidable in the context of a witness of \(p\), then \(\exists u:p\,q(u)\) is semidecidable.
\end{lemma}
\begin{proof}\leavevmode
\begin{enumerate}[titledsteps]
    \item \textbf{Combine and enumerate tests.} For conjunction, run both tests and wait for both successes. For disjunction, return on the first success. For a family of tests $T_n$, simulate the first $k$ tests for $k$ steps in round $k$. This procedure succeeds precisely when some $T_n$ does. For an internal family $p_n\in\Sigma$, the choice instance used here has domain $\N$ and codomain $2^{\N}$. It supplies the witnessing streams locally in the parameter context by \cite[Proposition 9.1 and its following remark]{hyland}. A stream is an effective object element with a program index (equivalently use the decidable-bit subobject of $\N^{\N}$); Hyland distinguishes this case from his proof for an arbitrary codomain, which uses choice in the base metatheory. The resulting truth value is independent of the choice of streams and thus defines the restriction of existential quantification to $\Sigma$.


    \item \textbf{Search under a witness.} For the dependent assertion eliminate a witnessing stream \(b\) for \(p\) into the proposition. Define \(q_n=\bot\) for \(b(n)=0\), and \(q_n=q(u_n)\) for \(b(n)=1\), using the witness \(u_n:p\) from the hit. Then
\[
 \exists n\,q_n
   \ \Longleftrightarrow\ \exists n\ \exists u:p\,
                 (b(n)=1\ \wedge\ q(u))
   \ \Longleftrightarrow\ \exists u:p\,q(u).
\]
The first equivalence uses subterminality of \(p\). The second uses \(p\Leftrightarrow\exists n\,b(n)=1\). Each \(q_n\) has a positive test by decidable bit elimination, so the countable dovetailing just constructed realizes the result.\qedhere
\end{enumerate}
\end{proof}

\paragraph{Evidence for a concept column.}
For a subobject $E_m\hookrightarrow H$, the column at a port $h:X\to H$ is the pullback $h^*E_m$. Its classifying map is $\chi_{E_m}h:X\to\Omega$. If the attribute has a positive test $H\to\Sigma$, composing with $h$ supplies a positive test for the column. Internally, the common-attribute assertion for $A\hookrightarrow X$ is the truth value $\forall x:X\, (A(x)\Rightarrow I_h(x,m))$. It can be formed in $\E$ even when the available observations only enumerate positive witnesses. The decision procedures of \cref{thm:concept-expressibility} use a finite table with decidable entries or elimination on rational linear presentations to settle this universal assertion. In particular, a missing positive witness at a finite stage does not certify an implication over all inputs.

\subsection{Representing lower weights and signed numerical values}
A positive test may eventually confirm that a row contributes weight, while an unfinished search still leaves room for further contributions. Lower bounds represent exactly the numerical information available during this search. Later, division and signed arithmetic will require values with two-sided approximations.
Write \(Q\) for the internal object of nonnegative extended lower reals. Concretely a member is a rounded downward cut of the rationals containing every negative rational. The negative rationals alone represent zero, and the full cut represents infinity. The order is inclusion. A supremum is represented by the union of the cuts together with the negative rationals. For a lower real \(m\), put \(B(m)=\{0\}\cup\{a\in\Q_{>0}:a<m\}\). Define
\[
 m+n=\sup_{a\in B(m),\,b\in B(n)}(a+b),\qquad
 mn=\sup_{a\in B(m),\,b\in B(n)} ab
\]
using rational embeddings. These operations include \(0\cdot\infty=0\). Each \(B(m)\) is directed by maximum and has supremum \(m\). A finite set of lower approximants to the same coordinate has a common approximant, so iterated suprema in an associative or distributive law can be taken over the same tuples of rational approximants on both sides. The rational laws then give the nonnegative semiring laws and preservation of directed suprema by finite sums and products. For \(p\in\Omega\) define \(\tau(p)=\sup_{u:p}1\). Its cut at \(0\le r<1\) says precisely \(p\). Thus indicators of truth values land in the lower reals even when \(p\) is undecidable.

For a computation on names we use $Q_\Sigma\subseteq Q$. Membership means that the strict rational bound $r<m$ has a test uniform in $r$. A name can instead enumerate the lower bounds of $m$, with silent stages. Dovetailing all rational tests produces such an enumeration. Waiting for a given rational to appear recovers its test.

The formulas above enumerate finite sums and products by ranging over tuples of approximants. To name an indicator, enumerate the bounds on zero from the start and the additional bounds on one once the test succeeds. Interleaving the enumerations computes the supremum of a uniformly named countable family. The semantic construction in $Q$ also takes suprema of arbitrary internal families. Only the presented countable construction supplies this enumeration.

Signed activities take values in $\Q$ or in the internal real line $\R$. In $\E$ the Cauchy and Dedekind real lines agree by the number-to-number choice instance \cite[\S9]{hyland}. Hyland's Proposition 9.6 presents the real line by indices of recursive Cauchy names with a prescribed rate \cite[\S9]{hyland}. A rate gives rational enclosures of width at most $2^{-n}$. Conversely the midpoints of such enclosures approximate the real with a known error bound. Interval arithmetic and refinement therefore compute finite sums and products.

These enclosures also explain when division is effective. For a positive located real, refine until the lower endpoint of an enclosure is positive. Interval division then computes the reciprocal. A lower name alone supplies only one side of such an enclosure, a distinction we use to identify the normalization domain in \cref{prop:normalization-unit-domain,prop:attention-normalization-boundary}.

\subsection{Foreign keys and typed observations in the schema collage}
A multisorted algebraic theory \(\mathcal T\) is a finite-product category of contexts of sorts. A map \(\Gamma\to t\) represents a term with free variables in \(\Gamma\), modulo the equations of the theory. A \(\mathcal T\)-algebra interprets contexts by products and terms by maps. A presentation constructs terms as finite well-sorted expressions and quotients by the least substitution-compatible congruence containing the stated equations. A map from the resulting algebra amounts to an interpretation of the generators satisfying those equations: extend by recursion on expressions, descend by the congruence, and use induction to prove uniqueness. This formulation retains the operations by which numerical attributes can subsequently be combined.

Fix $\mathcal T$. An algebraic database schema is $(C,O)$ with entity category $C$ and algebraic profunctor
\begin{equation}\label{eq:database-schema}
 O:C^{\mathrm{op}}\times\mathcal T\longrightarrow\E,
\end{equation}
whose transpose $O(c,-)$ preserves finite products. An element $u\in O(c,t)$ is a type-$t$ observable at the entity $c$. The collage $\widetilde C$ adjoins these observables as arrows from the entity category $C$ to the type category $\mathcal T$.
In this schema of Schultz, Spivak, Vasilakopoulou and Wisnesky \cite{algebraic-databases,algebraic-databases-errata}, the two kinds of mixed composition have concrete interpretations. Precomposing an observable with an entity arrow reads it through a foreign key. Postcomposing with a type arrow applies a term to its value. The collage has no generating arrows from types back to entities.

\begin{theorem}[Typed-observation presentation]\label{thm:schema-presentation}
Finite entity arrows, type operations, and observable generators and equations determine an internal algebraic schema with the usual generator-and-relation universal property, retaining the finite-product operations on observables.
\end{theorem}
\begin{proof}
Form \(C\) by quotienting entity paths and \(\mathcal T\) by quotienting typed terms by their respective equations. If \(G(d,t)\) is the family of observable generators, define
\[
 G_c(t)=\coprod_d C(c,d)\times G(d,t),\qquad
 O(c,-)=F_{\mathcal T}(G_c)/{\equiv_c}.
\]
Here \(F_{\mathcal T}\) forms the free type algebra, and \(\equiv_c\) is the algebra congruence generated by \((h^*a,h^*b)\) for every defining observable equation \(a=b\) at \(d\) and \(h:c\to d\). For \(g:c_0\to c\), set \(g^*(f,u)=(fg,u)\). Extension from generators sends \((h^*a,h^*b)\) to \(((hg)^*a,(hg)^*b)\). It descends and gives \(1^*=1,\ (fg)^*=g^*f^*\). Products in each \(O(c,-)\) are the type-algebra products.

For a target functor \(P:C^{\rm op}\to\mathcal T\text{-}\mathsf{Alg}\), the universal bijection at fixed \(C,\mathcal T\) is
\[
 \operatorname{Nat}(O,P)
 \ \cong\ \{(s_u\in P(d,t))_{u\in G(d,t)}:
            \llbracket a\rrbracket_s=\llbracket b\rrbracket_s
                         \text{ for every defining }a=b\}.
\]
Restriction evaluates at \((1_d,u)\). The inverse maps \((f,u)\) to \(P(f)s_u\) and extends by the type operations. The relation calculation makes the extension well defined. Agreement on these generators gives both inverse identities. Varying the entity and type presentations first uses their path and term universal properties. Equality of induced term maps is the quotient on presentation morphisms \cite{algebraic-databases-errata}.
\end{proof}

\begin{definition}[Instance and its comparison maps]
An instance consists of an entity action \(I_e:C\to\E\), a product-preserving type algebra \(I_t:\mathcal T\to\E\), and evaluations \(I_e(c)\times O(c,t)\to I_t(t)\). They satisfy
\[
 \ev_{c,t}(x,O(f,t)u)=\ev_{c',t}(I_e(f)x,u)\quad(f:c\to c')
\]
and commute with each type operation in the observable arguments. A morphism is a pair of natural transformations on the entity and type sides commuting with these evaluations. We can fix a numerical type algebra and take the type-side transformation to be its identity.
\end{definition}
\begin{lemma}[Collage interpretation]\label{lem:collage}
These data are equivalent to a functor from the collage into \(\E\) preserving finite products on the type side. The equivalence respects instance morphisms and holds fibrewise in each parameter slice.
\end{lemma}
\begin{proof}
For \(u:c\to t\) in the mixed hom, let \(J(u)(x)=\ev_{c,t}(x,u)\). The two functor laws to check at mixed arrows are
\[
 J(u f)=J(u)I_e(f),\qquad J(\alpha u)=I_t(\alpha)J(u).
\]
For \(u_i:c\to t_i\), applying these equations to the projections from \(\prod_i t_i\) gives \(J(\langle u_i\rangle)=\langle J(u_i)\rangle\) under the canonical product comparison. Hence the second law includes substitution of every multiary type operation. Conversely these equations extend \(I_e,I_t\) to the collage: each composable string contains at most one mixed arrow. Restriction and extension are inverse up to the chosen product identifications. For a transformation with components \(a_c,b_t\) the remaining square is \(b_tJ(u)=J'(u)a_c\). All equations are stable under reindexing in a parameter slice.
\end{proof}

\paragraph{Legal observable tables.}\phantomsection\label{sec:legal-attribute-tables}
Fix the shape \(A\) and type algebra \(S\). For each generating attribute \(u:c\to t\) choose a column \(b_u:A(c)\to S(t)\); let \(B=\prod_u S(t)^{A(c)}\) be the full table object. Each defining observable equation \(a=b\) with source \(c\) and target \(t\) induces two evaluation maps from \(B\) to \(S(t)^{A(c)}\), by following the specified entity paths and evaluating the scalar terms. Taking their product gives
\[
 \mathcal L\ \lhook\joinrel\longrightarrow\ B\
 \rightrightarrows\ \prod_{(a=b):c\to t} S(t)^{A(c)} .
\]
Here \(\mathcal L\) is the equalizer. An assignment over a test object \(T\) factors through it precisely when its rowwise interpretations satisfy every defining relation. Extension by paths and scalar operations then gives the unique interpretation of all observables by \cref{thm:schema-presentation,lem:collage}. This proves the assertion also for families of assignments.

If the columns split as parameter and data columns \(B=P\times X_0\), the legal input object at \(p\) is the fibre \(\mathcal L_p\) of \(\mathcal L\to P\). The gate composite on \(P\times X_0\) restricts along \(\mathcal L\hookrightarrow P\times X_0\), so its evaluation at \(p\) has domain \(\mathcal L_p\). A relation involving only parameters similarly cuts out a subobject of \(P\). This supplies the legal-domain interpretation of \cref{prop:relative-models}. The independent rational coordinate fields in \cref{thm:support-readback} are the full tables specified in that theorem.

\begin{construction}[A layer schema]\label{con:layer-schema}
Use entity sorts \(J,I,K,P,D,C_0\) for tokens, input and output channels, parameters, edge occurrences, and feature cells. The foreign keys include
\[
 C_0\to J,\quad C_0\to I,\qquad
 D\xrightarrow{s} I,\quad D\xrightarrow{t}K,\quad D\xrightarrow{\rho}P .
\]
A feature is a scalar attribute on \(C_0\). A parameter is a scalar attribute \(\theta:P\to R\), and there is a bias attribute \(b:K\to R\). In the dense instance construct \(C_0=J\times I\) with its two projections. The weight read by an edge is the interpretation of the observable \(\theta\circ\rho\). Distinct edges can read a single parameter, and distinct tokens can carry equal feature values. These are equations between keys or attributes with distinct mathematical consequences.
\end{construction}

A finite numerical instance constructs each key carrier from a finite ordinal $[n]$ and each foreign key from an array of numerical addresses. The interpretation of a declared path is the composite of its arrays. The layer schema of \cref{con:layer-schema} introduces no additional equations between independently stored numerical attributes, so the entity-category laws follow from composition and the type-side laws from the selected scalar algebra. For a general presented schema the quotient universal property identifies exactly which interpretations descend, without supplying a decision procedure for arbitrary equations between real attributes. Executing a join compares numerical row addresses, not arbitrary words in the free entity category; the latter word problem need not be decidable \cite{algebraic-databases-errata}.

To see what naturality says about a parameter lookup, let $u\in O(P,R)$ be the parameter attribute and let $A_R=I_t(R)$ be its value carrier. An edge can read $u$ after following its foreign key $\rho$, or it can evaluate the reindexed observable $u\circ\rho$ directly. The instance equation identifies these readings.
\[
\begin{tikzcd}[column sep=large]
 I_e(D)\arrow[r,"I_e(\rho)"]\arrow[d,"{\ev(-,u\circ\rho)}"']&
 I_e(P)\arrow[d,"{\ev(-,u)}"]\\
 A_R\arrow[r,equal]&A_R .
\end{tikzcd}
\]
On \(d\) both routes give \(\theta_{\rho(d)}\) by the instance equation. The composite is constant on each \(\rho\)-fibre even if \(I_e(\rho)\) is neither injective nor surjective. In particular this square distinguishes a parameter lookup from assigning a new parameter independently at every edge.

\subsection{Finite information and the witnesses of projection}
Let \(\Fin X\) be the image of finite lists into the subobject lattice of \(X\), ordered by inclusion. An ideal is an inhabited directed lower subobject. A relation \(U\) on \(X\) is specified by membership in \(\Omega^X\).
\begin{theorem}[Finite-information representation]\label{thm:finite-relations}
There is a dcpo isomorphism, natural under direct image,
\begin{equation}\label{eq:relations-ideals}
 \Idl(\Fin X)\xrightarrow{\cong}\Omega^X,
 \quad I\longmapsto\bigcup_{F\in I}F,
\end{equation}
with inverse $U\mapsto I_U=\{F\in\Fin X:F\subseteq U\}$. The compact elements of $\Omega^X$ are exactly finite relations.
\end{theorem}
\begin{proof}\leavevmode
\begin{enumerate}[titledsteps]
    \item \textbf{Recover the ideal from its union.} The empty finite relation belongs to $I_U$, and finite union makes it directed. Every $x\in U$ is contained in the singleton $\{x\}\in I_U$, so $\bigcup I_U=U$. Conversely, let $I$ be an ideal and $F\subseteq\bigcup I$, with $F$ represented by a list $x_1,\dots,x_n$. For each list entry, membership in the union supplies an $F_i\in I$ containing it. Finite induction using directedness gives one $G\in I$ containing every $F_i$. Since $F\subseteq G$, lower closure gives $F\in I$. This proves $I_{\bigcup I}=I$ without deciding equality in $X$ or choosing an infinite family of witnesses.


    \item \textbf{Identify the compact relations.} Every relation is the directed union of its finite subrelations. If $U$ is compact, apply compactness to $U=\bigvee_{F\in I_U}F$ to obtain an $F\subseteq U$ with $U\subseteq F$, hence $U=F$. Conversely, if a finite $F$ lies in a directed union $\bigvee_aU_a$, choose a witnessing index for each element of a representing list. Directedness and finite induction yield one $a$ with $F\subseteq U_a$, proving compactness.


    \item \textbf{Transport along direct image.} For $f:X\to Y$, direct image takes a finite list to the list of its images and preserves arbitrary unions. Thus $\exists_f(\bigcup_{F\in I}F)=\bigcup_{F\in I}f(F)$, which is the asserted naturality. Inverse image is not asserted to preserve finite relations.\qedhere
\end{enumerate}
\end{proof}

For row maps $L\xrightarrow{t}Y\xleftarrow{s}D$, the pullback $L\times_Y D$ retains each compatible pair of rows. Its projections agree over $Y$, and pairing gives the unique factorization of any compatible row maps. Projecting the join to a truth value keeps only the fact that some pair exists. Projecting by addition instead indexes the sum by the pairs themselves, so two pairs with equal payloads remain distinct.

\subsection{A uniform evaluator for guarded positive queries}
For each base relation sort, a partial recursive map $K(i)$ on numerical occurrence addresses returns the row's finite tuple of key fields. Its domain is the occurrence object. Once a key tuple is confirmed it is permanent. A second program enumerates nonnegative rational approximants to the annotation. Let $h_i^{(n)}$ be the maximum of zero and the approximants seen in $n$ steps, and put
\[
 a_i=\sup_n h_i^{(n)},\qquad a_i^{(n)}=\max(0,h_i^{(n)}-2^{-n}).
\]
These stage bounds increase to $a_i\in Q_\Sigma$. For $0\le r<a_i$, choose $n_0$ with $r<h_i^{(n_0)}$, then choose $n\ge n_0$ with $2^{-n}<h_i^{(n_0)}-r$. This gives $r<a_i^{(n)}$. Each positive stage bound is strictly below $a_i$, even if the enumerator has already emitted its supremum. The bound zero is available for every confirmed row.

A primitive guard has a partial recursive recognizer $G$ on the candidate's confirmed numerical fields $k$ and a finite certificate code $C$. A certificate names finitely many typed row slots, including the slots of the candidate that it uses. Its facts specify rows, comparisons of confirmed key fields, and strict rational bounds $r<a_i$ on their annotations. Additional slots are bound by existence of matching rows. A numerical address inspected by a guard is part of the declared key data. The slot and enumeration numbers used to implement a witness are private to the presentation. Write $C\preceq (U,b)$ when the slots can be matched in the instance $U$, with distinguished slots fixed by the candidate $b$, so that all the listed facts hold. The guard's truth value is
\begin{equation}\label{eq:guard-cert}
 g_U(b)\quad\Longleftrightarrow\quad
 \exists C\;\bigl(C\preceq(U,b)\ \wedge\ G(k,C)\!\downarrow\bigr).
\end{equation}
On a presentation, search for such a match and a finite halting trace by dovetailing. Confirmation of a strict bound uses $r<a_i^{(n)}$ at some stage. Consequently two streams naming the same annotation validate the same bounds. Extension of the evidence preserves each successful match and trace. In particular, a temporary zero maximum gives no certificate that a final annotation equals zero.

A guarded positive query is built from base scans, tagged union, natural join on numerical keys, restriction by these positive guards, and projection along an output key. A scan returns a base address, a union tags the branch, and a join pairs the two witness addresses. A restriction retains the address when its test succeeds. Projection changes the output key but retains the witness address. Repeated confirmations of one guard do not create new occurrences. Existential evaluation asks whether a witness reaches an output; weighted evaluation adds its annotation, multiplying those of joined rows. The empty query contributes zero and the unit query one nullary witness.

For a fixed query $q$, a finite fragment $D\Subset U$ records finitely many confirmed row IDs and fields, a rational bound for each such row (zero or a strict positive lower bound), and finitely many successful guard records. A record specifies the guard site, candidate witness ID, certificate, row match, and a finite accepting computation. Its field facts must be present, and every strict bound in it must be below the recorded bound for that row. Extension adds rows or records and raises bounds. Finite unions with coordinatewise maximum give common extensions. To evaluate $q$ on $D$, perform its finite joins and projections and pass a restriction when its candidate has a recorded success at that site. The witness ID omits the choice of certificate, match and trace.

For $D_n$, run the first $n$ addresses for $n$ steps and retain the confirmed fields and bounds $a_i^{(n)}$. Enumerate possible guard records by natural numbers. At round $n$ check the first $n$ codes against these fields and bounds, verifying their finite traces. Process the query from its leaves upward so that a record can enable a restriction once its candidate has been built. All the round's checks are finite. A later round still verifies the same trace and facts. Foreign-key arrays and interpreted type operations belong to the underlying schema instance, as in \cref{thm:schema-presentation}.

\begin{theorem}[Finite evidence and uniform recursion]\label{thm:evidence}
For each occurrence presentation $U$, guarded positive query code $q$ and output key $j$, the positive truth and lower-weight interpretations satisfy
\[
 \llbracket q\rrbracket(U)_j
    =\bigvee_{D\Subset U}\llbracket q\rrbracket(D)_j
    =\bigvee_n\llbracket q\rrbracket(D_n)_j .
\]
The weighted finite stages use their rational lower bounds. A single partial recursive interpreter, taking the query code, the underlying row and guard programs, the output key $j$ and a rational $r$, succeeds exactly when $r<\llbracket q\rrbracket(U)_j$. For existential interpretation it succeeds exactly when the output witness fibre is inhabited.
\end{theorem}
\begin{proof}\leavevmode
\begin{enumerate}[titledsteps]
    \item \textbf{Discover each finite witness.} Write \(W_q(D)\) for the witness IDs obtainable in a fragment \(D\). Scanning, tagging and pairing preserve IDs under extension. Restriction retains an existing ID once its record passes. Projection changes only the key. Induction on query syntax therefore gives
\[
 D\le D'\Longrightarrow W_q(D)\subseteq W_q(D'),\qquad
 b\in W_q(U)\Longleftrightarrow \exists n\ b\in W_q(D_n).
\]
For the reverse inclusion in the equality of witness sets, finite execution on a fragment is sound in \(U\). For eventual discovery of a semantic witness, take the finite collection of row confirmations, required strict bounds and guard records in its derivation. Choose a stage past their confirmation times, the stages at which all required bounds are exceeded, and the codes of all records and row addresses. Bottom-up evaluation builds the same ID at this stage. The same maximum construction for the finite data in any \(D\Subset U\) yields \(D\le D_n\). In particular repeated certificates proving one guard never multiply its witness weight.


    \item \textbf{Recover its contribution.} Write $W_{q,j}^{(n)}$ for the stage-$n$ witnesses, and $E_b$ for the finite input uses of witness $b$. The weight computed at that stage is the rational sum
\[
 S_{q,j}^{(n)}=\sum_{b\in W_{q,j}^{(n)}}
                      \prod_{e\in E_b}a_{s(e)}^{(n)}.
\]
Let \(a_b=\prod_{e\in E_b}a_{s(e)}\). For each finite set \(F\) of semantic witnesses, take a common discovery stage \(n_F\). Finite sums and products preserve directed suprema of the nonnegative lower approximants. Consequently
\[
 \sum_{b\in F}a_b
  =\sup_{n\ge n_F}\sum_{b\in F}\prod_{e\in E_b}a_{s(e)}^{(n)}
  \le\sup_n S^{(n)}_{q,j}
  \le\sup_{F'\ {\rm finite}}\sum_{b\in F'}a_b .
\]
The last inequality uses soundness of every stage and its bounds. Taking the supremum over \(F\) gives both inequalities as equalities, including an empty output fibre. The strict rational lower cut is
\[
 r<\llbracket q\rrbracket(U)_j\quad\Longleftrightarrow\quad
 \exists n\;r<S_{q,j}^{(n)}.
\]

    \item \textbf{Construct the uniform interpreter.} A universal interpreter enumerates $n=0,1,\ldots$, constructs these finite sums and halts on the first successful rational comparison. Structural recursion on the finite query code and bounded simulation of the supplied programs make this a single partial recursive procedure. For existential output it halts on the first completed output witness, without testing its weight. Neither procedure decides negative information or equality of denoted real attributes.\qedhere
\end{enumerate}
\end{proof}

For fixed $q$ and $j$, the theorem determines the query on an arbitrary ideal from its values on finite evidence. Let $\mathcal D$ be the poset of consistent finite evidence fragments and take $L=\Omega$ for truth or $L=Q$ for lower weights. These semantic codomains have all directed joins. The explicit stage enumeration above gives the corresponding positive names on presented inputs. The left arrow embeds a fragment as its principal ideal.
\[
\begin{tikzcd}[column sep=large]
 \mathcal D \arrow[r,"{\llbracket q\rrbracket_{\rm fin}}"] \arrow[d,"\downarrow"']&
 L \arrow[d,equal]\\
 \Idl(\mathcal D) \arrow[r,"{\overline{\llbracket q\rrbracket}}"']& L .
\end{tikzcd}
\]
Every ideal is the directed join of its principal ideals, so the finite-stage evaluator has the unique directed-supremum-preserving extension displayed by the lower arrow. Its witness addresses retain the input uses required by the polynomial elaboration below. This construction distinguishes the supplied database content from the derived properties of its execution: numerical addresses, finite-stage monotonicity, witness search and threshold tests are computed from the occurrence and query codes rather than introduced as separate admissibility conditions.

\section{Polynomial functors from query witnesses}\label{sec:witness}

The evaluator retains a witness ID as long as its contribution still has to be counted. On a finite table we can inspect that ID to see how the contribution was formed. A unary use gives a weighted span. A join can use several inputs in one contribution, so we must also record the positions at which those inputs occur. Polynomial diagrams supply those positions and describe what happens to them under query composition.

\subsection{Base change for additive queries}

For finite presented objects, let \(K_0\) be a commutative semiring. For a row map \(f:A\to B\), set \(f^*u(a)=u(fa)\) and
\((f_!v)(b)=\sum_{a:fa=b}v(a)\).
This \(f_!\) is additive aggregation of annotations. On a family of types the corresponding dependent-sum symbol constructs an object of witnesses. We connect the two by an interpretation of finite witnesses, rather than identify their recipients. Semiring-annotated relational algebra provides the finite query semantics \cite{green-provenance}.

An internal finite decidable fibre also has a canonical aggregate without a chosen order. Over its base form the object of pairs $(n,e)$ with $e:[n]\simeq F$ in the fibre. Its projection is regular epic: a finite covering list can be made repetition-free by scanning with decidable equality, and the resulting enumeration is bijective. On the kernel pair, two enumerations of the same fibre differ by their unique change of indices $e'^{-1}e$. Finite commutativity makes their sums and products agree. They descend uniquely along the regular epimorphism. Pulling back the enumeration object and its descended map gives the same construction after a base change. Numerical row addresses give one such enumeration in the concrete tables.

\begin{lemma}[Pullback and aggregation]\label{lem:bc}
Pulling back a row family before summing gives the same result as summing first and then reindexing the result. The right square states this identity for the pullback at left. Also \(f_!(v\cdot f^*u)=(f_!v)\cdot u\).
\[
\begin{tikzcd}
 A' \arrow[r,"g'"]\arrow[d,"f'"']& A\arrow[d,"f"]\\
 B'\arrow[r,"g"']&B
\end{tikzcd}
\qquad
\begin{tikzcd}
 K_0^{A}\arrow[r,"(g')^*"]\arrow[d,"f_!"']&
 K_0^{A'}\arrow[d,"f'_!"]\\
 K_0^{B}\arrow[r,"g^*"']&K_0^{B'}
\end{tikzcd}
\]
\end{lemma}
\begin{proof}
Present \(A'\) as pairs \((b',a)\) with \(g(b')=f(a)\). Its fibre over \(b'\) is isomorphic to the fibre of \(f\) over \(g(b')\) by projection to \(a\), with inverse \(a\mapsto(b',a)\). Reindex the finite sum along these inverse maps. For the second identity, \(u(fa)=u(b)\) on that fibre. Distribute the constant factor over the sum, including the empty sum. This proves both comparisons constructively.
\end{proof}

\begin{proposition}[Additive convolution]\label{prop:conv}
For finite middle indices, the action of a weighted relation \(M:I\times J\to K_0\) on columns is \((L_Mx)_j=\sum_i M(i,j)x_i\). Composing with \(N:J\times L\to K_0\) gives the relation \((N\circ_+ M)(i,k)=\sum_j M(i,j)N(j,k)\). The diagonal with coefficient one is the identity. Over nonnegative rational weights, the support of the composite is the existential composite of the supports.
\end{proposition}
\begin{proof}
Expand the two iterated actions by finite distributivity. Both sum \(x_iM(i,j)N(j,k)\) over \(I\times J\), differing only in the order of enumeration. Multiplication by the diagonal removes every summand except the matching index. A finite nonnegative rational sum is positive precisely when some term is positive, and a product is positive precisely when both factors are. This last argument has no signed analogue: opposite weights may cancel on projection.
\end{proof}
Finite-subsum counting also describes continuous matrices on a decidable-equality index object. Define $\sum_i a_i$ in $Q$ as the directed supremum of its finite subsums. Any finite subset of a product lies in the rectangle formed by its two projections. Rectangles are therefore cofinal, and counting pairs agrees with iterated counting.

For a decidable-equality index object \(I\), write \(e_i\) for its indicator column. For any \(u:I\to Q\) the finite subsets \(F\) form a directed family under union. Pointwise \(u=\sup_F\sum_{i\in F}u(i)e_i\), since at coordinate \(j\) the singleton \(\{j\}\) already attains \(u(j)\). Hence
\[
 (Tu)_j=\sup_F\sum_{i\in F}u(i)(Te_i)_j,\qquad
 T\longleftrightarrow\big((Te_i)_j\big)_{i,j}
\]
for Scott-continuous \(Q\)-linear maps. Conversely the finite-subsum formula defines the map for any such matrix: finite distributivity gives linearity, and finite sets of witnesses for a strict bound have a common stage in a directed family. On enumerated indices with named entries, enumerate finite subsums. Signed numerical execution uses finite free modules.

\subsection{From weighted incidences to linear maps}
For a constructed finite decidable object $I$, write $F_R(I)=R^{(I)}$, where $R$ is $\Q$ or the numerical $\R$. The finite basis identifies this free module with $R^I$. The construction below also applies to the finite additive $Q$-module realization. The later use of subtraction takes place over $R$.

\begin{definition}[Weighted incidence presentation]
A presentation from $I$ to $J$ is a finite span $I\xleftarrow{s}E\xrightarrow{t}J$ with a weight $w:E\to R$, taken up to weight-preserving isomorphism of spans. Its linear interpretation is
\[
 L_E=t_!\,m_w\,s^*:R^I\longrightarrow R^J.
\]
Here $s^*$ is reindexing, $m_w$ multiplies the coordinate at an occurrence $e$ by $w(e)$, and $t_!$ sums over the finite fibre. Distinct edge occurrences remain distinct until their weights are summed.
\end{definition}

\begin{theorem}[Functorial linearization]\label{neural:thm:linearization}
Compose weighted spans by pullback and multiplication of weights. Weighted-span interpretation preserves identities and composition, and sends disjoint union of port objects to direct sums of free modules. After quotienting presentations by equality of their aggregated coefficients, it identifies the resulting category with the category of finite free $R$-modules on these bases and their linear maps.
\end{theorem}
\begin{proof}
For $I\xleftarrow{s}E\xrightarrow{t}J$ and $J\xleftarrow{u}D\xrightarrow{v}K$, form $E\times_JD$. The composite occurrence $(e,d)$ has weight $w(e)z(d)$. Finite Beck--Chevalley for this pullback, together with the projection formula, gives
\[
 v_!m_z u^*t_!m_ws^* = L_{E\times_J D}.
\]
Explicitly, both sides at $k$ sum $z(d)w(e)a(s(e))$ over pairs with $t(e)=u(d)$ and $v(d)=k$. This is the finite Fubini law of \cref{prop:conv}. The diagonal span with weight $1$ gives the identity. Disjoint unions give block sums, including the empty-port unit. The coherence maps are the canonical finite-product and coproduct bijections.

The coefficient of $(i,j)$ is $\sum_{e:s(e)=i,\,t(e)=j}w(e)$. A linear map is determined uniquely by its images of the free generators. Conversely its coefficients are represented by the span with occurrence object $I\times J$ and those weights. Thus the quotient has exactly the stated hom objects. The quotient uses equality of the denoted coefficients. Constructing the matrix of a presented span requires no decision procedure for equality of located reals.
\end{proof}
The witness span displays the composition used in the proof.
\[
\begin{tikzcd}[column sep=small,row sep=normal]
 &&E\times_J D\arrow[dl,"p_E"']\arrow[dr,"p_D"]\\
 &E\arrow[dl,"s"']\arrow[dr,"t"]&&D\arrow[dl,"u"']\arrow[dr,"v"]\\
 I&&J&&K .
\end{tikzcd}
\]
The diamond with bottom vertex \(J\) is cartesian. Its upper vertex enumerates precisely the compatible wires. Pulling both weight fields to that vertex and multiplying gives the weight of the composite span. Linearization sends this pasting to composition of linear maps. Choosing bases writes its coefficients as matrix multiplication.
In \cref{con:layer-schema} let $j:C_0\to J$ and $i:C_0\to I$ be the feature keys, $\gamma=\langle j,i\rangle$, and $u:C_0\to R$ the stored feature. The field $x=\Sigma_\gamma u$ has the input type of the channel layer. The query is
\begin{equation}\label{eq:layer}
 x_{j,i}=\sum_{\gamma(c)=(j,i)}u_c,\qquad
 a_{j,k}=b_k+\sum_{t(d)=k}\theta_{\rho(d)}x_{j,s(d)},\qquad
 h_{j,k}=\phi(a_{j,k}).
\end{equation}
Indeed the join has rows $(c,d)\in C_0\times_I D$, output $(j(c),t(d))$ and contribution $u_c\theta_{\rho(d)}$. Grouping by $d$ gives the displayed inner fibre sum by distributivity. Empty fibres contribute zero, and repeated cells contribute once per row. In the dense instance $\gamma$ is the identity and the join is $J\times D$. Pullback identifies the same pairs fibrewise by \cref{lem:bc}, so this calculation commutes with parameter substitution. Figure~\ref{fig:feature-fibre} follows one fibre in both directions.

\begin{figure}[!htbp]
    \centering\small
    \begin{tikzpicture}[x=1cm,y=1cm]
        \begin{scope}
            \node at (1.9,1) {Forward sum};
            \node (a) at (0,.42) {$u_{c_1}$};
            \node (b) at (0,-.42) {$u_{c_2}$};
            \node[figbox] (sum) at (1.75,0) {$+$};
            \node (y) at (3.65,0) {$x_{j,i}$};
            \draw[figwire] (a) -- (.7,.42) -- (sum);
            \draw[figwire] (b) -- (.7,-.42) -- (sum);
            \draw[figwire] (sum) -- (y);
        \end{scope}
        \begin{scope}[xshift=6cm]
            \node at (1.9,1) {Reverse pullback};
            \node (a) at (0,.42) {$\bar u_{c_1}$};
            \node (b) at (0,-.42) {$\bar u_{c_2}$};
            \node[dot] (copy) at (1.75,0) {};
            \node (y) at (3.65,0) {$\bar x_{j,i}$};
            \draw[figwire] (y) -- (copy);
            \draw[figwire] (copy) -- (.7,.42) -- (a);
            \draw[figwire] (copy) -- (.7,-.42) -- (b);
        \end{scope}
    \end{tikzpicture}
    \caption{A fibre of $\gamma$ with two stored cells. Read the left picture from left to right. Its transpose, read from right to left, copies the same receipt $\bar x_{j,i}$ to both cells.}
    \label{fig:feature-fibre}
\end{figure}


\subsection{Foreign keys, operations, and input positions}
In the layer above, a single contribution reads both a feature and a weight. More generally, let $b\in B_j$ index a summand at output $j$, and let $e\in E_b$ index an input use in that summand. We call this position a socket. Each use has a source address, each use belongs to one summand, and each summand has an output address. Reading these three assignments gives the polynomial diagram \cite{poly:gambino-kock,poly:containers}
\[
\begin{tikzcd}[column sep=large]
 I & E \arrow[l,"s"'] \arrow[r,"p"] & B \arrow[r,"t"] & J
\end{tikzcd}
\]
which defines the internal polynomial functor
$\mathcal P=\Sigma_t\Pi_p\Delta_s:\E/I\to\E/J$.
Here $B_j$ is the object of operations with output type $j$, and
$E_b$ is the fibre of input positions of $b$. Internally,
\[
 (\mathcal P X)_j=\sum_{b\in B_j}\prod_{e\in E_b}X_{s(e)}.
\]
The sums and products in this formula are dependent type constructors.
Their arithmetic interpretations below are separate maps, constructed
by finite summation and multiplication in a semiring. The distinction
keeps a type of possible inputs separate from their numerical values.
In the database interpretation, two input occurrences may refer to the same input \(s(e)\). The product then uses that input twice. A nullary summand supplies the empty product \(1\). An output with no summands receives \(0\). In the type reading those conventions are the terminal object and the initial object of the corresponding fibres.
\begin{theorem}[Composition and finite arithmetic interpretation]
\label{poly:thm:interpretation}
Polynomial diagrams compose by substituting an operation into each
input position. For finite presented occurrence diagrams, an internal
commutative semiring $K$ gives an interpretation
\[
 \langle\mathcal P,w\rangle_K:K^I\to K^J,\qquad
 x\longmapsto
 \left(\sum_{b\in B_j}w_b\prod_{e\in E_b}x_{s(e)}\right)_{j\in J}.
\]
Composition is interpreted by substitution of these expressions, with
the weight of a composite operation equal to its outer weight times
the weights of its selected inner operations. The interpretation
commutes with semiring homomorphisms and substitution of external
parameters. Singleton input fibres give the existing weighted-span
interpretation. Empty input fibres give constants.
\end{theorem}
\begin{proof}
For an outer operation $c$ with positions $D_c$, a composite operation
consists of $c$ and, for every $d\in D_c$, an inner operation
$b_d$ of the required output type. Its positions are the dependent
sum $\sum_{d\in D_c}E_{b_d}$. These are the objects of operations and
positions of the composite diagram. The dependent distributivity
isomorphism constructs this substitution in $\E$. In the finite
arithmetic interpretation the corresponding identity is
\[
 v_c\prod_{d\in D_c}
       \left(\sum_{b\in B_{s(d)}}w_b\prod_{e\in E_b}x_{s(e)}\right)
 =\sum_{(b_d)}v_c\prod_dw_{b_d}
                  \prod_{(d,e)}x_{s(e)}.
\]
Finite distributivity proves the equality without selecting a summand.
The unit has one unary operation at each type, of weight $1$.
Both associations of a triple composite describe the same three-level
family of operations and positions. Dependent reassociation gives the
associator and its pentagon. Semiring maps preserve each displayed
finite sum and product. Parameter substitution preserves the diagram
and these operations, proving the two naturalities. For
$I\leftarrow E\xrightarrow{\id}E\to J$ the formula is
$t_!m_ws^*$, as in \cref{neural:thm:linearization}.
\end{proof}

\subsection{Elaborating a query into operations and input uses}
\begin{theorem}[Witness elaboration]\label{thm:elaboration}
In a finite key context a guarded query with confirmed guards and finitely presented witness fibres constructs a diagram $I\leftarrow E_q\to B_q\to J$. Equivalently one may retain the guard as a base subobject internally. Finite arithmetic execution then uses a presentation of the eligible rows. The finite evidence stages of \cref{thm:evidence} have such presentations. For every commutative semiring the diagram's arithmetic interpretation gives exactly the annotated query result. Substituting one eligible query presentation into another corresponds to polynomial substitution.
\end{theorem}
\begin{proof}
For each witness \(b\) keep its output \(t(b)\), socket fibre \(E_b\), and source function \(s_b:E_b\to I\). With the eligible guard table in the base, the recursion is
\[
\begin{array}{l|l|l}
 \text{constructor}&\text{witness at the output}&\text{socket fibre}\\ \hline
 \text{scan}& b=\text{eligible base row}&1\ \text{at its annotation address}\\
 1,\ 0 &\text{one nullary witness, or no witness}&\varnothing\ \text{for the nullary witness}\\
 q+q'&\mathrm{inl}(b)\ \text{or}\ \mathrm{inr}(c)&E_b\ \text{or}\ E_c\\
 q\Join q'&(b,c)\in B_q\times_{\rm key}B_{q'}&E_b+E_c\\
 q|_G& b\in B_q,\ G(b)&E_b\\
 \Sigma_h q& b\in B_q,\ \text{output }h(t(b))&E_b
\end{array}
\]
Set \(E=\sum_b E_b\), \(p(b,e)=b\), \(s(b,e)=s_b(e)\). For a composite with outer witness \(c\) and positions \(D_c\), the comparison of nested and flat witnesses is the bijection
\[
 (c,(b_d)_{d\in D_c},d,e)\ \longleftrightarrow\
 ((c,(b_d)_{d\in D_c}),(d,e)),\qquad e\in E_{b_d}.
\]
It preserves the source of \(e\) and the output of \(c\). Arithmetic evaluation sends a tagged union to addition and \(E_b+E_c\) to multiplication of the two socket products. Grouping witnesses over \(h(t(b))\) gives the projection sum. The distributivity identity in \cref{poly:thm:interpretation} gives substitution. For three levels both reassociations retain the same ordered chain of three socket addresses.
\end{proof}
\paragraph{What the coherent row language sees.}\phantomsection\label{sec:coherent-row-language}
The arithmetic step in witness elaboration has to interpret the eligible table. Even the graph of the sum of unit weights over one varying fibre has no uniform coherent definition in the bare entity--attribute language. Take one parent, one child of that parent in the first structure and two in the second, all of weight \(1\), with the same rational scalar algebra and no symbols for reductions. Include the single child into the two-child table and act identically on the parent and on scalars. This is a homomorphism. Every coherent formula is preserved by it: equality is preserved on terms, conjunction and disjunction preserve this implication, and an existential witness maps to an existential witness. A graph formula true of the parent and the scalar \(1\) in the first structure would therefore be true of that same pair in the second, where the sum is \(2\). In the field theory the symbol \(\Sigma_f\) has the finite execution of \cref{thm:elaboration}. Coherent completion of this enlarged signature retains that interpretation, as explained in \cref{app:field-presentations}.

For a finite candidate table \(B\) over the current shape, let \(g:S\to\Omega^B\) classify the guards and \(i:2\hookrightarrow\Omega\) classify decidable truth. The exact mask domain is characterized by
\[
\begin{tikzcd}[column sep=large]
 S_{\rm mask}\arrow[r,"m"]\arrow[d,hook,"j"']&2^B\arrow[d,hook,"i^B"]\\
 S\arrow[r,"g"']&\Omega^B
\end{tikzcd}
\qquad
 \Hom_S(Y,S_{\rm mask})\cong
 \{d:Y\to2^B:i^Bd=g h\}\quad(h:Y\to S).
\]
Given \(d\), the eligible rows are \(\{(y,b):d(y)(b)=1\}\). Apply this construction in the current candidate slice for each successive guard. At stage \(n\), the finite table decides which successes are recorded by that stage. New records can enlarge the next table. To see the effect on marking, the one-row guarded query \(a\,\mathbf1_{a>1}\) vanishes at every rational \(0\le a\le1\) and takes value \(2\) at \(a=2\). A polynomial vanishing on that interval is identically zero. The mask base thus determines where the indicated occurrence polynomial and its derivatives apply.

The two diagrams distinguish restriction from projection. Restricting the operation table to $B'$ pulls its socket table back to $E'=B'\times_B E$. Projecting a result along $h$ keeps the operation and socket tables and changes only the output address.
\[
\begin{tikzcd}
 E'\arrow[r]\arrow[d,"p'"']&E\arrow[d,"p"]\\
 B'\arrow[r,"g"']&B
\end{tikzcd}
\qquad
\begin{tikzcd}
 B\arrow[r,"t"]\arrow[dr,"ht"']&J\arrow[d,"h"]\\
 &J'.
\end{tikzcd}
\]
\begin{corollary}[Finite polynomial characterization]\label{thm:polycl}
The numerical interpretations of finite weighted occurrence diagrams are precisely tuples of presented coordinate polynomials. With the same scalar coefficients over $\Q$ or the located real line, affine--square circuits have precisely these polynomial interpretations.
\end{corollary}
\begin{proof}
Make one operation row for each monomial occurrence, labelled by its coefficient and output, and one input socket for each factor in that monomial. Reading back the diagram gives the initial polynomial, including constant monomials and empty output sums. In the reverse direction each operation contributes its coefficient times the monomial spelled by its labelled sockets. An affine--square circuit expands by induction on its gate list. Conversely use $uv=((u+v)^2-(u-v)^2)/4$ at each product, then form the finite weighted sums. The translations use finite syntax recursion and never compare two real coefficients for equality. The comparison identifies formal polynomial interpretations. A register presentation additionally records which intermediate results were shared.
\end{proof}

Consider a single product row with two socket IDs and the same foreign key $\theta$ at both positions. It evaluates to $\theta^2$. Its two uses can be varied independently in the witness table before lookup identifies their values. We follow them through the construction of jets and return to their shared parameter in \cref{thm:joint-jet}.

\section{Free algebras and parameterized evaluation}\label{sec:para}

A query for a single gate specifies its input positions. A network also specifies what fills those positions, how often a cell runs, and which calls read the same parameter. Trees describe the substitutions. To describe sharing, we evaluate a tree with a parameter held fixed throughout the recursion. This will let us ask when a comparison verified on one cell remains valid on the whole network.

\subsection{Operations and well-founded trees}

An operation grammar has a single object $I$ of all its input and output
sorts. Its diagram is therefore an endopolynomial
$F:\E/I\to\E/I$. A variable of sort $i$ is a leaf. A node labelled
$b$ has children indexed by $E_b$, with the prescribed sorts. Write
$\mathsf T_FX$ for the object of these well-founded trees whose free
leaves are labelled in $X\to I$.
\begin{theorem}[The free polynomial monad on a query signature]
\label{poly:thm:free}
The tree object has an isomorphism
\[
 \mathsf T_FX\cong X+F(\mathsf T_FX).
\]
Variable insertion $\eta$ and grafting $\mu$ make $\mathsf T_F$ the
algebraically free monad on $F$. Its underlying functor is polynomial:
its operations are operation-labelled shapes with undecorated free
leaves, and its input positions are those leaves. Its strength distributes an external parameter to
every free leaf. These constructions are stable under substitution of
external parameters.
\end{theorem}
\begin{proof}
Hyland's topos and each of its slices are elementary toposes with a natural numbers object. The existence of W-types in such a topos is recalled in \cite[\S4.3]{poly:gambino-kock}. The passage to indexed branching is the initial-algebra theorem for dependent polynomials \cite[Theorem 12]{gambino-hyland}. Apply it to \(X+F(-)\) in \(\E/I\), using variables of sort \(i\) as its leaf constructors. The free-polynomial-monad theorem \cite[Theorem 4.5]{poly:gambino-kock} identifies the resulting strong monad and its cartesian unit and multiplication.

In the present notation its polynomial is
\[
 (\mathsf T_FX)_i
   =\sum_{c\in C_i}\prod_{l\in L_c}X_{\operatorname{sort}(l)},
\]
where $C_i$ is the object of well-founded operation shapes and $L_c$
their free leaves. The unit inserts a variable leaf, multiplication
grafts its decorating tree, and strength copies the external
parameter to each leaf. Initiality identifies algebras with
interpretations of the generating operations. For an external parameter substitution let $r^*$ denote reindexing on the associated slices of sorted families. Put $H=X+F(-)$ and let $H'$ be its reindexed polynomial. The dependent sum and product comparisons give $r^*H\cong H'r^*$. If $b:H'B\to B$ is an algebra, give $\Pi_r B$ the $H$-algebra structure adjoint to
\[
 r^*H\Pi_r B\ \cong\ H'r^*\Pi_r B
      \xrightarrow{H'\varepsilon_B} H'B\xrightarrow{b}B .
\]
Under $r^*\dashv\Pi_r$, a map $r^*W\to B$ respects the algebra structures precisely when its transpose $W\to\Pi_r B$ does, by naturality of the counit. Initiality of $W$ now gives initiality of $r^*W$. The resulting comparison with the tree object in the new slice is unique as an algebra map, so these comparisons compose. This argument keeps the operation sorts fixed under external parameter substitution.
\end{proof}

\paragraph{The free--forgetful adjunction.}
Write \(U:\operatorname{Alg}(F)\to\E/I\) for the forgetful functor and \(L X=(\mathsf T_FX,\mathrm{node})\). The universal extension just constructed gives
\[
 \Hom_{\operatorname{Alg}(F)}(LX,(A,a))
 \ \cong\ \Hom_{\E/I}(X,U(A,a)).
\]
Restriction composes with \(\eta_X\). In the other direction an assignment on \(X\) labels the variables of a tree and the algebra evaluates its nodes. The two maps are inverse by induction, and substitution of assignments and postcomposition with an algebra morphism commute with evaluation by that same induction. Thus \(L\dashv U\), with counit evaluation at the identity assignment on \(A\). The induced monad \(UL\) has multiplication \(U\epsilon_{LX}\), precisely grafting a tree of trees. Working in the external parameter context gives the fixed-parameter evaluation below.

For an ordered finite presentation of the gate arities, encode a sorted variable by a leaf and a gate by its label and an ordered list of child codes. Recursion returns the assigned value at a leaf. At a gate it evaluates the child list and applies the gate to the resulting values. Induction gives uniqueness. Grafting fills each free leaf with its specified tree. Filling the leaves in two successive stages gives the same tree in either association, and grafting variables leaves the tree unchanged. These finite codes allow arbitrary finite depth.

An implementation can also retain a shared intermediate value in a register. Encode it by a list of assignments, each referring only to earlier registers. Expanding a reference into its defining expression yields a tree with the same value, by induction on the list. The list additionally records which intermediate values the program evaluates just once.

The sorts belong to the fixed base $I$. An external parameter $P$ in $\E$ acts on $\E/I$ by $(P\odot X)_i=P\times X_i$. In the parameterized formulas below $P\times X$ in a slice denotes this action. An operation can read a different sort from its output, so a single register is supplied to all sorts. For example, a unary operation of output sort $1$ reading sort $0$ would need to turn a sort-$1$ parameter into a sort-$0$ parameter if one used arbitrary sort-indexed registers. The external action constructs the required distribution directly.

\begin{construction}[Parameterized composition]\label{poly:con:para}\leavevmode
\begin{enumerate}[titledsteps]
    \item \textbf{Compose the implementations.} Let \(P\odot X\) be the external action on the slice fixed above. A parameterized arrow \((P,f):X\to Y\) has implementation \(f:P\odot X\to Y\). The two diagrams give serial composition and a change of parameter:
\[
\begin{tikzcd}[column sep=large]
 (P\times Q)\odot X\arrow[r,"\simeq"]\arrow[d,"g\star f"']&
 Q\odot(P\odot X)\arrow[d,"1\odot f"]\\
 Z&Q\odot Y\arrow[l,"g"]
\end{tikzcd}
\qquad
\begin{tikzcd}[column sep=small]
 P'\odot X\arrow[r,"r\odot1"]\arrow[dr,"f'"']&
 P\odot X\arrow[d,"f"]\\
 &Y .
\end{tikzcd}
\]
The first isomorphism uses the cartesian symmetry to list parameters in execution order. 
    \item \textbf{Change the parameter register.} In the convention of \cite[Definition G.1]{cdl}, \(r:P'\to P\) exhibiting \(f'=f(r\odot1)\) is a 2-cell \((P,f)\Rightarrow(P',f')\). The underlying map \(r\) always points from the new parameter register to the old one. Horizontal composition takes the product of these maps in the displayed order. The product associators and unitors give the bicategory coherence. Parallel composition applies \(f\) and \(g\) on separate inputs after permuting the parameter and input factors. To share \(P\) between two calls precompose their independent implementation with \(\Delta_P\odot1\).
\end{enumerate}
\end{construction}
The strength sends $(\theta,b,(x_e)_e)$ to $(b,((\theta,x_e))_e)$. Thus every position $e$ can read $\theta_{\rho(e)}$ from the same external table. Applying strength for two successive parameter objects pairs the same two parameters with every child, in the order specified by the product associator.
Let \(r:P\to Q\) compare \(f:P\odot X\to Y\) and \(g:Q\odot X\to Y\), so \(f_p=g_{r(p)}\). For a shape \(b\) with positions \(E_b\), the strength acts by
\[
 \widehat F(f)(p,b,(x_e)_e)=(b,(f_p x_e)_e)
                          =\widehat F(g)(r(p),b,(x_e)_e).
\]
For two composable maps the entry at \(e\) is \(g_q(f_p x_e)\) in either order. Write \(T=\mathsf T_F\). Compatibility with grafting is the commutative square
\[
\begin{tikzcd}[column sep=large]
 P\odot TTX \arrow[r,"\mathrm{st}_{P,TX}"] \arrow[d,"1\odot\mu_X"']&
 T(P\odot TX) \arrow[d,"{\mu_{P\odot X}\,\circ\,T(\mathrm{st}_{P,X})}"]\\
 P\odot TX \arrow[r,"\mathrm{st}_{P,X}"']&T(P\odot X).
\end{tikzcd}
\]
At a leaf of an inner tree both composites return the pair \((p,x)\) at exactly that leaf. The constructors and their order agree, which proves equality by induction on the outer tree.

In a concrete field model the composition rule of \(\mathbb F_\Theta\) in \cref{sec:model-interfaces} induces the same recursion on finite gate trees. At a variable, evaluation projects to that input. At a node with children \(t_i\), it gives \(a(p,(\llbracket t_i\rrbracket_p)_i)\), with the same parameter passed to every child. Induction therefore identifies evaluation of a finite register expression after expansion of its references with the free evaluation below. A repeated reference uses its compiled value at each occurrence; the parameter diagonal appearing in the free-algebra law is the one in parameterized composition.

\begin{theorem}[Parameterized free evaluation]
\label{poly:thm:parametric}
Every operation interpretation $a:\Theta\times FA\to A$ has a
unique extension $\bar a:\Theta\times\mathsf T_FA\to A$ satisfying
\begin{align}
 \bar a_\theta(\eta x)&=x,\label{poly:eq:param-unit}\\
 \bar a_\theta(\operatorname{node}(b,(u_e)))
   &=a_\theta(b,(\bar a_\theta(u_e))).\label{poly:eq:param-node}
\end{align}
For each internal parameter it satisfies
\begin{equation}\label{poly:eq:param-mult}
 \bar a_\theta\mu_A
       =\bar a_\theta\mathsf T_F(\bar a_\theta).
\end{equation}
In the $\para$ convention of \cref{poly:con:para} these equations give a lax algebra for the lifted free
monad, with coherence cells given by
$!_\Theta:\Theta\to1$ and
$\Delta_\Theta:\Theta\to\Theta\times\Theta$.

More explicitly, let morphisms between operation interpretations
$(\Theta,A,a)$ and $(\Lambda,B,b)$ be pairs
$(r:\Theta\to\Lambda,h:A\to B)$ satisfying
$h a_\theta=b_{r\theta}Fh$. The extension is an equivalence with
the category of parameterized $\mathsf T_F$-algebras satisfying
\eqref{poly:eq:param-unit} and \eqref{poly:eq:param-mult}, with
morphisms satisfying the corresponding extended equation.
\end{theorem}
\begin{proof}\leavevmode
\begin{enumerate}[titledsteps]
    \item \textbf{Extend evaluation to trees.} Perform the tree recursion of \cref{poly:thm:free} in the context
$\theta:\Theta$. This proves existence, uniqueness, and parameter
stability without choosing a global parameter. To prove
\eqref{poly:eq:param-mult}, induct on the outer tree. At a variable
leaf both sides evaluate its substituted tree. At a node, the root
operation on both sides is $a_\theta$. The induction equations at
its children identify the arguments. This proves the equation.


    \item \textbf{Verify the lax algebra laws.} In parameterized composition, the right-hand side initially carries
two parameters: one for the inner evaluations and one for the outer
evaluation. Precomposition by $\Delta_\Theta$ makes both equal to
$\theta$, yielding the lax multiplication cell. The unit equation
forgets an unused parameter through $!_\Theta$. Both three-level
comparisons copy $\theta$ to the same triple. Coassociativity of the
diagonal and its two counit equations prove the lax associativity and
unit laws, with the displayed product associators.


    \item \textbf{Extend comparisons.} For a pair $(r,h)$ as in the statement, induction gives
$h\bar a_\theta=\bar b_{r\theta}\mathsf T_Fh$: variables use $h$
and the node step uses the defining equation for $(r,h)$. Conversely
restrict that equation to one-node trees. Restriction and extension
are inverse on objects by uniqueness, and on these pairs of arrows
by their unchanged underlying maps. Identity and composition of the
pairs preserve the equations, proving the stated equivalence.\qedhere
\end{enumerate}
\end{proof}
For each $\theta$, the square identifies two ways to compare a tree evaluation. One can evaluate first and map the result by $h$, or apply $h$ at the variables and evaluate using the translated parameter $r\theta$.
\[
\begin{tikzcd}[column sep=large]
 \mathsf T_F A\arrow[r,"\bar a_\theta"]\arrow[d,"\mathsf T_F h"']&
 A\arrow[d,"h"]\\
 \mathsf T_F B\arrow[r,"\bar b_{r\theta}"']&B .
\end{tikzcd}
\]
The square reduces comparison of two evaluators to comparison of their generating gates. So far the carrier map $h$ in that square has been fixed. A parameterized carrier map raises a further question: can its own parameter change as it passes from a node to a child?

\subsection{Comparison cells for a whole architecture}
Throughout, a cell $(P,f)\Rightarrow(Q,g)$ is represented by $Q\to P$ as in \cite[Definition G.1]{cdl}. In the unit pasting the identity, with terminal register, maps to evaluation after variable insertion, with register $P$, by the underlying discard $P\to1$. In the multiplication pasting the two successive evaluations map to evaluation after grafting by the underlying diagonal $P\to P\times P$. These are the lax algebra cells in \cite[Definition F.2 and Theorem G.10]{cdl}.
For $T=\mathsf T_F$, the triangle factors that action into distributing a register to the leaves and applying $f$ there.
\[
\begin{tikzcd}[column sep=large]
 P\odot TX\arrow[r,"\mathrm{st}"]\arrow[dr,"{\widehat T(f)}"']&
 T(P\odot X)\arrow[d,"Tf"]\\
 &TY .
\end{tikzcd}
\]
We identify products by their canonical associators. In a cartesian parameter category, the counit $P\to1$ is unique. For a parameter comultiplication $d:P\to P\times P$, the two counit laws require $\pi_1d=\pi_2d=\id$. Hence $d=\Delta_P$, and the algebra laws reduce to the pointwise equations of \cref{poly:thm:parametric}. The remaining freedom lies in the morphisms between algebras.

\begin{theorem}[Extension of architecture comparisons]\label{thm:cell-ext}
Let $a:P\odot FA\to A$ and $b:Q\odot FB\to B$ be operation interpretations. A comparison with carrier $h:U\odot A\to B$ has a one-gate reparameterization
\[
 k:U\times P\longrightarrow Q\times U,\qquad
 k(u,p)=(r(u,p),v(u,p)),\qquad
 h_u a_p=b_{r(u,p)}F(h_{v(u,p)}).
\]
It extends, with the given carrier and register map, to a lax morphism of the lifted tree algebras if and only if $v=\pi_U$. Restriction and extension identify the category of such one-gate comparisons with the category of lax morphisms between these tree algebras.
\end{theorem}
\begin{proof}\leavevmode
\begin{enumerate}[titledsteps]
    \item \textbf{Apply the unit law.} Let \(k(u,p)=(r(u,p),v(u,p))\). The parameter map at a variable is independent of the algebra. The unit coherence equation is
\[
 (!_{Q}\times1_U)k=1_U\times!_P:U\times P\longrightarrow U
             \quad\Longleftrightarrow\quad v=\pi_U,
\]
using the canonical unitors. This is equality of register maps even when two register values induce the same map on the carrier.


    \item \textbf{Extend the comparison by induction.} With that equality, fix $u,p$ and induct on a tree. Both routes in the right square return $h_u(x)$ on a variable. At a node the one-gate equation identifies the root operations, and the inductive equations identify their arguments.
\[
\begin{tikzcd}[column sep=small]
 U\times P\arrow[r,"k"]\arrow[d,"\pi_U"']&Q\times U\arrow[d,"\pi_U"]\\
 U\arrow[r,equal]&U
\end{tikzcd}
\qquad
\begin{tikzcd}[column sep=small]
 TA\arrow[r,"Th_u"]\arrow[d,"\bar a_p"']&TB\arrow[d,"{\bar b_{r(u,p)}}"]\\
 A\arrow[r,"h_u"']&B .
\end{tikzcd}
\]

    \item \textbf{Check multiplication and composition.} For multiplication let $q=r(u,p)$. Both routes send $(u,p)$ to $(q,q,u)$. One copies $q$. The other translates two copies of $p$ with the same $u$. Restricting the tree equation to one-node trees recovers precisely the one-gate equation and the register map $k$. The identity uses the terminal register and the identity on $P$. The composite of two comparisons is determined by
\[
 (v,u)\longmapsto j_vh_u,\qquad
 (v,u,p)\longmapsto s(v,r(u,p)).
\]
Extension and restriction preserve these maps. For the ordinary category assertion take comparisons up to isomorphism of their carrier register: $\phi:U\simeq U'$ intertwines $h$ and $k$ when $h'_{\phi u}=h_u$ and $k'(\phi u,p)=(1_Q\times\phi)k(u,p)$. Composition respects these isomorphisms by their product. Its associator and unitors are the canonical product maps, which intertwine the displayed formulas. Equivalently use lists of registers, concatenation and the canonical flattening of tuples. At the level of individual representatives the restriction and extension formulas are mutually inverse before these identifications. The unsimplified parameter pastings are given in \cref{app:pasting}.\qedhere
\end{enumerate}
\end{proof}

\begin{example}[A one-gate comparison that changes state]\label{ex:bit-obstruction}
For the unpointed unary signature $F=\id$, one has $TX=\N\times X$. Take terminal algebra parameters, bit carriers, $a(x)=1-x$, $b(x)=x$, and $h_u(x)=x\mathbin{\oplus}u$ on a comparison bit $u$. The one-gate equation holds with $v(u)=1-u$. A comparison for arbitrary terms would need a single map $v:2\to2$. The variable term forces $v(u)=u$, whereas a one-node term forces $v(u)=1-u$. This exhibits the obstruction for unrestricted comparisons of an unpointed signature. The pointed transfinite formulation in \cite[Conjecture G.12]{cdl} has a different source construction.
\end{example}

\begin{theorem}[Traversal with changing parameters]\label{thm:traversal}
Write $S_P$ for the free monad on $P\odot F(-)$, with a parameter label at each node. Every $k:U\times P\to Q\times U$ induces a natural traversal $d_X:U\odot S_P X\to S_Q(U\odot X)$, a colax morphism of monads in this displayed orientation. For operations $a,b,h$ satisfying the one-gate comparison and their evaluations $e_a:S_P A\to A$, $e_b:S_Q B\to B$,
\[
 h_u(e_a(t))=e_b\big(S_Q(h)(d_A(u,t))\big).
\]
\end{theorem}
\begin{proof}
At a variable return $(u,x)$. At a node labelled $(p,c)$ compute $(q,v)=k(u,p)$, emit $(q,c)$, and traverse all children with state $v$. Relabelling variables commutes with this recursion. The variable clause is the unit law. The node equation for $h$ and induction at its children give the evaluation identity. The multiplication law is
\[
\begin{tikzcd}[column sep=normal]
 U\odot S_P^2X \arrow[r,"d_{S_PX}"]\arrow[d,"{1\odot\mu^P}"']&
 S_Q(U\odot S_PX) \arrow[r,"S_Q d_X"]&S_Q^2(U\odot X)\arrow[d,"\mu^Q"]\\
 U\odot S_PX \arrow[rr,"d_X"']&&S_Q(U\odot X).
\end{tikzcd}
\]
Induct on the outer tree. At a variable, each route enters the inserted tree with the state carried to that variable. At a node the two routes compute the same $(q,v)$ before entering any child. This also settles a nullary node, where there are no children. The square therefore commutes for grafting. In \cref{ex:bit-obstruction} the state arriving at a variable records the parity of its depth.
\end{proof}

The choice between the two evaluations can now be read from a single gate. If $v(u,p)=u$ and all source nodes have label $p$, the entire traversal uses one target parameter $r(u,p)$. Otherwise the traversal records the target parameter at each node, and the grafting law explains how those records compose.

\clearpage
\part{Representability, descent and feedback}\label{part:feedback}

\section{Marked sockets and substitution}\label{sec:variation}

A polynomial operation is presented by \(I\leftarrow E\xrightarrow{p}B\to J\). Its object of operations with one distinguished input position is \(E\) itself: a mark on operation \(b\) chooses precisely an element of its fibre. Relative to this chosen socket, the other positions are the complement of the diagonal in \(E\times_B E\), on the presented finite decidable fibres. The square
\[
\begin{tikzcd}[column sep=large]
 E\times_B E\arrow[r,"\pi_2"]\arrow[d,"\pi_1"']&E\arrow[d,"p"]\\
 E\arrow[r,"p"']&B
\end{tikzcd}
\]
is cartesian. Over \(e\) on the left its fibre is the entire input family of \(p(e)\). The diagonal section singles out the position to be varied. Labelling that position by an increment and its complement by ordinary inputs gives the one-mark evaluation below. The square-zero formula computes it in an annotation algebra.

\subsection{Square-zero computation and marked positions}
First-order variation is constructed in the annotation algebra before
choosing a neural realization. For a commutative semiring object $K$,
put $\operatorname{Sq}(K)=K\times K$, with
\[
 (a,u)+(b,v)=(a+b,u+v),\qquad
 (a,u)(b,v)=(ab,av+bu),
\]
and units $(0,0),(1,0)$. The notation $a+\epsilon u$ records this
pair and the multiplication law $\epsilon^2=0$. It does not introduce
infinitesimals into every object of the ambient topos.
\begin{theorem}[Square-zero evaluation and the chain rule]
\label{poly:thm:square-zero}
The displayed operations make $\operatorname{Sq}(K)$ a commutative
semiring. Every presented finite arithmetic polynomial $f:K^I\to K^J$
has an induced map
\[
 \operatorname{Sq}(f)(x,u)=(f(x),D f_x(u)),
\]
obtained by evaluating its constants as $(c,0)$ and its variables as
$(x_i,u_i)$. It satisfies
\[
 D(gf)_x=Dg_{f(x)}\circ Df_x,\qquad D\id_x=\id,
\]
and $Df_x$ is $K$-linear in $u$. The construction respects semiring
specialization. Coefficients or parameters that vary are additional
inputs, so their variation is included by the same construction.
\end{theorem}
\begin{proof}
For three factors the second component of either association is
$ab w+ac v+bc u$, and the first is $abc$. The other semiring laws follow
from the corresponding laws in $K$ and finite distributivity.
Induction on expressions gives $D c=0$, $D x_i(u)=u_i$,
$D(f+g)=Df+Dg$, and $D(fg)=fDg+gDf$. These also prove linearity
in $u$. Evaluation commutes with substitution by
\cref{poly:thm:interpretation} and tree substitution.
Evaluating a substituted expression in $\operatorname{Sq}(K)$
and comparing second components gives the chain rule. A semiring
homomorphism $h:K\to L$ induces $(h,h)$ on square-zero extensions;
the displayed operations show that this is again a homomorphism,
proving specialization naturality.
\end{proof}
\begin{proposition}[Marked positions and positive query variation]
\label{poly:prop:marked}
For a presented finite occurrence diagram, mark one input position
of each operation and retain the other positions. Its arithmetic
interpretation is exactly the first-order map
\begin{equation}\label{poly:eq:marked}
 (D\langle\mathcal P,w\rangle_{K,x}(u))_j
  =\sum_{b\in B_j}w_b\sum_{e\in E_b}
       u_{s(e)}\prod_{e'\in E_b\setminus\{e\}}x_{s(e')}.
\end{equation}
For guarded positive database queries, with typing and guards retained
in the base, finite union and projection differentiate by their finite
additive interpretations, and natural join satisfies
\[
 D(R\Join S)(\dot R,\dot S)
       =\dot R\Join S+R\Join\dot S.
\]
\end{proposition}
\begin{proof}
Finite decidable occurrence positions have constructed complements
of singletons, so the displayed punctured product is defined without
deciding equality of semantic atom labels. Multiplying the pairs
$(x_{s(e)},u_{s(e)})$ selects precisely one second component, since two
such selections multiply to zero. Summing over operations proves
\eqref{poly:eq:marked}. Under substitution a marked leaf lies in a
uniquely specified inner operation at an outer input position. This
gives the same chain rule as square-zero evaluation. In a guarded
query the guards have zero variation and source annotations carry
the input variations. Finite sums and products then give the query
formulas.
\end{proof}

The marking construction gives a row $(b,e)$ for each socket $e$ of operation $b$. Its marked input reads the tangent $u_{s(e)}$, and each other input $e'$ reads its original value $x_{s(e')}$. Testing whether $e'=e$ compares addresses in a finite decidable socket fibre. Two sockets that happen to carry equal numerical values remain distinguishable. This operation is the one-hole construction familiar from container derivatives \cite{poly:container-derivative}. Its numerical meaning follows here from multiplication in the square-zero extension. A variable coefficient gets a socket of its own, contributing $\sum_b\dot w_b\prod_{e\in E_b}x_{s(e)}$ at the corresponding output.

\subsection{Several marks and the order of a jet}
For a row $b$ of $I\xleftarrow{s}E\to B\to J$, a mark of colour $i$ retains the row, a position $e$ with $s(e)=i$, and values at its other positions. In the slice over an output key $j$, its one-hole contexts form the container
\[
 (\partial_i\mathcal P)(X)_j=
 \sum_{b\in B_j}\ \sum_{\substack{e\in E_b\\s(e)=i}}
                  \prod_{d\in E_b\setminus\{e\}}X_{s(d)} .
\]
Its direction is supplied by filling the hole. For several distinguishable marks, retain their position IDs before projecting back to the input keys. This distinction is visible even in a product with two equal factors.

\begin{theorem}[Marked substitution and truncated evaluation]\label{thm:marked-jets}
For a finite presented arithmetic query over a commutative semiring $K$, let $K_n$ be the length-$(n+1)$ coefficient tables, with convolution truncated above $n$. Evaluation over $K_n$ is a functorial forward order-$n$ jet. In the separate algebra $K[\epsilon_1,\ldots,\epsilon_n]/(\epsilon_i^2)$, feed each socket a base value plus a sum of linear marked inputs. The coefficient of $\epsilon_1\cdots\epsilon_n$ is the sum over injections of the labelled marks into the socket fibre. For a substituted query, the same coefficient is obtained by partitioning the marks by outer socket and passing each block to its inner query. Both rules commute with truncation, change of coefficients and relabelling of the displayed keys.
\end{theorem}
\begin{proof}
For \(S\subseteq[n]\), write \(\epsilon_S=\prod_{i\in S}\epsilon_i\). Coefficient tables multiply by
\[
 (ab)_S=\sum_{U\sqcup V=S}a_Ub_V,\qquad
 [\epsilon_S]\prod_{e\in E}y_e
     =\sum_{\bigsqcup_{e\in E}S_e=S}\prod_{e\in E}[\epsilon_{S_e}]y_e .
\]
If \(y_e=x_e+\sum_i\epsilon_i u_{e,i}\), each nonzero contribution assigns each colour to a different socket:
\[
 [\epsilon_S]\prod_e y_e
 =\sum_{j:S\hookrightarrow E}\ \prod_{i\in S}u_{j(i),i}
                         \prod_{e\notin j(S)}x_e .
\]
For nested sockets the indexing bijection sends a marking \(j:S\hookrightarrow\coprod_d E_d\) to the partition \(S_d=(\pi j)^{-1}(d)\) and injections \(S_d\hookrightarrow E_d\). Its inverse assembles the injections by coproduct. Unmarked outer sockets retain their whole inner witness. For a single increment,
\[
 [t^m]\prod_e y_e(t)=\sum_{\sum_e m_e=m}\prod_e[t^{m_e}]y_e(t).
\]
Nonnegative degrees give compatibility with truncation. Relabelling transports the index sets bijectively. Semiring morphisms preserve their finite products and sums. Orders zero and one give ordinary and square-zero evaluation.
\end{proof}
For the row reading $\theta$ twice, the degree-two term on the arc $\theta+ut$ is $u^2$. With two independently labelled inputs $u,v$, the mixed coefficient is $uv+vu$. Passing both positions through their common foreign key preserves both ordered witnesses. This comparison also explains why ordinary truncated coefficients require no factorial divisions in $K$.

\subsection{Reading a reverse program back as a query}
\begin{proposition}[Transpose instructions and the marking square]\label{thm:compiler-reverse}
On presented finite numerical fields, the reverse pass of a compiled arithmetic query has a typed source expression. For mixed increments of every finite order, marking before compilation and evaluating the compiled program on the same increments agree, including after parameter sharing.
\end{proposition}
\begin{proof}\leavevmode
\begin{enumerate}[titledsteps]
    \item \textbf{Transpose the elementary instructions.} Pair fields by the finite sum of coordinate products. A gather along $f$ sends a seed back by $\Sigma_f$ and a sum along $f$ sends it back by $f^*$, by grouping pairs with equal image. A product reduction sends an output seed $g$ to
\[
 h_e=g_{f(e)}
       \prod_{\substack{e'\ne e\\ f(e')=f(e)}}x_{e'} .
\]
The displayed index is a structural query over ordered pairs of rows. Expansion of the product with one marked factor proves the pairing identity also at zero-valued inputs. A scalar product sends $g$ to $(gy,gx)$ at operands $(x,y)$; addition copies the seed. At a gate with a declared derivative use its transposed derivative. Initialize adjoints by output seeds (sum them if outputs repeat), traverse registers in reverse topological order, and add each receipt at its operand. At each step the finite pairing of all pending adjoints with register variations is unchanged by the local identities. At the external boundary it equals the original output pairing. Reading these instructions back by \cref{thm:compiler-struct} gives the required query.


    \item \textbf{Transport the higher-order coefficients.} At order $n$ use the single algebra $R[h_1,\ldots,h_l]/(h_1,\ldots,h_l)^{n+1}$ for all increment coordinates. Each polynomial gate is interpreted by substitution in this algebra, and a gather or reduction has exactly the same row indices in both programs. Induction on the ordered register list proves equality of each coefficient, base included. Marking sockets enumerates precisely the terms in that substitution by distributivity (\cref{thm:marked-jets}). For an owner lookup $\rho$, the induced ring map substitutes $h_{\rho(e)}$ for each use increment. In particular $p q$ at $(p,q)=(\theta,\theta)$ gives $(\theta+h)^2$ with terms $2\theta h$ and $h^2$. Mixed contributions survive even when the two sockets have the same owner.\qedhere
\end{enumerate}
\end{proof}

\section{Representing derivations, arcs and pointed observables}\label{sec:charts}

\subsection{Universal characterization in a numerical module}
Marking computes a first variation from a chosen expression. To compare expressions in the numerical ring $R$, we ask what determines that variation after passing to polynomials. Let $A$ be an internal $R$-algebra. Form the free $A$-module on symbols $[a]$ for $a\in A$ and quotient by the submodule generated by
\[
 [a+b]-[a]-[b],\quad [ab]-a[b]-b[a],\quad [r]\ (r\in R).
\]
Write $\Omega^1_{A/R}$ for the quotient and $da$ for the class of $[a]$. The effective quotient exists by the free finite-sum module construction and exactness. The relations impose additivity, the product rule, and vanishing on scalars. Thus they specify precisely when an assignment on the symbols descends to a linear map from $\Omega^1_{A/R}$ \cite[Tag 00RM]{poly:stacks-differentials}.
\begin{proposition}[Represented derivations and arbitrary generators]
\label{poly:prop:universal}
For each $A$-module $M$ there is a natural bijection
\[
 \Hom_A(\Omega^1_{A/R},M)\cong\operatorname{Der}_R(A,M).
\]
Equivalently, derivations are $R$-algebra sections of
$A\ltimes M\to A$, where
$(a,u)(b,v)=(ab,av+bu)$. For the free algebra $R[X]$ on any internal
generator object $X$,
$\Omega^1_{R[X]/R}\cong R[X]^{(X)}$.
\end{proposition}
\begin{proof}
A linear map on the free symbols descends exactly when their images
satisfy the three displayed derivation equations. This proves the
representing bijection and its naturality in $M$. A section has the
form $a\mapsto(a,\delta a)$. The algebra-homomorphism equations are
exactly those derivation equations. Finally arbitrary images of the
generators in an $R[X]$-module extend uniquely by finite sums and the
Leibniz rule to a derivation on $R[X]$. This is the universal property
of the free $R[X]$-module on $X$, proving the last isomorphism.
\end{proof}
The representing map \(d:A\to\Omega^1_{A/R}\) and any derivation \(\delta:A\to M\) determine the triangle
\[
\begin{tikzcd}[column sep=large]
 A\arrow[r,"d"]\arrow[dr,"\delta"']&\Omega^1_{A/R}\arrow[d,dashed,"{\exists!\,\widetilde\delta}"]\\
 &M .
\end{tikzcd}
\]
It commutes on underlying objects. The diagonal and top arrows are derivations and the dashed map is \(A\)-linear. Naturality of the representing bijection makes change of coefficients postcomposition by a module map. Change of algebra gives the following pullback of differentials.
For a finite free module \(V\), set \(A_V=\Sym_R(V^*)\). Its elements are coordinate polynomials, with equality of coefficients as the algebraic quotient. A polynomial map \(f:V\to W\) here is the opposite of an algebra map \(f^*:A_W\to A_V\). Evaluation at a point gives a function on the carrier. Polynomial maps keep their algebraic presentation. The construction requires no decision procedure for extensional equality of numerical programs.
\begin{lemma}[Pullback of the universal differential]\label{neural:lem:kahler}
There is a natural isomorphism $\Omega^1_{A_V/R}\cong A_V\otimes_R V^*$. The algebra map $f^*$ induces
\[
 A_V\otimes_{A_W}\Omega^1_{A_W/R}\longrightarrow\Omega^1_{A_V/R},
 \qquad a\otimes db\longmapsto a\,d(f^*b),
\]
and these pullbacks compose contravariantly with polynomial maps.
\end{lemma}
\begin{proof}
Apply \cref{poly:prop:universal} to the free generators of
$A_V$. An $R$-linear assignment on $V^*$ extends uniquely to a
derivation, so the representing module is $A_V\otimes_RV^*$.
The displayed pullback is well-defined by the same universal property
after restricting scalars along $f^*$. On a generator, successive
pullbacks give $d(f^*g^*b)$. Uniqueness identifies this with the
pullback of the composite.
\end{proof}
\begin{theorem}[Comparison of occurrences, square-zero evaluation, and cotangents]
\label{poly:thm:differential-comparison}
For each finite expression or acyclic circuit in the polynomial
numerical interpretation, the marked-occurrence rule, square-zero evaluation, and evaluation of
the universal differential give the same first-order map. They commute
with circuit substitution and reparameterization. Consequently the
cotangent interpretation of a composite architecture is the composite
of its local cotangent interpretations, including accumulation at
shared inputs and shared parameters.
\end{theorem}
\begin{proof}\leavevmode
\begin{enumerate}[titledsteps]
    \item \textbf{Construct the derivation at a point.} Let $f:V\to W$ have algebra map $f^*:A_W\to A_V$. First fix a point $x$ and a direction $u$. Evaluation at $x$ makes $R$ an $A_V$-module. Sending each linear generator $\lambda$ to $\lambda(u)$ therefore extends to a derivation on $A_V$. By the universal property, this is evaluation of $d$ in direction $u$.


    \item \textbf{Evaluate it on output coordinates.} Now apply that derivation to each output coordinate $f^*\lambda$, for $\lambda\in W^*$. Its value is the second component of $f^*\lambda$ evaluated at $x+\epsilon u$. Indeed the two assignments agree on the input generators and both satisfy additivity and the product rule. By \cref{poly:prop:marked}, the same component is the sum over marked input occurrences. This identifies all three constructions at $(x,u)$.


    \item \textbf{Check substitution and parameter sharing.} Polynomial substitution commutes with square-zero evaluation by
\cref{poly:thm:square-zero}. Pullback of the universal differential gives the
same chain rule by \cref{neural:lem:kahler}. The same calculation applies when parameters are held in a register shared by several substitutions. Include
parameter ports among the inputs when taking this differential.
For $r:\Theta'\to\Theta$ in that polynomial interpretation, pullback
along $r\times\id$ is the reparameterization comparison. In
particular $d\Delta_\Theta(u)=(u,u)$, whose dual sends
$(\alpha,\beta)$ to $\alpha+\beta$. Together with the chain rule this proves the two cotangent claims.\qedhere
\end{enumerate}
\end{proof}
The following square compares an execution on square-zero inputs with the execution that returns a value and its first variation together. Let $\chi_I$ split the pair in each input coordinate, and put $\widetilde f(x,u)=(f(x),Df_xu)$.
\[
\begin{tikzcd}[column sep=large]
 \operatorname{Sq}(R)^I \arrow[r,"f_{\operatorname{Sq}(R)}"]
    \arrow[d,"\chi_I"',"\cong"] &
 \operatorname{Sq}(R)^J \arrow[d,"\chi_J","\cong"']\\
 R^I\times R^I \arrow[r,"\widetilde f"'] &
 R^J\times R^J .
\end{tikzcd}
\]
Projecting to the second component recovers $Df_xu$ for arbitrary $u$. Restricting the comparison to constant inputs would determine only the forward value $f(x)$.

\subsection{The free differential algebra of a numerical presentation}\label{sec:free-differential-algebra}

A symbolic equation such as $y=x^2$ already specifies what its differentiated equations must say once its variables are allowed to vary. We construct a single algebra containing all these consequences. Work over $k=\Q$, or a commutative $\Q$-algebra of constant scalars. Let $A=k[x_1,\ldots,x_d]/(f_a)_a$ be a numerical chart. Fix $r\ge1$ independent formal parameters. Write $\mathsf{DiffAlg}_{k,r}$ for commutative $k$-algebras with $r$ commuting $k$-derivations and maps preserving them, and let $U$ forget the derivations.

\begin{theorem}[Free differential extension and its jet evaluations]\label{thm:free-differential-extension}
The functor $U$ admits a left adjoint on the displayed presentations. Its value on $A$ is
\[
 \mathcal J_r A=
 k[x_{i,\alpha}\mid 1\le i\le d,\ \alpha\in\N^r]\big/
       \big(\partial^\alpha f_a(x_{1,0},\ldots,x_{d,0})\big)_{a,\alpha},
 \qquad \partial_j x_{i,\alpha}=x_{i,\alpha+e_j}.
\]
Here $e_j$ is the $j$th coordinate vector and $x_{i,0}$ is the original coordinate. For total order $n$, use just $|\alpha|\le n$ in both the generators and relations and call the resulting ordinary algebra $\mathcal J_{r,n} A$. For a $k$-algebra $B$ its $B$-points correspond naturally to maps $A\to B[t_1,\ldots,t_r]/(t_1,\ldots,t_r)^{n+1}$.
\end{theorem}
\begin{proof}\leavevmode
\begin{enumerate}[titledsteps]
    \item \textbf{Construct the differential extension.} Extend the shifts $\partial_j$ by the Leibniz rule on polynomials. They commute on the generators, hence everywhere, and preserve the ideal since they carry each indicated derivative of a relation to its next derivative. Given a map $u:A\to UB$ into a differential algebra $(B,\delta_1,\ldots,\delta_r)$, the only possible differential extension sends $x_{i,\alpha}$ to $\delta^\alpha u(x_i)$. It respects every relation because differentiating $f_a(u(x))=0$ gives zero. Conversely any differential map restricts along $x_i\mapsto x_{i,0}$ to such a $u$. Restriction and extension are inverse and commute with maps on both sides, proving the adjunction. In particular the result is independent of the chosen polynomial presentation. The same construction with a set of generators works for any presented commutative $k$-algebra.


    \item \textbf{Identify the truncated points.} For finite order put $X_i(t)=\sum_{|\alpha|\le n}x_{i,\alpha}t^\alpha/\alpha!$, where $\alpha!=\prod_j\alpha_j!$. For every polynomial $f$, the coefficient of $t^\alpha$ in $f(X(t))$ is $\partial^\alpha f(x_0)/\alpha!$. It holds on each generator, and the iterated Leibniz rule gives exactly series convolution on products. Thus the defining relations of $\mathcal J_{r,n} A$ are precisely the conditions for $x_i\mapsto X_i(t)$ to define a truncated arc. Conversely extract each coefficient and multiply by $\alpha!$ to recover $x_{i,\alpha}$. These operations are inverse and natural in $B$. This is the divided-coefficient description of jets \cite{vojta-jets}; the factorials convert it to the ordinary commuting derivations used here.\qedhere
\end{enumerate}
\end{proof}

For $y=x^2$ and $r=1$, putting $x_j=\partial^j x$ and $y_j=\partial^j y$ gives $y_1=2x_0x_1$ and $y_2=2x_1^2+2x_0x_2$. The second equation allows the input itself to have a second derivative. Substituting the straight variation $x(t)=a+ht$ gives $y_2=2h^2$. With two independent straight increments $h,k$, applying $\partial_1\partial_2$ instead gives $2hk$.

Each polynomial numerical query on a fixed finite key instance induces an algebra map $q^*:A_{\mathrm{out}}\to A_{\mathrm{in}}$. The extension sends an output coordinate $y_{i,\alpha}$ to $\partial^\alpha(q^*y_i)$. The uniqueness just proved makes these maps compose under query substitution. There is also a finite execution of the formula: store the divided coefficients in fields with the original keys. Lookup reindexes each coefficient, a grouped sum adds the coefficients in its fibre, and a product takes their convolution. These are precisely the instructions of the compiled circuit over the truncated polynomial ring. Thus the symbolic interpreter and the neural interpreter compute the same coefficients, including when several sockets read one parameter. At any requested order only finitely many new fields occur. Key equalities and guards selecting the presented fibre stay in the structural context. The construction here concerns polynomial scalar operations; the normalization and implicit-solution sections give the separate coefficient rules for their numerical primitives.

Two constructions meet in this comparison. The ordinary algebra underlying $\mathcal J_r A$ is the direct limit of the $\mathcal J_{r,n} A$ along the maps sending a generator to the generator with the same index. Each defining relation occurs at a finite stage. By contrast, for a $k$-valued point $a$ of a free affine register, the inverse limit of $A/\mathfrak m_a^{n+1}$ is $k[[h_1,\ldots,h_d]]$, where $h_i=x_i-a_i$. A compatible sequence in that tower picks one coefficient per increment monomial. An arc based at $a$ pairs the two constructions by the substitution $h_i\mapsto X_i(t)-a_i$, which has zero constant term and so respects truncation. Even for one coordinate the free differential algebra $k[x_0,x_1,x_2,\ldots]$ and the point completion $k[[h]]$ have different universal properties.

\subsection{Arcs, pointed observables and their evaluation}
Work internally in $\E$, or in the slice over the presented shape. For the relative geometry take $k=\Q$ internally; the constructions below commute with scalar extension, in particular to the located real ring. Write $X=\operatorname{Spec}A$ for the functor $B\mapsto\Hom_{k\text{-alg}}(A,B)$ on internal commutative $k$-algebras. In this convention an equation between charts can be checked in a generic algebra. The finitely presented charts have finite syntax even when an equality between numerical values is undecidable.

\begin{proposition}[The forward jet chart]\label{prop:jet-chart}
Write $B_n=B[t]/(t^{n+1})$ and $J_nX(B)=X(B_n)$. If $A=k[x_1,\dots,x_d]/(F_1,\dots,F_r)$, the coordinate algebra of $J_nX$ is the polynomial algebra on $x_{i,j}$ for $0\le j\le n$, modulo the coefficients through degree $n$ in the $F_a(\widetilde x_1,\ldots,\widetilde x_d)$, where $\widetilde x_i=\sum_{j=0}^n x_{i,j}t^j$. This presentation represents the construction independently of generators.
\end{proposition}
\begin{proof}
A point of the presented chart in $B$ gives one $B_n$-element per generator. Its coefficient equations say exactly that these elements kill each relation, so they induce an algebra map $A\to B_n$. Extracting coefficients on generators reverses this assignment. Both constructions commute with maps out of $B$. The represented functor therefore identifies charts obtained from different presentations. Reduction of coefficients induces $J_nX\to J_mX$ for $m\le n$ and defines the base point map to $J_0X=X$. On an affine register a query acts by the degree-grouping construction of \cref{thm:marked-jets}.
\end{proof}

Two nearby evaluations also define an algebra independently of any chosen arc. Let $I_\Delta$ be the kernel of multiplication $A\otimes_k A\to A$. The algebra of principal parts of order $n$ is
\[
       P^n_{A/k}=(A\otimes_k A)/I_\Delta^{\,n+1}.
\]
The first factor records the base evaluation and the second the varying one \cite{stacks:parts}. At $x:A\to B$, set $A_{B}=B\otimes_k A$ and $I_x=\ker(A_B\to B)$, where $b\otimes a\mapsto bx(a)$. Base change along the first structure map gives $A_B/I_x^{n+1}$: the images of $1\otimes a-a\otimes1$ generate the augmentation ideal after base change. Its augmentation splits by the constants. Denote the remaining $B$-module by
\[
             \mathcal D_x^{n}(X)=I_x/I_x^{\,n+1}.
\]
It describes a scalar observable through order $n$ with its constant removed. At order zero this module is zero.
\begin{theorem}[Evaluation of backward jets on forward arcs]\label{thm:jet-pairing}
A chart map $f:X\to Y$ sends arcs forward and induces
$f_x^{[n]}:\mathcal D_{f(x)}^n(Y)\to\mathcal D_x^n(X)$ by substitution. These maps compose contravariantly, commute with truncation and are natural under change of test algebra. For an arc $\gamma$ based at $x$ and a pointed observable $h$ at $f(x)$,
\[
       \langle f_x^{[n]}h,\gamma\rangle
        =\langle h,J_n f(\gamma)\rangle\quad\text{in }t B_n .
\]
More generally let $\mathcal A_{B,n}$ be the category of augmented $B$-algebras $C$ with $J(C)=\ker(C\to B)$ and $J(C)^{n+1}=0$. For $C$ in this category let $N_x(C)$ consist of the $C$-points of $X$ reducing to $x$, namely the $k$-algebra maps $A\to C$ over $x:A\to B$. Evaluation induces a $B$-module isomorphism
\[
  \mathcal D_x^n(X)\ \cong\ \operatorname{Nat}(N_x,J).
\]
At $n=1$ the backward module is $B\otimes_A\Omega^1_{A/k}$.
\end{theorem}
\begin{proof}
Write \(C_x=(B\otimes_k A)/I_x^{n+1}\). For an augmented \(B\)-algebra \(C\) with \(J(C)^{n+1}=0\), the maps
\[
 N_x(C)\ \underset{u\mapsto(a\mapsto u(1\otimes a))}{
                  \overset{y\mapsto(b\otimes a\mapsto by(a))}{\rightleftarrows}}\
 \Hom_{\mathcal A_{B,n}}(C_x,C)
\]
are well defined: \(I_x\) maps into \(J(C)\), whence \(I_x^{n+1}\) maps to zero. They are inverse by \(B\)-linearity and natural in \(C\). Evaluating a natural operation at \(1_{C_x}\) now gives
\[
 \operatorname{Nat}(N_x,J)\longrightarrow J(C_x),\quad
 \eta\longmapsto\eta_{C_x}(1_{C_x}),\qquad
 h\longmapsto(u\longmapsto u(h))\ \text{in reverse}.
\]
The inverse identities follow by naturality along \(u:C_x\to C\) and evaluation at the identity. This also realizes the internal operation object by \(J(C_x)\). A chart map induces an augmented algebra map \(C_{f(x)}\to C_x\). The arc pairing is composition
\(C_{f(x)}\to C_x\to B_n\) restricted to augmentation ideals, giving the stated equation, functoriality and truncation. Under scalar change \(B\to B'\), the presentation by the increments \(1\otimes a-x(a)\) and their products gives \(B'\otimes_B C_x=C_{x'}\) canonically up to isomorphism, and the split augmentation transports \(J(C_x)\). At order one the class of \(1\otimes a-x(a)\) satisfies the derivation rules. Every derivation at \(x\) extends \(B\)-linearly to \(B\otimes_k A\) and kills \(I_x^2\). The two constructions agree on generators in each direction, giving \(I_x/I_x^2\cong B\otimes_A\Omega^1_{A/k}\) by \cref{poly:prop:universal}.
\end{proof}
On a polynomial chart the increments \(d_i=1\otimes x_i-x(x_i)\) form polynomial coordinates over \(B\), so the degree-\(j\) graded piece is \(\Sym_B^j(B^d)\) for \(1\le j\le n\). On a chart with relations include \(F_a(x+d)\) in the quotient. Two computations at the rational origin exhibit the effect:
\[
\begin{array}{c|c|c}
 A &\mathcal D^2_0(\Spec A)&\text{evaluations on }\Q[t]/(t^3)\\ \hline
 \Q[x]/(x^2)&\Q\,x& x\mapsto b t^2\\
 \Q[x,y]/(x^2,y^2)&\Q\,x\oplus\Q\,y\oplus\Q\,xy
       &x\mapsto b t^2,\ y\mapsto c t^2 .
\end{array}
\]
Indeed \((at+bt^2)^2=a^2t^2=0\) forces \(a=0\). Every ordinary arc in the second row kills \(xy\). Evaluation at \(1_{C_0}\), the identity on the universal augmented test algebra \(C_0=A\) in this example, detects \(xy\). On a free affine register over a \(\Q\)-algebra, write \(h=\sum_{j=1}^nh_j\) in homogeneous components. Then \([t^j]h(tu)=h_j(u)\). Grid interpolation at \(u\in\{0,\ldots,n\}^d\) recovers every coefficient. This proves separation in the free case.

On the free register with shared substitution $(p,q)=(\theta,\theta)$, pulling back $pq$ and evaluating on $\theta+ut$ gives the same coefficient $u^2$ as applying the forward query first. In general the full filtered observable transports under a coordinate change; its list of homogeneous pieces uses the chosen affine coordinates.

\section{Localizations, \texorpdfstring{\'etale}{etale} lifts and descent}\label{sec:site}
A finite normalization or implicit equation can require an inverse. To identify where it is available, we use relative affines in the symmetric monoidal category of internal $k$-modules, with tensor product over $k$. Their coordinate objects are internal commutative $k$-algebras. We use the finitely presented affine charts and finite covering families below as a site. One can take presented codes and identify their represented charts. All tests, tensor products and covering conditions are interpreted in the current shape context. This gives an ordinary relative affine geometry in the monoidal approach of \cite{relative:tvq}; the homotopical setting of \cite{relative:hag} uses further structure. The algebraic site used here records where equations and infinitesimal lifts glue. The observation opens of \cref{sec:eff} record terminating numerical evidence.

\subsection{The covering families and the equalizer}
For $a\in A$, the principal open $D(a)$ is represented by $A[z]/(az-1)$. A map to a test algebra factors through it exactly when the image of $a$ has an inverse, which is unique. A finitely presented map $A\to C$ is \'etale when every $A$-algebra map $C\to D/J$ has a unique lift to $D$ for $J^2=0$ \cite{etale}. A finite family $\{A\to C_i\}$ in this class covers when its scalar extensions on modules are exact and jointly reflect isomorphisms. Equivalently $C=\prod_i C_i$ is faithfully flat over $A$, since finite products of modules are finite direct sums. These conditions concern the family. An individual chart can be \'etale over a proper open.

\begin{proposition}[Internal affine descent]\label{prop:affine-descent}
These families form a pretopology. For each cover, the diagrams for functions and, more generally, for any $A$-module $M$, are equalizers:
\[
 M\longrightarrow C\otimes_A M
       \ \substack{\longrightarrow\\[-4pt]\longrightarrow}\ 
                 C\otimes_A C\otimes_A M .
\]
The arrows insert a unit in the first and second tensor positions. Represented affine functors are sheaves on this site.
\end{proposition}
\begin{proof}
Identities satisfy both tests. In a pushout $A\to A'$, extension along $A'\to A'\otimes_A C_i$ on an $A'$-module has underlying module $C_i\otimes_A M$. Exactness and reflection thus persist. The lifting property follows from the pushout universal property, and the finite generators and relations base change. For successive families, lift the intermediate map first and then the upper map. Composite presentations are finite; exactness composes and the two successive tests for isomorphisms reflect the original one. This proves the pretopology laws.

Write \(i(m)=1\otimes m\) and \(\delta(c\otimes m)=1\otimes c\otimes m-c\otimes1\otimes m\). On tensoring with \(C\), set
\[
 s_0(a\otimes b\otimes m)=ab\otimes m,\qquad
 s_1(a\otimes b\otimes c\otimes m)=ab\otimes c\otimes m .
\]
On elementary tensors \(s_0(1\otimes i)=1\) and
\(s_1(1\otimes\delta)=1-(1\otimes i)s_0\). Thus \(1\otimes i\) identifies its source with \(\ker(1\otimes\delta)\). Flatness identifies this kernel with \(C\otimes_A\ker\delta\). Conservativity of \(C\otimes_A-\) reflects the comparison \(M\to\ker\delta\) as an isomorphism. For \(M=A\) this is an equalizer of algebras, and applying \(\Hom(D,-)\) gives the matching-family bijection for every represented affine \(D\), also in a slice.
\end{proof}
Finitely many principal opens give covers if $\sum_i a_i b_i=1$ in $A$. Localization is exact by clearing denominators in a finite relation. If an element vanishes after every localization, some power $a_i^{m_i}$ kills it at each key. A power of $\sum_i a_i b_i=1$ of degree greater than $\sum_i m_i$ expands into terms each divisible by one of these powers. Hence the element is zero. Applying this to kernels and cokernels proves joint conservativity. Conversely, if the localizations are jointly conservative, apply them to $A/(a_i)_i$; each gives zero, so the ideal is the unit ideal. For the empty family the statement means $A=0$. An inverse lifts across a square-zero ideal: if $xy=1-e$ with $e$ in the ideal, $y(1+e)$ is an inverse of $x$. This proves the lifting property for the principal opens.

\subsection{Transport of a chosen lift}
\begin{theorem}[\'Etale lift of an arc]\label{thm:etale-jets}
For an \'etale map $Y\to X$ and $n\ge0$, evaluation at the base point and projection induce an isomorphism
\[
             J_nY \ \cong\ Y\times_X J_nX .
\]
It is natural in the test algebra and agrees with every truncation.
\end{theorem}
\begin{proof}
The claimed natural isomorphism says precisely that the left square below is cartesian. For any test algebra \(B\) and \(B_j=B[t]/(t^{j+1})\), a lifting step is the right square of affine test objects:
\[
\begin{tikzcd}[column sep=small]
 J_nY\arrow[r,"J_nf"]\arrow[d,"\ev_0"']&J_nX\arrow[d,"\ev_0"]\\
 Y\arrow[r,"f"']&X
\end{tikzcd}
\qquad
\begin{tikzcd}[column sep=large]
 \Spec B_{j-1}\arrow[r,"y_{j-1}"]\arrow[d,hook]&Y\arrow[d,"f"]\\
 \Spec B_j\arrow[r,"x_j"']\arrow[ur,dashed,"{y_j}" description]&X .
\end{tikzcd}
\]
The diagonal is unique for each \(j\ge1\). Iteration supplies the inverse to the comparison with the fibre product.
In $B$-points, start with the chosen point in $Y(B)$ and the compatible arc in $X(B_n)$. For $j=1,\ldots,n$, the kernel of $B_j\to B_{j-1}$ is generated by $t^j$ and has square zero. Lift the previous $Y$-point uniquely over the prescribed $X$-point in $B_j$. The final lift is inverse to the displayed map. Truncating the construction or changing $B$ gives the same lift by uniqueness at each stage. For presented charts \cref{prop:jet-chart} represents these maps, so the pointwise inverse is an inverse of charts.
\end{proof}
\begin{example}[A double cover for a balanced table]\label{ex:root-cover}\leavevmode
\begin{enumerate}[titledsteps]
    \item \textbf{Construct the cover.} Let $A=k[a^{\pm1},b^{\pm1},c^{\pm1},d^{\pm1}]$ and $z=ad/(bc)$. For a two-by-two table with these input weights and unit row and column masses, the balanced rows are $(p,1-p)$ and $(1-p,p)$. Scaling rows and columns preserves $z$. Hence the odds $u=p/(1-p)$ obey $u^2=z$. Put $C=A[u]/(u^2-z)$. The basis $1,u$ makes extension faithfully flat and $2u$ is a unit. For a provisional lift $v$ across $J^2=0$, the unique correction is $-(v^2-z)/(2v)$, using the given image of $z$. Thus $\Spec C\to\Spec A$ is a cover. On $D(1+u)$ one recovers $p=u/(1+u)$. Indeed the resulting normalized table has cross ratio $z$; its row and column scales are reconstructed by matching the first row and the first column, and the cross-ratio equation forces the fourth entry. On positive located weights, interval refinement and rational square-root approximation select $u=\sqrt z>0$ and give the required inverses.
  

    \item \textbf{Test for a section and lift the coefficients.} There is no algebraic section over \(k=\Q\): specializing at \((2,1,1,1)\) would give a rational square root of two. The specialization \((-1,1,1,1)\) also obstructs an algebraic section after scalar extension to the located real ring, since a real square cannot equal \(-1\). This specialization tests the algebraic torus. The positive numerical locus has its selected root. Under the deck involution \(u\mapsto-u\), the logarithmic differential descends because \(2u\,du=dz\), hence \(du/u=dz/(2z)\). Write \(z(t)/z_0=1+\sum_{m\ge1}c_mt^m\) and \(u(t)=u_0(1+\sum_{m\ge1}v_mt^m)\). Then
\[
 v_m=\tfrac12\!\left(c_m-\sum_{i=1}^{m-1}v_i v_{m-i}\right).
\]
The relative coefficients depend only on the base arc, while the constant \(u_0\) specifies the sheet. The many-key lift in \cref{sec:balance} uses the same square-zero correction principle.
\end{enumerate}
\end{example}

\section{Functorial feedback through parametric lenses}\label{sec:optic-feedback}
The underlying forward map of a differentiable program already has its interpretation in the category of keyed programs. Sending each gate to its cotangent lens gives a second interpretation with the same forward part. The chain rule identifies the backward part of the composite, including the contributions from parameters shared between gates. At higher order the corresponding interface retains the joint class in the pointed jet algebra.

\subsection{Cotangent optics and compiled circuits}
Take the bimorphic lens of \cref{sec:intro} on pairs \((V,V^*)\). Restrict to backward maps linear in the incoming covector at each retained forward value. Identities, tensor and lens composition preserve that linearity. Write \(\lens\) for this category of cotangent optics. For each \(x\), the forward and backward paths for \(g f\) are displayed separately, with \(y=f(x)\) and \(z=g(y)\):
\[
\begin{tikzcd}[column sep=large]
 x\in V\arrow[r,mapsto,"f"]&y\in W\arrow[r,mapsto,"g"]&z\in U\\
 V^*& W^*\arrow[l,"f^\sharp_x"]&U^*\arrow[l,"g^\sharp_y"]
\end{tikzcd}
\]
The lower composite is evaluated at the retained points on the upper path. The formal cotangent functor selects a particular lens over each polynomial map, by the universal differential of \cref{neural:lem:kahler}.
\paragraph{Serial wiring of parametric lenses.}
In the diagram the upper horizontal wire is the value and the lower wire is the returned feedback. Each box is one lens. Its vertical wires are the parameter and the feedback to that parameter.
\begin{center}
\begin{tikzpicture}[x=1cm,y=1cm,>=stealth,every node/.style={font=\small}]
 \draw[rounded corners=2pt] (-2,-.5) rectangle (-.9,.5);
 \draw[rounded corners=2pt] (.9,-.5) rectangle (2,.5);
 \node at (-1.45,0) {\(L\)}; \node at (1.45,0) {\(K\)};
 \draw[->] (-3.45,.25) -- node[above]{\(X\)} (-2,.25);
 \draw[->] (-.9,.25) -- node[above]{\(Y\)} (.9,.25);
 \draw[->] (2,.25) -- node[above]{\(Z\)} (3.45,.25);
 \draw[->] (3.45,-.25) -- node[below]{\(Z'\)} (2,-.25);
 \draw[->] (.9,-.25) -- node[below]{\(Y'\)} (-.9,-.25);
 \draw[->] (-2,-.25) -- node[below]{\(X'\)} (-3.45,-.25);
 \draw[->] (-1.72,1.18) node[above]{\(P\)} -- (-1.72,.5);
 \draw[->] (-1.17,.5) -- (-1.17,1.18) node[above]{\(P'\)};
 \draw[->] (1.17,1.18) node[above]{\(Q\)} -- (1.17,.5);
 \draw[->] (1.72,.5) -- (1.72,1.18) node[above]{\(Q'\)};
\end{tikzpicture}
\end{center}
The retained \(Y\)-value determines how the second box returns its \(Y'\)-feedback. Reading the lower wire right to left gives the backward composite in \cref{sec:intro}. Drawing the two boxes on disjoint pairs of wires gives their tensor instead. Tying their parameters inserts the diagonal on the incoming wires. Its cotangent lens adds the outgoing receipts in \cref{neural:cor:sharing}.
\begin{theorem}[Differentiation is a compositional optic]\label{neural:thm:cotangent}
The assignment
\[
 \mathsf D^*f=\bigl(f,(x,\lambda)\mapsto df_x^*\lambda\bigr)
\]
is a functor from polynomial numerical maps to cotangent optics. Independent products give a strong symmetric monoidal structure, using the canonical isomorphism $(V\oplus W)^*\cong V^*\oplus W^*$.
\end{theorem}
\begin{proof}
\Cref{poly:thm:differential-comparison,neural:lem:kahler} gives
$d(gf)_x^*=df_x^*\,dg_{f(x)}^*$ and $d\id_x^*=\id$. These are exactly the optic composition and unit equations. The differential of an independent product is the direct sum of its component differentials, and dualization carries it to the direct sum of the cotangent maps. Canonical reassociation and symmetry of finite direct sums prove coherence. Copying has reverse map $(\alpha,\beta)\mapsto\alpha+\beta$, while discarding has reverse map $0$. This follows directly by dualizing the diagonal and terminal linear maps. Arbitrary paired circuits therefore accumulate sensitivities correctly even when they share inputs.
\end{proof}
\begin{corollary}[Shared parameters and cotangent accumulation]\label{neural:cor:sharing}
For $f:P\times X\to Y$ and $g:Q\times Y\to Z$, an output covector $\lambda$ pulls back to
\[
 \bigl((d_Pf)^*(d_Yg)^*\lambda,\ (d_Qg)^*\lambda,\
       (d_Xf)^*(d_Yg)^*\lambda\bigr),
\]
with derivatives evaluated at the forward values. Replacing $P\times P$ by a shared parameter through $\Delta_P$ adds the two parameter covectors. A finite recurrent unrolling therefore returns the sum of the cotangents contributed by every occurrence of a tied parameter.
\end{corollary}
\begin{proof}
Apply \cref{poly:thm:differential-comparison} to the
composition of \cref{poly:con:para}. Its diagonal comparison gives
$d\Delta_P^*(\alpha,\beta)=\alpha+\beta$. The shared evaluation
of \cref{poly:thm:parametric} and induction on its finite unrolling
give the last assertion. The same retained forward parameter is used
at every occurrence. Updating between occurrences defines a different
stateful computation.
\end{proof}

\subsection{Solving the adjoint equations of a circuit}

The optic composition law describes the backward pass through two successive gates. A register with several consumers must collect several such pullbacks. To see how their local equations determine the derivative of the whole circuit, execute a finite acyclic code in its given topological order and retain the forward values. Let $V_i$ be the activity module at node $i$, and let $A_{ij}:V_j\to V_i$ be the partial differential from parent $j$ at those values. Count repeated parent wires separately and sum their contributions. Write $B_i:P\to V_i$ for the partial differential with respect to the global parameter register.

\begin{theorem}[Finite local reverse propagation]\label{neural:thm:nilpotent}
On $D=\bigoplus_iV_i^*$ define
\[
 (N\alpha)_j=\sum_{i:j\to i}A_{ij}^*\alpha_i.
\]
Let $b$ contain the loss cotangents directly incident at each node. Write $k$ for the number of circuit nodes and take $d=k+1$ (or any positive bound exceeding all path lengths). Then $N^d=0$ and
\[
 \alpha=b+N\alpha,\qquad
 \alpha=(1+N+\cdots+N^{d-1})b
\]
has a unique solution. The synchronous iteration $\alpha_0=0$, $\alpha_{n+1}=b+N\alpha_n$ reaches it in at most $d$ rounds. The parameter cotangent is $\sum_i B_i^*\alpha_i$.
\end{theorem}
\begin{proof}
For \(m\ge1\) the block of \(N^m\) from output seed \(i_m\) to node \(i_0\) is
\[
 (N^m)_{i_0i_m}
    =\sum_{i_0\to i_1\to\cdots\to i_m}
             A_{i_1 i_0}^*\cdots A_{i_m i_{m-1}}^* .
\]
One multiplication by \(N\) appends one edge to the path formula. Since the DAG has no path of length \(d\), \(N^d=0\) and
\[
 (1-N)\sum_{m<d}N^m=1-N^d=1,\qquad
 \alpha_n=\sum_{m<n}N^m b,\qquad
 \bar p=\sum_iB_i^*\sum_{m<d}(N^m b)_i .
\]
The commuting reverse product gives the other inverse identity for \(1-N\). The chain rule identifies the path sum with the local input receipts. Transposing the shared parameter lookup adds all receipts at that parameter (\cref{neural:cor:sharing}). In particular two sockets using one parent contribute twice before grouping.
\end{proof}
\paragraph{Which occurrence tables come from an owner table?}
Let \(P_0=\operatorname{im}(\rho:D\to P)\) and \(K=D\times_P D\), with projections \(k_1,k_2\). Then
\[
 R^{P_0}\ \cong\ \operatorname{Eq}\bigl(R^D
             \overset{k_1^*}{\underset{k_2^*}{\rightrightarrows}} R^K\bigr).
\]
Indeed an occurrence table \(w\) in the equalizer has \(w_d=w_{d'}\) whenever \(\rho d=\rho d'\). The relation \(\{(\rho d,w_d):d\in D\}\) is total on \(P_0\) and single-valued there, so it is the graph of the unique descended table. Conversely lookup of any such table equalizes the two restrictions. The same graph argument over a test object proves the isomorphism for families. For finite presented owners, \(R^P\cong R^{P_0}\times R^{P\setminus P_0}\); forward evaluation ignores the second factor, and the reverse sum over an unused owner is empty. Thus the full table classified in \cref{prop:relative-models} and the observable tied occurrence table can differ.

Parameter sharing is a reparameterization of the occurrence register. For a key \(\rho:D\to P\), the free-module map \(F_R(\rho):F_R(D)\to F_R(P)\) sends each basis generator to its owner. The field register \(R^P=\Hom_R(F_R(P),R)\) is its dual presentation. Forward lookup \(r=\rho^*\) is precomposition with \(F_R(\rho)\). The reverse lookup is determined by the natural bidual isomorphism on the finite free modules.
\[
\begin{tikzcd}[column sep=large]
 F_R(D)\arrow[r,"F_R(\rho)"]\arrow[d,"\simeq"']&
 F_R(P)\arrow[d,"\simeq"]\\
 (R^D)^*\arrow[r,"r^*"']&(R^P)^* .
\end{tikzcd}
\]
On generators both routes are evaluation at the owner. The following formula computes this dual map in the occurrence and owner bases.
\begin{lemma}[Reverse lookup is aggregation]\label{lem:sharing}
With the canonical pairings on these finite free modules, the dual of \(r\) is the additive pushforward \(\rho_!\). In a composite with \(r\) the parameter covector is the pushforward of the occurrence covector.
\end{lemma}
\begin{proof}
For an arbitrary tangent \(u\) and occurrence covector \(g\), the defining pairing is
\[
 \langle g,r u\rangle=\sum_{d\in D}g_d u_{\rho(d)}
   =\sum_{p\in P}\left(\sum_{d:\rho(d)=p}g_d\right)u_p.
\]
This identifies the dual on each basis vector and hence as a linear map. Composition with the suffix pullback is the chain rule. When the same edge is read at several token rows its occurrence carrier is \(J\times D\) and the parameter key is \((j,d)\mapsto\rho(d)\), so both sums are accounted for.
\end{proof}

\subsection{From a keyed program to its parametric lens}

Start with polynomial gates on finite free registers, together with copies, discards, products and compositions. A numerical generator extends this language on its domain by a tracked value and the jet calculation derived for that generator. Write $\mathcal N$ for the category of finite based registers and maps denoted by total programs in a fixed such language. The domain of a partial gate is its value domain, such as the unit open for division. If $f:D_f\to Y$ is followed by a gate with domain $D_g\hookrightarrow Y$, their composite has domain $D_f\times_Y D_g$ over the original inputs. Independent products take products of domains. This constructs a domain from the syntax without a test deciding totality of arbitrary code. Identities, composites and independent products of total programs are total. We construct its first-order compiler before instantiating the extension at normalized lookup.

\begin{construction}[Syntax-directed cotangent compiler]\label{con:backward-compiler}\leavevmode
\begin{enumerate}[titledsteps]
    \item \textbf{Compile the gate pullbacks.} For a polynomial gate form one receipt per marked occurrence as in \cref{poly:prop:marked}, multiply by the output seed, and group by the source key. Pairing with an input increment identifies the result with the cotangent pullback. For a scalar gate with a tracked Taylor program multiply its incoming receipt by the first coefficient for a unit perturbation. Copy sends $u$ to $(u,u)$ and pulls $(\alpha,\beta)$ back to $\alpha+\beta$; discard pulls back to zero. A bilinear pairing sends $(\dot q,\dot k)$ to $\dot q(k)+q(\dot k)$ and uses the corresponding two evaluation-dual maps in reverse. 
    \item \textbf{Wire the reverse instructions.} At an independent product take the direct sum of backward maps. If $f$ is followed by $g$, retain $f(x)$ and define the reverse program recursively by
\[
 (gf)^\sharp(x,\lambda)
    =f^\sharp\bigl(x,g^\sharp(f(x),\lambda)\bigr).
\]

    \item \textbf{Verify the compiled map.} Finite syntax recursion constructs this compiler and its forward-value registers uniformly. Induction on the syntax, using the product and chain rules, proves that each compiled reverse map is $Df_x^*\lambda$. In particular a \emph{tracked pullback} is an output of the compiler, not an additional primitive property of a keyed program.
\end{enumerate}
\end{construction}

\begin{theorem}[Parametric cotangent comparison]\label{thm:para-lens}
For a finite keyed execution $c$ in the numerical language above, with interpretation $f:P\times X\to Y$, its syntax-compiled backward program has the parametric lens interpretation
\[
 f^\sharp(p,x,\lambda)=\big((D_p f)^*\lambda,(D_x f)^*\lambda\big).
\]
The assignment respects serial composition and independent products. A differentiable change of parameter acts on the interpreted covector by cotangent pullback. A change of parameter given in the numerical language has an executable pullback obtained by compiling its code. On $\mathcal N$ this construction gives the pseudofunctor on parameter bicategories induced by the strong monoidal cotangent functor.
\end{theorem}
\begin{proof}\leavevmode
\begin{enumerate}[titledsteps]
    \item \textbf{Compare generators and composites.} For a polynomial generator, the claim follows from \cref{poly:thm:differential-comparison}. For an analytic generator, the compiler of \cref{con:backward-compiler} computes its pullback by differentiating the implemented primitive. The product and composition clauses of the compiler then give the claim by induction on the whole program. For a subsequent $g:Q\times Y\to Z$, put $y=f(p,x)$ and $\beta=(D_y g)^*\lambda$. Pairing with independent variations of $(p,q,x)$ shows that the backward program for the composite returns
\[
 ((D_pf)^*\beta,\ (D_qg)^*\lambda,\ (D_xf)^*\beta).
\]
Lens composition returns the same triple of covectors. For independent products the differential is block diagonal. For a differentiable reparameterization $r:P'\to P$, the chain rule at $(p',x)$ gives
\[
  (f\circ(r\times\id))^\sharp(p',x,\lambda)
   =\bigl((Dr_{p'})^*\alpha,\beta\bigr),\qquad
  (\alpha,\beta)=f^\sharp(r(p'),x,\lambda).
\]
When $r$ is a numerical program, run its compiled backward map on $(p',\alpha)$ to compute the first component. This also proves that on total maps the result is independent of the program denoting $r$, by uniqueness of the differential. The identity, chain and product rules give the identity and the two compositions of reparameterizations. Reassociation of product registers supplies the pseudofunctorial coherence.


    \item \textbf{Locate the effective boundary.} The gate syntax matters for effectivity. Already at the external level, no index transformation extracts derivatives from value programs even when each individual denotation is a rational function analytic on the whole real line. Uniformly in $e$, at requested precision $2^{-N}$ simulate the diagonal computation through step $N$. If it has halted at step $t$, evaluate $\delta\varphi(x/\delta)$ to that precision, where $\delta=2^{-t}$ and $\varphi(v)=v/(1+v^2)$. Otherwise return zero. A later halt changes the output by at most $2^{-N-1}$ since $|\varphi|\le1/2$. At the origin the derivative of the eventual function is one in the halting case and zero in the nonhalting case. A uniform value-code-to-derivative compiler would decide halting. The compiler above uses the actual expression for the gate.


    \item \textbf{Accumulate at shared registers.} For a program with shared registers, apply induction in reverse execution order, using cotangent addition at any register with several consumers (\cref{neural:thm:nilpotent}). This gives the comparison square below. Its upper route differentiates the interpreted program. Its lower route interprets the backward program compiled from the keyed occurrences.
\[
\begin{tikzcd}[column sep=normal,row sep=normal]
 \text{keyed programs}\arrow[r,"\mathrm{execute}"]\arrow[d,"\mathrm{mark/pull\ back}"']&
 \para(\mathcal N)\arrow[d,"{\para(\mathsf D^*)}"]\\
 \text{forward/backward programs}\arrow[r,"\mathrm{interpret}"']&
 \para(\lens).
\end{tikzcd}
\]
Here the lower left denotes programs with their explicit backward pass. The right hand interpretation forgets which rows implemented equal numerical maps. Parametric lenses for gradient-based learning are developed in \cite{lens:deep}.\qedhere
\end{enumerate}
\end{proof}
\begin{comparison}[Differentiable descent at an interpreted port]\label{comp:descent}
A differentiably realized factorization through a port decoder lifts to a factorization of cotangent lenses. Every returned port covector annihilates the kernel of the decoder differential at the retained point.
\end{comparison}
\begin{proof}
Take a factorization \(C=\bar C(1_X\times d)\) from \cref{prop:response-quotient} in the total differentiable program language, with \(d:Z\to D\) and \(\bar C:X\times D\to O\). Put \(\mathsf L V=(V,V^*)\). The same triangle then commutes in cotangent optics:
\[
\begin{tikzcd}[column sep=small]
 \mathsf L X\otimes\mathsf L Z\arrow[rr,"{1\otimes\mathsf D^*d}"]\arrow[dr,"\mathsf D^*C"']&&
 \mathsf L X\otimes\mathsf L D\arrow[dl,"{\mathsf D^*\bar C}"]\\
 &\mathsf L O .
\end{tikzcd}
\]
Indeed \(\mathsf D^*\) preserves composition and its tensor comparison identifies the two input pairs. Thus if \((\alpha,\gamma)=\bar C^\sharp(x,d(z),\lambda)\), the pullback through \(C\) is \((\alpha,d_z^\sharp\gamma)\). In particular every \(u\in\ker Dd_z\) is annihilated by the port component of every returned covector, by evaluation on \(Dd_z u=0\).
\end{proof}
This conclusion uses the differentiable factorization itself. The logical quotients in \cref{sec:hidden-concepts} require only a predicate or a map of sets. The differential argument applies when their maps are realized in the differentiable language.

\subsection{Joint feedback for parameterized maps}\label{sec:joint-feedback}
For an algebraic register $X$ and a test algebra $B$, put
$\mathcal C_X^{(n)}(\mathsf y)=\sum_{x\in X(B)}\mathsf y^{\mathcal D_x^n(X)}$,
interpreted on a test algebra in its current context. A position is a retained value; a direction is a pointed scalar observable truncated at that value. A chart map $f:X\to Y$ acts forward on positions and backwards by $f_x^{[n]}$ on directions. For based numerical affine registers this same interface is computed by truncated polynomials in the coordinate increments. For the numerical language used below, run the syntax on the generic increment in $\R[h_1,\ldots,h_d]/(h)^{n+1}$, where $(h)$ is the ideal of all increments. Addition and multiplication are truncated arithmetic. For a unit $a$ and nilpotent $v$ use $(a+v)^{-1}=a^{-1}\sum_{j=0}^n(-v/a)^j$; for the exponential use $\exp(a)\sum_{j=0}^n v^j/j!$. The positive square root lifts one homogeneous degree at a time by division by twice its positive constant value. The base values use the procedures in \cref{sec:sto}. Compositions substitute centred polynomials, and a product evaluates both coordinates. Induction gives truncation and substitution compatibility and the Taylor remainder of each resulting program from the primitive Taylor formulas. Balanced and equilibrium gates obtain their coefficients in \cref{sec:balance,sec:equilibrium,app:limits}. Thus the syntax computes its Taylor table. Extension to a new primitive uses the orders constructed for that primitive. This permits the numerical gates of \cref{sec:sto} to be interpreted on their effective domains.

\begin{theorem}[Composition of joint parameter jets]\label{thm:joint-jet}
Let $f:P\times X\to Y$ and $g:Q\times Y\to Z$ be interpreted by chart maps or by such Taylor programs. Above a retained point $(q,p,x)$ the reverse of the composite is
\[
  \mathcal D_{g(q,y)}^n(Z)
  \xrightarrow{\ g^{[n]}_{(q,y)}\ }
  \mathcal D_{(q,y)}^n(Q\times Y)
  \xrightarrow{\ (\id_Q\times f)^{[n]}_{(q,p,x)}\ }
  \mathcal D_{(q,p,x)}^n(Q\times P\times X),\qquad y=f(p,x).
\]
It commutes with parameter reindexing, product wiring and truncation. At order one its parameter and input components give \cref{thm:para-lens}.
\end{theorem}
\begin{proof}
Write \(\Delta f=f(p+\delta p,x+\delta x)-f(p,x)\) in the order-\(n\) algebra, and similarly for \(g\). For a centred output polynomial \(H\), the composite acts by
\[
 H(\delta z)\longmapsto H(\Delta g(\delta q,\delta y))
 \longmapsto H(\Delta g(\delta q,\Delta f(\delta p,\delta x))).
\]
Every substituted increment belongs to the augmentation ideal \(J\). A degree-\(m\) product therefore lands in \(J^m\). This proves well-definedness modulo \(J^{n+1}\) and compatibility with truncation. For a chart the same computation descends along its relations by \cref{thm:jet-pairing}. Associativity of substitution gives triple composition. Substitution in both variable blocks gives the product law, including mixed monomials. Substitution for \(\delta p\) gives reparameterization. In degree one the mixed products vanish and the two summands are the usual parameter and input cotangents. For smooth maps equality on the numerical open implies equality of Taylor coefficients there. On a constrained chart the induced equality is tested in its pointed quotient.
\end{proof}
For a chart map \(f:X\to Y\), a retained point \(x\), and \(y=f(x)\), the order filtration gives a commutative square of feedback maps:
\[
\begin{tikzcd}[column sep=large]
 \mathcal D_y^{n+1}(Y)\arrow[r,"{f_x^{[n+1]}}"]\arrow[d,"\mathrm{tr}"']&
 \mathcal D_x^{n+1}(X)\arrow[d,"\mathrm{tr}"]\\
 \mathcal D_y^n(Y)\arrow[r,"{f_x^{[n]}}"']&\mathcal D_x^n(X).
\end{tikzcd}
\]
The pullback takes the pointed ideal into the pointed ideal, hence its \(j\)-th power into the \(j\)-th power for every \(j\). This induces all four maps and proves commutation on representatives. At first order the ideal of a product modulo its square splits into a parameter summand and an input summand. At higher order mixed classes survive. The joint feedback on the product is needed before a diagonal reparameterization.
If the parameter lookup is $\rho:D\to K$, its jet pullback replaces the increment at occurrence $d$ by the increment at $\rho(d)$ in \emph{every} monomial and groups the resulting coefficients. At first order this is $\rho_!$. At order two it preserves pairs of occurrences, including pairs that subsequently acquire the same key. The joint observable $pq$ at $(0,0)$ restricts to zero on both axes. Under $p=q=\theta$ it pulls back to $\theta^2$; its mixed variation in directions $u,v$ is $2uv$. Separate axis feedback would lose this term before it reached the shared owner.

For an affine numerical register over $\R$, write a second-order loss as $\lambda(\delta)+H(\delta,\delta)/2$. Filling two distinct marks gives
\[
 \lambda'(u)=\lambda(Df\,u),\qquad
 H'(u,v)=H(Df\,u,Df\,v)+\lambda(D^2f(u,v)).
\]
In the first term the marks enter two downstream positions. In the second they enter one downstream position and split in its argument. The identity follows by extracting the mixed coefficient in \cref{thm:marked-jets}; fixing $u$ gives a program for $H'(u,-)$ without returning a table with every pair of free keys. The full pulled-back jet is independent of coordinates. The displayed splitting into its linear and quadratic parts uses the affine register.

\subsection{How a defining equation becomes a bidirectional instruction}
A relation $R\hookrightarrow U\times Y$ specifies an equation for a module call. In the current slice define
\[
 D_R=\{x:U\mid\exists y\,[R(x,y)\wedge\forall y'(R(x,y')\Rightarrow y'=y)]\}.
\]
This constructs the input type on which the unrestricted equation itself determines an output. Fixing a sheet or branch changes the relation whose graph is read back.
\begin{proposition}[Composition of solved relations]\label{thm:solved-composition}
The relation restricted to $D_R$ is the graph of a canonical map $s_R:D_R\to Y$. Every base change on which the pulled-back relation is the graph of a total map factors uniquely through $D_R$. These calls commute with reindexing, parallel wiring and serial wiring.
\end{proposition}
\begin{proof}\leavevmode
\begin{enumerate}[titledsteps]
    \item \textbf{Invert the graph projection.} Let $G=R\times_U D_R$ and $p:G\to D_R$ be projection. Existence makes $p$ regular epic and uniqueness makes its kernel pair the diagonal. As the coequalizer of this pair $p$ is invertible. Explicitly $1_G$ coequalizes the pair and descends to $j:D_R\to G$ with $jp=1_G$; cancelling $p$ gives $pj=1$. Set $s_R=\pi_Yj$. This proves the unique-choice rule for this graph, in the sense of \cite{streicher-logic}. 
    \item \textbf{Identify the universal domain.} A base on which the relation is a total graph satisfies the defining formula for $D_R$, giving the asserted factorization. Inverse image between slices preserves the quantifiers in that formula by Beck--Chevalley. Products of the graph isomorphisms give the parallel comparison. 
    \item \textbf{Compose the solution graphs.} For a successive call with domain $D_T\hookrightarrow Y$ and output $t:D_T\to Z$, the execution domain is $D_R\times_YD_T$ using $s_R$. The graph pullback keeps the unique intermediate $s_R(x)$; insertion of this value and forgetting it are inverse between this pullback and the graph of the composite on that domain. These identifications commute with threefold composition because they insert the same successive outputs. Hiding an ambiguous intermediate in a relational composite can enlarge its unique-output domain: the full relation from one point to two points followed by the map back to one point has a unique final output. Serial execution of a solved call retains its intermediate, as the graph pullback specifies. The explicit numerical constructions below exhibit names on the stated input domains. Membership in the maximal semantic domain can itself be undecidable: the one-output relation true when a machine halts has exactly the halting inputs as its domain.\qedhere
\end{enumerate}
\end{proof}

Consequently a solved call in a presented architecture denotes \(\pi_Y(p^{-1}(-))\), with $p:G\to D_R$ the graph projection just inverted. For a node whose input is \(v:T\to U\) and whose execution domain is \(D\hookrightarrow U\), its context of definition is \(T\times_U D\). Iterating these pullbacks gives the domain of the glued interface in \cref{app:field-presentations}.
The serial comparison is a pullback statement about graphs. Write \(\Gamma_s\hookrightarrow D\times Y\) and \(\Gamma_t\hookrightarrow Y\times Z\). For total selected calls the square is cartesian:
\[
\begin{tikzcd}[column sep=large]
 \Gamma_s\times_Y\Gamma_t\arrow[r,"\pi_2"]\arrow[d,"\pi_1"']&
 \Gamma_t\arrow[d,"\mathrm{in}_t"]\\
 \Gamma_s\arrow[r,"\mathrm{out}_s"']&Y .
\end{tikzcd}
\]
Projection to the two external ports identifies its upper left object with \(\Gamma_{ts}\) by the inverse constructed in the proof. For a partial second call the same pullback restricts the allowed first inputs. The graph therefore carries both the domain calculation and the equation that can be read back at the composite port.

\clearpage
\part{Kernels, quotients and constrained solutions}\label{part:solutions}

\section{Stochastic rows and normalized attention}\label{sec:sto}

Which probability law does an attention head use when the input table is fixed? It weights eligible row occurrences within that table. A law on database instances describes a separate source of randomness. To construct the row law, let $\pi:E\to X$ give the query address of an eligible occurrence and $t:E\to J$ its target. For weights $w$ and target values $v$, the query first computes the pair $N=\pi_!(w\,t^*v)$ and $Z=\pi_!w$. The same witnesses contribute to both sums. Division gives the readout $Z^{-1}N$ on the domain carrying an inverse for $Z$. Lower integration will give the corresponding construction on represented row spaces.

\subsection{Valuations and a computable lower integral}
Internally, a row space \(S\) carries its specified topology \(\mathcal O(S)\). We use discrete spaces on finite indices. A continuous valuation is a monotone map \(\mu:\mathcal O(S)\to Q\) with
\[
 \mu(\varnothing)=0,\qquad
 \mu(U)+\mu(V)=\mu(U\cup V)+\mu(U\cap V),\qquad
 \mu(\bigcup_\lambda U_\lambda)=\sup_\lambda\mu(U_\lambda)
\]
for directed families of opens. Write $\V S$ for the valuations of total mass at most one. A valuation kernel assigns such a law to each input, with lower-continuous evaluation on every open.

There is an exact domain for computing masses of named opens. Write $j:Q_\Sigma\hookrightarrow Q$, let $T_R$ be the object of positive row tests denoting opens, and let $i:T_R\to\mathcal O(R)$ take denotations. For the state topology one takes $T_R=\Sigma^R$ and $i(u)=\iota u$. For the semantic valuation object $\mathsf{Val}(R)$ form the pullback
\[
 \mathsf{Val}_{\rm obs}(R)
 =\mathsf{Val}(R)\times_{Q^{T_R}}(Q_\Sigma)^{T_R},
 \qquad \mu\longmapsto\mu i .
\]
Its projection is monic since $j^{T_R}$ is monic. A family factors through it exactly when its evaluation on the row tests factors through $j$, by exponential adjunction. Evaluation of this unique factor is the mass computation on names; different program indices can realize the same factor. Finite sums of named weighted Dirac masses give examples by testing each addressed row and interleaving lower bounds. There is no uniform extraction for every semantic valuation. On one row the semantic subprobability of mass $\tau(\neg\operatorname{Halts}(e,e))$ would give a test of diagonal nonhalting by evaluation on the full row and comparison with $1/2$. For a lower-semicontinuous feature $g:S\to Q$, the relevant tests are the rational threshold opens $U_r=\{s:r<g(s)\}$. On represented inputs, the threshold preimages and their finite intersections are computable on names. This specifies the interface for the following integration procedure on both discrete rows and general represented spaces.

\begin{lemma}[Lower integration, change of rows, and names]\label{lem:eff-integral}
For a valuation \(\mu\) and nonnegative lower-semicontinuous \(g\), define
\[
 I_\mu(g)=\sup_{\ 0=r_0<\cdots<r_n\ \mathrm{rational}}
              \ \sum_{k=1}^n(r_k-r_{k-1})\mu(U_{r_k}).
\]
This is the lower integral of Vickers \cite[Definition 9]{vickers}. It is increasing, preserves directed suprema of integrands, and satisfies \(I_\mu(\tau U)=\mu(U)\) for an open indicator. With \(f_*\mu(V)=\mu(f^{-1}V)\), it satisfies \(I_{f_*\mu}(g)=I_\mu(g\circ f)\). On the named inputs just specified its value has a uniform \(Q_\Sigma\)-name.
\end{lemma}

The partition and naming proofs are in \cref{app:integration}.

For a finite weighted discrete measure, substitution into this formula gives
\(I_{\sum_i a_i\delta_i}(g)=\sum_i a_i g(i)\). Indeed, for each row the scalar partition sum of the indicators of its strict thresholds has supremum \(g(i)\). One common refinement suffices for finitely many row approximations. Moving the finite nonnegative weighted sum through this directed supremum proves the formula.

\begin{lemma}[Linearity of the lower integral]\label{lem:eff-linearity}
The functional \(I_\mu\) preserves addition of nonnegative lower-semicontinuous functions and multiplication by nonnegative rational constants.
\end{lemma}

\subsection{A law on tables and a law on rows}

Fix a query $q$ and an instance $I$ over a source-state object $Z$. We describe a success event at one output key. Allowing the key to vary amounts to working over $Z\times J$. Build the witness object $W_q\to Z$ by recursion in the slice. A scan reads the entity fibre, union takes tagged sums, join takes pullbacks over matching keys, and restriction selects the subobject on which the guard holds. In \eqref{eq:guard-cert}, matching a certificate means existentially quantifying its finitely many row slots in the fibre. Projection retains the witness and reindexes its output. This construction uses the declared key fields and annotations of the instance and applies in every parameter slice.

\paragraph{The event and its semidecidability domain.}
Let $\chi_q:Z\to\Omega$ classify the image of $W_q\to Z$. Write $\iota:\Sigma\hookrightarrow\Omega$ for the inclusion and form the pullback
\begin{equation}\label{eq:query-test-domain}
\begin{tikzcd}[column sep=large]
 Z_q\arrow[r,"u_q"]\arrow[d,hook,"j_q"'] & \Sigma\arrow[d,hook,"\iota"]\\
 Z\arrow[r,"\chi_q"'] & \Omega .
\end{tikzcd}
\end{equation}
For every $h:Y\to Z$, the success event after base change has classifying map $\chi_q h$, by stability of images under pullback. A factorization $\chi_q h=\iota v$ is equivalent to a unique lift $Y\to Z_q$ carrying $u_q$ back to $v$. Thus $Z_q$ is the exact domain of semidecidability of this event.

The occurrence interpreter constructs a map into this domain. For the finite input vocabulary read by $q$, let $c$ range over tuples of the row and weight program indices of \cref{thm:evidence}. Let $A(c)$ be the resulting family of numbered occurrences, with its interpreted key fields and annotations. Form
\[
 \widehat Z_q=\sum_{z:Z}\sum_{c:\mathrm{Code}_q}
     \operatorname{Iso}_q\bigl(A(c),I_z|_q\bigr),
 \qquad
 \widehat Z_q\xrightarrow{e_q}Z_{q,\mathrm{pres}}
            \xhookrightarrow{m_q} Z .
\]
The arrows factor the projection to $z$ through its regular image. The comparison object $\operatorname{Iso}_q$ consists of inverse row maps at each input sort commuting with all fields read by $q$, including any observable address and the annotations. It is a subobject of the product of the corresponding function objects, cut out by the inverse and field equations. It permits two presentations to use different private enumerations. The code and guard recognizers give a success test $\widehat u_q:\widehat Z_q\to\Sigma$ by \cref{thm:evidence}. Transport along the row comparison preserves scans, tagged and paired witnesses, their keys and each fact of a guard certificate. The inverse comparison gives transport back. Induction on $q$ therefore gives
\[
 \iota\widehat u_q=\chi_q m_q e_q .
\]
Put \(K=\widehat Z_q\times_{Z_{q,\mathrm{pres}}}\widehat Z_q\), with projections \(p_1,p_2\). The descent calculation is
\[
 \iota\widehat u_q p_1
   =\chi_qm_q e_qp_1=\chi_qm_q e_qp_2
   =\iota\widehat u_q p_2
 \ \Longrightarrow\ \widehat u_q p_1=\widehat u_q p_2 .
\]
Since \(e_q\) coequalizes \(p_1,p_2\), there is a unique \(v_q:Z_{q,\mathrm{pres}}\to\Sigma\) with \(v_qe_q=\widehat u_q\). Cancelling the epimorphism \(e_q\) gives \(\iota v_q=\chi_qm_q\), so \eqref{eq:query-test-domain} gives the inclusion \(Z_{q,\mathrm{pres}}\hookrightarrow Z_q\). A presentation over an epimorphic cover makes \(m_q\) an isomorphism.

For the necessity of this distinction, take one entity sort over $Z=\N$ whose fibre at $e$ is the subterminal $\{*\mid\neg\operatorname{Halts}(e,e)\}$, where $\operatorname{Halts}(e,e)$ denotes termination of $\varphi_e(e)$. Its existential scan has precisely this nonhalting predicate as its semantic event. A test $\N\to\Sigma$ for that event would semidecide diagonal nonhalting. Dovetailing with the halting test would decide the diagonal halting set. This also distinguishes the two domains. Take the tagged union of the unit query and the same scan, projected to a common output. Its success predicate is $\top$, hence $Z_q=Z$. An epimorphic cover presenting its scanned relation would, by descent for the scan alone, semidecide diagonal nonhalting on $Z$. Thus $Z_{q,\mathrm{pres}}\ne Z_q$ for this union.

\paragraph{Formal opens and subsets of source states.}
The bounded distributive lattice $B_Z=\Sigma^Z$ has pointwise finite meets and joins by \cref{lem:eff-search}. Its ideal completion defines a source locale $\mathcal Z_+$ with frame $\mathcal O(\mathcal Z_+)=\Idl(B_Z)$. The formal open generated by a test $u$ is its principal ideal $[u]=\mathord\downarrow u$. To interpret ideal opens on the actual source states, define
\begin{equation}\label{eq:state-open-comparison}
 \Phi_Z:\Idl(B_Z)\longrightarrow\Omega^Z,\qquad
 \Phi_Z(V)(z)\;\Longleftrightarrow\;\exists u\in V\;u(z).
\end{equation}
This is a frame homomorphism. A member of a join of ideals lies below a finite join of members of the generating ideals. If it succeeds at $z$, one of those members succeeds there. The converse follows from the inclusions into the ideal join, including the empty-join case. For a meet, successful witnesses $u\in V$ and $v\in W$ give $u\wedge v\in V\cap W$, and the reverse implication follows from the projections. The unit test witnesses preservation of top. Moreover $\Phi_Z([u])=\iota u$. It follows that $\mathcal O_{\rm st}(Z)=\operatorname{im}\Phi_Z$ is the subframe of $\Omega^Z$ generated by the positive predicates. Indeed every element of the image is their union, and every such union is the image of the join of their principal ideals.

For an arbitrary semantic event $A\in\Omega^Z$, its canonical ideal of sufficient tests is
\[
 J_A=\{u\in B_Z\mid\iota u\le A\}.
\]
It contains bottom and is closed downward and under finite joins. Directly from \eqref{eq:state-open-comparison},
$V\subseteq J_A$ if and only if $\Phi_Z V\le A$. It follows that $\Phi_Z J_A$ is the largest open in $\mathcal O_{\rm st}(Z)$ contained in $A$. In particular, $\Phi_Z J_A=A$ for every state open $A$, and $J_{\iota u}=[u]$. This gives a formal observable outcome for every semantic event.

The principal-ideal assignment need not preserve a pointwise countable join. On $Z=\N$ take $u_n(k)\Leftrightarrow k<n$. Its pointwise join is $\top$, while in the formal frame
\[
 \bigvee_n[u_n]=\bigcup_n\mathord\downarrow u_n
       \;\subsetneq\;[\top].
\]
The union is already an ideal since the family is increasing. Membership of $\top$ in it would supply an $n$ with $\top\le u_n$, contradicted at $k=n$. Under $\Phi_Z$ the two opens agree, since $u_{k+1}(k)$ holds. In particular a law $\mu_{\rm st}$ on $\mathcal O_{\rm st}(Z)$ induces the formal-open law $V\mapsto\mu_{\rm st}(\Phi_Z V)$, and this law identifies the two opens. For an independently specified formal-open law, directed continuity applies to $\bigvee_n[u_n]$ in its own frame.

\begin{definition}[Stochastic algebraic database and query outcome valuation]\label{def:gen-measured-instance}
For an algebraic schema $\mathcal S$, a stochastic database consists of a source-state object $Z$, an instance $I\in\operatorname{Inst}_{\E/Z}(\mathcal S)$, and a continuous subprobability valuation $\mu:\mathcal O(\mathcal Z_+)\to Q$. Its semantic event for $q$ is $\chi_q$. Put $V_q=J_{\chi_q}$. The observable outcome valuation is the direct image
\[
 \nu_q=(\gamma_{V_q})_*\mu ,
 \qquad \gamma_{V_q}:\mathcal Z_+\longrightarrow\mathbb S,
 \qquad \gamma_{V_q}^{-1}(a)=V_q ,
\]
where $a$ generates the frame of the Sierpi\'nski locale $\mathbb S$. Its mass at $a$ is $\mu(V_q)$. On source states $V_q$ denotes $\Phi_Z V_q$, the largest state open contained in the semantic event. When $\chi_q=\iota u$ on $Z$, the test $u$ is unique and $V_q=[u]$. The descent above constructs $u$ on a presented family, and its success mass is
\begin{equation}\label{eq:gen-query-law}
 \Pr_{\mathfrak D}(q)=\nu_q(a)=\mu([u])
       =\int_{\mathcal Z_+}\tau([u])\,d\mu.
\end{equation}
More generally, the pulled-back instance over $Z_q$ has the test $u_q$ of \eqref{eq:query-test-domain}. The same formula on that base uses a source law there. For a law on $\mathcal O_{\rm st}(Z)$ the mass of any semantic event belonging to this frame is directly $\mu_{\rm st}(\chi_q)$.
\end{definition}
\begin{proof}[Verification of the outcome valuation]
The frame of $\mathbb S$ is the ideal completion of the free bounded distributive lattice on $a$. Sending $a$ to $V_q$ extends to a lattice homomorphism and then to a frame homomorphism by taking joins over ideals. This gives the displayed $\gamma_{V_q}$ uniquely. Its inverse-image map preserves bottom, finite meets and arbitrary joins, so composing with $\mu$ preserves strictness, modularity and directed continuity. It preserves top as well, giving $\nu_q(\top)=\mu(\top)$. Evaluation at $a$ gives the observable mass $\mu(V_q)$. For a test $u$ the equality $V_q=[u]$ and the indicator identity of \cref{lem:eff-integral} give \eqref{eq:gen-query-law}. For an induced law $\mu=\mu_{\rm st}\Phi_Z$, the observable mass is $\mu_{\rm st}(\Phi_Z J_{\chi_q})$. This is $\mu_{\rm st}(\chi_q)$ on every state-open event, in particular $\mu([u])=\mu_{\rm st}(\iota u)$ for a test. Evaluation over $Z_q$ uses the source law specified on that base.
\end{proof}

For an explicitly finite source $Z=[n]$ with row masses $m_i\in Q$, define a law on state subsets by
\[
 \mu_{\rm st}(A)=\sum_{i<n}m_i\tau(A(i)),\qquad
 \mu(V)=\mu_{\rm st}(\Phi_Z V)
       =\sum_{i<n}m_i\tau\bigl(\exists u\in V\;u(i)\bigr).
\]
Each evaluation at $i$ preserves finite meets and arbitrary joins. Its Dirac valuation is strict, modular and directed-continuous, as is the finite weighted sum. The total mass is $\sum_i m_i$, giving a subprobability whenever this sum is at most one. In fact the first formula makes sense on all of $\Omega^{[n]}$ and assigns a semantic lower real even to an event without a uniform positive name. For a test $u$, it gives $\mu([u])=\sum_i m_i\tau(u(i))$.

The one-point source of mass one gives a deterministic table. A source with several states gives a law on tables, whereas \cref{thm:attention-kernel} constructs a law on the eligible rows \emph{within} each table. If $A_z(i,j)$ is the latter kernel, the joint finite law on $(z,j)$ has weight $m_zA_z(i,j)$. Two reads of the same random table share its source state.

For effective evaluation, use the same evaluation pullback with the named tests for formal opens; sending a named test to its principal ideal gives a name of that principal open in the principal-test representation. If only principal-open masses are requested, one can take $T=\Sigma^Z$ and $i(u)=[u]$ directly. An independent source law is necessary already on two states: the Dirac laws at opposite states give different masses to the same singleton event on the same instance family. On a presented source family \cref{thm:evidence} supplies that test and hence the name for $[u]$. Applying the valuation program yields the lower name of the success mass. The source law remains part of the statistical input.

\begin{proposition}[Best agreement available from an attribute profile]\label{prop:profile-risk}
Let \(\mu\) be a probability law on \(X\), \(Y:X\to2\) a decision, and \(b:X\to2^M\) the profile for finite \(M\). Write \(w_{S,i}=\mu\{b=S,Y=i\}\) for the masses of the measurable joint events and \(R(D)=\Pr_\mu[D(b)\ne Y]\). The optimal risk value and the predicate for attainment are
\[
\begin{aligned}
 R_*=\inf_D R(D)
    &=\sum_{S\subseteq M}\min(w_{S,0},w_{S,1})
      =\tfrac12\!\left(1-\sum_{S\subseteq M}|w_{S,1}-w_{S,0}|\right),\\
 R(D)=R_*&\quad\Longleftrightarrow\quad
         \forall S\quad w_{S,D(S)}\ge w_{S,1-D(S)} .
\end{aligned}
\]
Classically this gives an attained minimum and zero disagreement is attainable exactly when no profile has both labels of positive mass. Rational masses give a computable exact minimizer. For located mass data every positive rational tolerance gives a computable classifier within that tolerance. The internal constructive predicate for exact zero disagreement is the existence of a table \(D\) with \(w_{S,1-D(S)}=0\) for every \(S\).
\end{proposition}
\begin{proof}\leavevmode
\begin{enumerate}[titledsteps]
    \item \textbf{Compute the least error.} A Boolean formula can realize any table \(D:2^M\to2\) by disjoining complete profiles. The disjoint error events give
\[
 R(D)=\sum_S w_{S,1-D(S)},\qquad
 R(D)-R_*=\sum_S\bigl(w_{S,1-D(S)}-\min(w_{S,0},w_{S,1})\bigr).
\]
Here \(R_*\) denotes the displayed sum of minima for the moment. Each summand of the difference is nonnegative, so equality holds precisely under the stated pointwise inequalities. The absolute-value formula follows from \(2\min(a,b)=a+b-|a-b|\) and \(\sum_{S,i}w_{S,i}=1\). Finite normalized lower masses are located by the complementary-sum argument of \cref{thm:attention-kernel}. With names for the masses, take \(N=2^{|M|}\) and approximate each mass to within \(\varepsilon/(4N)\). Choosing the larger rational approximation at each profile incurs excess at most \(\varepsilon/(2N)\) there. Hence \(R_*\le R(D_\varepsilon)<R_*+\varepsilon\). This proves the infimum identity. Exact comparison gives the minimizing table for rational masses.


    \item \textbf{Test whether a selector can attain it uniformly.} For the obstruction to a uniform exact selector, let \(t\in[-1,1]\) be located, \(t^+=\max(t,0)\), \(t^-=\max(-t,0)\), and use two profiles \(A,B\) with masses
\[
 (w_{A,0},w_{A,1})=(t^+,t^-),\qquad
 (w_{B,0},w_{B,1})=(1-|t|,0).
\]
The value is identically zero. An attaining table gives \(t\ge0\) if \(D(A)=0\), and \(t\le0\) if \(D(A)=1\). To see the effective obstruction directly, present \(t_e\) by simulating \(\varphi_e(e)\). If its first halt at step \(k\) returns \(0\) or \(1\), give \(t_e\) the value \(+2^{-k-1}\) or \(-2^{-k-1}\), respectively. Use zero otherwise. Bounded simulation gives a uniform Cauchy name. A selector would compute a total bit \(f(e)=D_e(A)\) equal to \(\varphi_e(e)\) whenever that computation returns a bit. Taking an index \(j\) for the total bit program \(e\mapsto1-f(e)\) gives \(f(j)=1-f(j)\). This also rules out a uniform choice when the infimum is zero. On the uniform task grid the rounded firing bits attain risk \(6/49\) and the two tail names attain zero (\cref{ex:trained-discrepancy}). With a task label this measures accuracy. With the model decision it measures predictive fidelity \cite{concept-completeness}.\qedhere
\end{enumerate}
\end{proof}

\begin{lemma}[Witness composition of kernels]\label{lem:stoch}
Take a normalized law \(k_x\) on witnesses \(p\in P_x\) with target \(t(p)\), and a normalized law \(l_y\) on witnesses \(q\in D_y\) with target \(v(q)\). On finite presented fibres the joint witness weight is \(k_x(p)l_{t(p)}(q)\). Its output law is the Kleisli composite of the two output kernels. For a value attribute the corresponding expectation maps compose in reverse order.
\end{lemma}
\begin{proof}\leavevmode
\begin{enumerate}[titledsteps]
    \item \textbf{Sum the compatible witness pairs.} A composite witness is a pullback row $(p,q)$ with $q\in D_{t(p)}$. To find the mass at $z$, sum over all such rows with $v(q)=z$, grouping first by $y=t(p)$. The first witness law gives total weight $A(x,y)$ at $y$, and the second gives weight $B(y,z)$ from there to $z$. Hence the mass is $\sum_y A(x,y)B(y,z)$. Summing over $z$ uses normalization of each law. Multiplying instead by a value $u(z)$ and regrouping proves $(B\star A)^*u=A^*(B^*u)$.


    \item \textbf{Integrate the conditional test.} For general named kernels, a positive test $U$ on the pullback of compatible witness pairs gives at $p$ the test $U_p(q)=U(p,q)$ in the appropriate row fibre by dependent currying. The joint valuation is $M_x(U)=I_{k_x}(p\mapsto l_{t(p)}(U_p))$. Specialization of the pair test and evaluation of the second represented law give a positive lower name for the integrand uniformly in $p$. For a semantic state open use its directed union of positive tests; taking sections and evaluating $l_{t(p)}$ preserve this union, so the integrand remains lower-continuous. Linearity (\cref{lem:eff-linearity}) transports modularity through the outer integral. Directed continuity (\cref{lem:eff-integral}) transports directed joins. The integral of one recovers the mass of the first law because the second has mass one. Pulling the test back along the final projection gives the same iterated integral on the witness pairs.\qedhere
\end{enumerate}
\end{proof}

\subsection{Normalization as a pullback}
A lower bound on the total mass suffices to confirm that a row contributes. Division requires enough information to approximate the reciprocal as well. To identify the precise domain, let $\mathsf U(Q)$ be the image of the projection $(m,z)\mapsto m$ from $\{(m,z):mz=1\}$.
On the inverse-mass domain the pair \((m,z)\) is the graph data for a solved operation of \cref{thm:solved-composition}. The following calculation identifies that domain in the lower-real semantics, so the normalized port has a definite input type.

\begin{proposition}[Exact inverse-mass domain]\label{prop:normalization-unit-domain}
The projection is an isomorphism onto its image. A lower real \(m\) belongs to this image precisely when it is finite, strictly positive and located. For a mass map \(Z:X\to Q\) the domain carrying inverse mass is its pullback \(D=X\times_Q\mathsf U(Q)\).
\[
\begin{tikzcd}
 D\arrow[r]\arrow[d,hook]&\mathsf U(Q)\arrow[d,hook]\\
 X\arrow[r,"Z"']&Q
\end{tikzcd}
\]
\end{proposition}
\begin{proof}\leavevmode
\begin{enumerate}[titledsteps]
    \item \textbf{Determine the inverse graph.} If $zm=z'm=1$, then $z=z(mz')=(zm)z'=z'$. Thus the first projection of the inverse graph is monic, and regular-image factorization identifies the graph with its image.

Suppose $mz=1$. Positive rational approximants $r<m$ and $s<z$ satisfy $rs\le1$. The product's lower cut contains every rational below one, so such positive approximants exist for both factors. Each factor consequently has a positive rational lower bound, while a positive lower approximant to either factor gives a rational upper bound for the other. In particular both factors are finite and strictly positive.


    \item \textbf{Construct the upper cut.} Construct an upper cut for $m$ by
\[
 U_m=\{b\in\mathbb Q_{>0}:1/b<z\}.
\]
It is upward closed. If $b\in U_m$, roundedness supplies $s<z$ with $1/b<s$, and a rational $b'$ with $1/s<b'<b$; hence it is rounded. It is inhabited because $z$ has a positive lower approximant. If $r<m$ and $b\in U_m$, taking $s<z$ with $1/b<s$ gives $rs\le1$, and so $r<b$. Thus $U_m$ is disjoint from the lower cut of $m$.

For locatedness take $0<a<b$. Since $a/b<1=mz$, the product-cut definition supplies positive rationals $r<m$ and $s<z$ with $rs>a/b$. Rational comparison gives $r>a$ or $s>1/b$, since their simultaneous failure would imply $rs\le a/b$. The first case gives $a<m$; the second gives $b\in U_m$. For intervals meeting the nonpositive rationals the lower cut already supplies the required alternative. Hence the two cuts form a finite, strictly positive located real.


    \item \textbf{Recover the reciprocal from locatedness.} Conversely let $L_m,U_m$ be the lower and upper cuts of a positive located mass. Define a lower cut $L_z$ by including $s\le0$ and those $s>0$ for which $1/s\in U_m$. It is downward closed and rounded by upper roundedness; an inhabited upper cut gives positive members. Products of positive rational approximants $r<m$ and $s<z$ satisfy $rs<1$ by disjointness of the two cuts. Given $0\le q<1$, start with rationals $0<a<m<b$. Repeatedly trisect this bracket: locatedness at the two interior points permits keeping a bracket of width at most two thirds of the preceding width. Finitely many steps suffice for $a/b>q$, since the original lower endpoint stays positive. Now $a<m$, $1/b<z$, and $a/b>q$. Together these inclusions give $mz=1$ as lower cuts. This finite argument constructs the inverse cut without selecting a sequence of witnesses. On represented two-sided inputs interval division gives its numerical name. Pulling back the inverse graph along $Z$ gives the asserted base-change formula.\qedhere
\end{enumerate}
\end{proof}
\begin{definition}[Weighted occurrence lookup and its normalization data]
\label{def:quant-attention}
Let $I$ be a finite decidable row-occurrence object, $X$ a query object,
$K$ a key carrier, and $d:\N$ a finite feature dimension. A presented
query supplies $k:I\to K$, $v:I\to Q^d$, and a score
$s:X\times K\to Q$. Put
\begin{equation}\label{eq:attention-raw}
 w(x,i)=s(x,k(i)),\qquad Z(x)=\sum_{i\in I}w(x,i),\qquad
 N_j(x)=\sum_{i\in I}w(x,i)v_j(i).
\end{equation}
The normalization-data object is the internal subobject
\begin{equation}\label{eq:attention-domain}
 \operatorname{NormAtt}(w)=
 \{(x,z):X\times Q\mid zZ(x)=1\}.
\end{equation}
By \cref{prop:normalization-unit-domain}, this is the pullback
$X\times_Q\mathsf{U}(Q)$ along $Z$. Its second coordinate is
uniquely constructed on that domain, rather than an independent
choice. The domain consists exactly of queries with finite positive
located mass. Deciding that the mass is positive supplies only a lower bound.
\end{definition}
Figure~\ref{fig:lookup-mass} follows a weight through the two reductions.

\begin{figure}[!htbp]
    \centering\small
    \begin{tikzpicture}[x=1cm,y=1cm]
        \node (w) at (0,0) {$w_i$};
        \node[dot] (copy) at (.90,0) {};
        \node (v) at (.90,-1.15) {$v_i$};
        \node[figbox] (prod) at (2.28,-.52) {$\times$};
        \node[figbox] (mass) at (4.35,.60) {$\sum_i$};
        \node[figbox] (num) at (4.35,-.52) {$\sum_i$};
        \node[figbox] (div) at (7.65,0) {$zN$};
        \node[figbox] (inv) at (6.16,.60) {$(-)^{-1}$};
        \draw[figline] (w) -- (copy);
        \draw[figwire] (copy) -- (1.55,.60) -- (mass);
        \draw[figwire] (copy) -- (1.55,-.52) -- (prod);
        \draw[figwire] (v) -| (prod);
        \draw[figwire] (prod) -- (num);
        \draw[figwire] (mass) -- node[above] {$Z$} (inv);
        \draw[figwire] (inv) -- node[above] {$z$} (7.15,.60) -- (div);
        \draw[figwire] (num) -- node[below] {$N$} (7.15,-.52) -- (div);
        \draw[figwire] (div) -- (8.65,0);
    \end{tikzpicture}
    \caption{A normalized lookup on one occurrence fibre and one payload coordinate. Each weight enters both sums. The upper wire computes the reciprocal on the inverse-mass domain. The lower wire carries the weighted numerator.}
    \label{fig:lookup-mass}
\end{figure}

\begin{theorem}[Normalized attention as a finite probability kernel]
\label{thm:attention-kernel}
Give $D=\operatorname{NormAtt}(w)$ the state topology $\mathcal O_{\rm st}(D)$ constructed in \eqref{eq:state-open-comparison}. On $D$ define
\begin{equation}\label{eq:attention-formula}
 \alpha_{x,z}(i)=z w(x,i),\quad
 A(x,z)(U)=\sum_{i\in I}\alpha_{x,z}(i)\tau(U(i)),\quad
 \operatorname{Att}(x,z)=zN(x)\in Q^d.
\end{equation}
Then $\sum_i\alpha_{x,z}(i)=1$ and
$A:\operatorname{NormAtt}(w)\to\mathsf V_{\le1}I$ is a total-mass-one
valuation kernel. Coordinatewise,
\begin{equation}\label{eq:attention-expectation}
 \operatorname{Att}_j(x,z)=\int_Iv_j(i)\,dA(x,z)(i).
\end{equation}
The law, evaluated on named positive row tests and weights on the inverse-mass object, has uniform positive lower names.
\end{theorem}
\begin{proof}\leavevmode
\begin{enumerate}[titledsteps]
    \item \textbf{Construct the row law.} Since $zZ=1$, the nonnegative coefficients $\alpha_i$ sum to one. They therefore define the discrete valuation $\sum_i\alpha_i\delta_i$, where $\delta_i(U)=\tau(U(i))$. The indicator identities give strictness, modularity and preservation of directed joins for each Dirac valuation and hence for the finite weighted sum. Evaluating its lower integral on each coordinate gives the stated readout by the calculation following \cref{lem:eff-integral}.


    \item \textbf{Recover upper bounds on each weight.} For $c_i=\sum_{j\ne i}\alpha_j$, the upper cut for $\alpha_i$ is $b\mapsto[1-b<c_i]$. Indeed for every $\varepsilon>0$ the equality $\sum_i\alpha_i=1$ yields $l_i\in B(\alpha_i)$ with $\sum_i l_i>1-\varepsilon$. Then
\[
  l_i\le\alpha_i\le 1-\sum_{j\ne i}l_j,
  \qquad (1-\sum_{j\ne i}l_j)-l_i<\varepsilon .
\]
The upper cut is rounded and disjoint from the lower one by $c_i+\alpha_i=1$. Taking the displayed enclosure narrower than a given rational gap proves locatedness (enlarge the endpoints slightly for strict enclosures). Thus these coefficients give maps $D\to\R_{\ge0}$, whose strict rational comparisons are positive tests in $\E$. On named lower inputs the same enclosures are found by dovetailing.


    \item \textbf{Prove continuity on state opens.} For any $U\in\Omega^I$, the finite subrelations $F\subseteq U$ form an inhabited directed family with union $U$ by \cref{thm:finite-relations}. Directed continuity therefore gives
\[
 \{d:q<A_d(U)\}=
 \bigcup_{F\in\Fin I,\ F\subseteq U}
       \{d:q<\textstyle\sum_{i\in F}\alpha_i(d)\}.
\]
Every summand on the right is a state open, and the union is allowed even for an undecidable $U$. This proves kernel continuity. It also identifies the finite probability simplex with the normalized valuations on the discrete row object: the inverse takes singleton masses, disjoint finite additivity gives their sums on $F$, and the same directed union gives the value on $U$. On any separately specified source topology, the formula proves that kernel continuity is equivalent to lower continuity of all singleton masses. For a named test, run its finitely many row tests in parallel, including each confirmed row once in the enumeration of mass bounds.\qedhere
\end{enumerate}
\end{proof}

\begin{proposition}[Attention kernel composition and row migration]
\label{prop:attention-migration}
For a second finite probability kernel $B:I\to\mathsf V_{\le1}J$,
expectation of a fixed value function satisfies
\begin{equation}\label{eq:attention-kleisli}
 (B\star A)^*=A^*B^*.
\end{equation}
Let $f:I\to J$ be a finite row migration. Given $w:I\to Q$ put
$w'(j)=\sum_{i:f(i)=j}w(i)$, and suppose $v=v'f$. The total masses
of $w,w'$ agree, and on their common normalization-data object,
\begin{equation}\label{eq:attention-migration}
 \operatorname{Att}_{w'}(v')=\operatorname{Att}_w(v),
 \qquad A_{w'}=f_* A_w.
\end{equation}
\end{proposition}
\begin{proof}
The expectation law of \cref{lem:stoch} makes the top row below contravariant in row kernels. For a deterministic migration \(f:I\to J\), let \(\delta_f\) be its Dirac kernel. There are two descriptions of the same pasting:
\[
\begin{tikzcd}[column sep=large]
 Q^J\arrow[r,"{\delta_f^*=f^*}"]\arrow[dr,"{(\delta_f\star A)^*}"']&
 Q^I\arrow[d,"A^*"]\\
 &Q^D .
\end{tikzcd}
\]
Here \(D\) is the common normalization domain, and \(Q\) the lower-real weight object. The diagonal takes the expectation of a \(J\)-payload under the migrated kernel. The path through the upper right reads it at its original address and then takes expectation. Pushing the unnormalized weights forward gives the same normalized kernel since finite Fubini partitions the occurrences into their fibres:
\[
 \sum_{j\in J}w'(j)=\sum_{i\in I}w(i),\qquad
 \sum_{j\in J}w'(j)v'(j)
 =\sum_{i\in I}w(i)v'(f(i))=\sum_{i\in I}w(i)v(i).
\]
Therefore the normalization equation uses the same $z$ on both sides
and the normalized numerators agree. The same Fubini calculation on
characteristic predicates $U\subseteq J$ identifies the pushforward
valuation with the normalized weights $w'$. The proof counts each
original occurrence once, even when $f$ identifies distinct rows.
\end{proof}

\begin{proposition}[Splitting off total mass]\label{thm:mass-split}
For a finite occurrence fibre $J$, let $W^\times_J$ consist of weight tables with named inverse mass, and let $\mathsf{Prob}(J)=\{\alpha:J\to Q\mid\sum_j\alpha_j=1\}$. Internally there is a natural isomorphism for pushforward of finite rows
\[
 W^\times_J\ \cong\ \mathsf U(Q)\times\mathsf{Prob}(J),
 \qquad (w,z)\longmapsto(Z,zw).
\]
A map out of $W^\times_J$ invariant under common multiplication by a unit factors uniquely through the probability row. On presented inputs the maps use the inverse and weight names carried by the domain.
\end{proposition}
\begin{proof}
The inverse sends $(m,\alpha)$ to $(m\alpha,m^{-1})$. The identities $\sum_j\alpha_j=1$ and $zZ=1$ give the inverse laws. If $u$ is a unit, scaling sends $(w,z)$ to $(uw,z/u)$ and fixes the normalized row. Conversely $(w,z)$ is obtained by scaling $(zw,1)$ by $Z$. Restriction of an invariant map to this unit-mass section gives its unique factorization. For a finite row map $f:J\to J'$, pushforward preserves total mass by partition of the fibres and commutes with scalar multiplication. This proves the square
\[
\begin{tikzcd}
 W^\times_J\arrow[r,"\sim"]\arrow[d,"f_!"']&
 \mathsf U(Q)\times\mathsf{Prob}(J)\arrow[d,"{1\times f_*}"]\\
 W^\times_{J'}\arrow[r,"\sim"']&\mathsf U(Q)\times\mathsf{Prob}(J').
\end{tikzcd}
\]
On named inputs, the graph of inversion in \cref{prop:normalization-unit-domain} supplies the inverse. Multiplication, aggregation and projection then compute both sides of the factorization.
\end{proof}

\begin{proposition}[The constructive boundary of normalization]
\label{prop:attention-normalization-boundary}
Normalized attention is not monotone as a map on the coordinatewise
order of unnormalized weight vectors. Moreover, in the effective topos
there is no uniformly effective reciprocal of every strictly positive
lower real.
\end{proposition}
\begin{proof}\leavevmode
\begin{enumerate}[titledsteps]
    \item \textbf{Test monotonicity on two rows.} In the classical rational two-row case, take $v=(1,0)$,
$w=(1,0)$ and $w'=(1,1)$. Although $w\le w'$, their normalized
outputs are respectively $1$ and $1/2$. Hence no general
Scott-continuous raw-weight map can have this normalization as its
restriction to rational points.


    \item \textbf{Test the reciprocal on lower names.} For the effective obstruction, uniformly in a program index $e$ put
$Z_e=\tfrac12+\tfrac14\tau(\operatorname{Halts}(e))$.
It is always strictly positive. A uniform reciprocal would produce a
positive lower-real name $z_e=Z_e^{-1}$. The input enumeration initially supplies the lower bounds of \(1/2\) and adds those of \(3/4\) if the diagonal computation halts. Thus the construction is a single recursive transformation of \(e\). Its reciprocal is \(2\) in the nonhalting case and \(4/3\) in the halting case. A tracker returning an enumerable lower name of this reciprocal would semidecide nonhalting by waiting for the positive test \(3/2<z_e\). Dovetailing with the halting test would decide the diagonal halting set. If \(d(e)\) decided whether \(\varphi_e(e)\) halts, the program that halts on \(e\) exactly when \(d(e)\) says it does not halt gives a contradiction on its own index.

Thus \eqref{eq:attention-domain} is the exact domain for inversion of the supplied total mass. A particular family can determine its probability row on a larger domain. For a singleton with any strictly positive weight $m$, the row $(1)$ already satisfies $m\cdot 1=m$, and it agrees with the normalized row wherever the inverse exists.\qedhere
\end{enumerate}
\end{proof}

\subsection{Integral form}
\begin{theorem}[Integral attention and pushforward]
\label{thm:refine-integral-attention}
For topological row and parameter spaces take a continuous valuation $\mu$ and a continuous nonnegative lower-real score on their topological product. For the internal represented construction take the state opens generated by positive tests, use the represented product, and let the represented score and law range over $w:X\times R\to Q_\Sigma$ and $\mu\in\mathsf{Val}_{\rm obs}(R)$, respectively. Both constructions use the same formula on the underlying rows.
The lower integral of \cref{lem:eff-integral} constructs the weighted valuation
family and mass
\[
 \lambda_x(U)=\int_Rw(x,r)\tau(U(r))\,d\mu(r),\qquad
 Z(x)=\lambda_x(R).
\]
On the normalization-data object
$D=\{(x,z):zZ(x)=1\}$, the formula
\[
 A(x,z)(U)=z\lambda_x(U)
\]
defines a normalized valuation kernel. For every represented
nonnegative feature $v$, its readout is $\int_R v\,dA$.
For a continuous row map $f:R\to S$, put
$\lambda'_x=f_\#\lambda_x$. Then the masses and normalization-data
objects agree, $A'=f_\#A$, and
\[
 \int_S v'\,dA'=\int_R(v'\circ f)\,dA.
\]
The represented integration and normalization operations preserve their supplied names on the displayed domain.
\end{theorem}
\begin{proof}\leavevmode
\begin{enumerate}[titledsteps]
    \item \textbf{Weight the valuation.} We first show that weighting defines a valuation. In the partition formula of \cref{lem:eff-integral}, use the opens $U\cap\{w_x>r_k\}$. For each partition, the result is a finite nonnegative linear combination of restricted valuations $U\mapsto\mu(U\cap V)$. Restriction preserves strictness and modularity, and it preserves directed joins since intersection distributes over unions. Refinement increases the partition values for every $U$. Taking a common refinement on both sides of the modularity equation preserves that equation at the supremum. Exchanging the two suprema proves preservation of directed joins. Thus $\lambda_x$ is a valuation.


    \item \textbf{Establish continuity in the query.} First use the ordinary product topology on $X\times R$. To show lower semicontinuity in $x$, write a threshold open of $w$ in $X\times R$ as a union of rectangles. Its section at $x$ is the directed union of finite unions of row opens from rectangles whose first factor contains $x$. A strict rational bound on the mass of that section already holds on one finite union. Intersect the finitely many corresponding parameter opens. Throughout this neighbourhood the same row opens witness the same bound. Intersecting first with $U$ gives the statement for restricted sections, and taking finite sums and the supremum gives it for $\lambda_x(U)$.


    \item \textbf{Compute the represented masses.} For the represented product, instead fix a positive test $u:R\to\Sigma$ and a rational $a>0$. Curry the test $(x,r)\mapsto u(r)\wedge [a<w(x,r)]$ to obtain $X\to\Sigma^R$; apply the evaluator in $\mathsf{Val}_{\rm obs}(R)$. Finite rational sums over a partition and the enumerable supremum over partitions then give a morphism $X\to Q_\Sigma$, uniformly in $u$. This is exponential adjunction internally in $\E$; on program indices one specializes the test to the parameter name. The corresponding Type--2 operations are computed in \cite[Proposition~3]{pauly-represented}. For an arbitrary state open $U$, the positive tests contained in $U$ form the directed ideal $J_U$ and their union is $U$. Continuity of the valuation just proved gives $\lambda_x(U)=\sup_{u\in J_U}\lambda_x(\iota u)$. Rational lower bounds on this supremum are unions of state opens. On the inverse-mass graph $z$ is a positive located real by \cref{prop:normalization-unit-domain}; multiplication preserves lower continuity and computes the corresponding names. At the level of ordinary Type--2 represented spaces the same currying argument establishes continuity also for any fixed Scott-continuous valuation: the open-set representation is continuous to the Scott topology \cite[\S2]{hoyrup-represented}, and composition with the valuation is continuous into the lower-real topology. The latter has the admissible lower-cut representation, obtained by enumerating rational lower bounds. Thus the mass functional has a continuous realizer there; computable evaluation uniformly in a law uses its represented name. This external continuity statement does not identify Type--2 names with Kleene indices in $\E$. In the ordinary topological construction the graph carries its subspace topology from $X$ and the lower-real coordinate $z$. Thus both source topologies have the required kernel continuity.

All the partition thresholds here are positive, so these opens compute the lower integral of $w_x\tau U$. At $U=R$ that integral is $Z(x)$. Multiplication by its supplied inverse therefore normalizes $\lambda_x$.


    \item \textbf{Push the normalized law forward.} For a continuous \(f:R\to S\), preimages commute with the open tests and \(f^{-1}S=R\). Hence the mass stays fixed and scalar multiplication by the same \(z\) commutes with pushforward on every open. Change of rows in \cref{lem:eff-integral} proves the expectation identity. On names, intersection of open tests and the partition enumeration give a name for every required mass, uniformly in \(x\). With a name for the inverse mass the same finite arithmetic gives the normalized valuation. Finite counting on row occurrences recovers \cref{thm:attention-kernel}.\qedhere
\end{enumerate}
\end{proof}

On finite or represented rows, a weighted lookup thus exports two quantities: the weighted value and its mass. Their behaviour under change of rows is determined before normalization. The inverse-mass object identifies the queries for which those quantities also determine a probability law.

\subsection{Jets of a normalized lookup}
On the principal open \(D(Z)\), the scalar multiplication map by \(Z\) on the pulled-back value module is an automorphism. The normalized readout is its inverse applied to the numerator section \(N\). Pulling this module equation back along a first-order point gives \(Z\dot y+\dot Z\,y=\dot N\). The unit \(Z\) uniquely determines \(\dot y\). In coordinates for a finite located row, put \(Z=\sum_j w_j\), \(N=\sum_j w_jv_j\) and \(y=N/Z\). The module equation becomes
\begin{equation}\label{eq:unnormdiff}
 \dot y=\frac1Z\sum_j (v_j-y)\dot w_j+\sum_j\alpha_j\dot v_j,\qquad
 \alpha_j=w_j/Z.
\end{equation}
The centred coefficient of \(\dot w_j\) comes from subtracting \(y\dot Z\) from the numerator variation. Common rescaling acts through the same unit on both \(N\) and \(Z\), so the induced variation on the quotient vanishes. In row coordinates the cancellation is \(\sum_j w_j(v_j-y)=0\). This is the infinitesimal form of \cref{thm:mass-split}.

\begin{proposition}[Jets of finite normalized lookup]\label{prop:normalized-jets}
On finite signed weight coordinates the readout $Y=N/Z$ is a morphism on the principal open $D(Z)$. Its order-$n$ jet is the table $y_0,\ldots,y_n$ determined by
\[
 y_0=n_0/z_0,\qquad
 y_m=z_0^{-1}\left(n_m-\sum_{i=1}^m z_i y_{m-i}\right).
\]
It commutes with any row migration which transports the numerator and mass queries together.
\end{proposition}
\begin{proof}
Interpret both grouped queries in $B_n$. The input arc factors through $D(Z)$ exactly when $z_0$ is a unit, by the nilpotent inverse formula or \cref{thm:etale-jets}. Equating coefficients of $ZY=N$ now gives the recursion and its uniqueness. When $Z$ is a positive finite located mass, \cref{prop:normalization-unit-domain} constructs its real inverse. A migration preserving the two queries after substitution preserves all coefficients by \cref{thm:marked-jets}, hence preserves the unique quotient arc. At order one this is the marked derivative of the readout. Lower integrals on represented spaces have the change-of-rows statement proved above; the present signed jet calculation concerns the finite numerical registers.
\end{proof}

\subsection{Numerical primitives for a scored lookup}

\begin{lemma}[Effective analytic primitives]\label{neural:lem:analytic}
The exponential, its derivative, and the reciprocal on a positively witnessed located-real input have effective local approximation procedures. Their forward and cotangent operations extend \cref{neural:thm:cotangent} to finite circuits built from these primitives. The sigmoid $\sigma(t)=(1+\exp(-t))^{-1}$ is a total instance.
\end{lemma}
\begin{proof}
Given a rational bound $|t|\leq B$, the exponential series has a computable tail bound. Beyond an index greater than $2B$, the ratio of successive absolute terms is at most $1/2$, so the remaining tail is bounded by twice its first omitted term. The differentiated series is again the exponential series and has the same bound. Finite search on a located-real name provides $B$. The identity $\exp(t)\exp(-t)=1$ converts an upper bound for $\exp(-t)$ into a positive lower bound for $\exp(t)$.

For reciprocal, refine a positive located input until its lower endpoint is positive. Interval division then computes its reciprocal, whose derivative is $-t^{-2}$. The sigmoid denominator is at least one. The local estimates give tracked forward and reverse maps for these gates. Products and the chain rule extend the optic composition equations to their finite circuits.
\end{proof}
These primitive derivatives apply on their located-real numerical domains. In particular, inserting the gate for a positive test in a query and differentiating a real-valued scoring gate are distinct operations.

For a real input jet $a(t)$ the numerical gate has an effective extension in every finite order. Set $b_0=\exp(a_0)$ and
\[
           m b_m=\sum_{j=1}^m j\,a_j b_{m-j}\qquad(1\le m\le n).
\]
Indeed $b_0\sum_{h=0}^{n}(a(t)-a_0)^h/h!$ has just these coefficients, by formal differentiation through degree $n-1$. The finite expression is natural under substitution of nilpotent increments, including several marked directions. The analytic tail bound of \cref{neural:lem:analytic} computes the retained scalar; all higher coefficients use only finite arithmetic. Positive square roots on nonzero length registers likewise use the located square root at order zero and the successive corrections of \cref{ex:root-cover}. The weight coordinates in an affine chart can therefore be fed by real score programs, with their own tracked jets.

\subsection{Bilinear scoring and scaled dot-product attention}
The score port is the evaluation pairing \(H^*\otimes_R H\to R\). On query and key occurrences the two projections supply a dual field and a primal field. Pairing, applying the chosen positive weight map, and normalizing in the mass fibre gives the row kernel of \cref{thm:attention-kernel}. An isomorphism of feature registers acts contragrediently on the query field, as the following comparison shows.

\begin{construction}[Query into a dual, key into its primal]
Take a finite free key module \(H\). Produce \(q_i\in H^*\) and \(k_j\in H\) by two independently parameterized affine incidence queries of the token features. Produce \(v_j\in V\) by a third such query. The evaluation map \(H^*\times H\to R\) gives \(c_{ij}=q_i(k_j)\). In coordinates this is exactly \cref{ex:score-witness}. For \(d=\operatorname{rank}H>0\) and located-real values choose \(s_{ij}=c_{ij}/\sqrt d\), the scaled score of \cite{vaswani-attention}. The scale is an architectural choice, not a consequence of grouping rows.
\end{construction}
\begin{proposition}[Bilinear score and change of representation]\label{prop:pairing}
For free finite input modules \(U,W\), each bilinear score \(B:U\times W\to R\) corresponds to a linear map \(U\to W^*\), and equivalently to a linear functional on \(U\otimes_R W\). If a key basis is changed by an isomorphism \(S:H\to H'\), transform the key by \(S\) and the query by \((S^{-1})^*\). The score and the resulting row kernel are unchanged.
\end{proposition}
\begin{proof}
Send \(u\) to the functional \(w\mapsto B(u,w)\). Bilinearity supplies linearity both of this functional and of the assignment in \(u\). Evaluating it on \(w\) recovers \(B\), and these operations are inverse. The tensor presentation on pairs modulo additivity and scalar balancing identifies the same forms with linear functionals by quotient descent. Finally \(((S^{-1})^*q)(Sk)=q(k)\) slot by slot. Scores, weights, masses and normalized kernels agree. Thus the numerical dot product of two column representations uses their displayed dual pairing.
\end{proof}

\begin{example}[A pair score as a witness polynomial]\label{ex:score-witness}
Let $e\in M$ be an eligible occurrence with query address $\pi(e)$ and key address $t(e)$. With coordinate index $[d]$, the operation table is $M\times[d]$, projected to $M$. At $(e,h)$ the two sockets read $(\pi(e),h)$ in the query view and $(t(e),h)$ in the key view. Take the disjoint union of the views as input sorts. With coefficient one, the score before scaling at $e$ is $\sum_h q_{\pi(e),h}k_{t(e),h}$. Substituting projection queries for either view now has a specified composite witness object. Exponentiation, which will be applied to the resulting scalar, is a subsequent interpreted gate.
\end{example}

\begin{proposition}[Classical scaled dot-product recovery]
\label{prop:attention-softmax}
In the numerical located-real interpretation, let finite query and key
vectors have computable located coordinates with positive key
dimension $d_k\ge1$, and put
$c(x,i)=\langle q(x),k(i)\rangle/\sqrt{d_k}$, and take
$w(x,i)=\exp(c(x,i))$. Then every nonempty finite row collection has
strictly positive located normalizer and a computable inverse. On
located real-valued feature vectors the attention formula becomes
\begin{equation}\label{eq:attention-softmax}
 \operatorname{Attention}(Q,K,V)
 =\operatorname{softmax}\!\left(\frac{QK^\top}{\sqrt{d_k}}\right)V.
\end{equation}
This is a numerical specialization, not an assertion that signed dot
products are intrinsic operations of the semiring of nonnegative lower reals.
\end{proposition}
\begin{proof}\leavevmode
\begin{enumerate}[titledsteps]
    \item \textbf{Compute the exponential.} The rational exponential on $0\le y\le1$ is computed by the series $\sum_{n\ge0}y^n/n!$. After degree $N$ the tail is bounded by $2/(N+1)!$ using the bound on successive-term ratios. Integer range reduction, the identity $\exp(-t)\exp(t)=1$ and interval arithmetic extend this to all located-real inputs; the resulting enclosures calculate strictly positive located weights, as in \cref{neural:lem:analytic}.


    \item \textbf{Construct the inverse mass.} For a nonempty finite row, take a key $i_0$ directly from its enumerated list. Refining the exponential enclosure at $c(x,i_0)$ produces a positive rational $a<\exp(c(x,i_0))$. Hence
\[
 0<a<\exp(c(x,i_0))\le Z(x)=\sum_i\exp(c(x,i)).
\]
The sum is a located real, and its enclosures can be refined to $0<a\le l_n<Z(x)<u_n$ with widths tending to zero. The rational intervals $(1/u_n,1/l_n)$ have widths tending to zero, uniformly in the input names. They construct $Z(x)^{-1}$ with $Z(x)Z(x)^{-1}=1$. Thus the projection of the normalization-data graph is an isomorphism onto the finite nonempty softmax input object, with inverse $x\mapsto(x,Z(x)^{-1})$. No extra normalization datum is supplied in this numerical realization.


    \item \textbf{Recover the attention output.} Substitution into \eqref{eq:attention-formula} gives the displayed scaled dot-product formula. Signed value coordinates use the located real-vector structure and the same finite weighted sum; no identification of signed coordinates with $Q^d$ is made.\qedhere
\end{enumerate}
\end{proof}

\begin{construction}[A token-indexed circuit]\label{con:head}\leavevmode
\begin{enumerate}[titledsteps]
    \item \textbf{Read the projections.} For a finite token family $x_i\in R^m$ indexed by $J$, read the matrices $W_Q,W_K,W_V$ and biases once as parameter tables. Compute $q_i=W_Qx_i+b_Q$, $k_j=W_Kx_j+b_K$ and $v_j=W_Vx_j+b_V$. 
    \item \textbf{Form the masked score table.} Write the mask as an incidence span $J\xleftarrow{\pi}M\xrightarrow{t}J$. Each occurrence $e$ has query address $\pi(e)$ and key/value address $t(e)$. In a list of eligible keys the list positions are the occurrences, so multiple positions may read the same key. A binary mask has one occurrence $e=(i,j)$ for each eligible pair. Write $M_i=\pi^{-1}(i)$ and enumerate each nonempty $M_i$. Its first occurrence supplies the positive-mass witness of \cref{prop:attention-softmax}. The fields of the query are
\[
 s_e=\frac{q_{\pi(e)}(k_{t(e)})}{\sqrt d},\qquad w_e=e^{s_e},\qquad
 Z_i=\sum_{e\in M_i}w_e,\qquad \alpha_e=\frac{w_e}{Z_{\pi(e)}} .
\]

    \item \textbf{Normalize and return the values.} The two-socket query of \cref{ex:score-witness} supplies the pairing. Group the exponentiated scores by $\pi$ to compute $Z_i$, then join this table back along $\pi$ to normalize each occurrence. Joining the values along $t$ gives
\[
 y_i=\sum_{e\in M_i}\alpha_e v_{t(e)} .
\]
The active output carrier is $J_+=\pi(M)$, computed by grouping the finite occurrence list. A block returning every token row uses the shapes with $J_+=J$, as happens for full masks on nonempty token tables and for causal masks containing the diagonal. In the binary case $s_{(i,j)}=s_{ij}$ and $\alpha_{(i,j)}=\alpha_{ij}$ give the usual masked array formula.

\end{enumerate}
\end{construction}
\begin{example}[Scoring and grouping]\label{ex:group}
Let three eligible row IDs $a,b,c$ carry features $0,1,1$ and weights $1,2,2$. Their masses are $1/5,2/5,2/5$, so the expected feature is $4/5$. Grouping by feature value gives weights $1,4$ and the same expectation. Scoring the two distinct feature values afresh would give weights $1,2$ and expectation $2/3$. The two presentations in \cref{prop:attention-migration} agree because the migration pushes forward the original occurrence weights.
\end{example}

Take $x=(0,1,1)$, $q_i=u x_i$, and $k_j=v_j=x_j$, with full eligibility. At $u=\log 2$ the two parameter-dependent queries $b,c$ both generate the weights in \cref{ex:group}. Query $a$ has zero logits at all three keys and returns $2/3$. Queries $b$ and $c$ read one scalar register $u$. We will compute its update after deriving the pullback for a head.

\begin{example}[Two conventions at zero mass]
For an exact rational table the zero-mass test is decidable. Adding an always-present row with value zero and weight one returns $N/(1+Z)$ for every row. Returning zero only at $Z=0$ instead returns $N/Z$ on all positive rows. These choices define different operations already at positive mass. The nonempty softmax construction uses its computed inverse.
\end{example}

\subsection{Recovering the normalization equation and its adjoint}
On the inverse-mass domain of this construction, put $N=\sum_e w_ev_e$, $Z=\sum_e w_e$. Read back the result as the unique $y$ with $Zy=N$. Its computed section is $y=Z^{-1}N$, and multiplication by the inverse proves uniqueness. For a joint increment multiindex $\alpha\ne0$, comparing coefficients gives
\[
 y_\alpha=Z_0^{-1}\Bigl(N_\alpha-
       \sum_{0<\beta\le\alpha}Z_\beta y_{\alpha-\beta}\Bigr).
\]
The order-zero coefficient is $Z_0^{-1}N_0$. This is an algorithm in increasing total degree. For positive $w_e=\exp(s_e)$ put $\alpha_e=w_e/Z$. Differentiating the single equation gives
\[
 dy=\sum_e\alpha_e\{dv_e+(v_e-y)\,ds_e\}.
\]
For a covector $g$, the value receipt is $\alpha_e g$ and the logit receipt is $\alpha_e\langle g,v_e-y\rangle$. These are source queries, followed by grouping at the original owner keys. With a structured witness $b$, multiply its occurrence weights. Its log-weight differential is $H_b=\sum_{e\in E_b} ds_{\kappa(e)}$. The same derivation with $b$ as the row gives
$d\mathbb E[V]=\mathbb E[dV]+\mathbb E[VH]-\mathbb E[V]\mathbb E[H]$.
An isomorphism of witness tables over the output preserving weights and payloads transports both $Z$ and $N$ term by term; hence it transports the solution and every coefficient above.

For a law on \emph{input tables}, the corresponding comparison uses the source law specified in \cref{eq:gen-query-law}. A structural witness isomorphism over each source state gives the same image event there, and the retained semidecision test classifies that event on its domain. Pushing the same law through the equal tests gives equal masses. A finite two-state example with event $\{z_1\}$ and source masses $(1-p,p)$ gives mass $p$. Its value varies with the source law even if every within-table attention computation stays fixed.

Fix scores in one nonempty group \(E\), with independent scalar payload \(v\). Its value functional \(T_s\) satisfies
\[
 \begin{gathered}
 T_s(e_e)=\alpha_e,\qquad T_s=T_{s'}\ \Longleftrightarrow\ \alpha=\alpha'
      \ \Longleftrightarrow\ s-s'\in\R\mathbf1,\\
 s^0=\log\alpha-|E|^{-1}\Big(\sum_e\log\alpha_e\Big)\mathbf1,\qquad [s^0]=[s].
 \end{gathered}
\]
Here \(e_e\) is the basis payload. Positivity permits taking logarithms of \(\alpha_e=\exp(s_e)/Z\), which proves the middle equivalence and the inverse formula. If payloads are read through \(\kappa:E\to K\), the exposed operation is \(T_s\kappa^*\) and
\[
 (T_s\kappa^*)(e_a)=\sum_{\kappa(e)=a}\alpha_e,\qquad
 T_s\kappa^*=T_{s'}\kappa^*\ \Longleftrightarrow\
                   \kappa_!\alpha=\kappa_!\alpha'.
\]
This equality on basis inputs is both necessary and sufficient by linearity. If two occurrences share an address, a sufficiently small opposite perturbation of their probabilities preserves positivity, total mass and every exposed value, while changing the score class recovered from independent payloads.

\section{Attention on pointed witnesses}\label{sec:structured}

The interface selecting a pointed count first retains a socket address and then sums over complete witnesses. The two field operations below give its interpretation in the concrete model of \cref{sec:field-language}; the count vector specifies the observation used for the inverse problem.

\paragraph{Marking before marginalization.}
Let \(B\) carry the normalized law \(p\) on complete witnesses and let \(E\) be the table of pairs consisting of a witness and one of its sockets. Write \(s:E\to B\) for the projection and \(\kappa:E\to A\) for the factor address. On finite weight fields the pointed query is the linear interpretation of the span
\[
\begin{tikzcd}[column sep=large]
 &E\arrow[dl,"s"']\arrow[dr,"\kappa"]\\
 B&&A .
\end{tikzcd}
\qquad
 p\longmapsto \kappa_!s^*p .
\]
It takes a probability field on witnesses to a weight field on factor addresses. Its total mass is the expected number of sockets. This follows from \((!_A)_!\kappa_!=(!_E)_!\), where \(!_A:A\to1\) is the terminal map. For a payload depending on the marked socket, insert multiplication by that payload before summation. The product-weight model below instantiates this construction.
Choose a finite table $B$ of complete witnesses and a key $\kappa:E\to A$ on their factor sockets. Give witness $b$ weight $m_b=\prod_{e\in E_b}w_{\kappa(e)}$, with a row coefficient if desired, and put $Z=\sum_b m_b$. Work on the inverse-mass domain of \cref{prop:normalization-unit-domain}. Positive weights and any explicitly generated legal witness exhibit a positive summand. A one-mark query may read \emph{any} payload $v(b,e)$ available after the mark:
\[
     Y_v=Z^{-1}\sum_b\sum_{e\in E_b}m_b\,v(b,e).
\]
Construct it by taking the one-hole table, putting the removed weight back, joining the payload, and grouping. Grouping by $\kappa(e)=a$ gives $\mathbb E[N_a]$, where $N_a(b)$ counts occurrences of $a$. Normalization is over complete witnesses. A witness carrying three marks of interest can therefore contribute three times.

Here the ordinary forward table is built from the context object of \cref{thm:marked-jets}. Choosing the payload equal to the indicator of a factor key turns that object into a table of expected counts. The same construction in the chart $D(Z)$ supplies its variation by \cref{prop:normalized-jets}.

\subsection{Chain and projective-tree attention}

With one eligible row per witness and one factor socket this is ordinary attention. With witnesses equal to legal state strings or projective dependency trees and sockets equal to constituent factors, it is the structured attention of Kim et al.~\cite{structured}. A chain chart records $H_1(t)=a(t)$ and $H_{i+1}(t)=\sum_s H_i(s)w_{i,s,t}$. Joining successive endpoint keys gives a unique earlier prefix for each extension, so summing the last table gives $Z$.

For the segmentation experiment of \cite[\S4.2 and Appendix A.3]{structured}, the additional rescaling is also a marked query. Mark once at each selected state site, with payload either its feature $h_i$ or the unit. The two expectations are $V=\sum_i\Pr(z_i=1)h_i$ and $M=\sum_i\Pr(z_i=1)$. The reported context is $\lambda V/M$ with $\lambda=2$. On a nonempty binary chain with finite exponential potentials the all-selected string has positive weight, so these finite numerical queries construct a positive located $M$.

For trees, split each projective subtree at its root into a left and a right forest. Split each nonempty forest just after its first component and retain that component's root as the parent-arc socket. These cuts are unique and their inverse joins reconstruct the tree. The chart in \cref{app:treechart} implements this bijection with shared registers. The payload of the marked arc $i\to j$ can read the value $v_i$; grouping by $j$ and dividing by $Z$ gives the expected parent value of each dependent token.

\subsection{Two marked factors in the same structure}

For an incidence predicate $L(b,e,f)$ and a pair value $v(b,e,f)$ form
\[
 Z^{-1}\sum_b m_b\!\!\sum_{\substack{e,f\in E_b,\ e\ne f\\L(b,e,f)}}v(b,e,f).
\]
This head can attend to two child arcs of the same parent in one latent tree, or two linked records in a database witness. It measures their joint inclusion. If a probability law on all such marked pairs is wanted, its normalizer is the total marked-pair mass instead. In the located-real numerical chart, for log-weight perturbations $h$, marking $m_b$ gives $m_b H_b$ with $H_b=\sum_e h_{\kappa(e)}$. The marked unit equation for $Z^{-1}$ therefore gives
\[
 \dot{\mathbb E[N_a]}=\mathbb E[N_a H]-\mathbb E[N_a]\mathbb E[H].
\]
Write \(E_a(b)=\{e\in E_b:\kappa(e)=a\}\). The ordered pairs split into off-diagonal and diagonal parts:
\[
 N_a(b)N_c(b)=
  \#\{(e,f)\in E_a(b)\times E_c(b):e\ne f\}
           +\mathbf1_{a=c}\,N_a(b).
\]
For \(h\) at factor addresses, put \(H_b=\sum_cN_c(b)h_c\). After weighting complete witnesses this gives
\[
 D M_a(s)[h]=\sum_c\bigl(\mathbb E[N_aN_c]
                         -\mathbb E[N_a]\mathbb E[N_c]\bigr)h_c.
\]
At a single witness with monomial \(w^2\) there are two ordered distinct-position pairs and two diagonal pairs. The second moment is four and the variance is zero. For an input-dependent payload, the product rule also marks the payload. Expectation semirings for higher moments on packed forests are described in \cite{moments}.

\subsection{Recovering structured scores from expected factor counts}
A structured head often exposes only factor marginals. The one-mark query still gives a reverse construction. Let $b$ run over a nonempty finite presented witness table, let $n_b\in\N^l$ count its uses of the $l$ factor addresses, and let $q_b>0$ be its known located base coefficient. Put
\[
 p_b(s)=\frac{q_b e^{\langle s,n_b\rangle}}
                 {\sum_cq_c e^{\langle s,n_c\rangle}},\qquad
 M(s)=\sum_b p_b(s)n_b,\qquad
 U=\mathrm{span}\{n_b-n_{b_0}\}.
\]
The classical mean-parameter theorem is given in \cite[Theorem 3.3]{wj-mean}. Here the inverse must produce a located name from the marked query output. The counts are integer data, so rational elimination computes $U$ and its orthogonal projection.
\begin{theorem}[Constructive readback of expected counts]\label{thm:mean-inverse}
The factor score quotient $\R^l/U^\perp$ is in effective bijection with the means having a strictly positive barycentric presentation
$m=\sum_b p_b n_b$, $\sum_b p_b=1$, $p_b>0$.
The forward map is the one-mark query $M$. The reverse map computes its unique score in $U$, and thus computes every witness probability. This includes repeated count vectors and repeated uses of one address.
\end{theorem}
\begin{proof}\leavevmode
\begin{enumerate}[titledsteps]
    \item \textbf{Recover a positive barycentric presentation.} Adding a vector in $U^\perp$ changes every log weight by the same scalar and cancels in normalization. The forward probabilities give a positive presentation. To invert, work in the affine hull of the count vectors, whose equations have rational coefficients. Row-reduce the rational matrix with columns $(1,n_b)$. Enumerate rational values for its free probability coordinates; solve for the pivot coordinates using $(1,m)$, and dovetail tests of strict positivity. If a strictly positive presentation exists, positivity persists in a neighbourhood of its free coordinates. Density of the rationals gives a successful trial. Conversely a successful trial gives the presentation exactly on this affine-hull base. Choose a positive rational $\delta$ strictly below every returned coordinate. For $U=0$, all $n_b$ coincide and score zero is the unique representative.


    \item \textbf{Bound the search region.} Otherwise choose a rational basis matrix $B$ with column span $U$, write $s=Bz$, and set $v_b=B^{\mathsf{T}} n_b$. The rational differences $v_b-v_{b_0}$ span the full coordinate space since $B$ has independent columns in $U$. Invert an independent square minor to compute $c>0$ rational so that
\[
 \max_b \langle z,v_b\rangle-\min_b\langle z,v_b\rangle
      \ \ge\ c\|z\|_2 .
\]
For instance take a reciprocal rational upper bound on $\sqrt{\dim U}$ times the infinity operator norm of the inverse minor with differences as rows. Define
\[
 J(z)=\log\sum_b q_b\exp\big(\langle z,v_b-B^{\mathsf{T}}m\rangle\big).
\]
Write \(\alpha_b> \delta\) for the feasible barycentric coefficients, \(l_b=\langle z,v_b-B^{\mathsf T}m\rangle\), \(L=\max_b l_b\), \(l=\min_b l_b\). Since \(\sum_b\alpha_b l_b=0\),
\[
 L=\sum_b\alpha_b(L-l_b)\ge\delta\sum_b(L-l_b)
       \ge\delta(L-l)\ge\delta c\|z\|_2 .
\]
The penultimate bound follows because the sum dominates each gap. These are inequalities of lattice values and require no selection of an extremizing mass index. Take a rational \(0<q_-<\min_bq_b\). Then \(J(z)\ge\log q_-+L\). Choose a rational \(R\) such that \(\log q_-+\delta cR>J(0)+1\). The boundary of the cube of radius \(R\) lies above that level.


    \item \textbf{Bound the Hessian below.} Indeed on that cube a positive rational $\epsilon$ below every normalized probability is obtained by bounding the finitely many linear exponents and the $q_b$ above and below. The gradient is $B^{\mathsf{T}}(M(Bz)-m)$. For a direction $u$, the Hessian quadratic form is
\[
 \frac12\sum_{b,c'}p_b p_{c'}\langle u,v_b-v_{c'}\rangle^2
 \ \ge\ \epsilon^2 c^2\|u\|_2^2.
\]
For rational \(u\), select the two extreme values of the rational projections \(\langle u,v_b\rangle\). Their two ordered terms give the Hessian bound. Continuity in \(u\) extends it to located directions. Set \(\kappa=\epsilon^2c^2\) and let \(L_0\) be a positive rational upper bound for \(\max_b\|v_b-B^{\mathsf T}m\|_2\). 
    \item \textbf{Compute the minimizing score.} For each error \(e_j>0\) minimize rational objective approximations of error at most \(e_j/4\) on a rational grid of covering radius \(e_j/(4(1+L_0))\). At its chosen point \(z_j\), \(J(z_j)\le J(z)+e_j\) throughout the cube. Applying this inequality for \(z=(z_i+z_j)/2\) and each index gives
\[
 \frac{\kappa}{8}\|z_i-z_j\|_2^2
 \le \tfrac12(J(z_i)+J(z_j))-J((z_i+z_j)/2)
 \le \tfrac12(e_i+e_j).
\]
Take \(e_j=\kappa2^{-2j-4}\). The limit is a computed minimizer, lies in the interior by the coercive bound, and satisfies \(B^{\mathsf T}(M(Bz)-m)=0\). Since \(M(Bz)-m\in U\) and \(B^{\mathsf T}|_U\) is injective, \(M(Bz)=m\). Conversely any solution in \(U\) is a critical point and hence a minimizer by convexity. Coercivity and the midpoint inequality give uniqueness. Orthogonal projection onto \(U\) then gives both score-quotient identities.\qedhere
\end{enumerate}
\end{proof}
This mean domain is the relative interior of the convex hull of the $n_b$. At any strictly positive barycentric vector a rational right inverse to the map from sum-zero probability increments onto $U$ lifts a small ball of mean increments while keeping all entries positive. Conversely let $\bar n$ be the average of all $n_b$. At a relative interior point $m$, choose $0<\epsilon<1$ sufficiently small that $(m-\epsilon\bar n)/(1-\epsilon)$ lies in the convex hull. A barycentric presentation of this last point, mixed with the uniform one with weight $\epsilon$, yields a strictly positive presentation of $m$.

For example, witnesses with count vectors $(2,0)$ and $(0,1)$ and unit base weights give means $(2p,1-p)$ with $0<p<1$. The recovered score is specified by $2s_1-s_2=\log(p/(1-p))$, with representative $\frac15(2,-1)\log(p/(1-p))$ in $U$. The factor of two is the count of occurrences of the first address. This inverse uses the retained witness table and the base coefficients; a scalar head exposing a single constant payload could give the same output for every score.

The inverse also has a differential program. In the basis above let $u=B^{\mathsf{T}} m$ and $\bar v=\sum_b p_b v_b$. The covariance $C=\sum_b p_b(v_b-\bar v)(v_b-\bar v)^{\mathsf{T}}$ is the positive definite Hessian used in the proof. Marking the softmax gives $dp_b=p_b\langle v_b-\bar v,dz\rangle$, whence $du=C\,dz$. A second mark at directions $a,d$ gives the vector
\[
 T(a,d)=\sum_b p_b(v_b-\bar v)
       \langle v_b-\bar v,a\rangle\langle v_b-\bar v,d\rangle .
\]
The derivatives of the centring vectors have zero contribution after summing \(p_b(v_b-\bar v)\). For mean coordinates \(u=B^{\mathsf T}m\), two independent mean directions \(h,k\), and \(u(z(u))=u\), the two coefficient equations are
\[
 C z_h=h,\qquad C z_{hk}+T(z_h,z_k)=0 .
\]
Thus \(z_h=C^{-1}h\), \(z_{hk}=-C^{-1}T(C^{-1}h,C^{-1}k)\). Positive definiteness gives positive leading pivots for elimination. The feasible probability trial and the strict cube bounds persist locally. Pairing the gradient difference with \(\widetilde z-z\) on this cube yields
\(\kappa\|\widetilde z-z\|_2\le\|\nabla J(\widetilde z)\|_2\).
An approximate inverse with residual \(o(\|t\|^m)\) consequently has error of the same order. This identifies the computed coefficients with derivatives of the selected inverse.

\section{Balancing two foreign keys}\label{sec:balance}

Here the selected output is a table with two prescribed margin fields. The fibre below presents the constraint. The construction in \cref{thm:balance-exact} supplies its balanced representative for positive dense inputs, and \cref{thm:global-balance-inverse} computes the observable score class on its stated support.

The span of the two keys induces a marginal map \(m=(\pi_!,\tau_!)\) on edge weights. Write \(W_+(S)\) for the positive located weight fields on a finite carrier \(S\). At compatible fixed targets its positive solution space is the fibre
\[
\begin{tikzcd}[column sep=large]
 W_{\rm bal}\arrow[r,hook]\arrow[d]\arrow[dr,phantom,"\lrcorner" very near start]&
 W_+(E)\arrow[d,"m"]\\
 1\arrow[r,"{(\mu,\nu)}"']&W(I)\times W(J).
\end{tikzcd}
\]
Here \(W(I),W(J)\) are the ambient real weight fields, so summing an empty fibre is allowed. Balancing chooses the representative of a score orbit lying in this fibre. The existence and effective construction of that representative on the stated support domains are the content of \cref{thm:balance-exact,thm:global-balance-inverse,sec:sparse-solve}.
Use nonempty finite row and column tables $I,J$, a finite edge-occurrence span $I\xleftarrow{\pi}E\xrightarrow{\tau}J$, and positive target masses $\mu_i,\nu_j$ of the same total. This includes rectangular joins and repeated foreign keys. For the exact forward construction generate at least one positive row at every pair, with positive weights $w_e$. Set $K_{ij}$ to the sum of weights at the pair. Equal total is forced by summing the two marginal equations. An unrestricted real register generates all numerical inputs of this kind. Take $w_e=\exp(s_e)$, $M=\exp(m)$, $\mu=M\operatorname{softmax}(a)$ and $\nu=M\operatorname{softmax}(b)$. This map is a split epimorphism onto the positive located equal-total tables, with section $(\log w,\log\sum\mu,\log\mu,\log\nu)$. The positive logarithm is computed in \cref{sec:sparse-solve}. In the sequel the targets are held fixed during differentiation. Define two grouped lookups
\[
 Ac=\pi_!(w\,\tau^*c),\qquad Br=\tau_!(w\,\pi^*r).
\]
Read the target masses back at their own keys and alternate $r=\mu/(Ac)$, $F(c)=\nu/(Br)$. Divisions are pointwise. Every bounded number of passes is a finite query with tracked units. A fixed point gives $P_e=r_{\pi e} w_e c_{\tau e}$ with both marginals prescribed. Choosing the same token table at both ends, one edge per pair, unit marginals and $w_{ij}=\exp(s_{ij})$ recovers dense Sinkformer attention~\cite{sink}. For general targets, $\mu_i^{-1}\sum_{\pi e=i}P_e v_e$ is the corresponding conditional readout.

\begin{theorem}[Effective balancing on positive dense tables]\label{thm:balance-exact}
For positive located-real affinities and targets, the alternating grouped lookups determine a unique positive balanced table of the form $P_e=r_{\pi e}w_e c_{\tau e}$. Fixing one column scale to one makes both scales unique. The scales and the balanced table have uniformly computed Cauchy names on this input domain. More precisely, let $B$ be the positive dense lower-real input object with equal target totals. Let $S:B\to Q$ sum all edge weights and all entries of both target tables. The graph of positive located anchored scales and located balanced entries over $B$ is isomorphic over $B$ to $B\times_Q\mathsf U(Q)$.
\end{theorem}

\begin{proof}\leavevmode
\begin{enumerate}[titledsteps]
    \item \textbf{Recover located inputs.} On the pullback $S$ is a positive located unit. Normalize the finite concatenated list of inputs by $S^{-1}$. The complementary-cut proof of \cref{thm:attention-kernel} makes every normalized entry located; multiplication by $S$ recovers every input as a located real. Conversely, from the located outputs one gets $w_e=P_e/(r_{\pi e}c_{\tau e})$ and both target tables by summing $P$. Hence $S$ has a located reciprocal. 
    \item \textbf{Contract the anchored sweep.} For the forward construction choose strict rational bounds $0<a<K_{ij}<b$ by finite search. For positive column fields compare their largest and smallest coordinate ratios. Each row lookup averages these ratios and contains the same subdistribution of mass $a/b$ in every row. After two lookups their relative spread contracts by $(1-a/b)^2$; rescaling at the anchor does not change it. The explicit tail and the verification of both marginal equations are given in \cref{app:limits}. This constructs the output uniformly and uniquely. Uniqueness of the anchored output and of $S^{-1}$ proves both inverse laws over $B$ and stability under base change. For fixed located targets the sum of just the affinities gives the same pullback. Every exponential-score table on a generated dense mask lands there by finite real arithmetic. 
    \item \textbf{Test the effective boundary.} A lower name alone cannot be promoted by a promise about its denotation. For unit marginals take the $2$-by-$2$ table $(m,1;1,1)$ with $m=1+\tau(\operatorname{Halts}(e,e))$. By \cref{ex:root-cover} its numerical diagonal is $1/2$ or $\sqrt2/(1+\sqrt2)$ according to nonhalting or halting. Comparing a purported uniform enclosure against $11/20$ would decide diagonal halting. This example even has known entry bounds.\qedhere
\end{enumerate}
\end{proof}

\subsection{Marked marginal equations and elimination}

Orient an edge from its row to its column. On vertex fields set $d\phi(e)=\phi_{\pi e}-\phi_{\tau e}$. On edge fields let $\partial$ group by row with sign $+$ and by column with sign $-$. Write $h_e=\dot w_e/w_e$ and anchor $\phi_o=0$. Marking the three sockets of $P_e$ gives
\[
 \dot P_e=P_e(h_e-d\phi(e)),\qquad
 \partial(P\,d\phi)=b,\qquad b=\partial(Ph).
\]
Here $\phi_i=-\dot r_i/r_i$ and $\phi_j=\dot c_j/c_j$ for row and column keys respectively. The second equation just says that marking either fixed target mass returns zero.

Aggregate the positive edge entries into symmetric conductances \(c_{uz}\). At a nonanchor vertex \(v\) let \(D_v=\sum_u c_{uv}\), \(c=(c_{uv})_{u\ne v}\). List the other vertices first, leaving the anchor equation in the system for now. In block form one deletion is
\[
 \begin{pmatrix}A&-c\\-c^{\mathsf T}&D_v\end{pmatrix}
 \binom{\phi_{\rm rem}}{\phi_v}=\binom{b_{\rm rem}}{b_v}
 \ \Longleftrightarrow\
 \begin{cases}
 (A-cc^{\mathsf T}/D_v)\phi_{\rm rem}= b_{\rm rem}+cb_v/D_v,\\
 \phi_v=(b_v+c^{\mathsf T}\phi_{\rm rem})/D_v .
 \end{cases}
\]
Thus for \(u\ne z\) the new conductance is \(c'_{uz}=c_{uz}+c_{uv}c_{vz}/D_v\). The new right side sums to \(\sum_{u\ne v} b_u+b_v=0\). A path through \(v\) gives a positive shortcut between its distinct neighbors, so the remaining graph stays connected. Each deleted vertex has a positive pivot. At the last vertex the zero right side is the anchor equation, and \(\phi_o=0\) initializes reverse substitution. Each displayed equivalence recovers the saved coordinate uniquely. This computes the marked tangent, including on connected sparse support at the positive solution constructed in \cref{sec:sparse-solve}.

\subsection{Reverse receipts and independent cycle directions}

Write $\mathcal R(h)=h-d\phi(h)$. The push--pull identity
$\sum_e P_e \mathcal R(h)_e d\psi(e)=\sum_v\psi_v\partial(P \mathcal R(h))_v=0$
is regrouping by endpoints. Replacing $g$ by $\mathcal R(g)$ and then $\mathcal R(h)$ by $h$ in the weighted pairing proves
\[
 \sum_e g_e\,\dot P_e=\sum_e P_e \mathcal R(g)_e h_e.
\]
Thus the same elimination run on the seed gives the receipts $P_e\mathcal R(g)_e$ at log-weight sockets; divide by $w_e$ at raw-weight sockets. Final delivery sums tied addresses. Endpoint-potential perturbations disappear, and perturbations with $\partial(Ph)=0$ pass through as $Ph$.

\paragraph{Which changes can balancing express?}
At a positive balanced solution with independently variable edge weights, the possible variations are exactly the signed flows $f$ with $\partial f=0$. One inclusion is the marked constraint. Conversely choose $h_e=f_e/P_e$. The solve then has right side zero, so its unique anchored output is zero and $\dot P=f$. The kernel consists exactly of $d$ of vertex potentials, since every $P_e$ is a unit. Thus marking induces a canonical isomorphism
\[
       R^E/d(R^{I+J})\ \longrightarrow\ \ker(\partial:R^E\longrightarrow R^{I+J}),
       \qquad [h]\longmapsto P \mathcal R(h),\quad f\longmapsto[f/P]\ \text{in reverse},
\]
where $R$ is the located real ring and $\mathcal R(h)$ denotes the residual defined above. Choose a spanning tree by scanning the enumerated edges and keeping each edge that joins distinct components. Each edge outside this tree closes a unique cycle. Orient that cycle to have coefficient $+1$ on the extra edge, and use the signs of traversal on the tree path. Subtract these cycle flows with coefficients equal to the respective extra-edge values of $f$. What remains lives on a tree and has zero signed sum at each vertex. Removing leaves forces its edge values to vanish one by one. This gives a basis indexed by the extra edge IDs, including length-two cycles for parallel rows. In particular a tree support has zero attention-weight freedom at fixed marginals; tied scores can further reduce the attainable directions on any support.

The same proof gives a finite layer with \emph{exact} marginal preservation. Start from a balanced table $P^0$. On dense rectangular data one is constructed by splitting $\mu_i\nu_j/\sum_i\mu_i$ equally among the occurrences at $(i,j)$. For each extra edge $a$ let $\sigma^a$ be its signed cycle and let a trainable cell produce a coefficient $t_a$. Return
\[
        P(t)=P^0+\sum_a t_a\sigma^a.
\]
This is a single affine occurrence query. Since $\partial\sigma^a=0$, both budgets hold for every coefficient and every input. Conversely the preceding leaf argument applied to $P-P^0$ gives unique coefficients for every table with those budgets. The inequalities $P(t)_e>0$ specify exactly the positive part of this parameter space. To construct an everywhere-positive trainable example, find a rational $0<\gamma<\min_e P_e^0$; with $m$ extra edges set $t_a=\gamma\,\varphi(u_a)/(1+m)$, where $\varphi(u)=u/(1+u^2)$ has absolute value at most $1/2$. Every entry then stays above $\gamma/2$. The bounded trainable subfamily need not cover the whole positive region. The balancing score layer and the cycle-coordinate layer parameterize the output differently: the former takes independent affinity scores and computes the balancing response, while the latter takes admissible circulation coordinates and evaluates by a finite sum.

On the $2$-by-$2$ table of \cref{ex:root-cover} write $\sigma=(+1,-1,-1,+1)$ in row order. Every tangent is $\dot P=t\sigma$. Since $\sum_e\sigma_e d\phi(e)=0$, the marked formula yields $t=(\sum_e\sigma_e h_e)/(\sum_e1/P_e)$. Thus there is one scalar degree of freedom and the sign of every compensating change is determined. On a general connected support the corresponding count is $|E|-|I|-|J|+1$ at independent logits. Figure~\ref{fig:balanced-cycle} gives the signed incidence pattern.

\begin{figure}[!htbp]
    \centering\small
    \begin{tikzpicture}[x=1cm,y=1cm]
        \node[figbox] (i1) at (0,.65) {$i_1$};
        \node[figbox] (j1) at (2.7,.65) {$j_1$};
        \node[figbox] (i2) at (2.7,-.70) {$i_2$};
        \node[figbox] (j2) at (0,-.70) {$j_2$};
        \draw[figwire] (i1) -- node[above] {$+t$} (j1);
        \draw[figwire] (i1) -- node[left] {$-t$} (j2);
        \draw[figwire] (i2) -- node[right] {$-t$} (j1);
        \draw[figwire] (i2) -- node[below] {$+t$} (j2);
        \node at (6.5,0) {$\displaystyle
          \dot P=\begin{array}{c|rr} &j_1&j_2\\\hline
             i_1&+t&-t\\ i_2&-t&+t\end{array}$};
    \end{tikzpicture}
    \caption{The four edges of a $2$-by-$2$ support and their signed variations in row order. At each row vertex the outgoing changes sum to zero. At each column vertex the incoming changes also sum to zero, so the same cycle preserves both marginals.}
    \label{fig:balanced-cycle}
\end{figure}


\begin{proposition}[The anchored balancing chart]\label{prop:balanced-lift}
On the open defined by the pivot units of a vertex elimination, the anchored scale equations define a finitely presented \'etale morphism over the affine weight chart. Every positive dense solution lies on this open. For a fixed initial scaling, jets of the scales over an affinity arc exist uniquely in every finite order.
\end{proposition}

\begin{proof}\leavevmode
\begin{enumerate}[titledsteps]
    \item \textbf{Present the anchored chart.} Use universal targets with the equal-total relation, then base change to the required targets. Adjoin unit row and column scales, fix one column scale to one, and impose all marginal equations except the redundant anchor equation. Invert the successive pivot elements of the chosen deletion order. This uses only finitely many generators, relations and inverses.


    \item \textbf{Lift the scales through a square-zero extension.} Let \(B'\twoheadrightarrow B\) have square-zero kernel \(J\). Lift the finitely many retained scale units provisionally, with the anchor still one. Write \(\widetilde P\) for the resulting edge entries, \(D\) for the anchored incidence, and \(e=D\widetilde P-\widehat b\in J^{V\setminus\{o\}}\) for the error against the given lifted targets. Relative corrections multiply a row scale by \(1-\delta_i\) and a column scale by \(1+\delta_j\), with \(\delta_o=0\). Since \(J^2=0\),
\[
 P'=\widetilde P-\operatorname{diag}(\widetilde P)D^{\mathsf T}\delta,\qquad
 DP'=\widehat b
 \quad\Longleftrightarrow\quad
 D\operatorname{diag}(\widetilde P)D^{\mathsf T}\delta=e .
\]
The pivot units in \(B\) lift to units in \(B'\). Elimination therefore gives the unique \(\delta\) in the \(J\)-module, and the inverted elements have unique inverses. This proves the square-zero lifting property and finite presentation gives the claimed \'etale morphism \cite{etale}. At a positive solution all pivots are positive. Successive lifts through \(B[t]/(t^{j+1})\twoheadrightarrow B[t]/(t^j)\) give each jet by \cref{thm:etale-jets}. Different pivot orders on an overlap give the unique lift of the same point. Reanchoring at \(o'\) sends \((r,c)\mapsto(r c_{o'},c/c_{o'})\) and keeps every entry \(r_i w_e c_j\) fixed. The two-by-two case gives \cref{ex:root-cover}. The residual estimates of \cref{app:limits} identify the positive numerical jets.\qedhere
\end{enumerate}
\end{proof}

\subsection{Recovering the scaling equations through mixed order}
Let \(V=I+J\) be the vertex table and \(o\) the anchor. Use the same orientations in the global equations, the jet solver and the sparse forward solve (\cref{sec:sparse-solve}):
\[
 \partial:\R^E\to\R^V,\quad
 D:\R^E\to\R^{V\setminus\{o\}},\quad
 D^{\mathsf T}:\R^{V\setminus\{o\}}\to\R^E .
\]
Here \(D\) is \(\partial\) with its anchor row removed. Extend a potential by zero at \(o\). With \(s=\log w\), the block call is the unique anchored solution of
\[
 \partial P=(\mu,-\nu),\qquad
 \log P=s-D^{\mathsf T}\phi,\qquad \phi\in\R^{V\setminus\{o\}} .
\]
Indeed \(\phi_i=-\log r_i,\ \phi_j=\log c_j\) give the row and column scales of the dense or sparse solve. Marginals alone leave the continuum \((p,1-p;1-p,p)\) on a two-by-two unit-budget table. The orbit equation locates the computed section, whose inverse to graph projection follows from \cref{thm:solved-composition}. Anchor each component on its equal-total base.

For two independent logit directions $h,k$ keep targets fixed. Set
\[
 W=\operatorname{diag}(P),\qquad L=DWD^{\mathsf{T}},\qquad
 Q_P=1-D^{\mathsf{T}}L^{-1}DW .
\]
For an anchored vertex vector $v$, $v^{\mathsf{T}}Lv=\sum_eP_e((D^{\mathsf{T}}v)_e)^2$. Vanishing forces agreement along every edge and hence zero by the anchor. The finite positive-pivot elimination in this section computes $L^{-1}$. Marking the log equation in direction $h$ and imposing zero marginal derivative gives $L\phi_h=DWh$, so $P_h=WQ_Ph$. Mark it again with two distinct variables. The derivative of the exponential gives
\[
 P_{hk}=W\big((Q_Ph)\odot(Q_Pk)-D^{\mathsf{T}}\phi_{hk}\big),\qquad
 P_{hk}=WQ_P\big((Q_Ph)\odot(Q_Pk)\big).
\]
Indeed $DP_{hk}=0$ determines $\phi_{hk}$ by the same solve. Also $WQ_P=W-WD^{\mathsf{T}}L^{-1}DW$ is symmetric, giving the logit adjoint from the very same field computation. At each higher total degree the unknown potential coefficient enters linearly via $-WD^{\mathsf{T}}$. Exponentiating the already known lower-degree potential and input coefficients gives the source term, and applying $WQ_P$ to its unweighted value enforces the marginals. This gives a finite recursion from the selected base and agrees with the jets of the selected numerical section by the residual estimates in \cref{app:limits}.

On two rows and columns take affinity matrix $(e^s,1;1,1)$ with unit marginals. Its cross ratio is preserved by scaling, hence $(p/(1-p))^2=e^s$ for either diagonal entry $p$. Positivity selects $p=e^{s/2}/(1+e^{s/2})$. Direct differentiation gives $p'=p(1-p)/2$, $p''=p(1-p)(1-2p)/4$. At $s=\log4$ these give $2/3$, $1/9$, and $-1/54$. The mixed calculation above includes this repeated-direction second derivative by taking $k=h$.

\begin{proposition}[Behavioral inverse for balanced attention]\label{thm:global-balance-inverse}
On the fixed feasible support and targets, the balanced value operation with an independently variable edge payload determines the orbit of its log affinities modulo endpoint potentials. Positive feasible tables $P$ and these score orbits are in effective bijection.
\end{proposition}
\begin{proof}
With \(D\) the incidence from edges to unanchored vertices, \(W_{\rm bal}\) the positive tables with the specified marginals, and \(T_P\) the conditional value operator, the inverse identities are
\[
\begin{aligned}
 P_e&=\mu_{\pi e}(T_P e_e)_{\pi e},\\
 \operatorname{Bal}(\exp s)=\operatorname{Bal}(\exp s')
   &\ \Longleftrightarrow\ s-s'\in\operatorname{im}D^{\mathsf T},\\
 \R^E/\operatorname{im}D^{\mathsf T}
   &\ \underset{P\mapsto[\log P]}{\overset{[s]\mapsto\operatorname{Bal}(\exp s)}{\rightleftarrows}}\ W_{\rm bal}.
\end{aligned}
\]
The first line probes the basis payload at occurrence \(e\). For the forward implication in the second line subtract the two equations \(\log P=s-D^{\mathsf T}\phi=s'-D^{\mathsf T}\phi'\). Conversely, translate the anchored potential by the potential difference and apply uniqueness of the balanced table. For \(P\in W_{\rm bal}\), scores \(\log P\) have unit scales, proving the inverse identities. To compute a representative, solve \(D_T^{\mathsf T}\psi=(\log P)|_T\) along a spanning tree \(T\), taking anchor value zero. Then \(\log P-D^{\mathsf T}\psi\) vanishes on \(T\). Each chord value is its signed fundamental-cycle sum, including length-two cycles of parallel occurrences. Apply this construction separately on each feasible component.
\end{proof}

\paragraph{Observability with grouped edge payloads.}
On the direction space of the marginal fibre, grouping payloads is the restriction of a linear map. Write \(V=\R^E\), let \(D\) be the reduced signed marginal map, and write \(G\) for grouping by row and payload key. Its invisible directions form the upper left corner:
\[
\begin{tikzcd}[column sep=large]
 \ker D\times_V\ker G\arrow[r,hook]\arrow[d,hook]&\ker G\arrow[d,hook]\\
 \ker D\arrow[r,hook]&V .
\end{tikzcd}
\]
The restricted observation \(G|_{\ker D}\) is monic exactly when this pullback is zero. Over a positive feasible table every rational direction in the pullback is realized by a short two-sided line segment in the fibre. The rank test and the two indistinguishable tables below compute this criterion in coordinates.
The independent-payload requirement admits an exact test. Suppose an edge payload is read through \(k:E\to K\), with the row of the output retained. Put
\(G:\R^E\to\R^{I\times K}\), \((Gf)_{i,a}=\sum_{\pi e=i,\,k(e)=a}f_e\). Basis probes and the known \(\mu_i\) recover \(GP\). For two positive balanced tables of the same marginals,
\[
 T_P k^*=T_{P'} k^*\quad\Longleftrightarrow\quad GP=GP',
 \qquad P-P'\in\ker D.
\]
Consequently these probes recover the score orbit for every positive balanced table precisely when \(\ker D\cap\ker G=0\). For the converse direction, row reduction of the integer matrix \(\bigl(\begin{smallmatrix}D\\G\end{smallmatrix}\bigr)\) yields a nonzero rational kernel vector \(f\) when its column rank is smaller than the number of edges. At a feasible \(P^0\), take a rational
\(0<\delta<\min_e P^0_e/(1+|f_e|)\). Then \(P^\pm=P^0\pm\delta f\) are distinct positive tables with the same probes and marginals. Their scores \(\log P^\pm\) represent distinct orbits by \cref{thm:global-balance-inverse}. Two parallel occurrences at the same payload key give the particular invisible flow \((1,-1)\) on those edges. This test uses only the finite incidence and payload-key tables.
\section{Transport of graph messages and tensor features}\label{sec:arch}

The orbit coequalizer in \cref{sec:fixed-shape-orbits} classifies the parameter keys for equivariant linear fields. The computations below transport edge messages through that action and through maps of incidence tables. Lookup along each incidence key and summation at each receiving row interpret the corresponding field terms of \cref{sec:field-language}. The row kernel is the one constructed in \cref{thm:attention-kernel}.

\subsection{Message passing and graph attention}
Let the graph consist of finite vertex and edge carriers and foreign keys $s,t:E\to V$. At $e$ a shared cell reads $x_{s(e)},x_{t(e)}$, edge attributes, and parameters to produce a message $m_e$ and a nonnegative weight $w_e$. In the raw lower-weight interpretation messages have nonnegative coordinates. A finite numerical realization also allows located signed values. Incoming aggregation constructs
\[
 (N_j,Z_j)=\sum_{e:t(e)=j}(w_em_e,w_e).
\]
Ordinary additive message passing reads $N_j$ at unit weights. Graph attention reads $y_j=Z_j^{-1}N_j$ on its inverse-mass object. A vertex cell may then combine $y_j$ with $x_j$. A linear graph convolution is obtained by taking $m_e=W x_{s(e)}$ and prescribed edge coefficients. Token attention uses the very same formula on the graph of eligible key--query pairs, with each edge directed from key to query.

\begin{proposition}[Transport and coarsening of messages]\label{thm:graph}
Pullback of an incoming-edge fibre transports its pair $(N,Z)$. Coalescing occurrences by an address map over $V$ also preserves $(N,Z)$ when each new row carries the two grouped contributions. Normalization in either presentation gives the same readout on their common inverse-mass domain.
\end{proposition}
\begin{proof}
For transport apply \cref{lem:bc} to the two additive coordinates. For coalescing, factor $t=t'c$ and use $t_!=t'_!c_!$ on the pair $(wm,w)$. The grouped numerator and mass require no division inside a group of zero mass. The outer mass is unchanged, so its inverse is the same. A graph isomorphism transporting all fields is a fibrewise bijection and gives equivariance by the first assertion. In particular this proof counts each parallel edge once.
\end{proof}
How much information must a grouped edge retain? Let $M$ be the additive carrier of messages and $K$ the mass carrier, with scalar multiplication as above. If $c:E\to C$ groups edges over vertices, the two paths below sum the same pair in $M\times K$, either directly at the vertex or through the group address.
\[
\begin{tikzcd}[column sep=large]
 (M\times K)^E\arrow[r,"c_!"]\arrow[dr,"t_!"']&
 (M\times K)^C\arrow[d,"t'_!"]\\
 &(M\times K)^V .
\end{tikzcd}
\]
At a vertex, apply the inverse of its total mass after this diagram. Thus vertices with equal transported data have equal outputs even if the grouped rows had different individual messages.

Successive layers can read the same graph $\Gamma$. Their composite $g(\Gamma,f(\Gamma,x))$ uses the product-environment comonad to supply the adjacency data to both evaluations, as in the coKleisli account of graph convolutions \cite{graph:cokleisli}. Polynomial spans likewise describe graph message passing \cite{graph:dp}. Here the edge rows also identify the occurrences that contributed to each normalization and the addresses of the parameters they read. We will use both pieces of information when running the computation backwards.

For the orbit key of \cref{thm:orbit}, \cref{lem:sharing} gives $\bar\theta_o=\sum_{(j,i)\in o}\lambda_j x_i$. For a graph row with $t(e)=j$, pairing \cref{eq:unnormdiff} with $\lambda_j$ gives $\bar m_e=(w_e/Z_j)\lambda_j$ and $\bar w_e=Z_j^{-1}\lambda_j(m_e-y_j)$. Higher variations of finite normalized graph rows follow \cref{prop:normalized-jets} on the same incidence.

\subsection{Continuous symmetry on the feature register}

For a feature representation \(\rho\), a pair of representations \(\rho_U,\rho_V\), and an output representation \(\rho_W\), the scalar and feature gates used here obey
\[
 \langle \rho_{U^*}(g)q,\rho_U(g)k\rangle=\langle q,k\rangle,\qquad
 C\bigl(\rho_U(g)u,\rho_V(g)v\bigr)=\rho_W(g)C(u,v).
\]
Transport the edge keys by \(e\mapsto ge\). Relative data and invariant scoring give \(s_{ge}(gx)=s_e(x)\), hence \(Z_{gi}(gx)=Z_i(x)\) by fibrewise reindexing. If the value obeys \(v_{ge}(gx)=\rho_W(g)v_e(x)\), the normalized sum satisfies
\[
 y_{gi}(gx)=\sum_{\pi e=i}\alpha_e(x)\rho_W(g)v_e(x)
            =\rho_W(g)y_i(x).
\]
Realize each contraction by joining its repeated index and summing the fibre. The following construction computes the required coupling tables for rotations in three dimensions.

For three-dimensional features take equality of coordinate keys for the metric $\delta$, and signed permutations of three keys for volume. Form $H_l$ by imposing permutation symmetry and vanishing trace on words of length $l$. The rational projection $\STF_l$ and a basis of all equivariant couplings $H_k\otimes H_J\to H_l$ are computed in \cref{app:angular}. The algorithm tests rotations represented by pairs of reflections and reduces the resulting finite coefficient equations over $\Q$. Thus the equivariant channels come with finite keyed contraction tables.

For an edge displacement $r=x_j-x_i$ of positive length, the angular feature in degree $J$ is $\STF_J(r^{\otimes J})/\|r\|^J$. Pairing with a trace-free symmetric tensor returns its homogeneous harmonic polynomial on the unit sphere. Indeed taking two formal marks in a homogeneous polynomial and summing equal marked directions extracts $J(J-1)$ times the trace; symmetry identifies coefficients with symmetric tensors over $\Q$. This identifies the angular functions up to basis. Apply the compiled $C$ to a type-$k$ feature and this angular feature, multiply by a scalar radial gate, and group the desired output type and channel. The Clebsch--Gordan channels of Fuchs et al.~\cite[Appendix B]{se3} are such intertwiners; the completeness argument expresses them in the compiled basis after identifying their harmonic bases. Conversely each compiled channel satisfies their kernel equivariance equation. Allowing the indicated degrees, radial functions and channel multiplicities, forming invariant query--key pairings (using the metric in the chosen basis), and summing the weighted values recovers the finite angular attention mechanism of the $\mathrm{SE}(3)$-Transformer up to basis. Scalar radial networks use their numerical gates on their ordinary domains. Translations cancel in $r$. An angular direction requires a nonzero displacement; a self-interaction or a regular solid-harmonic kernel can be specified at zero. A nearest-neighbour realization uses a presented incidence or the numerical domain deciding the neighbour set.
For a fixed group element the forward identity is \(f\rho_X(g)=\rho_Y(g)f\). Substituting a generic increment gives \(J_nf\,J_n\rho_X(g)=J_n\rho_Y(g)\,J_nf\). Functoriality in \cref{thm:jet-pairing} reverses the same square on pointed scalar observables, with their base points transported by \(g\).
\paragraph{A concrete polynomial head.}
A value channel which is regular even at coincident points is
\[
 S(r)=r\otimes r-\tfrac13\langle r,r\rangle\delta,\quad
 v_{ij}=h_j S(r),\quad
 k_{ij}=\psi(h_j,\langle r,r\rangle)r.
\]
Here $h_j$ and the polynomial $\psi$ are scalar and a vector query is paired with $k_{ij}$. The two free value keys are symmetric, and tracing gives zero because $\delta$ has trace three. Mark the two $r$ sockets. In direction $u$ the receipts read $u\otimes r+r\otimes u-\tfrac23\langle r,u\rangle\delta$. Both this table and the original one transform by the two free output keys. This is a type-two head with a concrete first jet. More generally the same pair-of-marks head from the preceding sections can carry intertwiner-valued payloads on structured witnesses whenever the incidence and weights transport with the symmetry.

\section{Initial algebras and equalizers for recurrent cells}\label{sec:equilibrium}

An unfolded recurrent cell uses the parameter table of \cref{sec:model-interfaces} at every iteration. An equilibrium call instead designates the solution value; its construction below supplies the graph section on the specified convergence domain, to which \cref{thm:solved-composition} applies.

\subsection{Recurrent and recursive networks}
\begin{proposition}[Folding a learned cell]\label{thm:fold}
For the signature $F(X)=1+X\times L$, a parameterized algebra on $H$ is exactly an initial state $h_0:P\to H$ and a cell $c:P\times H\times L\to H$. Its extension on closed terms is the unique recurrent evaluation on finite snoc-lists. For $F(X)=L+X^2$, the corresponding data are a leaf encoder $P\times L\to H$ and a binary combiner $P\times H^2\to H$. Its extension is the unique evaluation of finite leaf-labelled binary trees.
\end{proposition}
\begin{proof}
Distribute the external product over each displayed coproduct. In the tree case
$P\times(L+H^2)\cong(P\times L)+(P\times H^2)$. The label thus belongs to the leaf constructor (compare \cite[Examples I.2 and J.3]{cdl}). In both cases $\mathsf T_F0$ is the initial $F$-algebra. Attaching a root gives its structure map $\mathsf{in}:F(\mathsf T_F0)\to\mathsf T_F0$. Initiality provides a unique homomorphism from closed terms to the algebra $a_p$.

For a list, the homomorphism sends the empty list to $h_0(p)$ and satisfies $h_p(\operatorname{snoc}(u,x))=c(p,h_p(u),x)$. For a tree it applies the encoder to each leaf label and the combiner to the two child states at a node. For both recursions, the square identifies evaluation after constructing a term with application of the cell after evaluating the children.
\[
\begin{tikzcd}
 F(\mathsf T_F0)\arrow[r,"Fh_p"]\arrow[d,"\mathsf{in}"']&FH\arrow[d,"a_p"]\\
 \mathsf T_F0\arrow[r,"h_p"']&H .
\end{tikzcd}
\]
Each recursive call reads the same $p$. The finite constructors produce a register list by structural recursion, uniformly in the presented input shape.
\end{proof}

For a finite list write $h_k=c(p,h_{k-1},x_k)$, $A_k=D_hc$ and $B_k=D_pc$ at the retained values. \Cref{neural:thm:nilpotent} gives $\lambda_{k-1}=A_k^*\lambda_k$ and $\bar p=Dh_0(p)^*\lambda_0+\sum_k B_k^*\lambda_k$. A binary node sends the corresponding two receipts to its children. We now allow the marked hidden path to return to a previous key.

\subsection{A trainable contraction on a hidden incidence}

For a parameterized cell \(F:S\times Z\to Z\), the object of fixed points over its external data is the equalizer
\[
\begin{tikzcd}[column sep=large]
 E_F\arrow[r,hook,"e"]\arrow[dr,"p"']&S\times Z\arrow[r,shift left=.65ex,"F"]\arrow[r,shift right=.65ex,"\pi_Z"']\arrow[d,"\pi_S"]&Z\\
 &S .
\end{tikzcd}
\]
A solver over \(S\) gives a section of \(p\). When it selects the unique fixed point in each fibre, it exhibits \(p\) as an isomorphism by the graph argument of \cref{thm:solved-composition}. For \(r:S'\to S\), pullback gives \(E_{F(r\times1)}\cong S'\times_S E_F\) by preservation of the defining equalizer in the slice. Thus an equilibrium port with its solver can be reparameterized before or after the solve.

\paragraph{A contraction realizing the equalizer.}

\begin{enumerate}[titledsteps]
    \item \textbf{Construct a contractive cell.} Choose finite directed hidden incidences with parallel edges, external input and parameter cells \(u_i(s)\), and a rational \(0<\lambda<1\). For a row \(i\) of nonzero in-degree \(d_i\), put
\[
 \varphi(t)=t/(1+t^2),\quad
 a_e^{\rm eff}=2\lambda\,\varphi(a_e)/d_i,\quad
 F_i(s,z)=\varphi\!\left(u_i(s)+\sum_{e:\,t(e)=i}a_e^{\rm eff}z_{s(e)}\right).
\]
Here $s$ outside a subscript denotes the external data; $s(e)$ is the source key. For an empty row the sum is empty and no coefficient is divided by zero. For first variations the $u_i$ may be any of the differentiable finite cells on their domains. For the smooth numerical syntax of this paper its generic-increment compiler constructs their coefficients in every finite order. These two statements differ even for an empty hidden incidence: $u(s)=s|s|$ gives the equilibrium $\varphi(u(s))=s|s|/(1+s^4)$, with derivative zero at the origin and different one-sided limits of $\varphi(u(s))/s^2$. Parallel uses may have the same parameter foreign key. The identities
\[
 |\varphi|\le\tfrac12,\quad
 \varphi(t)-\varphi(v)=\frac{(t-v)(1-tv)}{(1+t^2)(1+v^2)},\qquad
 (1+t^2)(1+v^2)-(1-tv)^2=(t+v)^2
\]
give $|\varphi(t)-\varphi(v)|\le |t-v|$. Each row has total absolute effective weight at most $\lambda$. Thus iterating from zero has tail bounded by $\lambda^N\|z_1\|_\infty/(1-\lambda)$ after $N$ steps and constructs the unique equilibrium. A complete hidden incidence, $\lambda=1/2$, and affine $u_i$ recover a dense nonlinear trainable instance of the equilibrium architecture of Bai et al.~\cite{deq}. All choices of its real trainable coefficients meet the computed bound.


    \item \textbf{Solve the adjoint equation.} At the selected fixed point \(\zeta(s)\), differentiating \(F\langle1,\zeta\rangle=\zeta\) gives
\[
 (1-H)\,D\zeta=A,\qquad
 D\zeta^*=A^*(1-H^*)^{-1},
 \qquad H=D_zF,\quad A=D_sF .
\]
The inverse is realized by the contraction bound on this cell. This describes the reverse lens of the selected section without choosing a basis. In the displayed registers put \(L=H^{\mathsf T}\), \(B=A^{\mathsf T}\). Marking \(y(1+t^2)=t\) gives \(\varphi'(t)=(1-t^2)/(1+t^2)^2\), hence \(\|H\|_\infty=\|L\|_1\le\lambda\). Let \(\beta_N=\sum_{k=0}^{N-1}L^k g\). The path and convergence calculations are
\[
 \beta_{N+1}=g+L\beta_N,\qquad
 \|\beta-\beta_N\|_1\le\frac{\lambda^N}{1-\lambda}\|g\|_1,\qquad
 \beta=g+L\beta,\qquad \bar s=B\beta .
\]
An entry of \(BL^k\) is a sum of products along \(k\) hidden-to-hidden edges followed by one external exit. Induction follows by matrix multiplication at the shared vertex. The external tail is at most \(\|B\|_1\lambda^N\|g\|_1/(1-\lambda)\). If \(\beta,\beta'\) both solve the adjoint equation then \(\|\beta-\beta'\|_1\le\lambda\|\beta-\beta'\|_1\), proving uniqueness.


    \item \textbf{Recover the higher coefficients.} The finite compiler in \cref{neural:thm:nilpotent} sums paths up to the circuit depth. Here the path sum converges with a rate constructed from the cell. Through each order $m$ compiled for the external cells, let $b_m$ be the coefficient found by running the body with its unknown hidden coefficient set to zero. The marked substitution rule gives $z_m=b_m+Hz_m$ with the same one-mark hidden action $H$ at every order. \Cref{app:limits} identifies these formal coefficients with the jets of the selected real equilibrium and proves the reverse formula by pairing finite path sums before taking their limits.
\end{enumerate}

\subsection{Reading back an equilibrium and its mixed query}
Consider the scalar cells $F_0(s,z)=s$ and $F_1(s,z)=(s+z)/2$. Their Lipschitz constants in the hidden state are zero and $1/2$. Solving either fixed-point equation gives $z=s$ for every input, although the cells have different hidden-state derivatives. Even exact knowledge of the equilibrium map at every input therefore leaves the body undetermined.

For the presented cell the symbolic readback is the relation $z=F(s,z)$. The displayed convergence rate computes its section, and uniqueness gives the inverse to projection in \cref{thm:solved-composition}. Write $A=D_sF$ and $H=D_zF$ at that section. For two external directions $h,k$ the source queries are
\[
 z_h=(1-H)^{-1}Ah,\qquad
 z_{hk}=(1-H)^{-1}D^2F[(h,z_h),(k,z_k)] .
\]
Twice marking $F(s,z(s))$ leaves one further term $Hz_{hk}$, which moves to the left. The inverse is the computed geometric sum, bounded by $(1-\lambda)^{-1}$. Pairing its finite partial sums against a seed $g$ gives $A^{\mathsf{T}}\sum_{j\ge0}(H^{\mathsf{T}})^jg$ by the tail bound. In any higher degree, evaluate the body on all previously computed coefficients and put the unknown top coefficient equal to zero. Inserting that coefficient $v$ afterwards adds exactly $Hv$, since its product with any positive-degree increment vanishes at this order. Solve with the same inverse. Finally $\|z-\tilde z\|\le\|F(s,\tilde z)-\tilde z\|/(1-\lambda)$ by the triangle inequality. A polynomial lift whose residual is $o(\|t\|^m)$ therefore represents the actual selected solution through order $m$.

\subsection{Partial iteration of a discrete neural transition}

Extend the structural language by an exact natural-number state sort and a partial first-order loop. Its operational meaning is iteration until a decidable halt flag is set. In the effective topos the graph of a terminating computation is witnessed by a finite execution trace. Each finite prefix can be checked effectively. To see that the target can implement this control numerically, use a deterministic two-counter machine with increment-and-jump, conditional decrement with two successors, and halt \cite{counter}. Give the recurrent field register two nonnegative integer coordinates $n_0,n_1$ and a one-hot bit $b_j$ for each program location. These are represented exactly as rational numbers. For a decrement instruction $j$ acting on counter $k$, compute
\[
 z_k=\operatorname{ReLU}(1-n_k),\qquad
 u_j=\operatorname{ReLU}(b_j+z_k-1),\qquad v_j=b_j-u_j .
\]
On the encoded states these are the zero test, the active zero branch, and the active positive branch. Add $u_j,v_j$ to their respective successor bits and subtract $v_j$ from $n_k$. An increment contributes $b_j$ both to its successor bit and to its counter increment. A halt bit feeds itself. Compute the new counters by adding all indicated changes to the old counters and compute each new bit by summing its incoming signals. All weights and biases are rational integers. ReLU and affine wiring implement one fixed cell per program.

Let \(\mathrm{enc}\) encode a machine configuration by the two integer coordinates and its one-hot instruction. The gates above give
\[
 F\mathrm{enc}(c)=\mathrm{enc}(\mathrm{step}(c)),\qquad
 F^n\mathrm{enc}(c)=\mathrm{enc}(\mathrm{step}^n(c)).
\]
Indeed exactly one branch signal is one, and \(v_j=1\) requires a positive counter. The halt bit is absorbing. For an encoded input \(x\), write \(h_n(x)\) for the bit after \(n\) steps. Put
\[
 D_n=\{x:h_n(x)=1,\ \forall j<n\ h_j(x)=0\},\qquad
 D=\{x:\exists n\ h_n(x)=1\}\cong\coprod_{n\ge0}D_n .
\]
Each \(D_n\) is decidable on these encodings. Search for the first hit computes the inverse over \(D\). On \(D_n\) the result is the decoded \(F^n(x)\), giving step-for-step agreement of the partial programs. A universal two-counter machine gives Turing completeness with unbounded exact counters and time \cite{counter}. Conversely each rational affine--ReLU step and halt test can be simulated on a Turing machine, so integer-encoded partial results are partial recursive. A terminating trace exposes a finite computation. Continuous subcomputations on fixed differentiable branches have their usual readback.

\section{Coproducts of guarded expert branches}\label{sec:experts}

Start with a finite family of polynomial programs $P_b$ with a common output interface and branch-indexed expert parameters. Select a finite subtable of a token's expert occurrences by positive evidence or decidable codes. Pull the expert family back along the selected keys and evaluate the resulting diagram. On each domain where the selected key subtable and capacity decisions are fixed, substitution and marking give its derivatives to every order supported by the branch gates. Input-dependent incidence is the new operation. The gate values used to score the experts can still depend on all the logits.
For a fixed strict branch, let \(j_A:U_A\hookrightarrow U\) be its inclusion. Write \(Y\) for the output register and \(e_A:U_A\to Y\) for the selected computation (including token and parameter data in the base). The strict-winner cover below exhibits \(U\cong\coprod_A U_A\). The universal property gives one dispatch arrow \(e:U\to Y\) and its restrictions:
\[
\begin{tikzcd}[column sep=large]
 U_A\arrow[r,hook,"j_A"]\arrow[dr,"e_A"']&
 U\arrow[d,dashed,"e"]\\
 &Y .
\end{tikzcd}
\]
For branch data depending on the base the same copairing is taken in its slice. Restricting along a new input reindexes each summand by pullback. On each strict open, the returned cotangent is therefore that of the chosen branch by restriction to that open. The following comparisons construct the coproduct decomposition and its effective domain.

For a finite expert table \(I\) and a router selecting \(k\) experts, define the open
\[
 U_A=\bigcap_{i\in A,\,j\notin A}\{s:s_i>s_j\},
 \quad A\subseteq I,\ |A|=k,\qquad U=\bigcup_{|A|=k}U_A .
\]
For \(A\ne A'\) of size \(k\), take \(i\in A\setminus A'\) and \(j\in A'\setminus A\). The joint test includes \(s_i>s_j>s_i\), so \(U_A\cap U_{A'}=\varnothing\). Dovetailing strict rational-enclosure comparisons computes the inverse of \(\coprod_A U_A\to U\). Ties within \(A\) do not affect these tests. For \(k=0\) or \(k=|I|\) the sole test is vacuous. With tokens ordered by their presented IDs, the zero-based rank and admission bit at expert \(i\in A_n\) are
\[
 R_{n,i}=\sum_{m<n}\mathbf1_{i\in A_m},\qquad
 \operatorname{admit}_{n,i}=\mathbf1_{R_{n,i}<\operatorname{cap}_i}.
\]
The overflow branch uses the complementary bit. A finite ordered score-code table gives a total route by lexicographic comparison of score and tie priority.

With $k=1$, select an expert $i$, let $p_j=\exp(s_j)/\sum_d\exp(s_d)$ and return the expert contribution $p_i E_i(x)$ for an admitted token. Adding the block's residual once and dropping the expert contribution on overflow gives Switch dispatch~\cite{switch}. Marking its numerator and its common mass gives
\[
 \dot p_i=p_i\Big(\dot s_i-\sum_j p_j\dot s_j\Big),\quad
 \bar E_i=p_i g,\qquad
 \bar s_j=p_i\langle g,E_i(x)\rangle\big(\mathbf1_{i=j}-p_j\big).
\]
Here $g$ seeds the expert contribution and the indicator compares finite keys. Reverse the selected expert body and group parameter receipts at their stored keys. This also explains exactly how an unselected logit receives a receipt. A finite-code execution has a discrete branch key; these numerical derivatives describe variations of continuous gate and expert values on a fixed branch.

There is an internal obstruction to an everywhere-total exact deterministic top-one tie rule on located reals. For program $e$ use scores $(0,h_e)$, where a name for $h_e$ simulates $N$ steps at precision $N$, returning $2^{-t-1}$ if halting is observed at step $t$ and zero otherwise. Its error is at most $2^{-N}$. A total router choosing the first expert on equality would select the second exactly on the diagonal halting set. Relabelling gives the same obstruction for the opposite tie policy. Thus a strict real branch and a total finite-code dispatch have different internal domains. These real observation opens use order; a principal algebraic unit open alone does not choose a sign.

Equivalently, restrict the numerical program to the witness of a fixed route, form its pointed jet by \cref{thm:joint-jet}, and deliver its mixed parameter monomials to the selected expert and the router keys they involve. Restriction to another input in the same strict-winner domain commutes with this construction. The integer batch order and capacity remain discrete parameters in this comparison.

Read back dispatch as the tagged sum of the guarded solution graphs, with tag given by the selected subset. Over the union of strict-winner domains the guards are disjoint, and the finite parallel search above supplies their unique tag. The tag and its selected expert value give an inverse to graph projection there. On the finite ordered-code sort, the same construction uses its decidable tie rule. Changing the real logits inside one winning region transports the readback of the whole guarded computation, including its softmax over the full denominator.

Consequently that graph inverse interprets the guarded call in the architecture model on the stated domain. Pulling back its selected-tag graph along an upstream field term computes the legal domain of the composite by \cref{thm:solved-composition}.

\clearpage
\part{Coalgebras and learning dynamics}\label{part:learners}

\section{Learner coalgebras and parameter dynamics}\label{sec:learning}
A backward pass ends with a covector at each parameter address. To begin a new pass we need a state transition that turns this information into new parameter values. We first give that transition an interface accepting covectors directly. A learner in the sense of Fong, Spivak and Tuy\'eras passes requested targets between components. We then construct the change of interface and show that a request transmits exactly the covector computed by the downstream component.

\paragraph{The parameter fibre specified by observations.}\phantomsection\label{sec:fit-fibre}
Fix a field model, and write \(P=M\Theta\), \(X=MX_{\rm in}\), \(Y=MY_{\rm out}\) for the interpreted records. The boundary selected in \cref{prop:relative-models} evaluates as \(F:P\times X\to Y\). A finite list of observations with index set \(D\) consists of \(a:D\to X\) and a target tuple \(y:1\to Y^D\). Evaluation on the inputs gives
\[
 e_a:P\longrightarrow Y^D,\qquad e_a(p)_i=F(p,a_i),
 \qquad P_{a,y}=P\times_{Y^D}1 ,
\]
with the pullback taken along \(y\). For every \(T\), maps \(T\to P_{a,y}\) are precisely the parameter families \(p:T\to P\) satisfying \(F(p(t),a_i)=y_i\) for every \(i\), by the pullback and finite-product properties. On a mixed legal domain first restrict \(P\) to the intersection of the pullbacks of \(\mathcal L\hookrightarrow P\times X\) along all \((1_P,a_i!)\). This makes each component of \(e_a\) well typed. For empty \(D\) the product is terminal and the fibre is \(P\).

Over the rationals, a two-address gate \(F((p,q),x)=(p+q)x\) observed at \(x=1\) with output \(2\) gives \(p+q=2\). Sharing the addresses is the diagonal \(k(\theta)=(\theta,\theta)\). The fibre of the shared gate is the pullback along \(k\), namely \(2\theta=2\) with unique solution \(\theta=1\). In general \(e_a^{F\circ(k\times1)}=e_a^F k\) proves the same base-change formula. The observation specifies the exact-fit tables. To compare all tables on the same data, we next retain the observed target in the comparison of losses.

\subsection{Fitting as a Kan extension over the observed target}\label{sec:fit-kan}
The coherent reduct in \cref{prop:coherent-port}, evaluated on the sample \(a\), gives the actual map \(e_a:P\to B\), where \(B=Y^D\). To compare approximate fits, equip observations with a prediction error \(E:B\times B\to\mathbb R_{\ge0}\) of zero diagonal, in the sense of \cite[Definition 3.1]{pugh-kan}. For a target \(y\), give \(B\) and \(P\) the following thin categories, denoted by \(\mathcal D_y\) and \(\mathcal M_y\):
\[
\begin{aligned}
 z\longrightarrow z'\text{ in }\mathcal D_y
       &\quad\Longleftrightarrow\quad E(y,z)\le E(y,z'),\\
 p\longrightarrow p'\text{ in }\mathcal M_y
       &\quad\Longleftrightarrow\quad E(y,e_a(p))\le E(y,e_a(p')).
\end{aligned}
\]
Composition and identities are supplied by transitivity and reflexivity. The second order is pulled back from the first, so evaluation gives a fully faithful functor \(e_a:\mathcal M_y\to\mathcal D_y\). In particular distinct parameter tables remain distinct objects. Its arrows record performance on the specified observations.

\begin{theorem}[Universal fitting tables over a target]\label{thm:fit-kan}
Over each target \(y\), the universal arrows from \(y\) to \(e_a:\mathcal M_y\to\mathcal D_y\) are precisely the global minimisers, including all ties. They are also precisely the left Kan extensions of \(y\) along the inference interface \(\iota\) in the following 2-category:
\[
 \mathcal T_y(0,1)=\{\iota\},\qquad
 \mathcal T_y(1,2)=\mathcal M_y,\qquad
 \mathcal T_y(0,2)=\mathcal D_y,\qquad p\iota=e_a(p).
\]
Here \(\{\iota\}\) denotes the terminal category, each endomorphism category of the three objects is terminal, and the remaining hom categories are empty. In the parameter fibre of \cref{sec:indexed-readback}, all universal tables are classified by the subobject
\[
 \mathsf{Opt}_a\ \lhook\joinrel\longrightarrow\ B\times P,\qquad
 (y,p)\in\mathsf{Opt}_a\ \Longleftrightarrow\
       \forall r:P\quad E(y,e_a(p))\le E(y,e_a(r)).
\]
For any \(b:T\to B\), its sections over \(T\) classify the families of Kan extensions for the target \(b\). This classification commutes with substitution.
\end{theorem}
\begin{proof}
View \(y\) as a 1-cell \(0\to2\) and \(p\) as a 1-cell \(1\to2\). Horizontal composition along \(\iota\) is the functor \(e_a\) on both objects and arrows. All other compositions are forced by the identities or empty hom categories. Thus functoriality of \(e_a\) gives interchange, and composition is unital and associative. For every \(r\) there is a unique 2-cell \(y\to e_a(r)\), since \(E(y,y)=0\le E(y,e_a(r))\). Write \(\eta_p\) for this cell at \(p\). The universal-arrow and the left-Kan-extension comparisons are the very same map
\[
 \mathcal M_y(p,r)\longrightarrow\mathcal D_y(y,e_a(r)),\qquad
 \alpha\longmapsto e_a(\alpha)\,\eta_p .
\]
The right-hand hom is terminal. The left-hand hom is the subterminal truth value of the comparison of the two losses. Hence this map is invertible for every \(r\) exactly when \(p\) has globally least loss. The comparisons are natural in \(r\) by composition. Two minimisers admit a unique arrow each way, and thinness makes the composites identities. This also proves that failure of attainment leaves the extension without a representing table.

For the indexed statement, form the two categories internally over \(B\), using \(B\times P\) and \(B\times B\) for their objects. Their arrow objects are the subobjects defined by the two displayed order relations with common first coordinate \(y\). Reflexivity, transitivity and the evaluation functor reindex. The condition that the hom comparison be invertible is the universal image of the comparison relation along \(B\times P\times P\to B\times P\). Pulling this projection back along a family of targets preserves its universal image by Beck--Chevalley. Explicitly, a section is a family \(p:T\to P\) over \(b\) for which the inequality holds against every candidate in the internal logic of the slice over \(T\). The unique units assemble there. Conversely any such Kan extension has that inequality by the same comparison, proving the asserted bijection of choices. The subobject therefore describes attainment even in the absence of a global selection. Inputs can also vary. In the slice over \(X^D\) the evaluation map is \((a,p)\mapsto(a,(F(p,a_i))_{i\in D})\), on the domain where all these calls are legal. The same arrow objects and projection give the classification there.
\end{proof}

\begin{example}[Tied predictions and the dependence on the target]\label{ex:fit-kan-tie}
Take the two legal constant tables with outputs \(-1,+1\) and scalar squared error. At target \(0\) both losses are \(1\). In \(\mathcal M_0\) both tables give universal arrows and are uniquely isomorphic. At target \(-1\) the losses are \(0,4\), so only the negative table gives a universal arrow. The construction in \cite[proof of Theorem 4.1]{pugh-kan} identifies models by equality of predictions and gives a nonidentity comparison from \(y\) only to the prediction of a selected minimiser. At target $0$ the two predictions are distinct. Whichever minimiser is selected, the other has no unit in that construction. More generally, a single fixed hom category on these tables cannot represent the full minimiser sets at both targets by isomorphism to a Kan extension. Universality at \(0\) would force the two objects to be isomorphic. Transport of the unit along that isomorphism at \(-1\) would make the positive table universal there as well. The indexed arrow relations of \cref{thm:fit-kan} give exactly the two required comparisons.
\end{example}

\begin{corollary}[Exact fit and parameter sharing]\label{cor:fit-kan-exact}
For finite rational output records with summed squared error, the restriction of \(\mathsf{Opt}_a\) to \(\operatorname{im}(e_a)\) is the exact-fit relation \(e_a(p)=y\). For a parameter map \(k:P'\to P\), inverse image sends every optimal table lying over \(k\) to an optimal table for \(e_a k\). The inclusion of these two minimiser relations is an equality over the image in \(B\) of the relation \( (y,k(p'))\in\mathsf{Opt}_a\).
\end{corollary}
\begin{proof}
Each squared coordinate error is nonnegative and the finite sum is zero exactly when every coordinate agrees. Over the image of \(e_a\) there locally exists a witness \(r\) with zero error. Comparison against \(r\) forces any optimal \(p\) to have zero error. Conversely zero error gives every required comparison. Eliminating the image witness into this equality of predicates proves the assertion for families, including an empty sample. The fibre over \(y\) is therefore the pullback \(P_{a,y}\) above.

For the parameter map, the universal inequality against \(P\) implies the one against \(P'\) by substituting \(k(r')\). On the indicated image choose a local witness \(s'\) with \(k(s')\) globally optimal. A minimiser \(p'\) for \(e_a k\) has loss at most that of \(k(s')\), which is at most that of each \(r:P\). Transitivity makes \(k(p')\) globally optimal, and the conclusion descends from the witness. For example, take the output \((p,q)\), target \((0,2)\), and diagonal \(k(t)=(t,t)\). The unrestricted perfect fit is \((0,2)\). On the diagonal the loss is \(t^2+(t-2)^2=2(t-1)^2+2\), with sole minimiser \(1\). Thus a restriction of parameter space can create an optimum even when the original optimum has empty inverse image.
\end{proof}

\begin{example}[Perfect predictions with different grounded concepts]\label{ex:fit-ground-distinction}
Let the reached states be bit pairs \((b,c)\), with reset \((0,0)\) and two actions flipping the respective bits. The environmental sensor returns \(c\). Two states have equal future sensor responses exactly when their \(c\) bits agree. Indeed a word flips \(c\) according to the parity of its second-bit flips, and the empty word already separates opposite values. Thus \(Q=2\). In a single rational architecture take a legal parameter \(p\in\{0,1\}\), expose the port
\[
 h_p(b,c)=(b+pc,pc),
 \qquad F_p(b,c)=\pi_1 h_p(b,c)-\pi_2 h_p(b,c)=b .
\]
Training on all four inputs with target \(b\) gives zero error to both tables. They have identical predictions everywhere and are isomorphic in the loss category. At \(p=1\) the second component reads the environmental bit, so neither quotient makes any identification in \(Q\). At \(p=0\) the histories \((0,0)\) and \((0,1)\) have the same port value and account for both elements of \(Q\). Thus both grounding quotients are terminal. The environmental test \(c=1\) is therefore readable exactly at the table \(p=1\).

Consequently even on this exact-fit fibre there can be no functor out of the loss category assigning to each table its grounding quotient up to isomorphism. It would have to send the isomorphism of the two tables to an isomorphism \(2\cong1\). The positive comparison is instead the factorization of observed ports in \cref{prop:ground-port-change}. Here \((u,v)\mapsto(u-v,0)\) takes the \(p=1\) port to the \(p=0\) port. The induced quotient map \(2\twoheadrightarrow1\) computes exactly the loss of the sensor predicate.
\end{example}

\paragraph{General objectives and attainment.}
The profile risk of \cref{prop:profile-risk} is given directly as an objective on decoder tables. For any objective \(L:P\to W\) valued in an internal preorder, let \(K_L\) have objects \(P\) and an arrow \(p\to r\) when \(L(r)\le L(p)\). Let \(J:P_{\rm disc}\to K_L\) be the identity on objects and \(!:P_{\rm disc}\to1\). A cocone from \(J\) to \(m\) is precisely the family \(L(m)\le L(p)\) for every \(p\). Given another cocone to \(n\), its component at \(m\) is the unique arrow \(m\to n\), and all required triangles commute by thinness. Thus every cocone is a colimit and its vertex gives \(\operatorname{Lan}_{!}J\). Conversely a Kan extension includes its cocone. Applied to profile risk, comparison against the decoder obtained by flipping a single entry \(S\) gives \(w_{S,m(S)}\ge w_{S,1-m(S)}\) by cancelling the other summands. Conversely these inequalities minimize every summand. This recovers precisely the attainment predicate of \cref{prop:profile-risk}. Its two-profile family has computable infimum zero and admits no uniform effective choice of a vertex, by the diagonal argument there.

The hom comparison can be tested on the output of a training step. For the scalar gate \(F(p,1)=p\) with target \(0\), loss \(p^2/2\), and step size \(1/2\), the gradient step from \(p=1\) gives \(p^+=1/2\). Its loss \(1/8\) exceeds the loss of \(p=0\), so the hom comparison in \cref{thm:fit-kan} at competitor \(0\) has empty domain. The learner coalgebra below describes the step. The Kan comparison tests the resulting table against the candidates.

\subsection{A polynomial interface for a training interaction}
For objects \(A,A'\) of \(\E\), let \(J(A,A')=A\mathsf y^{A'}\). Here a polynomial map is an internal container morphism, equivalently a strong natural transformation of extensions \cite[Proposition 2.12]{poly:gambino-kock}. A map \(J(A,A')\to J(B,B')\) chooses a position \(f(a):B\) and a backwards direction \(B'\to A'\) at each \(a\). Thus the hom-set is canonically the set of bimorphic lenses. Composition of these position and direction maps gives the lens composition law, and polynomial tensor satisfies \(J(A,A')\otimes J(B,B')\cong J(A\times B,A'\times B')\). This identifies the lens interface with the corresponding full subcategory of polynomial interfaces.

For the based registers take \(\mathcal C_X=J(X,X^*)\). Spivak's internal hom for the tensor of polynomial interfaces \cite{learn:languages} gives
\begin{equation}\label{eq:learn-poly}
 [\mathcal C_X,\mathcal C_Y](S)
 \cong Y^X\times (X^*)^{X\times Y^*}\times S^{X\times Y^*}.
\end{equation}
Here is the representing calculation in $\E$ for the based-register form of \cite[Proposition 2.8]{learn:languages}. Put $A_0=X\times Y^*$ and $B_0=Y^X\times (X^*)^{A_0}$. For any container $K=\sum_{i:I}\mathsf y^{D_i}$, a map $K\otimes\mathcal C_X\to\mathcal C_Y$ chooses $f_i:X\to Y$ and, at $i,x$, a map $Y^*\to D_i\times X^*$. Splitting this pair and currying gives $i\mapsto(f_i,\beta_i)\in B_0$ and $A_0\to D_i$ at each $i$. These are the position and direction maps $K\to B_0\mathsf y^{A_0}$. Uncurrying and pairing are inverse; substitution of a container map acts by the same precomposition on positions and postcomposition on directions on either side. This proves the internal-hom formula naturally in $K$.

\paragraph{An optimizer as a parameter lens.}
A state with readout \(w:S\to P\) and update \(u:S\times P^*\to S\) gives a bimorphic lens $W=(w,u):(S,S)\to(P,P^*)$ as in \cite[\S3]{lens:deep}. Write $j$ for the inclusion of the cotangent lenses of \cref{sec:optic-feedback} into all bimorphic lenses of \cref{sec:intro}. Identities and serial composition have the same formulas on both sides. The tensor comparison uses $(V\times V')^*\simeq V^*\times(V')^*$, with inverse given by adding the covectors on the two projections. On three factors both pastings add the same three evaluations, and on a zero factor the summand is zero. This gives the strong symmetric monoidal inclusion. The basis-paired update has $u(s,0)=s$, so an optimizer can lie outside the cotangent subcategory. For the implementation $f:P\times X\to Y$ set $L=\mathsf D^*f$. Apply $j$ and its tensor comparison to $L$, and form the composite in bimorphic lenses. The diagram abbreviates $j(L)$ by $L$. The source of the dashed composite distinguishes the state register from the ordinary input:
\[
\begin{tikzcd}[column sep=small]
 (S,S)\otimes(X,X^*)\arrow[rr,"W\otimes1"]\arrow[dr,dashed,"L(W\otimes1)"']&&
 (P,P^*)\otimes(X,X^*)\arrow[dl,"L"]\\
 &(Y,Y^*) .
\end{tikzcd}
\]
If \(L^\flat(w(s),x,\lambda)=(\alpha,\beta)\), the backward map of the composite returns \((u(s,\alpha),\beta)\). The first component is a next state. The second is feedback to the preceding port. Currying this lens at the state gives the three coalgebra components in the following theorem. This construction uses only the lens laws until \(L\) is instantiated by the compiled differential. Figure~\ref{fig:optimizer-wires} follows the two receipts to their destinations.

\begin{figure}[!htbp]
    \centering\small
    \begin{tikzpicture}[x=1cm,y=1cm]
        \node[figbox,minimum height=16mm,minimum width=13mm] (W) at (1.65,.55) {$W$};
        \node[figbox,minimum height=28mm,minimum width=15mm] (L) at (5.65,0) {$\mathsf D^*f$};
        \draw[figwire] (.05,.96) node[left] {$s$} -- ([yshift=4.1mm]W.west);
        \draw[figwire] ([yshift=-4.1mm]W.west) -- (.05,.14) node[left] {$s'$};
        \draw[figwire] ([yshift=4.1mm]W.east) -- node[above] {$p=w(s)$} ([yshift=9.6mm]L.west);
        \draw[figwire] ([yshift=1.4mm]L.west) -- node[below] {$\alpha$} ([yshift=-4.1mm]W.east);
        \draw[figwire] (.05,-.55) node[left] {$x$} -- ([yshift=-5.5mm]L.west);
        \draw[figwire] ([yshift=-11mm]L.west) -- (.05,-1.1) node[left] {$\beta$};
        \draw[figwire] ([yshift=7mm]L.east) -- (7.8,.70) node[right] {$y$};
        \draw[figwire] (7.8,-.7) node[right] {$\lambda$} -- ([yshift=-7mm]L.east);
    \end{tikzpicture}
    \caption{One training interaction. Follow $s$ and $x$ to the right to obtain $y=f(w(s),x)$. A covector $\lambda$ travels back through the retained execution. The input receipt $\beta$ leaves the learner, while the parameter receipt $\alpha$ returns to the optimizer and determines $s'=u(s,\alpha)$.}
    \label{fig:optimizer-wires}
\end{figure}

\begin{theorem}[A keyed execution as a learner coalgebra]\label{thm:feedback}
Let $f:P\times X\to Y$ be a total numerical query whose backward program is constructed by \cref{con:backward-compiler}. Let $S$ be a state register, $w:S\to P$ its parameter readout, and $u:S\times P^*\to S$ its chosen update. These maps construct a coalgebra $S\to[\mathcal C_X,\mathcal C_Y](S)$. At state $s$ put $p=w(s)$. Its three components are
\[
 x\mapsto f(p,x),\quad
 (x,\lambda)\mapsto (D_xf)^*\lambda,\quad
 (x,\lambda)\mapsto u\big(s,(D_pf)^*\lambda\big).
\]
Independent state registers compose by taking their product and passing the downstream input covector to the upstream query. A terminal loss supplies its output differential as the feedback seed.
\end{theorem}
\begin{proof}\leavevmode
\begin{enumerate}[titledsteps]
    \item \textbf{Construct the coalgebra.} At state $s$ each displayed map is an internal composite. Currying the three maps gives exactly the data in \cref{eq:learn-poly}. The readout $w$ can select the current parameters from a state that also stores optimizer data.


    \item \textbf{Compose the interactions.} For composition, let $f$ and $g$ have independent states $(s,t)$ and retain the intermediate value $y=f(p,x)$. Given output feedback $\lambda$, the second stage returns $\beta=(D_y g)^*\lambda$ at $y$. Its own state receives $(D_qg)^*\lambda$. The first stage uses $\beta$ to return $(D_x f)^*\beta$ at the input and $(D_p f)^*\beta$ to its state. By \cref{thm:para-lens}, these are exactly the input and parameter covectors for the product-state execution. Both updates read the original forward values. A third stage gives the same nested pullback in either association, and independent products act componentwise.

At an output with target $t$ and differentiable loss $\ell$, take $\lambda=d_1\ell_{(y,t)}$. A tracked loss differential makes this boundary operation effective.\qedhere
\end{enumerate}
\end{proof}
Apply \cref{lem:ground-finite-play} to the learner coalgebra with state $S$, move object $A_0=X\times Y^*$ and observed positions $B_0=Y^X\times(X^*)^{A_0}$. The behaviour of a state records its position after each finite sequence of input--feedback moves. Its image identifies exactly the distinctions observable in repeated training interactions. Given a named field of the learner state, \cref{thm:ground-common,thm:ground-stable} identify which of these predicates can be read at that field, both at the current state and under continuation.

A change of state representation must preserve the entire interaction. For $H=[\mathcal C_X,\mathcal C_Y]$ and transition maps $c,c'$, this amounts to the coalgebra morphism square
\[
\begin{tikzcd}
 S\arrow[r,"c"]\arrow[d,"k"']&H(S)\arrow[d,"H(k)"]\\
 S'\arrow[r,"c'"']&H(S').
\end{tikzcd}
\]
Following the upper route compares states after observing the current maps and taking a step. The lower route changes the state representation first. Equality requires the same observable feedback and corresponding next states. A differentiable reparameterization of the forward program need only supply its cotangent pullback.

\subsection{Target-valued requests and the category of learners}
How can a component expecting a target receive the covector computed by the next layer? In $\mathsf{Learn}(\E)$ a learner $X\to Y$ carries a state object $P$ and three maps. Its implementation $I:P\times X\to Y$ produces an output. On receiving a target, $U:P\times X\times Y\to P$ updates the state and $r:P\times X\times Y\to X$ requests an input target. We identify representatives under state isomorphisms preserving all three maps.

Suppose a learner $(Q,J,V,s)$ follows $(P,I,U,r)$. For output target $t$, the first learner supplies the intermediate value $y=I(p,x)$. The second returns $y'=s(q,y,t)$ as the target for the first. The composite therefore has implementation $J(q,y)$, request $r(p,x,y')$, and state update
\[
 (p,q)\longmapsto \big(U(p,x,y'),V(q,y,t)\big).
\]
The identity has terminal state and passes the target back to its input. Parallel composition takes pairs of states, outputs, requests and updates. Substitution into the formulas gives associativity and the unit laws up to the canonical product identifications. This is the category of learners of Fong, Spivak and Tuy\'eras, with the data taken in $\E$ \cite{neural:fong}. The pair $(U,r)$ forms the backward map of a parameterized simple lens. Fong and Johnson relate this description to symmetric lenses \cite{learn:johnson}.

To build the request, the upstream learner must be able to recover a covector from the target it receives. For a finite based register write $\flat_X:X\to X^*$ for the basis pairing and $\sharp_X$ for its inverse. At the retained value $x$, the following maps encode a target as a covector and recover a target from that covector:
\[
 e_x(t)=\flat_X(x-t),\qquad d_x(\alpha)=x-\sharp_X\alpha .
\]
The maps are inverse at each $x$. Write $\mathcal T_X=X\mathsf y^X$ for the target-valued interface. The isomorphism $\iota_X:\mathcal C_X\to\mathcal T_X$ fixes the position $x$ and acts backwards on directions by $e_x:X\to X^*$. The triangle transports the learner coalgebra across these interface isomorphisms while keeping the parameter state fixed.
\[
\begin{tikzcd}[column sep=normal]
 P\arrow[r,"c"]\arrow[dr,"c_{\rm target}"']&
 {[\mathcal C_X,\mathcal C_Y](P)}\arrow[d,"{[\iota_X^{-1},\iota_Y](P)}"]\\
 &{[\mathcal T_X,\mathcal T_Y](P)}.
\end{tikzcd}
\]
At the output, the transported coalgebra turns the target $t$ into $e_y(t)$ before calling the optimizer. At the input, it turns the returned covector into a requested target by $d_x$. A terminal loss can thus determine the output feedback while every internal wire uses the same target encoding.

\begin{theorem}[Compositional learner interpretation]\label{thm:learner}
Fix the based numerical registers of the query language and a rational step $\eta>0$. For total keyed programs in the constructed differentiable numerical language, composition with independent parameter states and parallel product gives a strong symmetric monoidal interpretation in $\mathsf{Learn}(\E)$. The source identifies canonical parameter-register reindexings preserving the based structure. A program $f$ has learner
\[
 \begin{aligned}
 I(p,x)&=f(p,x)=y, &\lambda&=e_y(t),\\
 U(p,x,t)&=p-\eta\sharp_P (D_p f)^*\lambda,&
 r(p,x,t)&=d_x\big((D_x f)^*\lambda\big).
 \end{aligned}
\]
\end{theorem}
\begin{proof}
Fix $p,x$ and let $y=f(p,x)$. The square expresses the compatibility required for composition. Encoding the input request gives the same covector as encoding the output target and pulling it back.
\[
\begin{tikzcd}[column sep=large]
 Y\arrow[r,"{r(p,x,-)}"]\arrow[d,"e_y"']&X\arrow[d,"e_x"]\\
 Y^*\arrow[r,"{(D_x f)^*}"']&X^* .
\end{tikzcd}
\]
For \(f_p:X\to Y\), \(g_q:Y\to Z\), set \(y=f_p(x)\), \(z=g_q(y)\), \(\lambda=e_z(t)\), \(\beta=(D_y g_q)^*\lambda\). The downstream request is \(y'=d_y(\beta)\), so \(e_y(y')=\beta\). Substitution into learner composition gives
\[
\begin{aligned}
 p^+&=p-\eta\sharp_P(D_pf)^*\beta,&
 q^+&=q-\eta\sharp_Q(D_qg)^*\lambda,\\
 x'&=d_x\big((D_xf_p)^*\beta\big)
    =d_x\big(D(g_qf_p)_x^*\lambda\big).
\end{aligned}
\]
The same calculation at three stages is associative by associativity of cotangent composition. At the identity \(d_x e_x(t)=t\). For parallel components \(e_{(x,y)}=e_x\times e_y\), \(d_{(x,y)}=d_x\times d_y\), and the basis pairing and update split by factors. The product associators and coordinate permutations give the symmetric monoidal coherence. An external differentiable loss supplies \(t=d_z(d_1\ell_{(z,t_0)})\) for the desired loss target \(t_0\), so the boundary encoder returns its loss covector.
\end{proof}

A parameter isomorphism identifies two learners when it preserves the implementation, request, and update. For the Euclidean step, an isometric linear isomorphism of parameter registers transporting the implementation also transports its request and update. Permutations of parameter keys are examples. General differentiable bijections still have cotangent pullbacks, but the resulting steps depend on the pairing \cite{lens:deep}. The rescaling after \cref{neural:con:optimizer} gives a scalar example.

\subsection{One state transition per shared table}
Serial composition assigns independent states to independent learners. For gates that share a parameter, one component stores it and each gate reads it through a stateless lookup. The reverse-lookup lemma determines what that component receives when the gates request different targets.
\begin{proposition}[Factorization through a stateful parameter table]\label{thm:state-table}
Let $\rho:D\to K$ be the foreign key from numerical occurrences to stored parameters. A single stateful learner $1\to R^K$, followed by stateless lookup and the downstream computation, implements the tied parameter update.
\end{proposition}
\begin{proof}
Write $P=R^K$. The table outputs its present state $\theta$ and upon a request $\theta'$ sets its next state to $\theta-\eta(\theta-\theta')$. Its request to the terminal input is unique. For lookup $L=\rho^*:P\to R^D$, the stateless request at a downstream occurrence covector $\gamma$ is $\theta-\sharp_P\rho_!\gamma$, by \cref{lem:sharing}. Substitution into the table update yields $\theta^+=\theta-\eta\sharp_P\rho_!\gamma$ exactly once. More explicitly, forward lookup and its reverse are related by
\[
\begin{tikzcd}[column sep=large]
 (R^D)^*\arrow[r,"L^*"]\arrow[d,"\sharp_D"']&(R^K)^*\arrow[d,"\sharp_K"]\\
 R^D\arrow[r,"\rho_!"']& R^K .
\end{tikzcd}
\]
For two uses of one scalar $\theta$ with targets $t_1,t_2$, the combined request is $t_1+t_2-\theta$. The table thus subtracts $\eta((\theta-t_1)+(\theta-t_2))$ in a single step. A general optimizer uses the transition $u(s,e_{w(s)}(\theta'))$ at the table. The stateless lookup wires perform exactly the same aggregation.
\end{proof}

The parameter key now has a training interpretation. Its fibre specifies all the terms in a parameter covector, and the table that owns the key takes the resulting step. In particular two tables with equal values can evolve independently, while two uses of one address must contribute to the same update.

\subsection{The update is an additional map}
Transporting an update along a parameter change requires a square in addition to the cotangent chain rule. For \(r:P'\to P\), an old-state update \(u\) and a new-state update \(u'\), both evaluated at the retained parameters, the required equation is
\[
\begin{tikzcd}[column sep=large]
 P'\times P^*\arrow[r,"{(p',\alpha)\mapsto(p',dr_{p'}^*\alpha)}"]\arrow[d,"r\times1"']&
 P'\times(P')^*\arrow[r,"u'"]&P'\arrow[d,"r"]\\
 P\times P^*\arrow[rr,"u"']&&P .
\end{tikzcd}
\]
The left-hand input supplies a covector at \(r(p')\). Commutation states exactly that updating after its pullback gives a state over the updated old state. This is a criterion on the chosen update. The basis-paired case below computes the two routes and identifies their discrepancy for a change of scale.
\begin{construction}[The basis-paired gradient step]\label{neural:con:optimizer}
For the presented free parameter module $P=R^{(I)}$, its basis defines
$\flat_I:P\to P^*$ by $\flat_I(p)(q)=\sum_i p_iq_i$ and an inverse $\sharp_I$. For a rational step code $\eta>0$, define
\[
 u_\eta:P\times P^*\longrightarrow P,
 \qquad u_\eta(p,\alpha)=p-\eta\sharp_I(\alpha).
\]
A loss $\ell:Y\times Y\to R$ and a target $y$ give the update
$u_\eta(p,d_Pf_{(p,x)}^*(d_1\ell_{(f(p,x),y)}))$.
In this realization the basis identifies a parameter covector with the vector used for the step. Other optimizers can retain additional state in the register of \cref{thm:feedback}.
\end{construction}
The evaluation of a covector on a tangent is intrinsic. Turning the covector into a displacement uses $\sharp_I$, the inverse of the register's basis pairing. For a parameter change $r:P'\to P$, the chain rule always gives the covector $dr^*\alpha$. For a linear isometric embedding $r:P'\to P$, let $r^\dagger=\sharp_{P'}r^*\flat_P$. Then $r^\dagger r=\id$ and the displacement computed in $P'$ pushes forward to $-\eta\,r r^\dagger\sharp_P\alpha$. Here $r r^\dagger$ is the orthogonal projection onto $r(P')$. For an isometric isomorphism it is the identity, so the whole Euclidean step is transported. For $r(z)=2z$, however, the covector in the $z$ coordinate is $2\alpha$. Taking the step there and mapping it to $p$ gives displacement $-4\eta\alpha$, whereas the step computed in $p$ gives $-\eta\alpha$.

\subsection{Transitions driven by a loss jet}
Write $\mathcal D^n(P)$ for the family of pointed jets over the parameter register. A state with readout $w:S\to P$ accepts feedback in the fibre at its present readout. An order-$n$ transition therefore has type
\[
              U_n:S\times_P\mathcal D^n(P)\longrightarrow S .
\]
For a loss $\ell$ and a retained input $x$, take the class of $\ell-\ell(y)$ at $y=f(w(s),x)$, pull it back by $f^{[n]}_{(w(s),x)}$, and restrict to the parameter fibre by holding $x$ fixed. Apply $U_n$ once. Shared parameters are substituted in the joint program before this restriction by \cref{thm:joint-jet}. Serial components with an optimizer using cross terms keep their joint state through this operation. For an executable loss in the stated syntax, the same generic-increment evaluation computes its seed. At first order the joint loss jet at $(p,x)$ splits as $(\alpha,\beta)\in P^*\oplus X^*$. Restriction to the parameter fibre yields $\alpha$ and gives the transition $u(s,\alpha)$. Restriction to the input fibre, with $p$ fixed, yields the input receipt $\beta$. Together with the retained output these two restrictions give all three components of \cref{thm:feedback}. The target-valued input request is $d_x(\beta)$ by \cref{thm:learner}. For example, if $f(p,x)=x$ on a scalar input and the output cotangent is $1$, these restrictions give $\alpha=0$ and $\beta=1$. A first-order transition extends to higher feedback by composing with truncation.

On a finite affine register the quadratic coefficient gives a symmetric map $H:P\to P^*$. On the principal open $D(\det H)$ the adjugate identity constructs $v=(\det H)^{-1}\operatorname{adj}(H)\alpha$. The step $p\mapsto p-v$ uses the linear part $\alpha$ and the quadratic part $H/2$ of the same loss jet. Conversely an inverse matrix forces $\det H$ to be a unit by multiplicativity. Thus the full Newton solve has its inverse domain as a pullback. The gradient transition constructed above gives a different response to the same feedback. On a scalar unit chart this is a single division. The numerical instance in \cref{ex:update} exhibits that inverse explicitly. On registers without a chosen affine identification an observable still pulls back as a whole jet, as in \cref{thm:jet-pairing}.

\begin{proposition}[An update on balanced attention tables]\label{prop:balanced-step}
Let the state be a dense positive balanced table $P$ with fixed marginals $\mu,\nu$. An output receipt $g_e$ at each edge determines the update
\[
 U_\eta(P,g)=\operatorname{Bal}_{\mu,\nu}\big((P_e\exp(-\eta g_e))_e\big).
\]
This is an effective state transition with those same marginals. Its velocity at $\eta=0$ is $-v$, where $v$ is the metric gradient of the receipt on cycle flows for
$G_P(v,k)=\sum_e v_e k_e/P_e$.
\end{proposition}
\begin{proof}
Every tilted entry is positive; the constructive balancing program of \cref{thm:balance-exact} therefore computes the next table. At $\eta=0$ the initial state already satisfies both marginals. Its marked equations give $v=P(g-d\phi)$ with $\partial v=0$ and the anchored $\phi$ determined by \cref{prop:balanced-lift}. For every tangent $k$ we have $\partial k=0$ and hence
\[
    G_P(v,k)=\sum_e g_e k_e-\sum_a\phi_a(\partial k)_a=\sum_e g_e k_e .
\]
This identifies the metric gradient uniquely: a vector orthogonal to every cycle including itself has zero sum of positive weighted squares, hence vanishes. Substituting the log-weight direction $-g$ gives the negative velocity asserted. A different representative $g+d\psi$ of the same receipt multiplies the tilted weights by row and column factors; absorbing these into the scales leaves the balanced output unchanged by uniqueness. For marginals of total one $G_P$ is the restriction of the Fisher pairing. A state that stores upstream affinity parameters instead receives the loss jet pulled back to those parameter keys.
\end{proof}

\section{Numerical realizations and an executed training step}\label{sec:realizations}

The \hyperref[sec:fit-fibre]{exact-fit fibre} describes the tables satisfying a prescribed data set. In the execution below the optimizer changes the parameter component of its state after reading the current datum. Its one-step comparison therefore concerns a map of states. The test for global optimality of an updated table is the hom comparison of \cref{thm:fit-kan}.

A fixed training datum closes the input and output boundaries of the learner and gives an endomorphism \(T_n:S\to S\). For two presentations with a state comparison \(k:S\to S'\), the comparison for a whole run is the horizontal pasting of the one-step squares:
\[
\begin{tikzcd}[column sep=large]
 S\arrow[r,"T_0"]\arrow[d,"k"']&S\arrow[r,"T_1"]\arrow[d,"k"']&S\arrow[d,"k"]\\
 S'\arrow[r,"T'_0"']&S'\arrow[r,"T'_1"']&S' .
\end{tikzcd}
\]
In the present pair of executions the state object and update \(u\) coincide, so \(k=1_S\). The joint-jet comparison identifies the feedback at each shared parameter owner. Applying \(u\) gives \(T'_n k=kT_n\) for each datum, and pasting gives the comparison for every finite run, including the empty run. At first order its coordinate form is
\[
 g_n(s)=D_p[\ell(f(p,x_n),t_n)]_{p=w(s)},\qquad T_n(s)=u(s,g_n(s)).
\]
An update consuming a higher loss jet uses the same square with the corresponding feedback. The block below evaluates these morphisms on specified registers.

\subsection{Assembly of a block and equivariance}

For heads \(a=1,\ldots,H_0\), take the finite product of the head programs on one shared token table. Concatenate their values in a finite direct sum at each row. With compatible dimensions form
\[
 u_i=x_i+W_O\!\left(\bigoplus_a y_i^{(a)}\right)+b_O,\qquad
 B_\Theta(x)_i=u_i+W_2\,\sigma(W_1u_i+b_1)+b_2 .
\]
The sigmoid $\sigma(t)=(1+e^{-t})^{-1}$ acts coordinatewise in the second residual term. Both residual branches read the same retained input, as do the different heads. Each affine map has the incidence realization from \cref{neural:thm:linearization}.

The positive integer $d$, obtained as the enumerated key-coordinate count, supplies an effective scale. Bisect $[1,d]$ by comparing rational squares to enclose $\sqrt d$, then apply reciprocal away from zero. The remaining reciprocals come from the head normalizers and the sigmoid denominator, which is at least one in the latter case. Finite composition therefore gives a program on the stated numerical names.

\begin{proposition}[Equivariance of the constructed block]
A bijection \(\gamma:J\to J'\) of the token indices, transporting the mask, the feature rows and any positional attributes, intertwines the two forward evaluations. With full eligibility and position-independent projections this holds for every permutation of the token indices.
\[
\begin{tikzcd}
 X^J\arrow[r,"B_{\Theta,M}"]\arrow[d,"\gamma_*"']&
 X^J\arrow[d,"\gamma_*"]\\
 X^{J'}\arrow[r,"B_{\Theta,\gamma M}"']&X^{J'}
\end{tikzcd}
\]
\end{proposition}
\begin{proof}
The reindexed features satisfy \(x'_{\gamma i}=x_i\). Shared projections give \(q'_{\gamma i}=q_i\), and similarly for keys and values. Transporting the span replaces $(\pi,t)$ by $(\gamma\pi,\gamma t)$ on the same occurrence carrier. The identity on occurrences therefore bijects the old fibre over $i$ with the new fibre over $\gamma i$, preserving both the weight and the value read at each occurrence. Both its normalizer and its value sum consequently agree. Direct sums, output projections, positionwise nonlinearities and residual sums all commute with \(\gamma_*\). This proves the square. A mask or positional field kept fixed admits the bijections preserving that field.
\end{proof}

\subsection{Coordinate pullbacks for the layer and head}

\begin{theorem}[Differential of a keyed neural layer]\label{thm:layer-back}
For a scalar numerical activation $\phi$, let $D_\phi$ be the image of its ordinary derivative graph $(a,l)$: for each rational $\epsilon>0$ there is a rational $\delta>0$ with $|\phi(a+t)-\phi(a)-lt|\le\epsilon |t|$ for $|t|<\delta$. The joint differentiability domain of the layer in \cref{eq:layer}, on its full numerical feature, bias and parameter registers, is the intersection of the inverse images of $D_\phi$ under its preactivations. There let \(\bar h_{j,k}\) be output covector coordinates and set \(g_{j,k}=\phi'(a_{j,k})\bar h_{j,k}\). Then the parameter, bias and input ports receive
\begin{align}
 \bar\theta_p&=\sum_{\substack{j,d\\\rho(d)=p}} g_{j,t(d)}x_{j,s(d)},&
 \bar b_k&=\sum_j g_{j,k},\label{eq:keygrad}\\
 \bar x_{j,i}&=\sum_{d:s(d)=i}\theta_{\rho(d)}g_{j,t(d)},& \bar u_c&=\bar x_{\gamma(c)}.
\end{align}
\end{theorem}
\begin{proof}\leavevmode
\begin{enumerate}[titledsteps]
    \item \textbf{Determine the differentiability domain.} The slope in the scalar graph is unique: compare two slopes at one nonzero rational increment inside both error radii and let $\epsilon$ tend to zero. It therefore descends to $D_\phi$ by \cref{thm:solved-composition}. On the displayed intersection, the product and chain rules give the layer differential. Conversely vary only $b_k$ by $t$ and project the output to $(j,k)$. This restriction of the layer is $t\mapsto\phi(a_{j,k}+t)$, so differentiability of the joint map supplies that scalar derivative at every output. For a rectifier the two strict chambers give slopes zero and one; at zero the two bias directions give incompatible one-sided slopes. 
    \item \textbf{Pull the receipts back to their owners.} Marking a weight or a feature in the occurrence product now gives
\[
 \dot a_{j,k}=\dot b_k+
 \sum_{t(d)=k}(\dot\theta_{\rho(d)}x_{j,s(d)}+
                        \theta_{\rho(d)}\dot x_{j,s(d)}).
\]
Pair with \(g\), expand the sum over the pullback rows, and group respectively by parameter key, bias key and feature key. Each output in \cref{eq:keygrad} is the coefficient of the corresponding independently variable input. For the stored feature, the remaining pairing is

\[
 \sum_{j,i}\bar x_{j,i}\sum_{\gamma(c)=(j,i)}\dot u_c
=\sum_c\bar x_{\gamma(c)}\dot u_c.
\]

Thus the transpose of the feature aggregation is reindexing along $\gamma$. In the dense instance with one parameter per matrix entry, the first equation is a sum of rank-one weight covectors. For a repeated parameter key, \cref{lem:sharing} sums the contributions after all its uses have been counted.\qedhere
\end{enumerate}
\end{proof}

\begin{theorem}[Full backward query of a head]\label{thm:head-back}
Let $\lambda_i\in V^*$ be the loss covector at each active output $i\in J_+$ in \cref{con:head}. Empty sums at unused ports are zero. The score, value, query and key ports receive
\begin{align}
 g_e&=\alpha_e\lambda_{\pi(e)}(v_{t(e)}-y_{\pi(e)}),&
 \bar v_j&=\sum_{e:\,t(e)=j}\alpha_e\lambda_{\pi(e)},\\
 \bar q_i&=\frac{1}{\sqrt d}\sum_{e:\,\pi(e)=i}g_e k_{t(e)},&
 \bar k_j&=\frac{1}{\sqrt d}\sum_{e:\,t(e)=j}g_e q_{\pi(e)} .
\end{align}

Here \(\bar q_i\in H^{**}\cong H\) and \(\bar k_j\in H^*\). These are the cotangents at the projected ports. At each of the three affine projections \(z_j=W_z x_j+b_z\), the parameter covectors in the displayed bases are
\[
 \bar W_z=\sum_j\bar z_j x_j^{\mathsf T},\qquad \bar b_z=\sum_j\bar z_j .
\]
The input feature receives \(\bar x_j=W_Q^{\mathsf T}\bar q_j+W_K^{\mathsf T}\bar k_j+W_V^{\mathsf T}\bar v_j\) in these coordinates.
\end{theorem}
\begin{proof}
In \cref{eq:unnormdiff} set $\dot w_e=w_e\dot s_e$ on each query fibre. Pair with $\lambda_i$ and sum over $i$. Before grouping by any address the resulting expression is
\[
 \sum_{i\in J_+}\lambda_i(\dot y_i)=
 \sum_{e\in M}\bigl(g_e\dot s_e+
               \alpha_e\lambda_{\pi(e)}(\dot v_{t(e)})\bigr).
\]
Grouping the second terms along $t$ yields $\bar v_j$. In the first terms substitute
$\sqrt d\,\dot s_e=\dot q_{\pi(e)}(k_{t(e)})+
q_{\pi(e)}(\dot k_{t(e)})$. Group along $\pi$ for $\bar q_i$ and along $t$ for $\bar k_j$. Finally $\dot z_j=\dot W_z x_j+W_z\dot x_j+\dot b_z$. Pairing and grouping by parameter key gives the matrix and bias covectors by \cref{thm:layer-back}. The dual of the triple diagonal adds the three contributions at $x_j$.

For a binary mask, the fibre $t^{-1}(j)$ is indexed by the queries $i$ with $(i,j)\in M$, recovering the usual column sum. For a single query with two equally weighted occurrences reading the same value $v_j$, each has $\alpha_e=1/2$. The value pullback gives $\bar v_j=(1/2+1/2)\lambda_i=\lambda_i$, as direct differentiation of $y_i=v_j$ requires.
\end{proof}

\begin{corollary}[Forward and backward comparison, including the update]\label{thm:training}
Construct a finite program by composing the keyed affine queries, the polynomial gates, the normalized heads of \cref{con:head}, the analytic gates of \cref{neural:lem:analytic}, products and residual addition. Retain all shared parameter keys. At every finite numerical instance its forward table equals the array interpretation after enumerating the presented ports. For a differentiable loss on outputs and fixed targets, its backward query returns the parameter covector of the array composite. With the displayed pairing on the same parameter register and the same step code, one update gives the same new parameter table in both presentations.
\end{corollary}
\begin{proof}
For each forward port \(i\), let \(E_i\) enumerate the keyed field as a coordinate vector. For a gate \(g\) and its array implementation \(\widetilde g\), put \(E_{\rm in}=\prod E_i\) on its ordered operands and \(E_{\rm out}\) on its result. The two local identities are
\[
 E_{\rm out}g=\widetilde g E_{\rm in},\qquad
 Dg_x^* E_{\rm out}^*
       =E_{\rm in}^*D\widetilde g_{E_{\rm in}x}^* .
\]
At affine and polynomial gates these follow from \cref{neural:thm:linearization,poly:thm:differential-comparison}. For a masked head enumerate \emph{occurrences} before gathering by \(\pi,t\). The two identities are the grouped forward sum and \cref{thm:head-back}. Binary masks give \cref{prop:attention-softmax}. Analytic gates use \cref{neural:lem:analytic}. Products and residual addition use their tuple and sum laws. Topological induction gives the first identity for the circuit. Reverse induction sums incoming receipts at each shared register and gives the second identity. At the owner address it yields
\[
 \bar\theta_{\rm keyed}=E_P^*\bar\theta_{\rm array},\qquad
 E_P\,u_\eta(\theta,\bar\theta_{\rm keyed})
      =u_\eta^{\rm array}(E_P\theta,\bar\theta_{\rm array}).
\]
The second equality uses that enumeration permutes the specified basis. Iterating the same equality on successive common batches compares the entire finite trajectory.
\end{proof}
Writing \(u_\eta\) for that update on \(P\), an isometric linear isomorphism \(T:P\to P'\) gives the final comparison explicitly.
\[
\begin{tikzcd}[column sep=huge]
 P\times P^* \arrow[r,"u_\eta"]\arrow[d,"T\times (T^{-1})^*"']&P\arrow[d,"T"]\\
 P'\times(P')^*\arrow[r,"u'_\eta"'] &P'
\end{tikzcd}
\]
Indeed \(\sharp' (T^{-1})^*=T\sharp\) by the isometry identity. The identity of register orderings and a permutation of finite parameter keys are both instances.

\begin{proposition}[An exact finite descent comparison]
For one affine scalar-output row \(f_w(x)=w^{\mathsf{T}} x\) with a fixed nonzero input vector of rational coordinates, target \(t\), and quadratic loss \(L=\tfrac12(f_w(x)-t)^2\), the basis-paired update gives
\[
 L(w^+)=(1-\eta\|x\|^2)^2 L(w).
\]
Consequently a rational step \(0<\eta<2/\|x\|^2\) strictly decreases any positive loss.
\end{proposition}
\begin{proof}
With \(e=w^{\mathsf{T}} x-t\), \cref{thm:layer-back} gives parameter gradient \(ex\). The new residual is \(e-\eta e\,x^{\mathsf{T}} x\). Squaring gives the identity. The displayed interval yields \(|1-\eta\|x\|^2|<1\). A concrete admissible step is \(1/\|x\|^2\), computed from the nonzero rational input by finite arithmetic. For the zero vector the parameter covector vanishes. This computation separates a proved finite change of objective from a first-order signal.
\end{proof}

\begin{example}[An executed update at a shared query parameter]\label{ex:update}\leavevmode
\begin{enumerate}[titledsteps]
    \item \textbf{Evaluate the shared query.} For the three-token head all three masked rows are active. Give the two parameter-dependent positions $b,c$ target one and define
\(L(u)=\tfrac12\big((y_b-1)^2+(y_c-1)^2\big)\).
There are two sums to check. Within each loss-bearing query, two eligible keys have value one. Across the two queries, both projections $q_b=u x_b$ and $q_c=u x_c$ read the same parameter address.

At $u=\log 2$ each of these two rows returns $4/5$ and receives covector $-1/5$. Take $\dot u=1$. Each unit-valued key has $\dot s=1$ and contributes $(2/5)(1-4/5)=2/25$ to the row derivative. The zero-valued key has $\dot s=0$. The row derivative is thus $4/25$, and its contribution to the loss derivative is $(-1/5)(4/25)=-4/125$. Summing at the shared address gives
\[
 L'(\log2)=-8/125,\qquad u^+=\log2+8\eta/125 .
\]

    \item \textbf{Check the derivative directly.} Direct evaluation gives an independent check. Both loss-bearing outputs equal $2e^u/(1+2e^u)$, so $L(u)=(1+2e^u)^{-2}$. Differentiating this scalar function gives $-8/125$ at $\log 2$. Moreover, $L$ decreases as $u$ increases. In this example every positive step size in the displayed update therefore decreases the exact loss.

The downstream rows express the same derivative as a request to the parameter table. Together their contributions give the target $u'=u+8/125$. The table takes one step, $u-\eta(u-u')=u+8\eta/125$. Thus the request transmits the sum over the two query occurrences, and the state transition applies the step size once at the address they share.


    \item \textbf{Compute the second-order transition.} For the same data put $u=\log 2+t$ in the ring with $t^3=0$. The denominator of $L$ before squaring is $1+4(1+t+t^2/2)=5+4t+2t^2$. Squaring and solving for the reciprocal by \cref{prop:normalized-jets} gives
\[
             L(\log 2+t)=\tfrac1{25}-\tfrac8{125}t+\tfrac{28}{625}t^2 .
\]
Its curvature in the retained affine coordinate is $56/625$. The scalar Newton transition therefore adds $(8/125)/(56/625)=5/7$ to $u$, while the first-order transition above adds $8\eta/125$. Both use the one parameter table on which the two query positions agree.

\end{enumerate}
\end{example}

\section*{What the reverse translation recovers}
\phantomsection\addcontentsline{toc}{section}{What the reverse translation recovers}
The inverse in \cref{thm:support-readback} reconstructs across shapes: exact evaluations with valid bounds determine the coefficient of every decorated embedding, even a pattern whose extra rows occur only as witnesses. Concatenation and substitution make this reconstruction a finite-product isomorphism. A named interface in \cref{prop:relative-models} fixes which arrow is observed. Its structural presentation retains the sharing and owner keys recovered in \cref{thm:compiler-struct}. On reachable states the profile quotient carries exactly those decisions that descend along it. On a common replacement domain the response quotient is universal among sufficient decoders. Predicate feedback pulls observations back along either map. In the differentiable language a realized factorization also transports its cotangent optics, and sharing a parameter acts by precomposition with its lookup. These universal properties specify what is invariant under a change of presentation. The rounded port of \cref{ex:trained-discrepancy} instantiates them as follows:
\[
\begin{array}{l|l}
 \text{domain and observable}&\text{determining data}\\ \hline
 \text{49 inputs, decision}&(h'_2=0)\ \vee\ (h'_2\ge39/4)\\
 \text{all reachable rational states, score}&
          6-5v+\tfrac{21}{4}[2v-10]_+,\quad v=h'_2\\
 \text{all nonnegative replacements, score}&d(z)=21z_1-20z_2\\
 \text{all nonnegative replacements, decision}&d(z)\ge-24\\

 \text{box }[0,2]^2,\ \text{score distance}&
       D_C((0,2),(1,2))=21/4
\end{array}
\]
The firing-bit vocabulary on the grid incurs six errors at its optimum. The upper and lower profile approximations of \cref{prop:profile-adjoints} locate exactly where that vocabulary mixes decisions. For a proposed intervention decoder on a common domain, \cref{cor:repair-decoder} constructs its least sufficient refinement. A target extent, its Boolean profile saturation, the infimum risk under a source law, and an interventional quotient can each be tested against the relevant construction in \cref{thm:concept-expressibility,prop:profile-risk,prop:response-quotient}. With varying admissibility, a single coarsest decoder can fail to exist.

Operational grounding gives a different observation of a reached history. In \cref{ex:ground-pulse} the hidden test $z_1=1$ holds precisely when token $9$ occurs now, after $a$, or after $aa$. The action sequence and the observed token give this hidden predicate an independent experimental test. Calling the token a measurement of a particular physical quantity would also require a calibration of the sensor. The four classes readable at both interfaces reduce to two when the quotient must support every next action. The pushout and the coequalizer under word actions explain the two counts (\cref{thm:ground-common,thm:ground-stable}).

For a frozen rational rectifier score on a compact rational box, eliminating the bypass coordinates computes a term for the distance on the response image (\cref{cor:dmv-response}). At fixed shape, equivariant polynomials have monomial-orbit coordinates for the finite group action. At normalized, structured, and balanced value ports the corresponding identifiable data are respectively the probed probability coefficients, the projection of scores onto the span of count differences, and the score orbit modulo endpoint potentials. Grouped payloads on a balanced support have the additional kernel test following \cref{thm:global-balance-inverse}.

For an implicit or guarded call the graph projection must have its exhibited inverse on the execution domain. Its equations then transport forward and backward jets at the selected point. At first order a shared lookup has adjoint \(\rho_!\). At higher order it substitutes shared increments in every monomial before the loss jet is sent to the owner. Consequently the two executions in \cref{ex:update} give the same transition for either of the specified optimizers.

\paragraph{Obstructions and their witnesses.}
\begin{center}
\small
\begin{tabular}{@{}p{.36\linewidth}p{.59\linewidth}@{}}
\toprule
Proposed inference & Witness or test\\
\midrule
Infer finite support from probes alone &
The zero transcript also fits \(\binom{|V|}{N}\) for sufficiently large \(N\) (\cref{sec:semantic}).\\[2pt]
Certify complete environmental recovery from an unbounded class of finite devices &
A unary device can agree with a constant sensor up to the longest probe and then change its output (\cref{ex:ground-delayed}).\\[2pt]
Select an exact classifier from located masses &
The two-profile law has risk infimum zero. Selecting its label at \(A\) selects a weak sign (\cref{prop:profile-risk}).\\[2pt]
Use one quotient for varying replacement domains &
The sufficient partitions \(ab\mid c\) and \(a\mid bc\) have an insufficient common join (\cref{prop:response-quotient} and the following calculation).\\[2pt]
Recover all pointed jets from ordinary arcs &
Every rational order-two arc at the origin of \(x^2=y^2=0\) kills the nonzero class \(xy\) (\cref{thm:jet-pairing}).\\[2pt]
Recover balanced edge weights from grouped payloads &
Any nonzero vector in \(\ker D\cap\ker G\) yields distinct positive tables with identical probes (\cref{thm:global-balance-inverse} and the following calculation).\\
\bottomrule
\end{tabular}
\end{center}
\clearpage
\part*{Appendices}
\phantomsection\addcontentsline{toc}{part}{Appendices}
\appendix

\section{Constructive lower integration on represented rows}\label{app:integration}

The attention kernel uses three facts about lower integration: refinement preserves lower bounds, change of rows acts by preimage on tests, and the integral is additive on nonnegative features. The partition construction gives the first two directly and yields a uniform enumeration on named inputs. For additivity we first handle finite sums of indicators, then pass to increasing approximations of the integrands.

\subsection{Partition evaluation and computable lower names}

\begin{proof}[Proof of \cref{lem:eff-integral}]\leavevmode
\begin{enumerate}[titledsteps]
    \item \textbf{Refine the partition.} For a partition \(P\) write \(S_P(g)=\sum_k\Delta r_k\,\mu\{g>r_k\}\). Inserting \(s\in(a,b)\) into one step replaces \((b-a)\mu(U_b)\) by \((s-a)\mu(U_s)+(b-s)\mu(U_b)\), which dominates the old term since \(U_b\subseteq U_s\). Sorted unions therefore give common refinements. If \(g=\sup_i g_i\) for a directed family, \(\{g>r\}=\bigcup_i\{g_i>r\}\). Using valuation continuity and a common index for the finitely many bounds yields
\[
 I_\mu(\sup_i g_i)=\sup_P\sup_i S_P(g_i)
                  =\sup_i\sup_P S_P(g_i)=\sup_i I_\mu(g_i).
\]


    \item \textbf{Evaluate an indicator.} For an indicator of $U$, positive thresholds below one give $U$ and thresholds at or above one give the empty open. Letting a rational endpoint approach one from below gives $\mu(U)$, including when the mass is unbounded. For change of rows, the threshold opens of $g\circ f$ are the preimages of those of $g$, so the two partition sums agree at each partition.


    \item \textbf{Enumerate the lower name.} Finally enumerate all finite increasing rational partitions. For each partition, run the valuation program on its threshold opens and enumerate lower approximants to the sum. Interleave the enumerations. By roundedness every strict rational bound on the supremum appears, which gives the required lower name.\qedhere

\end{enumerate}
\end{proof}

\subsection{Additivity and rational scaling}

\begin{proof}[Proof of \cref{lem:eff-linearity}]\leavevmode
\begin{enumerate}[titledsteps]
    \item \textbf{Sort the indicator levels.} Fix a positive rational mesh $\delta$ and a finite list of opens $U_i$. Put $h=\delta\sum_i\tau U_i$. Replacing two opens by their union and intersection preserves the pointwise sum of their indicators. Modularity preserves the sum of their valuations. Apply these replacements in the order specified by a finite sorting network. Its outputs form a nested list $V_j$. By induction on insertion, using distributivity of meets over joins, $V_j$ is the union of all $j$-fold intersections of distinct input occurrences. It is therefore the open on which at least $j$ tests have succeeded.


    \item \textbf{Evaluate the mesh integral.} For \(V_j\) the open of at least \(j\) successes, \(\{h>r\}=V_j\) when \((j-1)\delta\le r<j\delta\). Subdividing a partition step at mesh boundaries bounds its contribution by the corresponding lengths times \(\mu(V_j)\). For \(0<\eta<\delta\) use endpoints \(r_j=j\delta-\eta\). Their first increment is \(\delta-\eta\) and all others are \(\delta\). Thus, for a nonempty list of length \(m\),
\[
 (\delta-\eta)\mu(V_1)+\delta\sum_{j=2}^m\mu(V_j)
 \le I_\mu(h)\le\delta\sum_{j=1}^m\mu(V_j).
\]
Taking the supremum as \(\eta\downarrow0\) uses directed continuity of nonnegative lower multiplication even for infinite mass. Modularity through the sorting network gives \(\sum_j\mu(V_j)=\sum_i\mu(U_i)\). The empty list has integral zero.


    \item \textbf{Approximate the general integrand.} Put \(g_n=2^{-n}\sum_{k=1}^{n2^n}\tau\{g>k2^{-n}\}\). At the next stage the two half-mesh terms indexed \(2k-1,2k\) dominate the old \(k\)-term, giving \(g_n\le g_{n+1}\). If \(0\le r<g(x)\), take \(r<s<g(x)\) rational and a fine mesh below \(s-r\) with cutoff \(n>s\). The confirmed mesh levels through \(\lfloor 2^n s\rfloor\) give \(g_n(x)>r\). Thus \(\sup_n g_n=g\). For a second integrand \(f\) use the same mesh, concatenate the two lists, and obtain
\[
 I_\mu(f+g)=\sup_n I_\mu(f_n+g_n)
           =\sup_n\bigl(I_\mu(f_n)+I_\mu(g_n)\bigr)= I_\mu(f)+I_\mu(g).
\]
Multiplying the rational endpoints of each partition by the same positive rational gives the scaling law. Zero gives the zero integrand.\qedhere

\end{enumerate}
\end{proof}

\section{Coherence for parameterized algebra morphisms}\label{app:pasting}
In \cref{thm:cell-ext}, a single-gate comparison extends to trees precisely when its own parameter survives the passage through a node. The unit and multiplication pastings show where that constraint enters. Use the convention of \cite[Definitions F.2, F.3 and G.1]{cdl} and suppress product unitors. The source and target use the canonical comonoids on $P$ and $Q$. Passing the comparison register $U$ across a gate uses $k:U\times P\to Q\times U$. In the first diagram both routes discard the algebra parameter. In the second both routes distribute it across two successive gates.
\[
\begin{tikzcd}[column sep=large]
 U\times P\arrow[r,"k"]\arrow[d,"{1\times !_P}"']&
 Q\times U\arrow[d,"{!_Q\times1}"]\\
 U\times1\arrow[r,"\simeq"']&1\times U
\end{tikzcd}
\]
and
\[
\begin{tikzcd}[row sep=large,column sep=small]
 U\times P\arrow[r,"k"]\arrow[d,"{1\times\Delta_P}"']&
 Q\times U\arrow[r,"{\Delta_Q\times1}"]&Q\times Q\times U\arrow[d,equal]\\
 U\times P\times P\arrow[r,"{k\times1}"']&
 Q\times U\times P\arrow[r,"{1\times k}"']&Q\times Q\times U .
\end{tikzcd}
\]
Define \((q,v)=k(u,p)\), \((q',v')=k(v,p)\). Evaluating the two pastings at an arbitrary generalized element gives
\[
\begin{aligned}
 \text{unit:}\quad &u=v,\\
 \text{multiplication:}\quad &(q,q,v)=(q,q',v').
\end{aligned}
\]
The first equation implies the second by substitution into \(k\). The tree induction in \cref{thm:cell-ext} supplies the carrier equation. In \cref{ex:bit-obstruction}, \(v=1-u\). Evaluating the unit equation at \(u=0\) refutes extension even though the one-gate carrier equation holds.

\section{A projective-tree chart with unique witnesses}\label{app:treechart}


\begin{enumerate}[titledsteps]
    \item \textbf{Write the interval recurrences.} Here is a finite tree chart with the same uniqueness property. Number a nonempty sentence $1,\ldots,n$. Use a single artificial-root arc and the projective convention that each subtree occupies an interval. Let $T(l,h;k)$ count the total weight of subtrees on $[l,h]$ rooted at $k$. Let $F_r(l,h)$ be the weight of an ordered forest on an interval entirely on one side of its exterior parent $r$, with each component attached directly to $r$. Set $F_r(l,l-1)=1$. Compute
\begin{align*}
 T(l,h;k)&=F_k(l,k-1)F_k(k+1,h),\\
 F_r(l,h)&=\sum_{t=l}^h\sum_{k=l}^t T(l,t;k)w_{r k}F_r(t+1,h),\\
 Z&=\sum_k w_{0k}\,T(1,n;k).
\end{align*}

    \item \textbf{Recover the tree from its first component.} Before assigning weights, let \(\mathcal T(l,h;k)\) and \(\mathcal F_r(l,h)\) be the corresponding sets of trees and attached forests. Projectivity makes each component an interval. For \(l\le h\), cutting at the root and cutting after the leftmost component give the bijections
\[
\begin{aligned}
 \mathcal T(l,h;k)&\cong\mathcal F_k(l,k-1)\times\mathcal F_k(k+1,h),\\
 \mathcal F_r(l,h)&\cong
 \coprod_{l\le t\le h}\ \coprod_{l\le k\le t}
   \mathcal T(l,t;k)\times\{r\to k\}\times\mathcal F_r(t+1,h).
\end{aligned}
\]
The inverse of the second map attaches the first component at \(k\) and appends the remaining ordered forest. Its first interval ends at \(t\), so recutting recovers both indices uniquely. The empty interval has one empty forest. Thus products of arc weights transport these bijections to the displayed recurrences. Evaluate by increasing length, \(T\) before \(F\) at a given length. For a single root child sum the root arc times its subtree. For a boundary root allowing several children use \(F_0(1,n)\), the convention of the cited inside algorithm. For a nonzero partition function with its inverse, the expected multiplicity at arc \(r\to k\) is \(w_{rk}(\partial Z/\partial w_{rk})/Z\), obtained by reversing the shared chart.
\end{enumerate}

\section{Compiling the angular coupling tables}\label{app:angular}


\begin{enumerate}[titledsteps]
    \item \textbf{Construct the trace-free projector.} Let \(V\) be the three-coordinate space. Equality of coordinate keys gives \(\delta\), and signed permutations give the alternating volume tensor. Write \(K_l\) for the rational matrix whose rows impose adjacent transpositions \(x-\tau_i x=0\) and all trace equations on \(V^{\otimes l}\). Thus \(H_l=\ker K_l\). Row-reduce \(K_l\) over \(\Q\) and put a basis of its kernel in the columns of \(B_l\). With the standard word-coordinate pairing,
\[
 \STF_l=B_l(B_l^{\mathsf T}B_l)^{-1}B_l^{\mathsf T},\qquad
 \STF_l^2=\STF_l,\quad \operatorname{im}\STF_l=H_l .
\]
The Gram matrix is invertible because \(a^{\mathsf T}B_l^{\mathsf T}B_l a=\|B_la\|^2\) is positive for a nonzero rational column \(a\). If the kernel is zero use the zero projection. The finite rational row operations give the same kernel basis after scalar extension. Rotations preserve the symmetry and trace equations and the pairing, so they preserve both \(H_l\) and its orthogonal complement and commute with \(\STF_l\).


    \item \textbf{Compute the coupling tables.} To generate the couplings $C:H_k\otimes H_J\to H_l$, one can use the following finite rational algorithm, including for degrees where the coupling space is zero. For a formal three-vector $u$ use the reflection
\[
   H(u)x=x-2u\,\langle u,x\rangle/\langle u,u\rangle .
\]
In the chosen rational bases, write the intertwining equation
\[
 C\bigl(g^{\otimes k}x,\ g^{\otimes J}y\bigr)
       =g^{\otimes l} C(x,y),\qquad g=H(u)H(v),
\]
on every pair of basis inputs. Write the resulting matrix difference as
\[
 A(u,v;C)=C(\rho_k(g)\otimes\rho_J(g))-\rho_l(g)C,\qquad
 q(u,v)=\langle u,u\rangle\langle v,v\rangle .
\]
Here \(\rho_j(g)\) is the action on \(H_j\). With \(N=\max(k+J,l)\), all entries of \(q^N A\) are polynomials: each occurrence of \(g=H(u)H(v)\) contributes at most one factor \(q\) to the denominator. Stack the coefficients of \(q^NA\) as rows of a finite rational linear system \(K\,\mathrm{vec}(C)=0\). Rational elimination produces a basis matrix for \(\ker K\), hence a finite keyed contraction table for each basis element. Row reduction uses only rational presentation data and the same pivot units after extension to the real register.


    \item \textbf{Prove completeness on the real register.} Here is a completeness argument on the real register, so no choice of tabulated angular coefficients is hidden in the construction. Each reflection preserves the metric and reverses the alternating volume: expand the volume on the three reflected basis vectors; terms with two copies of $u$ vanish, and the one-copy corrections sum to $-2\langle u,u\rangle/\langle u,u\rangle$. Every real isometry is a product of these reflections constructively. To fix a basis vector $e$ whose current unit image is $a$, the squared norms of $a-e$ and $a+e$ sum to four. Search for one exceeding one. Reflection in $a-e$ sends $a$ to $e$; reflection in $a+e$ sends it to $-e$, followed by reflection in $e$. Earlier fixed basis vectors are orthogonal to the chosen normals and stay fixed. Induction exhausts the three keys; a volume-preserving product has even parity. Thus all rotations are products of pairs tested by the compiler. Conversely any real intertwiner satisfies the cleared equations on all pairs of nonzero rational normals. Evaluate on a rectangular grid of strictly positive integers with more entries in each coordinate than the degree in that coordinate. Repeated univariate interpolation forces every coefficient to vanish. Rational elimination expresses every real solution in the compiled basis. This proves completeness for real multilinear channels of the indicated types.
\end{enumerate}

\section{Effective limits and their variations}\label{app:limits}

\subsection{Tail and local variation of the balanced scaling}


\begin{enumerate}[titledsteps]
    \item \textbf{Contract the normalized sweep.} Choose rational bounds \(0<a<K_{ij}<b\), fix an anchor \(o\) and normalize the full sweep \(F\) by \(T(c)=F(c)/F(c)_o\). For \(r_j=c_j/d_j\) put \(m=\min r_j,\ M=\max r_j,\ \Delta(c,d)=M/m-1\). With \(\xi=a/b\), split the probabilities in a row of \(A\) by
\[
 \omega_{ij}=\frac{K_{ij}d_j}{\sum_kK_{ik}d_k}
           =\xi\,\frac{d_j}{\sum_k d_k}+(1-\xi)\widetilde\omega_{ij}.
\]
The rational bounds imply \(\widetilde\omega_{ij}\ge0\) and \(\sum_j\widetilde\omega_{ij}=1\). The common first term cancels between rows. Hence the spread of \((Ac)_i/(Ad)_i=\sum_j\omega_{ij}r_j\) is at most \((1-\xi)(M-m)\), while its minimum is at least \(m\). Thus \(\Delta(Ac,Ad)\le(1-\xi)\Delta(c,d)\). The transpose lookup has the identical bound. Common coordinatewise multiplication by a positive vector, simultaneous inversion and separate positive scalar rescalings each preserve \(\Delta\). For the two grouped lookups in a sweep one obtains
\[
 \Delta(Tc,Td)\le q\,\Delta(c,d),\qquad q=(1-a/b)^2<1 .
\]

    \item \textbf{Compute the tail and identify the fixed point.} Comparison of the two column sums in a sweep gives, for every \(j\),
\[
 \frac ab\frac{\nu_j}{\nu_o}\le T(c)_j\le\frac ba\frac{\nu_j}{\nu_o}.
\]
Choose \(C>1\) rational above all upper bounds. For bounded anchored \(c,d\), \(m\le1\le M\) implies \(\|c-d\|_\infty\le C\Delta(c,d)\). Starting with \(c_0=\mathbf1\), put \(c_{n+1}=T(c_n)\). Telescoping \(\|c_{n+1}-c_n\|_\infty\le Cq^n\Delta(c_1,c_0)\) gives the explicit tail \(Cq^N\Delta(c_1,c_0)/(1-q)\). The positive lower bound survives at the limit \(c=Tc\). Hence \(F(c)=\lambda c\) with \(\lambda=F(c)_o>0\). For \(r=\mu/(Ac)\) the computed table gives
\[
 \sum_e P_e=\sum_i\mu_i=\lambda^{-1}\sum_j\nu_j .
\]
The common positive target total forces \(\lambda=1\), proving both marginals. For another anchored fixed point \(d\), \(\Delta(c,d)\le q\Delta(c,d)\) gives \(\Delta(c,d)=0\), hence \(c=d\). All rational bounds used in the tail can be found from the located inputs.


    \item \textbf{Control the first variation.} For completeness the square-zero tangent is also the ordinary derivative of the selected real solution. Nearby weights have common strict bounds, hence a common $q<1$. Put $c,c'$ for the two anchored fixed points and $\eta=\Delta(T'c,Tc)$. The finite rational sweep gives $\eta=O(\|w'-w\|)$ locally. Since $1+\Delta(u,z)\le(1+\Delta(u,v))(1+\Delta(v,z))$, the quantity $D=\Delta(c',c)$ satisfies $D\le (1+qD)(1+\eta)-1$. Shrinking the numerical neighbourhood gives $D\le\eta/[1-q(1+\eta)]$. Thus both scales vary locally by $O(\|w'-w\|)$. Expanding the finite products in their defining equations leaves a quadratic error; the bounded finite deletion at the base point sends that error to a quadratic correction. This identifies the derivative with the marked solve above.


    \item \textbf{Identify the higher jets.} For higher orders take a realized weight arc and the polynomial scale approximants $(\tilde r(t),\tilde c(t))$ produced by \cref{prop:balanced-lift} through degree $m$. They are positive near the retained point and $\tilde c_o=1$. Substitution in the row and column equations has residual $o(|t|^m)$, since all coefficients through $m$ vanish. On this neighbourhood the finite divisions have denominators bounded away from zero, hence the anchored sweep at the perturbed weights satisfies $\eta(t)=\Delta(T_t\tilde c(t),\tilde c(t))=o(|t|^m)$. For the exact fixed point $c(t)$ use the same multiplicative triangle inequality to obtain
\[
 \Delta(c(t),\tilde c(t))
 \le \frac{\eta(t)}{1-q(1+\eta(t))}.
\]
The denominator has a positive lower bound for small $t$. Bounded anchored coordinates turn this into $o(|t|^m)$ error in the column scales. The row equation gives the same accuracy for the row scales. Thus the lifts compute the Taylor coefficients of the selected solution. On the union of the pivot-unit opens, their equality on intersections is also descent of the coefficient functions by \cref{prop:affine-descent}.
\end{enumerate}

\subsection{Constructing the starting scale on a sparse support}\label{sec:sparse-solve}


\begin{enumerate}[titledsteps]
    \item \textbf{Find a strictly positive feasible flow.} For a connected finite bipartite support orient each edge from row to column. Write \(D:\R^E\to\R^{V\setminus\{o\}}\) for the signed incidence with one anchor row removed, and \(\widehat b\) for \(b=(\mu,-\nu)\) with that entry removed. Equal total makes \(Dp=\widehat b\) equivalent to \(\partial p=b\). Choose a spanning tree \(T\). Let \(J=E\setminus T\). Leaf deletion inverts \(D_T\), so every marginal solution has the form
\[
 p_J=y,\qquad p_T=D_T^{-1}(\widehat b-D_J y).
\]
Enumerate rational \(y\) and dovetail the strict positivity tests on both blocks. Given a strictly positive feasible flow, positivity persists in a neighborhood of its finitely many off-tree coordinates, which contains a rational trial. Thus search returns a located feasible \(p\) precisely on the image of the positive edge cone, within the equal-total base. On the sparse support \((1,1),(1,2),(2,2)\) with unit budgets the row and column equations force \(p_{12}=0\). For disconnected supports impose the already necessary equalities of totals componentwise and run the same solve on each nontrivial component. An isolated vertex cannot carry a positive budget.


    \item \textbf{Bound the scaling objective.} Extend vertex potentials by \(z_o=0\). The edge difference is \((D^{\mathsf T}z)_e=z_{\pi e}-z_{\tau e}\). The convex scaling objective of \cite{rothblum-scale}, its gradient and Hessian are
\[
\begin{aligned}
 P_e(z)&=w_e\exp((D^{\mathsf T}z)_e),&
 \Phi(z)&=\sum_e P_e(z)-\langle\widehat b,z\rangle,\\
 \nabla\Phi(z)&=DP(z)-\widehat b,&
 \nabla^2\Phi(z)&=D\,\operatorname{diag}(P(z))D^{\mathsf T}.
\end{aligned}
\]
Choose positive rationals \(l<\min w_e,\ h<\min p_e,\ H>\max p_e\). Put \(C=(H+h)^2/(2l)\). For \(t\ge0\), \(w_e e^t-p_et-h|t|\ge lt^2/2-(H+h)t\ge-C\). For \(t\le0\), \(w_e e^t-p_et\ge h|t|\). The bound extends from rational to located \(t\) by continuity. Writing \(t_e=(D^{\mathsf T}z)_e\), the path to the anchor and \(Dp=\widehat b\) give
\[
 \|z\|_\infty\le\sum_e|t_e|,\qquad
 \Phi(z)=\sum_e(w_e e^{t_e}-p_et_e)\ge h\|z\|_\infty-|E|C .
\]

    \item \textbf{Compute its minimizer.} Choose \(R>0\) rational with \(hR-|E|C>\Phi(0)+1\). On the cube of radius \(R\), write \(d=|V|-1\). Cauchy--Schwarz along each tree path gives
\(\|v\|_2^2\le d^2\sum_e((D^{\mathsf T}v)_e)^2\). Thus any rational \(0<\kappa<le^{-2R}/d^2\) bounds the Hessian below by \(\kappa\,\id\). Take \(L>2e^{2R}\sum_e w_e+\sum_v|b_v|\) as a rational Lipschitz bound in the sup norm. Approximate \(\Phi\) to error \(\epsilon/4\) on a grid of radius \(<\epsilon/(4L)\) and choose the least rational approximation. Its point \(z_\epsilon\) satisfies \(\Phi(z_\epsilon)\le\Phi(y)+\epsilon\) everywhere on the cube.

For \(\epsilon_n=\kappa 2^{-2n-4}\), midpoint strong convexity as in \cref{thm:mean-inverse} gives

\[
 \|z_n-z_m\|_2^2\le4(\epsilon_n+\epsilon_m)/\kappa.
\]
 The limit minimizes \(\Phi\), lies in the interior by coercivity, and satisfies \(DP(z)=\widehat b\). Its scales are $r_i=\exp z_i,\ c_j=\exp(-z_j)$. In the notation of \cref{thm:global-balance-inverse}, the anchored potential is $\phi=-z$. Conversely positive anchored scales have logarithms and give a critical point. To compute $\log x$ for $x>0$, search for $n$ with $e^{-n}<x<e^n$ and start with the bracket $[-n,n]$ for $\log x$. At its trisection points $a<b'$ run the tests \(e^a<x\) and \(x<e^{b'}\) in parallel. Locatedness and \(e^a<e^{b'}\) ensure a hit. On the first success retain the interval from $a$ to the old upper endpoint, and on the second retain the interval from the old lower endpoint to $b'$. Each interval has two thirds of the previous width. Continuity and strict increase of the exponential identify the limit. Convexity makes a critical point minimizing. Coercivity places it in the cube and strong convexity gives uniqueness. A positive scaling supplies its own feasible flow, so graph projection is invertible on exactly the positive-feasibility image over located affinities.


    \item \textbf{Control the lifted scales.} Near a feasible point the same strictly feasible off-tree trial and strict bounds work on a neighbourhood in the corresponding total-mass base. Enlarge the cube to contain the scale approximants there and choose common strict $l,\kappa$. The Hessian bound on the segment from an approximate anchored log scale $\tilde z$ to the exact $z$ gives $\|\tilde z-z\|_2\le\|\nabla\Phi(\tilde z)\|_2/\kappa$ by pairing their difference with the gradient difference. Both this bound and the midpoint inequality follow by integration of the displayed Hessian along a segment; the finite exponential sums and their derivatives have uniformly convergent series on the bounded segment. For a lifted scale polynomial through degree $m$, take logarithms of its positive values near the base point. Its gradient residual is exactly its nonanchor marginal residual, hence $o(|t|^m)$ along a realized order-$m$ input arc. Re-exponentiation on the bounded cube preserves this estimate. This identifies the formal sparse lift with the jet of the computed section.
\end{enumerate}

\subsection{First variation at a selected equilibrium}

Put \(\Delta s=s'-s\), \(\Delta z=z(s')-z(s)\), \(A=D_sF\), \(H=D_zF\) at the base solution. With the common contraction bound,
\[
\begin{aligned}
 (1-\lambda)\|\Delta z\|_\infty
    &\le\|F(s',z(s))-F(s,z(s))\|_\infty=O(\|\Delta s\|),\\
 (1-H)\Delta z&=A\Delta s+o(\|\Delta s\|),\\
 Dz(s)&=\sum_{j\ge0}H^jA,\qquad
 Dz(s)^*g=A^*\sum_{j\ge0}(H^*)^jg .
\end{aligned}
\]
The first line bounds the joint Taylor remainder by \(o(\|\Delta s\|)\) in the second. The Neumann inverse has norm at most \((1-\lambda)^{-1}\), preserving this remainder. Taking adjoints at each finite partial sum and passing through the geometric tail bound proves the last identity.

\subsection{Higher coefficients at an equilibrium}
For each order compiled for the finite body in \cref{sec:equilibrium}, let $z_*$ be the selected real fixed point and $H$ its hidden one-mark action there. Given an external arc $s(t)$, define a candidate $z_{<m}(t)=\sum_{j<m}z_jt^j$ with $z_0=z_*$. In the coefficient of degree $m$ of $F(s(t),z_{<m}(t)+vt^m)$ the only contributions involving $v$ come from one hidden mark with all other degrees zero. Consequently that coefficient is $b_m+Hv$, where $b_m=[t^m]F(s(t),z_{<m}(t))$. The equation
\[
             z_m=b_m+H z_m
\]
has unique solution $\sum_{j\ge0}H^j b_m$, with tail at most $\lambda^N\|b_m\|_\infty/(1-\lambda)$ after $N$ terms. The coefficients at earlier orders do not change. This proves compatible lifting through all $B_m$ at the real base with the displayed bounded inverse (and coefficientwise for finite free real perturbation registers). For a realized smooth input arc, let $z_{\le m}(t)$ be the resulting polynomial. Taylor's formula for the finite gates implies its residual in the fixed-point equation is $o(|t|^m)$. The contraction in the hidden coordinate bounds the distance to the actual fixed point by that residual divided by $1-\lambda$. Hence these are jets of the selected equilibrium. For mixed increments use the quotient of the $m$th power of the increment ideal by its $(m+1)$st power. Multiplication by any increment vanishes on this quotient, so each degree-$m$ monomial obeys the same equation with linear part $H$. Its inverse acts coefficientwise. Centred substitution respects the filtration, and uniqueness in each degree makes the construction commute with reindexing and truncation. A general cyclic wiring such as $z=z$ does not select a point or give this bound.

\section{Specializing and composing presented field programs}\label{app:field-presentations}

\subsection{A finite shape gives scalar terms}
Fix a finite presented shape \(A\). For the independent-field language of \cref{sec:field-language}, choose an enumeration of each key context \(U(A)\) and set
\[
 e_A[U:t]=\prod_{u\in U(A)}t .
\]
An empty product is the terminal scalar context. Use the scalar term theory \(\mathcal T\), with names for the entries of any fixed background table used by the program and its declared relations interpreted. At each output coordinate, lookup is a projection, a pointwise operation is its scalar symbol, and \(\Sigma_f,\Pi_f\) are the finite sums and products in \cref{sec:field-language}. These assignments extend by pairing and substitution to \(e_A:\mathbb F\to\mathcal T\) on this language. The coordinate equations prove functoriality. Iterated reductions flatten by decomposing a composite fibre into its intermediate fibres. A pullback along a key induces a bijection on the corresponding fibres. For distributivity, expanding a product of finite sums indexes the result by one choice of summand in each factor, with the empty choice contributing the empty product. These identities and the scalar laws validate the arithmetic equations used here. Changing an enumeration conjugates the terms by the induced product permutations, since the reductions are commutative. The same coordinate bijections prove naturality under isomorphisms of shapes, with composition following from composition of the permutations. Interpreting the scalar terms in \(S\) gives \(M_A\). Related input attributes are evaluated on the \hyperref[sec:legal-attribute-tables]{legal table object}.

Parameter extension commutes with this specialization. Working with chosen products and strict maps of presentations, its pushout is
\[
\begin{tikzcd}
 \mathbb F\arrow[r,"J"]\arrow[d,"e_A"']&
       \mathbb F_\Theta\arrow[d]\\
 \mathcal T\arrow[r,"J"']&\mathcal T_{e_A\Theta}.
\end{tikzcd}
\]
Indeed a compatible pair of interpretations into a target consists of an interpretation \(L\) of \(\mathcal T\) and an extension of \(Le_A\) by a point of \(L(e_A\Theta)\). By \cref{prop:relative-models} this is exactly one interpretation of the lower right corner extending \(L\); the parameter-point condition there also gives the comparison on transformations. With products preserved up to their canonical isomorphisms this describes the corresponding universal equivalence. The right-hand functor takes a representative \(f:\Theta\times U\to V\) to \(e_A f\). Induction on the register list identifies its value at each named gate with that gate's scalar expansion.

\subsection{Gluing along a named boundary}
For programs \(N:X\to Y\) and \(K:Y\to Z\) over the same field theory, use disjoint parameter names \(\Theta,\Lambda\). Let \(B=\mathbb F[b:X\to Y]\). Interpret \(B\) in \(|P_N|\) by \(b\mapsto F_N\); on the other side include \(B\) in \(|P_K|[b:X\to Y]\). Form the presentation containing both parameter names, both lists of register definitions and the equation \(b=F_N\). This is the pushout of the two interpretations: a pair of interpretations agreeing on \(B\) gives values for every generator, each defining equation holds in the target, and substitution yields the unique extension. Restriction recovers the pair. Transformations agreeing on the boundary extend by naturality on those same generators. In this pushout \(F_K b\) selects the output and, writing \(F_N^\circ,F_K^\circ\) for the parameter-explicit representatives, it denotes
\[
  (p,q,x)\longmapsto F_K^\circ(q,F_N^\circ(p,x)).
\]
The two parameters classify independent choices. Pullback of values along the diagonal \(p\mapsto(p,p)\) gives the shared-table interpretation when the two table sorts coincide. More generally, a base term \(k:\Lambda'\to\Theta\) induces \(\mathbb F_\Theta\to\mathbb F_{\Lambda'}\) by precomposing representatives with \(k\times1\). The composition formula of \cref{sec:model-interfaces} gives preservation of composition and products immediately, and induction shows that all designated outputs and cuts specialize along it. Parameter families over a test object are interpreted in its slice.

\subsection{Coherent interpretation of a composite presentation}
Take the underlying many-sorted equational presentation of any of the field or interface theories above. In its coherent syntactic category objects are coherent formulas in context modulo provable equivalence. Arrows are equivalence classes of provably total, single-valued relations. A term interpretation translates equality, finite conjunction and disjunction, and existential quantification by substitution, renaming bound variables. It carries the proofs of totality and functionality to such proofs in the target. Conjunction followed by existential quantification composes relations, and the equality relation gives the identity. This constructs a coherent functor on syntactic categories.

If \(j:T_0\to T_1\) and \(k:T_1\to T_2\) present successive calls, the translation of a formula along \(kj\) is the successive translation along \(j\) and \(k\), up to renaming bound variables. For a total functional relation this identifies the two translated graphs, hence their interpreted arrows. Using \cref{prop:coherent-port}, the reductions on models therefore compose in the order \(j^*k^*\), including after a parameter specialization. Choosing \(M_A\) gives the concrete finite reductions used in execution. Their dependence on the field signature is exhibited in the \hyperref[sec:coherent-row-language]{row-language comparison}.

\clearpage
\phantomsection
\addcontentsline{toc}{section}{References}
\begin{thebibliography}{99}
\setlength{\itemsep}{1.2pt plus .3pt minus .3pt}
\bibitem{hyland} J. M. E. Hyland. The effective topos. In A. S. Troelstra and D. van Dalen (eds.), \emph{The L. E. J. Brouwer Centenary Symposium}, Studies in Logic and the Foundations of Mathematics 110, 165--216. North-Holland, 1982. \url{https://doi.org/10.1016/S0049-237X(09)70129-6}.

\bibitem{streicher-logic} T. Streicher. \emph{Introduction to Category Theory and Categorical Logic}, discussion of functional relations and unique choice. \url{https://www2.mathematik.tu-darmstadt.de/~streicher/CTCL.pdf}.
\bibitem{rutten-systems} J. J. M. M. Rutten. Universal coalgebra: a theory of systems. \emph{Theoretical Computer Science} 249 (2000), 3--80. \url{https://doi.org/10.1016/S0304-3975(00)00056-6}.
\bibitem{ground-psr} M.~L. Littman, R.~S. Sutton and S.~Singh. Predictive representations of state. \emph{Advances in Neural Information Processing Systems} 14 (2001), 1555--1561.
\bibitem{rota} G.-C. Rota. On the foundations of combinatorial theory. I. Theory of M\"obius functions. \emph{Zeitschrift f\"ur Wahrscheinlichkeitstheorie und Verwandte Gebiete} 2 (1964), 340--368. \url{https://doi.org/10.1007/BF00531932}.
\bibitem{moore} E. F. Moore. Gedanken-experiments on sequential machines. In C. E. Shannon and J. McCarthy (eds.), \emph{Automata Studies}. Princeton University Press, 1956. \url{https://doi.org/10.1515/9781400882618-006}.
\bibitem{ground-common-information} P.~G\'acs and J.~K\"orner. Common information is far less than mutual information. \emph{Problems of Control and Information Theory} 2 (1973), 149--162.
\bibitem{van-oosten} Jaap van Oosten. \emph{Realizability: An Introduction to Its Categorical Side}. Studies in Logic and the Foundations of Mathematics 152. Elsevier, 2008.

\bibitem{bauer} Andrej Bauer. On fixed-point theorems in synthetic computability. \emph{Tbilisi Mathematical Journal} 10(3), 167--181, 2017. \url{https://math.andrej.com/asset/data/recursion-theorem.pdf}.

\bibitem{gambino-hyland} Nicola Gambino and Martin Hyland. Wellfounded trees and dependent polynomial functors. In \emph{Types for Proofs and Programs}, Lecture Notes in Computer Science 3085, 210--225. Springer, 2004. \url{https://eprints.whiterose.ac.uk/id/eprint/113158/}.

\bibitem{pauly-represented} A. Pauly. On the topological aspects of the theory of represented spaces. \emph{Computability} 5 (2016), 159--180. \url{https://arxiv.org/abs/1204.3763}.
\bibitem{hoyrup-represented} M. Hoyrup. Results in descriptive set theory on some represented spaces, \S2. \url{https://arxiv.org/abs/1712.03680}.
\bibitem{rothblum-scale} U. G. Rothblum and H. Schneider. Scalings of matrices which have prespecified row sums and column sums via optimization. \emph{Linear Algebra Appl.} 114--115 (1989), 737--764. \url{https://doi.org/10.1016/0024-3795(89)90491-6}.

\bibitem{vickers} Steven Vickers. A localic theory of lower and upper integrals. \emph{Mathematical Logic Quarterly} 54(1), 109--123, 2008. \url{https://doi.org/10.1002/malq.200710028}.

\bibitem{algebraic-databases}
Patrick Schultz, David I. Spivak, Christina Vasilakopoulou, and Ryan Wisnesky.
\newblock Algebraic databases.
\newblock \textbf{Theory and Applications of Categories}, 32(16):547--619, 2017.
\newblock \url{https://tac.mta.ca/tac/volumes/32/16/32-16abs.html}.

\bibitem{algebraic-databases-errata}
Patrick Schultz, David I. Spivak, Christina Vasilakopoulou, and Ryan Wisnesky.
\newblock Algebraic databases, revised arXiv version 3, January 2025, Appendix B (Errata).
\newblock \url{https://arxiv.org/abs/1602.03501}.

\bibitem{green-provenance}
Todd J. Green, Grigoris Karvounarakis, and Val Tannen.
\newblock Provenance semirings.
\newblock In \textbf{Proceedings of PODS 2007}, pages 31--40, 2007.
\newblock \url{https://doi.org/10.1145/1265530.1265535}.

\bibitem{poly:gambino-kock}
Nicola Gambino and Joachim Kock.
\newblock Polynomial functors and polynomial monads.
\newblock \emph{Mathematical Proceedings of the Cambridge Philosophical Society},
154(1):153--192, 2013.
\newblock \url{https://doi.org/10.1017/S0305004112000394}.

\bibitem{poly:containers}
Michael Abbott, Thorsten Altenkirch, and Neil Ghani.
\newblock Containers: Constructing strictly positive types.
\newblock \emph{Theoretical Computer Science}, 342(1):3--27, 2005.
\newblock \url{https://doi.org/10.1016/j.tcs.2005.06.002}.

\bibitem{vaswani-attention}
Ashish Vaswani, Noam Shazeer, Niki Parmar, Jakob Uszkoreit,
Llion Jones, Aidan N. Gomez, \L ukasz Kaiser, and Illia Polosukhin.
\newblock Attention is all you need.
\newblock In \textbf{Advances in Neural Information Processing Systems 30},
2017.
\newblock \url{https://arxiv.org/abs/1706.03762}.

\bibitem{poly:container-derivative}
Michael Abbott, Thorsten Altenkirch, Neil Ghani, and Conor McBride.
\newblock Derivatives of containers.
\newblock In \emph{Typed Lambda Calculi and Applications (TLCA 2003)},
\emph{Lecture Notes in Computer Science} 2701, pages 16--30, Springer, 2003.
\newblock \url{https://people.cs.nott.ac.uk/psztxa/publ/tlca03.pdf}.

\bibitem{vojta-jets}
P. Vojta. \emph{Jets via Hasse--Schmidt derivations}.
\url{https://arxiv.org/abs/math/0407113}.

\bibitem{poly:stacks-differentials}
The Stacks Project Authors.
\newblock Differentials, Section 10.131, Tag 00RM.
\newblock \emph{The Stacks Project}.
\newblock \url{https://stacks.math.columbia.edu/tag/00RM}.

\bibitem{jacobs-fibrations} B. Jacobs.
\emph{Categorical Logic and Type Theory}. Studies in Logic and the Foundations of Mathematics 141, North-Holland, 1999.
\url{https://www.cs.ru.nl/B.Jacobs/CLT/bookinfo.html}.
\bibitem{pratt-chu} V. Pratt. \emph{Chu Spaces}. Notes for the School on Category Theory and Applications, University of Coimbra, 1999.
\href{https://ncatlab.org/nlab/files/Pratt-ChuSpaces.pdf}{Lecture notes}, Chapter~1.

\bibitem{neural:fong}
B. Fong, D. I. Spivak, and R. Tuy\'eras.
\emph{Backprop as Functor: A Compositional Perspective on Supervised Learning}.
LICS, 2019. \url{https://arxiv.org/abs/1711.10455}.

\bibitem{cdl}
B. Gavranovi\'c, P. Lessard, A. Dudzik, T. von Glehn, J. G. M. Ara\'ujo, and P. Veli\v{c}kovi\'c.
Position: Categorical Deep Learning is an Algebraic Theory of All Architectures.
\emph{ICML}, PMLR 235, 15209--15241, 2024.
\url{https://proceedings.mlr.press/v235/gavranovic24a.html}.

\bibitem{orbit:sharing}
S. Ravanbakhsh, J. Schneider, and B. P\'oczos.
Equivariance through parameter-sharing. \emph{ICML}, PMLR 70, 2892--2901, 2017.
\url{https://proceedings.mlr.press/v70/ravanbakhsh17a.html}.

\bibitem{graph:dp}
A. J. Dudzik and P. Veli\v{c}kovi\'c. Graph neural networks are dynamic programmers.
\emph{NeurIPS}, 2022. \url{https://arxiv.org/abs/2203.15544}.

\bibitem{graph:cokleisli}
B. Gavranovi\'c and M. Villani. Graph convolutional neural networks as parametric coKleisli morphisms. 2022. \url{https://arxiv.org/abs/2212.00542}.

\bibitem{learn:languages}
D. I. Spivak. Learners' languages. \emph{EPTCS} 372, 14--28, 2022.
\url{https://arxiv.org/abs/2103.01189}.

\bibitem{learn:johnson}
B. Fong and M. Johnson. Lenses and learners. 2019.
\url{https://arxiv.org/abs/1903.03671}.

\bibitem{gavranovic-thesis}
B. Gavranovi\'c. \emph{Fundamental Components of Deep Learning: A Category-Theoretic Approach}.
PhD thesis, University of Strathclyde, 2024. \url{https://arxiv.org/abs/2403.13001}.
\bibitem{lens:deep}
G. S. H. Cruttwell, B. Gavranovi\'c, N. Ghani, P. Wilson, and F. Zanasi.
Deep learning with parametric lenses. 2024. \url{https://arxiv.org/abs/2404.00408}.

\bibitem{moments} Z. Li and J. Eisner. First- and Second-Order Expectation Semirings with Applications to Minimum-Risk Training on Translation Forests. EMNLP, 2009. \href{https://aclanthology.org/D09-1005/}{Paper}.
\bibitem{structured} Y. Kim, C. Denton, L. Hoang and A. M. Rush. Structured Attention Networks. ICLR, 2017. \href{https://arxiv.org/abs/1702.00887}{Paper}.
\bibitem{sink} M. E. Sander, P. Ablin, M. Blondel and G. Peyr\'e. Sinkformers: Transformers with Doubly Stochastic Attention. AISTATS, 2022. \href{https://proceedings.mlr.press/v151/sander22a.html}{Paper}.
\bibitem{se3} F. B. Fuchs, D. E. Worrall, V. Fischer and M. Welling. SE(3)-Transformers: 3D Roto-Translation Equivariant Attention Networks. NeurIPS, 2020. \href{https://arxiv.org/abs/2006.10503}{Paper and appendix}.
\bibitem{deq} S. Bai, J. Z. Kolter and V. Koltun. Deep Equilibrium Models. NeurIPS, 2019. \href{https://arxiv.org/abs/1909.01377}{Paper}.
\bibitem{switch} W. Fedus, B. Zoph and N. Shazeer. Switch Transformers: Scaling to Trillion Parameter Models with Simple and Efficient Sparsity. JMLR 23(120), 2022. \href{https://jmlr.org/papers/v23/21-0998.html}{Paper}.
\bibitem{etale} The Stacks Project Authors. Formally \'etale maps, Tag 00UP. \href{https://stacks.math.columbia.edu/tag/00UP}{Section}.

\bibitem{relative:tvq} B. To\"en and M. Vaqui\'e. Au-dessous de $\operatorname{Spec}\mathbb Z$. \emph{Journal of K-Theory} 3 (2009), 437--500. \url{https://arxiv.org/abs/math/0509684}.
\bibitem{relative:hag} B. To\"en and G. Vezzosi. Homotopical algebraic geometry II: Geometric stacks and applications. \emph{Memoirs of the American Mathematical Society} 193 (2008), no.~902. \url{https://arxiv.org/abs/math/0404373}.
\bibitem{stacks:parts} The Stacks Project Authors. Finite order differential operators, Tag 09CH. \url{https://stacks.math.columbia.edu/tag/09CH}.

\bibitem{nrp} A. Soeteman, B. ten Cate, M. Funk, B. Kimelfeld, C. Lutz and M. Sch\"onherr. \emph{Neuro-Relational Programs: Unifying Queries and Neural Computation over Structured Data} (2026). \href{https://arxiv.org/abs/2606.11946}{Paper}.
\bibitem{pugh-kan} M. Pugh, J. Grundy, C. Cirstea and N. Harris. \emph{Learning Is a Kan Extension}. Transactions on Machine Learning Research (2025). \href{https://openreview.net/forum?id=xWKtKdeefL}{Published article}.

\bibitem{pugh} M. Pugh, J. Grundy, C. Cirstea and N. Harris. \emph{Presenting Neural Networks via Coherent Functors} (2026). \href{https://arxiv.org/abs/2604.15100}{Preprint}.

\bibitem{lambek-scott} J. Lambek and P. J. Scott. \emph{Introduction to Higher Order Categorical Logic}. Cambridge University Press, 1986. Part I, Section 6.

\bibitem{corradini-gadducci} A. Corradini and F. Gadducci. An algebraic presentation of term graphs, via gs-monoidal categories. \emph{Applied Categorical Structures} 7 (1999), 299--331. \url{https://doi.org/10.1023/A:1008647417502}.

\bibitem{makkai-reyes} M. Makkai and G. E. Reyes. \emph{First Order Categorical Logic}. Lecture Notes in Mathematics 611, Springer, 1977.

\bibitem{zhang-bolcskei-mv} Y. Zhang and H. B\"olcskei. Complete identification of deep ReLU networks through \L{}ukasiewicz logic. arXiv:2602.00266v2, 2026. \url{https://arxiv.org/abs/2602.00266v2}.

\bibitem{gerla-dmv} B. Gerla. Rational \L{}ukasiewicz logic and DMV-algebras. \emph{Neural Network World} 11 (2001), 579--594. \url{https://arxiv.org/abs/1211.5485}.

\bibitem{pkhakadze-qe} N. Pkhakadze. \emph{Quantifier Elimination in Quantified Propositional \L{}ukasiewicz Logic}. Master's thesis, 2010, Theorems 5.3.28--5.3.29. \href{https://citeseerx.ist.psu.edu/document?doi=58b23a00fbfc599d6a070f0359b9577b71e0e3c2&repid=rep1&type=pdf}{Full text}.
\bibitem{gnn} D. Tena Cucala, B. Cuenca Grau, B. Motik and E. V. Kostylev. \emph{On the Correspondence Between Monotonic Max-Sum GNNs and Datalog}. KR (2023). \href{https://arxiv.org/abs/2305.18015}{Paper}.
\bibitem{wz} P. Wilson and F. Zanasi. \emph{An axiomatic approach to differentiation of polynomial circuits}. \emph{Journal of Logical and Algebraic Methods in Programming} 135 (2023), 100892. \href{https://doi.org/10.1016/j.jlamp.2023.100892}{DOI}.
\bibitem{graph-poly-basis} O. Puny, D. Lim, B. T. Kiani, H. Maron and Y. Lipman. \emph{Equivariant Polynomials for Graph Neural Networks}. ICML (2023). \href{https://proceedings.mlr.press/v202/puny23a.html}{Paper}.
\bibitem{motif-basis} R. Curticapean, H. Dell and D. Marx. \emph{Homomorphisms Are a Good Basis for Counting Small Subgraphs}. STOC (2017). \href{https://arxiv.org/abs/1705.01595}{Paper}.
\bibitem{counter} A. Dudenhefner. \emph{Certified Decision Procedures for Two-Counter Machines}. FSCD (2022). \href{https://doi.org/10.4230/LIPIcs.FSCD.2022.16}{DOI}.
\bibitem{aggregate-equivalence} S. Cohen, W. Nutt and Y. Sagiv. \emph{Deciding Equivalences among Conjunctive Aggregate Queries}. Journal of the ACM 54(2), 2007.
\href{https://www.inf.unibz.it/~nutt/Teaching/FDBs1718/FDBsPapers/cohenEtAl-JACM-07.pdf}{Paper}.
\bibitem{wj-mean} M. J. Wainwright and M. I. Jordan. \emph{Graphical Models, Exponential Families, and Variational Inference}. Foundations and Trends in Machine Learning 1(1--2), 2008, Theorem~3.3.
\href{https://people.eecs.berkeley.edu/~jordan/sail/readings/wainwright-jordan-fnt.pdf}{Paper}.
\bibitem{cil2p} A. S. d'Avila Garcez and G. Zaverucha. \emph{The Connectionist Inductive Learning and Logic Programming System}. Applied Intelligence 11, 59--77 (1999). \url{https://doi.org/10.1023/A:1008328630915}.
\bibitem{tensorlog} W. W. Cohen, F. Yang and K. R. Mazaitis. \emph{TensorLog: A Probabilistic Database Implemented Using Deep-Learning Infrastructure}. JAIR 67 (2020). \url{https://doi.org/10.1613/jair.1.11944}.
\bibitem{deepproblog} R. Manhaeve, S. Duman\v{c}i\'c, A. Kimmig, T. Demeester and L. De Raedt. \emph{DeepProbLog: Neural Probabilistic Logic Programming}. NeurIPS (2018). \href{https://proceedings.neurips.cc/paper_files/paper/2018/hash/dc5d637ed5e62c36ecb73b654b05ba2a-Abstract.html}{Paper}.
\bibitem{hard-attn-ltl} A. Yang, D. Chiang and D. Angluin. \emph{Masked Hard-Attention Transformers Recognize Exactly the Star-Free Languages}. NeurIPS (2024). \url{https://arxiv.org/abs/2310.13897}.
\bibitem{cq-learning} B. ten Cate and V. Dalmau. \emph{Conjunctive Queries: Unique Characterizations and Exact Learnability}. ICDT (2021). \url{https://doi.org/10.4230/LIPIcs.ICDT.2021.9}.
\bibitem{tracr} D. Lindner, J. Kram\'ar, S. Farquhar, M. Rahtz, T. McGrath and V. Mikulik. \emph{Tracr: Compiled Transformers as a Laboratory for Interpretability}. NeurIPS (2023). \url{https://arxiv.org/abs/2301.05062}.
\bibitem{interpbench} R. Gupta, I. Arcuschin, T. Kwa and A. Garriga-Alonso. \emph{InterpBench: Semi-Synthetic Transformers for Evaluating Mechanistic Interpretability Techniques}. NeurIPS (2024). \href{https://proceedings.neurips.cc/paper_files/paper/2024/hash/a8f7d43ae092d9a5295775eb17f3f4f7-Abstract-Datasets_and_Benchmarks_Track.html}{Paper}.
\bibitem{transformer-programs} D. Friedman, A. Wettig and D. Chen. \emph{Learning Transformer Programs}. NeurIPS (2023). \href{https://proceedings.neurips.cc/paper_files/paper/2023/hash/995f693b73050f90977ed2828202645c-Abstract-Conference.html}{Paper}.
\bibitem{huang-rasp} X. Huang, A. Bakalova, S. Bhattamishra, W. Merrill and M. Hahn. Discovering Interpretable Algorithms by Decompiling Transformers to RASP. \emph{ICML}, PMLR 306 (2026), 46309--46412. \href{https://proceedings.mlr.press/v306/huang26b.html}{Paper}.
\bibitem{concept-neural-views} J. Hirth and T. Hanika. \emph{Formal Conceptual Views in Neural Networks} (2022). \url{https://arxiv.org/abs/2209.13517}.
\bibitem{concept-completeness} C.-K. Yeh, B. Kim, S. Arik, C.-L. Li, T. Pfister and P. K. Ravikumar. \emph{On Completeness-aware Concept-Based Explanations in Deep Neural Networks}. NeurIPS 33 (2020). \href{https://proceedings.neurips.cc/paper/2020/hash/ecb287ff763c169694f682af52c1f309-Abstract.html}{Paper}.
\bibitem{fca-foundations} B. Ganter and R. Wille. \emph{Formal Concept Analysis: Mathematical Foundations}. Springer, 1999. \url{https://doi.org/10.1007/978-3-642-59830-2}.
\bibitem{isbell-distributors} L. Shen and D. Zhang. \emph{Categories enriched over a quantaloid: Isbell adjunctions and Kan adjunctions}. Theory and Applications of Categories 28 (2013), 577--615. \url{https://arxiv.org/abs/1307.5625}.
\bibitem{fuzzy-galois} R. B\v{e}lohl\'avek. \emph{Fuzzy Galois Connections}. Mathematical Logic Quarterly 45 (1999), 497--504. \url{https://doi.org/10.1002/malq.19990450408}.
\bibitem{multi-concept-lattices} H. Lai and L. Shen. \emph{Multi-adjoint concept lattices via quantaloid-enriched categories}. Fuzzy Sets and Systems 405 (2021), 74--87. \url{https://doi.org/10.1016/j.fss.2020.03.007}.
\bibitem{ganter-exploration} B. Ganter. \emph{Attribute exploration with background knowledge}. Theoretical Computer Science 217 (1999), 215--233. \url{https://doi.org/10.1016/S0304-3975(98)00271-0}.
\bibitem{diophantine-undecidability} Yu. V. Matiyasevich. \emph{Diophantine representation of enumerable predicates}. Mathematics of the USSR-Izvestiya 5 (1971), 1--28. \url{https://doi.org/10.1070/IM1971v005n01ABEH001004}.
\bibitem{causal-abstraction} A. Geiger, D. Ibeling, A. Zur, M. Chaudhary, S. Chauhan, J. Huang, A. Arora, Z. Wu, N. Goodman, C. Potts and T. Icard. \emph{Causal Abstraction: A Theoretical Foundation for Mechanistic Interpretability}. JMLR 26(83), 1--64 (2025). \href{https://www.jmlr.org/papers/v26/23-0058.html}{Paper}.
\bibitem{attn-identifiability} G. Brunner, Y. Liu, D. Pascual, O. Richter, M. Ciaramita and R. Wattenhofer. \emph{On Identifiability in Transformers}. ICLR (2020). \href{https://research.google/pubs/on-identifiability-in-transformers/}{Paper}.
\bibitem{reluplex} G. Katz, C. Barrett, D. Dill, K. Julian and M. Kochenderfer. \emph{Reluplex: An Efficient SMT Solver for Verifying Deep Neural Networks}. CAV (2017). \url{https://arxiv.org/abs/1702.01135}.
\bibitem{relu-geometry} L. Zhang, G. Naitzat and L.-H. Lim. \emph{Tropical Geometry of Deep Neural Networks}. ICML (2018). \href{https://proceedings.mlr.press/v80/zhang18i.html}{Paper}.

\end{thebibliography}
\end{document}
